% ============================================================================
% Fichier     : GCI_main.tex
% Titre       : IT Governance, COBIT et ITIL
% Auteur      : MINKA MI NGUIDJOÏ Thierry Emmanuel
% Institution : Chaire UNESCO d'Afrique Centrale Accès TIC — ENSP / LIMSI
% Version     : 0.1 -- octobre 2026 (squelette)
% ============================================================================

\documentclass[12pt, twoside, openright, a4paper]{book}

% --- Feuille de style unifiée ---
\usepackage{Style_GCI}

% --- Bibliographie ---
\addbibresource{bibliographie/GCI_references.bib}

% ============================================================================
% MÉTADONNÉES PDF
% ============================================================================
\hypersetup{
    pdftitle={IT Governance, COBIT et ITIL},
    pdfauthor={MINKA MI NGUIDJOÏ Thierry Emmanuel},
    pdfsubject={MSN-2S3-MSI03 · Master Professionnel en Management des Systèmes d'Information · Chaire UNESCO d'Afrique Centrale Accès TIC},
    pdfkeywords={gouvernance des SI, COBIT 2019, ITIL 4, ISO/IEC 38500, ISO/IEC 20000, GCI}
}

\begin{document}

% --------------------------------------------------------------------------
% ÉLÉMENTS PRÉLIMINAIRES (pages romaines)
% --------------------------------------------------------------------------
\frontmatter

    % ==============================================================================
% Fichier : front/00_page_titre.tex
% Titre   : Page de titre - IT Governance, COBIT et ITIL
% Auteur  : MINKA MI NGUIDJOÏ Thierry Emmanuel
% Version : 1.0 -- octobre 2026
% ==============================================================================

\begin{titlepage}
\thispagestyle{empty}

% --- Fond décoratif : bandeaux bleu marine et filets or ---
\begin{tikzpicture}[remember picture, overlay]
    \fill[CouleurPrimaire]
        (current page.north west) rectangle
        ([yshift=-3.6cm] current page.north east);
    \fill[CouleurAccent]
        ([yshift=-3.6cm] current page.north west) rectangle
        ([yshift=-4.0cm] current page.north east);
    \fill[CouleurPrimaire]
        (current page.south west) rectangle
        ([yshift=2.5cm] current page.south east);
    \fill[CouleurAccent]
        ([yshift=2.5cm] current page.south west) rectangle
        ([yshift=2.9cm] current page.south east);
    % En-tête institutionnel, posé dans le bandeau supérieur
    \node[align=center, text=white] at ([yshift=-1.8cm] current page.north) {%
        {\fontsize{11}{13}\selectfont
        \textsc{République du Cameroun - Paix, Travail, Patrie}}\\[0.3ex]
        {\color{CouleurAccent!60!white}\rule{0.4\paperwidth}{0.4pt}}\\[0.6ex]
        {\fontsize{10}{12}\selectfont\textbf{\institutionprimaire}}\\[0.2ex]
        {\fontsize{8}{10}\selectfont \tutelle}\\[0.2ex]
        {\fontsize{9}{11}\selectfont \unitepedago}};
    % Logos : ENSPY à gauche, Chaire UNESCO à droite (cours de la Chaire)
    \node[fill=white, rounded corners=4pt, inner sep=2pt]
        at ([xshift=1.75cm, yshift=-1.8cm] current page.north west)
        {\includegraphics[height=2.6cm]{figures/logos/enspy_logo.jpg}};
    \node[fill=white, rounded corners=4pt, inner sep=2pt]
        at ([xshift=-1.75cm, yshift=-1.8cm] current page.north east)
        {\includegraphics[height=2.6cm]{figures/logos/unesco_logo.jpeg}};
    % Année et version, posées dans le bandeau inférieur
    \node[align=center] at ([yshift=1.25cm] current page.south) {%
        {\color{white}\bfseries\large Année académique \anneescolaire}\\[0.3ex]
        {\color{CouleurAccent!70!white}\small
        Version 1.0 | Document pédagogique de l'UE \codeUE}};
\end{tikzpicture}

\vspace*{0.6cm}

% --- Titre principal (sur deux lignes) ---
\begin{center}
    {\color{CouleurPrimaire}\rule{0.85\textwidth}{1.2pt}}\\[1.2ex]
    {\fontsize{26}{32}\selectfont\bfseries\color{CouleurPrimaire}
        IT Governance\\[0.4ex]
        COBIT et ITIL}\\[0.8ex]
    {\color{CouleurAccent}\rule{0.85\textwidth}{2pt}}\\[0.5ex]
    {\color{CouleurPrimaire}\rule{0.85\textwidth}{0.5pt}}
\end{center}

\vspace{0.8cm}

\begin{center}
    {\fontsize{13}{16}\selectfont\itshape\color{CouleurSecondaire}
        Gouverner, piloter et servir~:\\
        la gouvernance des systèmes d'information en contexte africain}
\end{center}

\vspace{0.5cm}

% --- Cartouche sigle et code UE ---
\begin{center}
    \begin{tikzpicture}
        \node[draw=CouleurPrimaire, fill=CouleurPrimaire, rounded corners=6pt,
              inner xsep=14pt, inner ysep=6pt]{%
            \color{white}\bfseries\large \sigleCours{} - UE~\codeUE
        };
    \end{tikzpicture}
\end{center}

\vspace{1.0cm}

% --- Schéma de la discipline : la gouvernance (EDM) dirige le management,
%     qui lui rend compte ; ITIL outille le domaine DSS ---
\begin{center}
\begin{tikzpicture}[
    gouv/.style = {
        draw=CouleurPrimaire, fill=CouleurPrimaire, rounded corners=4pt,
        minimum width=11.4cm, minimum height=0.8cm,
        font=\small\bfseries\color{white}, inner sep=4pt, align=center
    },
    box/.style = {
        draw=CouleurBordDef, fill=CouleurDef, rounded corners=4pt,
        text width=2.45cm, minimum height=1.1cm,
        font=\scriptsize\bfseries\color{CouleurPrimaire},
        inner sep=3pt, align=center
    },
    arrow/.style = {-{Stealth[length=5pt]}, thick, color=CouleurSecondaire}
]
\node[gouv] (edm) at (0,0) {Gouvernance~: évaluer · diriger · surveiller (EDM)};
\node[box] (apo) at (-4.35,-2.0) {Aligner, planifier,\\ organiser (APO)};
\node[box] (bai) at (-1.45,-2.0) {Bâtir, acquérir,\\ implémenter (BAI)};
\node[box] (dss) at (1.45,-2.0)  {Délivrer, servir,\\ soutenir (DSS)};
\node[box] (mea) at (4.35,-2.0)  {Surveiller, évaluer,\\ mesurer (MEA)};
\draw[arrow] ([xshift=-2.5cm]edm.south) -- node[left, font=\scriptsize\itshape,
    color=CouleurSecondaire]{diriger} ([xshift=-2.5cm, yshift=-0.85cm]edm.south);
\draw[arrow] ([xshift=2.5cm, yshift=-0.85cm]edm.south) -- node[right,
    font=\scriptsize\itshape, color=CouleurSecondaire]{rendre compte}
    ([xshift=2.5cm]edm.south);
\draw[arrow] (apo) -- (bai);
\draw[arrow] (bai) -- (dss);
\draw[arrow] (dss) -- (mea);
\node[font=\scriptsize\bfseries\color{CouleurAccent}, below=0.15cm of dss]
    {ITIL 4};
\node[font=\scriptsize\itshape\color{CouleurSecondaire}, below=0.15cm of bai]
    {Management (COBIT 2019)};
\end{tikzpicture}
\end{center}

\vspace{0.6cm}

\begin{center}
    {\large\bfseries\color{CouleurPrimaire}
        \auteurcours}\\[0.5ex]
    {\normalsize\color{CouleurSecondaire}
       Enseignant Chercheur | \laboratoire{} | ENSPY}\\[0.2ex]
    {\small\color{CouleurBordInfo}
        \institutionprimaire}\\[0.2ex]
    {\small\color{CouleurBordInfo}
        \href{mailto:maletyon@proton.me}{maletyon@proton.me}}
\end{center}

\end{titlepage}

    \cleardoublepage

    % ==============================================================================
% Fichier : front/01_avant_propos.tex
% Titre   : Avant-propos - GCI
% Auteur  : MINKA MI NGUIDJOÏ Thierry Emmanuel
% Version : 1.0 -- octobre 2026
% ==============================================================================

\chapter*{Avant-propos}
\addcontentsline{toc}{chapter}{Avant-propos}
\markboth{Avant-propos}{Avant-propos}

\epigraph{%
    « Spécifier les droits de décision et le cadre de responsabilité qui
    encouragent les comportements souhaitables dans l'usage des technologies
    de l'information. »%
}{Peter \textsc{Weill} et Jeanne \textsc{Ross}, définition de la
gouvernance des SI (2004, notre traduction)}

\vspace{0.5ex}

Ce cours repose sur une double conviction. La première est que la
gouvernance des systèmes d'information n'est plus une affaire de
spécialistes. Au Cameroun comme dans toute l'Afrique centrale, les banques,
les opérateurs, les administrations et les jeunes entreprises du numérique
confient désormais leur activité à des systèmes dont la défaillance arrête
tout : les paiements, les salaires, les ports, les soins. Quand un système
s'arrête, la question qui suit n'est jamais seulement technique. Elle est
toujours la même : \emph{qui avait décidé, et qui devait rendre compte~?}

La seconde conviction est que cette discipline ne s'importe pas telle
quelle. COBIT, ITIL et les normes ISO sont des acquis précieux, mais ils
ont été pensés dans des contextes où les régulateurs, les marchés, les
compétences et les contraintes budgétaires ne sont pas les nôtres. Les
appliquer en Afrique suppose de les relire à la lumière du droit
camerounais et communautaire (CEMAC, OHADA), des attentes de nos
régulateurs et de la réalité de nos organisations.

\bigskip

\textbf{Ce que ce cours n'est pas.} Il n'est ni un résumé des publications
d'ISACA ou de PeopleCert, ni une préparation à une certification. Les
référentiels y sont présentés pour ce qu'ils font bien, cités avec
exactitude, et discutés là où leur transposition pose problème.

\bigskip

\textbf{Ce que ce cours est.} Un traité en six parties, précédées d'un
prologue, qui va du cas concret vers la théorie, de la règle vers la
pratique, et du cas simple vers la recherche. La Partie~I pose les
fondements de la gouvernance des SI. La Partie~II est consacrée à
COBIT~2019, la Partie~III au cadre normatif et juridique, la Partie~IV à
ITIL~4 et aux processus, la Partie~V aux défis de l'agilité, du nuage et de
l'intelligence artificielle, et la Partie~VI à la recherche. Quatre
organisations fictives, une banque, un opérateur portuaire, une
administration et une jeune entreprise de paiement mobile, traversent le
cours et servent de terrain aux exercices.

\bigskip

\textbf{Comment utiliser ce cours.} Chaque chapitre s'ouvre sur une \emph{boîte
objectifs} précisant les compétences attendues, et se ferme sur une \emph{boîte
résumé}. Les exercices sont gradués en trois niveaux : \niveaubase\ fondamental,
\niveaumoyen\ intermédiaire, \niveauexpert\ expert. Chaque chapitre est
accompagné d'une fiche d'auto-évaluation sur NDAP SUKULU, et dix
laboratoires pratiques d'une heure complètent les 30~heures de présentiel
comme travail personnel.

\bigskip

\begin{flushright}
    \textit{\auteurcours}\\
    \textit{Yaoundé, octobre 2026}
\end{flushright}

    \cleardoublepage

    % ==============================================================================
% Fichier : front/02_mode_emploi.tex
% Titre   : Mode d'emploi — GCI
% Auteur  : MINKA MI NGUIDJOÏ Thierry Emmanuel
% Version : 1.0 -- octobre 2026
% ==============================================================================

\chapter*{Mode d'emploi de ce cours}
\addcontentsline{toc}{chapter}{Mode d'emploi}
\markboth{Mode d'emploi}{Mode d'emploi}

\section*{Structure générale}

Le cours est organisé en \textbf{6 parties précédées d'un prologue}, soit
21~chapitres numérotés de façon continue. Il accompagne une unité
d'enseignement de 30~heures en présentiel, réparties en quinze séances de
deux heures qui couvrent l'ensemble des parties, du prologue à la
partie~VI. Les ateliers de questions-réponses par chapitre et les
laboratoires de NDAP SUKULU sont des travaux à faire à la maison, hors de
ces 30~heures. Les parties~V et~VI servent en outre de point de départ aux
mémoires.

\begin{tableaucours}{lp{3.8cm}p{3.8cm}c}{%
    \tableauHeader{Partie} &
    \tableauHeader{Thème} &
    \tableauHeader{Compétences visées} &
    \tableauHeader{Niveau}
}
0 & Prologue opérationnel & Distinguer une cause technique d'une cause de gouvernance & \niveaubase \\
I & Fondements de la gouvernance des SI & Définir, situer et justifier la gouvernance des SI & \niveaubase \\
II & COBIT 2019 & Concevoir et évaluer un système de gouvernance & \niveaubase\,--\,\niveaumoyen \\
III & Cadre normatif et juridique & Identifier les obligations et auditer la conformité & \niveaubase \\
IV & ITIL 4 et processus & Organiser et mesurer la gestion des services & \niveaumoyen\,--\,\niveauexpert \\
V & Défis avancés & Gouverner l'agilité, le nuage, les données et l'IA & \niveauexpert \\
VI & Recherche & Lire et produire la recherche sur la gouvernance des SI & \niveauexpert \\
\end{tableaucours}

\section*{Les environnements visuels}

\begin{boxdef}
\textbf{Définition}. Introduit un terme technique avec sa définition formelle.
Tous les termes définis sont indexés en fin d'ouvrage.
\end{boxdef}

\begin{boxobjectifs}
Objectifs du chapitre. Liste les compétences mesurables attendues à l'issue
du chapitre.
\end{boxobjectifs}

\begin{boxscenario}{Titre du scénario}
Scénario opérationnel. Présente un cas fictif ancré dans la réalité
camerounaise ou africaine, qui revient d'un chapitre à l'autre (opérations
SANAGA, KRIBI, MEFOU et NOUN).
\end{boxscenario}

\begin{boxmetier}{DSI}
Application métier. Traduit le concept académique en démarche
opérationnelle pour l'un des cinq métiers visés : DSI, Responsable GRC,
Auditeur SI, Responsable des services IT et RSSI.
\end{boxmetier}

\begin{boxprocedure}
\libelleProcedure. Indique ce qui doit être documenté et prouvé pour qu'une
pratique soit auditable : la pièce attendue, son responsable, sa fréquence.
\end{boxprocedure}

\begin{boxobjet}{Code et nom}
\libelleObjet. Présente un objectif COBIT ou une pratique ITIL sous une forme
constante : finalité, pratiques, entrées et sorties, matrice RACI,
indicateurs.
\end{boxobjet}

\begin{boxjuris}
Jurisprudence \& Droit applicable. Référence les textes de loi et articles
pertinents, avec accent sur le droit camerounais et africain.
\end{boxjuris}

\begin{boxalert}
Attention. Signale une erreur fréquente de gouvernance, de conformité ou de
mise en œuvre.
\end{boxalert}

\begin{boxresume}
À retenir. Synthèse des points essentiels du chapitre.
\end{boxresume}

\begin{boxexo}
Exercice. Travaux gradués : \niveaubase\ fondamental,
\niveaumoyen\ intermédiaire, \niveauexpert\ expert.
\end{boxexo}

\section*{Conventions typographiques}

\begin{itemize}
    \item Les \termecle{termes clés} apparaissent en \textbf{gras bleu marine}
          et sont indexés automatiquement.
    \item Les \sigle{ACRONYMES} utilisent des petites capitales.
    \item Les références aux normes utilisent la commande \norme{ISO~XXXX}.
    \item Les références aux lois camerounaises utilisent \loicam{AAAA/NNN}.
    \item Les niveaux de capacité COBIT s'écrivent \capacite{2}{DSS02}.
\end{itemize}

\section*{Matériel d'accompagnement}

\begin{itemize}
    \item \textbf{NDAP SUKULU}, code de cours \codeUE~: une fiche
          d'auto-évaluation par chapitre et dix laboratoires pratiques de
          60~minutes.
    \item \textbf{Dossiers de cas}~: pièces des quatre organisations fictives
          (tableurs, extraits de tickets, contrats, procès-verbaux).
    \item \textbf{Modèles de documents} (Annexe~\ref{ann:templates})~: charte de gouvernance,
          matrice RACI, accord de niveau de service, grille d'évaluation de
          capacité.
\end{itemize}

    \cleardoublepage

    % ==============================================================================
% Fichier : front/03_ndap_sukulu.tex
% Titre   : NDAP SUKULU, l'environnement compagnon du cours
% Auteur  : MINKA MI NGUIDJOÏ Thierry Emmanuel
% Version : 1.0 -- novembre 2026
% ==============================================================================
% Texte commun aux cours de l'auteur ; seul le tableau des labs est propre au
% cours (fichier front/03b_labs.tex).
% ==============================================================================

\chapter*{NDAP SUKULU, la salle de classe compagnon}
\addcontentsline{toc}{chapter}{NDAP SUKULU, la salle de classe compagnon}
\markboth{NDAP SUKULU}{NDAP SUKULU}

\begin{center}
    \includegraphics[width=0.62\textwidth]{figures/logos/logo_ndap_sukulu.png}
\end{center}

\section*{Un nom venu du pays bassa}

\begin{quote}\itshape
Le soir tombe sur Éséka. Un enfant revient des champs avec son grand-père et
lui demande pourquoi il doit marcher si loin, chaque matin, pour aller à
l'école. Le vieil homme s'arrête au bord du chemin et répond, en bassa~:
« \textup{Ndap sukulu}, la salle de classe, ce ne sont pas les murs. Ce sont
les murs qui viennent après. La salle de classe, c'est l'endroit où quelqu'un
t'attend pour t'apprendre ce qu'il sait. » L'enfant insiste~: « Et qui
m'attend~? » Le grand-père sourit~: « \textup{Malet yon}. Ton enseignant. »
\end{quote}

En bassa, la langue paternelle de l'auteur, \textbf{\textit{ndap sukulu}}
désigne la salle de classe et \textbf{\textit{malet yon}} signifie « ton
enseignant ». L'environnement compagnon de ce cours porte ce nom parce qu'il
prolonge la salle de classe au-delà des heures de cours~: il attend
l'étudiant le soir, le week-end, à distance, avec des questions, des pièces à
examiner et une correction immédiate. Il est édité par \textbf{MaletYon},
« ton enseignant », et mis à la disposition des étudiants et des praticiens
qui souhaitent se perfectionner (\href{mailto:maletyon@proton.me}{maletyon@proton.me}).

\section*{Ce que l'on y trouve}

\begin{itemize}
    \item \textbf{Une fiche d'auto-évaluation par chapitre.} Chaque fiche
          puise vingt questions dans une banque de quarante (QCM, vrai ou
          faux, textes à trous). Elle dure trente minutes~; la copie validée
          est verrouillée et notée aussitôt. Le seuil de réussite est de
          60~\%.
    \item \textbf{Des laboratoires pratiques.} Chaque lab place l'étudiant
          devant les pièces d'une affaire du fil rouge (tableurs, journaux,
          contrats, procès-verbaux) dans un bureau d'analyste virtuel,
          ouvert dans le navigateur. Il dure soixante minutes, se termine
          par des points de contrôle notés automatiquement et renvoie aux
          sections du livre qu'il mobilise.
    \item \textbf{Un suivi pour l'enseignant}~: rapport de classe, notes par
          fiche et par lab, journal d'activité.
\end{itemize}

Les fiches et les labs sont des travaux à faire hors des heures de cours.
Ils ne remplacent ni le livre ni les exercices de fin de chapitre~: ils en
vérifient la maîtrise.

\section*{Comment y travailler}

\begin{enumerate}
    \item Se connecter au portail NDAP SUKULU avec ses identifiants et le
          code de cours \textbf{\codeUE}.
    \item Pour un lab, commencer par vérifier l'empreinte SHA-256 des pièces
          (\texttt{sha256sum -c MANIFESTE\_SHA256.txt}), puis travailler sur
          la copie du dossier \texttt{travail}~: les pièces originales sont
          en lecture seule et l'activité est journalisée, comme dans une
          vraie mission.
    \item Lire le rappel de cours affiché par le lab et relire les sections
          citées avant de répondre.
\end{enumerate}

\begin{boxinfo}
\textbf{Si le portail n'est pas encore ouvert.} En attendant le déploiement
de NDAP SUKULU, les fiches de questions sont fournies au format texte dans le
dossier \texttt{fiches} du cours, prêtes à être distribuées. Les corrections
des fiches figurent dans le livret de correction, à la suite des corrigés des
exercices de chaque chapitre.
\end{boxinfo}

\section*{Les laboratoires du cours}

% Tableau des labs du cours GCI dans NDAP SUKULU (propre au cours).
Le cours compte dix laboratoires, disponibles sous le code \codeUE. Le dernier
réunit les quatre organisations du fil rouge et sert de synthèse.

\begin{tableaucours}{lL{6.2cm}L{1.6cm}L{2.0cm}}{%
    \tableauHeader{Lab} &
    \tableauHeader{Sujet} &
    \tableauHeader{Chapitre} &
    \tableauHeader{Opération}
}
GCI-LAB-01 & Arbitrer le portefeuille de projets de l'ANPP & \ref{chap:alignement}
& MEFOU \\
GCI-LAB-02 & Fixer les capacités cibles de la BANQUE SANAGA &
\ref{chap:cobit-conception} & SANAGA \\
GCI-LAB-03 & Dérouler la cascade des objectifs de la BANQUE SANAGA &
\ref{chap:cobit-conception} & SANAGA \\
GCI-LAB-04 & Évaluer la capacité de BAI06 et de DSS02 & \ref{chap:cobit-capacite}
& SANAGA \\
GCI-LAB-05 & Tester les contrôles généraux informatiques de l'ANPP &
\ref{chap:audit} & MEFOU \\
GCI-LAB-06 & Prioriser et analyser les tickets de KRIBILOG &
\ref{chap:itil-incidents} & KRIBI \\
GCI-LAB-07 & Évaluer et décider un changement de la BANQUE SANAGA &
\ref{chap:itil-changements} & SANAGA \\
GCI-LAB-08 & Disponibilité, pénalités et reprise du suivi des conteneurs &
\ref{chap:itil-niveaux} & KRIBI \\
GCI-LAB-09 & Construire la matrice RACI du modèle de NOUN PAY &
\ref{chap:processus} & NOUN \\
GCI-LAB-10 & Synthèse des quatre cas~: feuilles de route chiffrées &
\ref{chap:cas-integrateurs} & les quatre \\
\end{tableaucours}
\vspace{-0.8em}
\begin{center}{\small\itshape Les laboratoires du cours dans NDAP SUKULU}\end{center}


    \cleardoublepage

    \tableofcontents
    \cleardoublepage

    \listoffigures
    \listoftables
    \cleardoublepage

% --------------------------------------------------------------------------
% CORPS DU COURS (pages arabes)
% Numérotation des chapitres CONTINUE sur tout le cours : le prologue est le
% chapitre 0, puis C01 à C21, un numéro par fichier.
% --------------------------------------------------------------------------
\mainmatter

    % ======================================================================
    % PARTIE 0 — PROLOGUE OPÉRATIONNEL
    % « Avant toute théorie, une banque qui s'arrête trois jours. »
    % But : montrer qu'une panne technique cache une défaillance de décision.
    % Niveau : introductif — aucun prérequis.
    % ======================================================================
    % Partie non numérotée : la Partie I reste « Fondements », comme dans
    % le plan (avec \part numéroté, COBIT devenait la Partie III).
    \part*{Prologue Opérationnel}
    \addcontentsline{toc}{part}{Prologue Opérationnel}

        \setcounter{chapter}{-1}
        % ============================================================================
% Fichier  : partie0/P00_trois_jours_sans_banque.tex
% Titre    : Prologue -- Trois Jours sans Banque : quand Personne n'a Décidé
% Auteur   : MINKA MI NGUIDJOÏ Thierry Emmanuel
% Version  : 1.0 -- octobre 2026
% ============================================================================

\chapter[Prologue : Trois Jours sans Banque]{Prologue : Trois Jours sans Banque, quand Personne n'a Décidé}
\label{chap:prologue}
\chaptermark{Prologue}

\epigraph{%
    « Une panne se répare en quelques heures.\\
    Une décision que personne n'a prise se paie pendant des années. »%
}{Thierry Minka, auditeur des systèmes d'information}

\niveauChapitre{Introductif}{%
    Aucun prérequis. Ce prologue se lit avant d'ouvrir le moindre
    référentiel~: il ne suppose aucune connaissance de COBIT ni d'ITIL.%
}

\begin{boxobjectifs}
\begin{itemize}
    \item Reconstituer la chronologie d'une défaillance informatique majeure
          (opération SANAGA).
    \item Distinguer une cause technique d'une cause de gouvernance.
    \item Poser les trois questions réflexes~: qui décide, qui exécute, qui
          rend compte.
    \item Distinguer gouvernance et management des systèmes d'information.
    \item Identifier les parties prenantes qui demandent des comptes
          (conseil d'administration, régulateur, clients).
    \item Situer COBIT et ITIL dans la réponse à la défaillance.
\end{itemize}
\end{boxobjectifs}

\bigskip

% ============================================================================
\section*{Avant de commencer : une panne qui n'était pas technique}
% ============================================================================

Douala, lundi 9~mars~2026, 7~h~30. Les 38~agences de la BANQUE
SANAGA\footnote{La BANQUE SANAGA, ses dirigeants, ses prestataires, les
dates et les montants de ce cas sont fictifs. Les mécanismes de
gouvernance, les référentiels et les textes cités sont exacts.} ouvrent
leurs portes. Aucun guichet ne fonctionne. Les écrans affichent un message
d'erreur du nouveau système bancaire central, mis en service pendant le
week-end. Les distributeurs automatiques passent en mode dégradé, les
virements et les paiements par carte sont bloqués. À Yaoundé, à Bafoussam,
à Garoua, les files d'attente s'allongent devant des agences qui ne peuvent
ni verser ni retirer.

La paralysie dure \textbf{trois jours}. Le service ne revient complètement
que le jeudi 12~mars au matin. Le coût, chiffré plus tard par la direction
financière, atteint \textbf{1,9~milliard de francs CFA}. Le 16~mars, la
Commission bancaire de l'Afrique centrale (\sigle{COBAC}) demande à la banque
un rapport circonstancié sous trente jours.

\medskip

Le conseil d'administration réunit une séance extraordinaire. Le directeur
général, M.~FOTSO, ouvre la séance par une phrase qui résume ce que chacun
pense alors~: « Le prestataire nous a livré un logiciel défaillant. » Le
rapport d'enquête interne, remis trois semaines plus tard, dit tout autre
chose.

\begin{boxalert}
\textbf{Ce que révèle l'enquête.} Le logiciel n'était pas en cause. Il
avait été recetté à 63~\% seulement quand la date de bascule a été
maintenue. La date elle-même avait été fixée pour coïncider avec le
lancement commercial d'une offre mobile, déjà annoncée à la presse. Le
comité de pilotage du projet ne s'était plus réuni depuis quatre mois.
Aucun comité n'avait autorisé la mise en production. Aucun retour arrière
n'avait été testé, et la conversion des données de l'ancien système était
irréversible. Le contrat de l'intégrateur ne prévoyait aucune pénalité. Un
rapport d'audit interne avait signalé ces faiblesses cinq mois plus tôt~:
ses recommandations étaient restées « en cours ».

\smallskip
Chacune de ces défaillances est une \textbf{décision} qui n'a pas été
prise, ou qui a été prise par la mauvaise personne, sans trace écrite et
sans que personne en rende compte.
\end{boxalert}

\medskip

Ce prologue n'enseigne ni COBIT ni ITIL. Ces référentiels occupent les
chapitres qui suivent. Il enseigne \textbf{le réflexe fondateur de la
discipline}~: devant toute défaillance d'un système d'information, chercher
la décision derrière l'incident. Une panne se répare. Une organisation qui
ne sait pas qui décide reproduira la même panne.

\separateur

% ============================================================================
\section{Chronologie de la Défaillance}
\label{sec:pro-chronologie}
% ============================================================================

\subsection{Reconstituer avant de juger}

La première tâche de quiconque analyse une défaillance (le DSI qui rédige le
rapport au régulateur, l'auditeur mandaté par le conseil, le responsable des
risques) est de \textbf{reconstituer la chronologie}. Elle fixe les faits
avant toute interprétation. Elle se construit à partir de pièces~:
procès-verbaux, courriels, tickets, journaux d'exploitation, contrats.

\begin{boxdef}
\termecle{Chronologie de défaillance}\index{chronologie de défaillance}
(\textit{incident timeline})~: suite datée des faits, des décisions et des
absences de décision qui ont précédé, accompagné et suivi une défaillance.
Chaque entrée est rattachée à une pièce justificative.
\end{boxdef}

\begin{tableaucours}{L{2.6cm}L{6.4cm}L{3.2cm}}{%
    \tableauHeader{Date} &
    \tableauHeader{Fait établi} &
    \tableauHeader{Pièce}
}
Juin 2025 & Le conseil d'administration approuve le projet HORIZON
(remplacement du système bancaire central), budget de 3,2~milliards FCFA.
Intégrateur retenu~: DIGICORE AFRIQUE. & Procès-verbal du conseil \\
Oct. 2025 & L'audit interne signale l'absence de procédure de gestion des
changements et de plan de continuité testé. Recommandations notées « en
cours ». & Rapport d'audit n°~2025-07 \\
Nov. 2025 & Dernière réunion du comité de pilotage du projet. & Compte rendu \\
Janv. 2026 & Le directeur général fixe la bascule au week-end du 7--8~mars,
date du lancement de l'offre « SANAGA Mobile ». & Courriel du DG \\
20 févr. 2026 & Recette à 63~\% des scénarios. La DSI demande un report par
note interne. Aucune réponse écrite. & Note de la DSI \\
7 mars, 22 h & Début de la migration. & Journal d'exploitation \\
8 mars, 14 h & Conversion des comptes~: écarts de solde sur 11~\% des
comptes. Le retour à l'ancien système n'est plus possible. & Journal
d'exploitation \\
9 mars, 7 h 30 & Ouverture des agences~: système indisponible. & Tickets
du centre de services \\
12 mars, 8 h & Retour complet du service. & Journal d'exploitation \\
16 mars & La COBAC demande un rapport circonstancié sous 30~jours. &
Lettre de la COBAC \\
\end{tableaucours}
\vspace{-0.8em}
\begin{center}{\small\itshape Tableau~P.1. Chronologie de l'opération SANAGA}\end{center}

\subsection{Constatation et analyse}

Une chronologie ne contient que des \textbf{constatations}~: des faits
vérifiables, rattachés à une pièce. L'interprétation vient ensuite, dans une
partie distincte, et elle s'accompagne d'un niveau de confiance. Mélanger
les deux est l'erreur la plus fréquente des rapports de défaillance, et la
plus facile à exploiter par celui qui veut contester le rapport.

\begin{boxalert}
\textbf{Violation typique.} Le premier projet de rapport de la BANQUE
SANAGA écrit~: « Le prestataire a livré un logiciel non conforme, ce qui a
provoqué la panne. » La phrase mêle un fait non établi (la non-conformité)
et une causalité non démontrée. DIGICORE AFRIQUE produit aussitôt le
procès-verbal de recette, qui montre que la banque a elle-même décidé de
basculer avec 37~\% des scénarios non testés.

\smallskip
\textbf{Règle~:} écrire d'abord la constatation (« La recette couvrait
63~\% des scénarios le 20~février~; la bascule a eu lieu le 7~mars »), puis
l'analyse, avec son niveau de confiance (« Il est très probable que les
écarts de solde proviennent de scénarios de conversion non testés »).
\end{boxalert}

% ============================================================================
\section{Cause Technique, Cause de Gouvernance}
\label{sec:pro-causes}
% ============================================================================

\subsection{Deux niveaux d'explication}

Toute défaillance d'un système d'information s'explique à deux niveaux. Le
\textbf{niveau technique} répond à la question \emph{comment}~: quel
composant a cédé, quelle donnée a été corrompue. Le \textbf{niveau de la
gouvernance} répond à la question \emph{pourquoi cela a été possible}~:
quelle décision, quel contrôle, quelle responsabilité a manqué.

\begin{boxdef}
\termecle{Défaillance de gouvernance}\index{défaillance de gouvernance}
(\textit{governance failure})~: absence, inadéquation ou contournement d'un
mécanisme de décision, de contrôle ou de reddition de comptes, qui permet
qu'un risque connu ou connaissable se réalise.
\end{boxdef}

\begin{tableaucours}{L{5.6cm}L{6.6cm}}{%
    \tableauHeader{Cause technique (comment~?)} &
    \tableauHeader{Cause de gouvernance (pourquoi était-ce possible~?)}
}
Écarts de solde lors de la conversion des comptes & Bascule maintenue avec
une recette incomplète, sans décision formelle d'acceptation du risque \\
Retour arrière impossible & Aucun plan de retour arrière exigé ni testé
avant l'autorisation de mise en production \\
Indisponibilité prolongée & Aucun plan de continuité testé, malgré la
recommandation de l'audit interne \\
Intervention lente de l'intégrateur & Contrat sans engagement de niveau de
service ni pénalité \\
Communication de crise improvisée & Aucun rôle de crise défini, aucune
information préalable du conseil sur le risque du projet \\
\end{tableaucours}
\vspace{-0.8em}
\begin{center}{\small\itshape Tableau~P.2. Les deux niveaux d'explication de l'opération SANAGA}\end{center}

\subsection{Remonter de l'incident à la décision}

Une méthode simple permet de passer d'un niveau à l'autre~: demander
« pourquoi~? » jusqu'à atteindre une décision, ou une absence de décision.
Appliquée à la BANQUE SANAGA, elle donne~:

\begin{enumerate}
    \item \emph{Pourquoi les agences étaient-elles paralysées~?} Parce que
          le nouveau système produisait des soldes erronés.
    \item \emph{Pourquoi les soldes étaient-ils erronés~?} Parce que des
          scénarios de conversion n'avaient pas été testés.
    \item \emph{Pourquoi a-t-on basculé sans les tester~?} Parce que la date
          était fixée par le lancement commercial.
    \item \emph{Pourquoi la date l'a-t-elle emporté sur la recette~?} Parce
          que personne n'avait l'autorité, ni l'obligation, d'arbitrer entre
          le risque technique et l'enjeu commercial.
\end{enumerate}

La quatrième réponse n'est plus technique. Elle décrit une
\textbf{organisation}~: elle sera au cœur des Parties~\ref{part:fondements}
et~\ref{part:cobit} du cours.

\begin{boxinfo}
\textbf{La règle d'or~:} une analyse de défaillance est terminée quand elle
a atteint une décision, ou une absence de décision, et qu'elle a nommé le
rôle qui aurait dû la prendre. Tant qu'elle s'arrête à un composant, elle
n'a expliqué que la moitié de l'incident.
\end{boxinfo}

% ============================================================================
\section{Les Trois Questions Réflexes}
\label{sec:pro-trois-questions}
% ============================================================================

\subsection{Décider, exécuter, rendre compte}

Devant tout projet, tout changement ou toute défaillance d'un système
d'information, trois questions suffisent à révéler l'état de la gouvernance.

\begin{enumerate}
    \item \textbf{Qui décide~?} Quelle instance ou quelle personne avait
          l'autorité de trancher, et l'a-t-elle fait, par écrit~?
    \item \textbf{Qui exécute~?} Qui a mis en œuvre la décision, avec quels
          moyens, et selon quelle procédure~?
    \item \textbf{Qui rend compte~?} Qui doit expliquer le résultat, à qui,
          et sur la base de quelles informations~?
\end{enumerate}

\begin{boxdef}
\termecle{Droits de décision}\index{droits de décision} (\textit{decision
rights})~: répartition formelle, entre les instances et les personnes d'une
organisation, du pouvoir de décider sur une question donnée. Weill et Ross
en font le cœur de la gouvernance des SI~\cite{weill2004itgovernance}.

\medskip
\termecle{Redevabilité}\index{redevabilité} (\textit{accountability})~:
obligation, pour celui qui a reçu une responsabilité, d'en rendre compte à
l'instance qui la lui a confiée, et d'en assumer les conséquences.
\end{boxdef}

\subsection{Application à la BANQUE SANAGA}

\begin{tableaucours}{L{3.4cm}L{2.9cm}L{3.1cm}L{2.6cm}}{%
    \tableauHeader{Décision} &
    \tableauHeader{Prise par} &
    \tableauHeader{Aurait dû l'être par} &
    \tableauHeader{Trace écrite}
}
Date de bascule & Directeur général seul & Comité de pilotage, sur avis de
la DSI & Courriel \\
Bascule avec recette incomplète & Personne (statu quo) & Comité de
pilotage, avec acceptation formelle du risque & Aucune \\
Autorisation de mise en production & Personne & Comité des changements &
Aucune \\
Suite des recommandations d'audit & Personne & Direction générale, suivie
par le comité d'audit & « En cours » \\
Information du conseil sur le risque & Personne & Directeur général &
Aucune \\
\end{tableaucours}
\vspace{-0.8em}
\begin{center}{\small\itshape Tableau~P.3. Les droits de décision dans l'opération SANAGA}\end{center}

Le tableau parle de lui-même. Sur cinq décisions critiques, une seule a été
prise, par une personne seule, sans l'avis de ceux qui portaient le risque.
Les quatre autres n'ont été prises par personne. C'est la signature d'une
organisation sans gouvernance des SI~: non pas des dirigeants incompétents,
mais une absence de règles sur \emph{qui décide de quoi}.

\begin{boxmetier}{DSI}
\textbf{Protéger la décision, pas seulement le système}
\begin{enumerate}
    \item Tout désaccord sur un risque majeur (date, périmètre, recette) se
          formule \textbf{par écrit}, adressé à l'instance qui a le droit de
          décider, et non au seul supérieur hiérarchique.
    \item La note expose les faits (taux de recette, scénarios non testés),
          le risque (ce qui peut arriver, son coût estimé) et les options
          (reporter, réduire le périmètre, accepter le risque).
    \item Elle demande une \textbf{décision explicite}~: le silence n'est
          pas une décision.
    \item Le DSI conserve la note et la réponse, ou l'absence de réponse,
          dans le dossier du projet.
\end{enumerate}
La note du 20~février de la DSI de la BANQUE SANAGA était juste sur le fond.
Elle a échoué sur la forme~: adressée au seul directeur général, elle ne
demandait aucune décision et n'a jamais atteint le comité de pilotage.
\end{boxmetier}

% ============================================================================
\section{Gouvernance et Management}
\label{sec:pro-gouv-management}
% ============================================================================

\subsection{Une distinction fondatrice}

Les trois questions réflexes font apparaître deux familles d'activités. Les
unes fixent la direction et contrôlent le résultat~; les autres
réalisent. COBIT~2019 en fait une distinction structurante, chacune
relevant d'activités, de structures et de finalités
différentes~\cite{isaca2018cobitintro}.

\begin{boxdef}
\termecle{Gouvernance des SI}\index{gouvernance des SI} (\textit{governance
of enterprise I\&T})~: ensemble des activités par lesquelles l'organe de
gouvernance \emph{évalue} les besoins des parties prenantes pour fixer des
objectifs équilibrés, \emph{dirige} par la priorisation et la décision, et
\emph{surveille} la performance et la conformité au regard de cette
direction. Elle relève, dans la plupart des organisations, du conseil
d'administration.

\medskip
\termecle{Management des SI}\index{management des SI} (\textit{management
of enterprise I\&T})~: ensemble des activités qui \emph{planifient},
\emph{bâtissent}, \emph{exploitent} et \emph{surveillent} les systèmes
d'information, conformément à la direction fixée par la gouvernance. Il
relève de la direction générale.
\end{boxdef}

\begin{tableaucours}{L{3.0cm}L{4.6cm}L{4.6cm}}{%
    \tableauHeader{Critère} &
    \tableauHeader{Gouvernance} &
    \tableauHeader{Management}
}
Question & Fait-on les bonnes choses~? & Fait-on bien les choses~? \\
Responsable & Conseil d'administration et ses comités & Direction générale,
DSI, responsables opérationnels \\
Verbes (COBIT) & Évaluer, diriger, surveiller & Planifier, bâtir,
exploiter, surveiller \\
Horizon & Stratégique, pluriannuel & Opérationnel, annuel et quotidien \\
Exemple SANAGA & Décider si le lancement commercial justifie le risque de
bascule & Conduire la recette, préparer le retour arrière \\
\end{tableaucours}
\vspace{-0.8em}
\begin{center}{\small\itshape Tableau~P.4. Gouvernance et management des SI}\end{center}

\begin{boiteRemarque}{Le DSI gouverne-t-il~?}
La question revient à chaque promotion, et la réponse est nuancée. Le DSI
\emph{participe} à la gouvernance~: il éclaire le conseil, propose des
arbitrages, alerte sur les risques. Mais il n'en est pas l'organe. Un DSI
qui décide seul de la stratégie numérique gouverne par défaut, faute
d'instance compétente. Un conseil qui délègue tout au DSI renonce à sa
propre responsabilité. Le Chapitre~\ref{chap:gouv-entreprise} reprendra
cette question à partir de la théorie de l'agence.
\end{boiteRemarque}

% ============================================================================
\section{Ceux qui Demandent des Comptes}
\label{sec:pro-parties-prenantes}
% ============================================================================

\subsection{Les parties prenantes de la BANQUE SANAGA}

Une défaillance majeure fait apparaître, d'un coup, tous ceux à qui
l'organisation doit des comptes. Chacun pose une question différente et
attend une réponse documentée.

\begin{tableaucours}{L{3.6cm}L{4.4cm}L{4.2cm}}{%
    \tableauHeader{Partie prenante} &
    \tableauHeader{Sa question} &
    \tableauHeader{Ce qu'elle exige}
}
Conseil d'administration & Pourquoi n'avons-nous pas été informés du
risque~? & Un rapport sur les causes et la responsabilité de chacun \\
Actionnaires & Combien cela coûte-t-il, et cela peut-il se reproduire~? &
Le chiffrage des pertes et un plan d'action \\
COBAC & Le contrôle interne de la banque fonctionne-t-il~? & Un rapport
circonstancié sous 30~jours \\
Clients et commerçants & Nos opérations sont-elles justes~? & La
correction des soldes et une indemnisation \\
Commissaire aux comptes & Les comptes de l'exercice restent-ils
fiables~? & La preuve de l'exactitude des soldes convertis \\
\end{tableaucours}
\vspace{-0.8em}
\begin{center}{\small\itshape Tableau~P.5. Les parties prenantes et leurs exigences}\end{center}

\subsection{Le coût d'une décision manquée}

Le chiffrage de la direction financière donne la mesure de l'enjeu.

\begin{tableaucours}{p{8.2cm}r}{%
    \tableauHeader{Poste} &
    \tableauHeader{Montant (millions FCFA)}
}
Commissions et frais non perçus pendant trois jours & 410 \\
Indemnisation des clients et des commerçants & 380 \\
Prestations d'urgence et heures supplémentaires & 260 \\
Retraits hors ligne au-delà du solde, non recouvrés & 330 \\
Correction des données et audit externe imposé & 520 \\
\textbf{Total} & \textbf{1~900} \\
\end{tableaucours}
\vspace{-0.8em}
\begin{center}{\small\itshape Tableau~P.6. Coût de l'opération SANAGA}\end{center}

Le report de la bascule de six semaines, demandé par la DSI, aurait coûté
selon le même service environ 85~millions FCFA (prolongation de
l'intégrateur et décalage de la campagne commerciale). Le rapport entre les
deux montants, plus de vingt pour un, est l'argument que la gouvernance
aurait dû voir avant la bascule, et qu'elle n'a vu qu'après.

\subsection{Le cadre légal camerounais}

La défaillance de la BANQUE SANAGA ne relève pas seulement de la
réglementation bancaire. La loi camerounaise sur la cybersécurité soumet
les systèmes d'information à des obligations d'audit.

\begin{boxjuris}
\textbf{\loicam{n°~2010/012} du 21~décembre~2010 relative à la
cybersécurité et la cybercriminalité au Cameroun, article~32} (version
anglaise officielle)~:

\medskip
\textit{« (1) Electronic communication networks and information systems
shall be subject to a regime of compulsory and periodic auditing of their
security systems by the Agency. (2) Security audit and severity scale
rating shall be undertaken each year or as required by the prevailing
circumstances. »}

\medskip
Soit, en substance~: les réseaux de communications électroniques et les
systèmes d'information sont soumis à un audit obligatoire et périodique de
leur sécurité par l'Agence (l'\sigle{ANTIC}), conduit chaque année ou
selon les circonstances. L'article~13 de la même loi soumet en outre les
systèmes d'information des opérateurs à un audit de sécurité obligatoire,
dont les modalités sont fixées par voie réglementaire.

\medskip
\textbf{Portée.} La loi fait de l'audit des systèmes d'information une
obligation, et non une bonne pratique facultative. Le
Chapitre~\ref{chap:droit} précisera son champ d'application exact, ses
textes d'application et son articulation avec la réglementation de la
COBAC sur le contrôle interne des établissements de crédit, ainsi qu'avec
le droit OHADA de la responsabilité des dirigeants.
\end{boxjuris}

% ============================================================================
\section{Ce que COBIT et ITIL Auraient Changé}
\label{sec:pro-cobit-itil}
% ============================================================================

\subsection{Deux référentiels, deux niveaux}

Les deux référentiels au programme de ce cours répondent chacun à un niveau
d'explication de la défaillance. \textbf{COBIT~2019} (Partie~\ref{part:cobit})
organise la gouvernance et le management de l'ensemble des systèmes
d'information~: qui décide, selon quels objectifs, avec quels contrôles.
\textbf{ITIL~4} (Partie~\ref{part:pratique}) organise la gestion des
services informatiques au quotidien~: comment autoriser un changement,
traiter un incident, tenir un engagement de service. COBIT dit \emph{ce qui
doit être maîtrisé}~; ITIL décrit en détail \emph{comment le faire} pour
les services.

\begin{tableaucours}{L{4.2cm}L{4.2cm}L{3.8cm}}{%
    \tableauHeader{Défaillance SANAGA} &
    \tableauHeader{Objectif COBIT 2019} &
    \tableauHeader{Pratique ITIL 4}
}
Aucune instance n'arbitre le risque du projet & EDM01, EDM03 (cadre de
gouvernance, optimisation des risques) & · \\
Comité de pilotage en sommeil & BAI01, BAI11 (programmes, projets) &
Gestion de projet \\
Mise en production non autorisée, sans retour arrière & BAI06, BAI07
(changements, acceptation et transition) & Habilitation des changements,
gestion des mises en production \\
Contrat sans engagement de service & APO09, APO10 (accords de service,
fournisseurs) & Gestion des niveaux de service, gestion des
fournisseurs \\
Aucun plan de continuité testé & DSS04 (continuité) & Gestion de la
continuité de service \\
Recommandations d'audit sans suite & MEA02 (système de contrôle
interne) & · \\
\end{tableaucours}
\vspace{-0.8em}
\begin{center}{\small\itshape Tableau~P.7. Défaillances SANAGA, objectifs COBIT et
pratiques ITIL (intitulés traduits par nos soins)}\end{center}

\subsection{Un exemple~: la gestion des changements}

La mise en production du nouveau système était un \termecle{changement}
\index{changement} au sens des deux référentiels~: une modification d'un
élément du système d'information susceptible d'affecter les services.
COBIT en fait un objectif de management, que le cours présentera toujours
sous la même forme.

\begin{boxobjet}{BAI06 · Gérer les changements informatiques}
\textbf{Finalité~:} permettre des changements rapides et fiables, en
limitant le risque d'affecter la stabilité ou l'intégrité de
l'environnement modifié.

\textbf{Pratiques clés~:} évaluer, prioriser et autoriser les demandes de
changement~; gérer les changements d'urgence~; suivre l'état des
changements et en rendre compte~; clore et documenter les changements.

\textbf{Entrées et sorties~:} demande de changement, évaluation d'impact et
de risque, plan de retour arrière~; décision d'autorisation, calendrier des
changements, rapport de clôture.

\textbf{RACI (extrait)~:} autorisation d'un changement majeur~:
\emph{A} au comité des changements, \emph{R} au responsable du changement,
\emph{C} aux métiers concernés et au RSSI, \emph{I} au centre de services.

\textbf{Indicateurs~:} taux de changements en échec, nombre de changements
non autorisés, part des changements majeurs dotés d'un retour arrière
testé.
\end{boxobjet}

\begin{boxprocedure}
\textbf{Mise en production d'un changement majeur.} Pour qu'un auditeur
puisse établir que le changement était maîtrisé, le dossier doit contenir~:
\begin{itemize}
    \item le procès-verbal de recette, avec le taux de couverture des
          scénarios et la liste des anomalies ouvertes~;
    \item l'évaluation du risque et la décision écrite d'acceptation des
          risques résiduels, signée par l'instance compétente~;
    \item le plan de retour arrière et la preuve qu'il a été testé~;
    \item le procès-verbal d'autorisation du comité des changements, daté
          et antérieur à la mise en production~;
    \item le rapport de clôture, avec les incidents survenus.
\end{itemize}
Dans l'opération SANAGA, seule la première pièce existait.
\end{boxprocedure}

\begin{boxmetier}{Auditeur SI}
\textbf{Les cinq pièces à demander après une défaillance majeure}
\begin{enumerate}
    \item Les procès-verbaux des instances de gouvernance (conseil, comité
          d'audit, comité de pilotage) sur la période du projet.
    \item Le dossier du changement (voir l'exigence de contrôle ci-dessus).
    \item Les contrats avec les prestataires et leurs annexes de niveau de
          service.
    \item Le plan de continuité et les comptes rendus de ses derniers
          tests.
    \item Le suivi des recommandations d'audit antérieures, avec leurs
          dates d'échéance et leurs responsables.
\end{enumerate}
L'absence d'une pièce est elle-même une constatation~: elle se consigne
comme telle.
\end{boxmetier}

% ============================================================================
\section{Retour sur la BANQUE SANAGA : ce qui Aurait Dû se Passer}
\label{sec:pro-retour-sanaga}
% ============================================================================

Voici la même bascule, conduite cette fois par une organisation gouvernée.

\begin{boxscenario}{Opération SANAGA. La bascule gouvernée}
\begin{enumerate}
    \item \textbf{Juin 2025.} En approuvant le projet HORIZON, le conseil
          d'administration fixe ses règles~: un comité de pilotage mensuel
          présidé par le directeur général adjoint, un point de risque
          trimestriel au comité d'audit, et l'obligation d'une décision
          formelle avant toute bascule.
    \item \textbf{Octobre 2025.} Les recommandations de l'audit interne sur
          la gestion des changements et la continuité reçoivent chacune un
          responsable et une échéance. Le comité d'audit en suit l'avancement.
    \item \textbf{20 février 2026.} La note de la DSI (recette à 63~\%) est
          inscrite à l'ordre du jour du comité de pilotage. Le comité
          compare le coût d'un report (85~millions FCFA) et le risque d'une
          bascule prématurée.
    \item \textbf{Décision.} Le comité reporte la bascule de six semaines et
          dissocie le lancement de « SANAGA Mobile », qui démarre sur
          l'ancien système. La décision est consignée au procès-verbal.
    \item \textbf{Avril 2026.} Recette à 98~\%, retour arrière testé à
          blanc, comité des changements réuni. Le contrat de DIGICORE
          AFRIQUE a été amendé~: présence sur site pendant la bascule et
          pénalités de retard d'intervention.
    \item \textbf{Bascule.} Un écart de solde est détecté le dimanche à
          15~h. Le retour arrière est déclenché à 16~h. Les agences ouvrent
          le lundi sur l'ancien système. La bascule réussit deux semaines
          plus tard.
\end{enumerate}
\textbf{Résultat~:} le même incident technique s'est produit. Il est resté
un incident d'exploitation, sans conséquence pour les clients, au lieu de
devenir une crise à 1,9~milliard FCFA.
\end{boxscenario}

Le point essentiel est là~: \textbf{la gouvernance n'empêche pas les
incidents techniques}. Elle fait en sorte qu'ils soient anticipés, décidés
et contenus. La suite du cours donne les outils pour construire une telle
organisation.

% ============================================================================
\separateur
% ============================================================================

\begin{boxresume}
\textbf{Ce qu'il faut retenir de ce prologue~:}

\begin{itemize}
    \item \textbf{Une défaillance a deux niveaux}~: la cause technique dit
          comment le système a cédé, la cause de gouvernance dit pourquoi
          c'était possible. Une analyse qui s'arrête au composant est
          incomplète.
    \item \textbf{La chronologie d'abord}~: on reconstitue les faits à partir
          de pièces, et on sépare toujours la constatation de l'analyse.
    \item \textbf{Trois questions réflexes}~: qui décide, qui exécute, qui
          rend compte. Une décision sans trace écrite, ou prise par
          personne, est un signal d'alerte.
    \item \textbf{Gouvernance et management sont distincts}~: la gouvernance
          évalue, dirige et surveille~; le management planifie, bâtit,
          exploite et surveille.
    \item \textbf{Les parties prenantes demandent des comptes}~: conseil,
          actionnaires, régulateur, clients, commissaire aux comptes. Au
          Cameroun, la loi n°~2010/012 fait de l'audit des systèmes
          d'information une obligation.
    \item \textbf{COBIT et ITIL se complètent}~: COBIT organise ce qui doit
          être maîtrisé à l'échelle de l'organisation, ITIL décrit comment
          gérer les services au quotidien.
    \item \textbf{La gouvernance n'empêche pas l'incident}~: elle le rend
          anticipé, décidé et contenu.
\end{itemize}
\end{boxresume}

\bigskip

\begin{boxexo}[title={\ding{45}\quad Exercice~P.1. Technique ou gouvernance~? \niveaubase}]
Pour chacun des constats suivants, relevés dans d'autres organisations,
indiquez s'il s'agit d'une cause technique ou d'une cause de gouvernance,
et justifiez en une phrase.
\begin{enumerate}[(a)]
    \item Le disque d'un serveur de messagerie tombe en panne.
    \item Aucune sauvegarde de ce serveur n'avait été vérifiée depuis
          deux ans, sans que personne en soit chargé.
    \item Un administrateur applique un correctif un vendredi soir sans
          prévenir personne, conformément à l'usage du service.
    \item Le budget de renouvellement du matériel a été supprimé par le
          directeur financier sans consulter la DSI.
\end{enumerate}
\end{boxexo}

\begin{boxexo}[title={\ding{45}\quad Exercice~P.2. Les droits de décision de KRIBILOG \niveaumoyen}]
La société KRIBILOG SA, opérateur logistique à Kribi et à Douala, change
d'infogérant en trois semaines, sur décision de son directeur des achats,
pour une économie annuelle de 48~millions FCFA. Le responsable des services
informatiques l'apprend par le nouveau prestataire. Deux mois plus tard, le
logiciel de suivi des conteneurs connaît trois arrêts majeurs.
\begin{enumerate}[(a)]
    \item Construisez, sur le modèle du Tableau~P.3, le tableau des
          décisions critiques de ce changement de prestataire.
    \item Pour chaque décision, indiquez qui l'a prise, qui aurait dû la
          prendre, et quelle trace écrite devrait exister.
    \item Rédigez la constatation que l'auditeur SI consignerait, puis son
          analyse, avec un niveau de confiance.
\end{enumerate}
\end{boxexo}

\begin{boxexo}[title={\ding{45}\quad Exercice~P.3. Le rapport à la COBAC \niveauexpert}]
Vous êtes responsable GRC de la BANQUE SANAGA. Le directeur général vous
confie la rédaction du rapport circonstancié demandé par la COBAC.
\begin{enumerate}[(a)]
    \item Proposez le plan du rapport, en séparant la chronologie, les
          constatations, l'analyse des causes et le plan d'action.
    \item Le directeur général vous demande d'écrire que la panne est due
          au seul prestataire. Expliquez, pièces à l'appui, pourquoi cette
          affirmation ne résisterait pas à l'examen du régulateur.
    \item Pour chacune des six défaillances du Tableau~P.7, proposez une
          mesure corrective, son responsable et son échéance.
    \item Quelle place donnez-vous, dans ce rapport, à la recommandation de
          l'audit interne d'octobre~2025 restée sans suite~? Justifiez.
\end{enumerate}
\end{boxexo}

\bigskip

\begin{boxinfo}
\textbf{La suite du cours.} Ce prologue a posé le réflexe fondateur~:
chercher la décision derrière l'incident. La Partie~\ref{part:fondements}
donne à ce réflexe ses fondements, en commençant par la gouvernance
d'entreprise et la théorie de l'agence (Chapitre~\ref{chap:gouv-entreprise}).
La Partie~\ref{part:cobit} construit avec COBIT~2019 le système de
gouvernance qui a manqué à la BANQUE SANAGA, et la
Partie~\ref{part:pratique} en organise, avec ITIL~4, la pratique
quotidienne.
\end{boxinfo}


    % ======================================================================
    % PARTIE I — FONDEMENTS DE LA GOUVERNANCE DES SI
    % « Gouverner, c'est décider qui décide. »
    % Standards de référence : ISO/IEC 38500, principes de gouvernance
    % d'entreprise de l'OCDE, droit OHADA des sociétés.
    % Niveau : fondamental
    % ======================================================================
    \part{Fondements de la Gouvernance des Systèmes d'Information}
    \label{part:fondements}

        % ============================================================================
% Fichier  : partie1/P01_C01_gouvernance_entreprise_si.tex
% Titre    : De la Gouvernance d'Entreprise à la Gouvernance des SI
% Auteur   : MINKA MI NGUIDJOÏ Thierry Emmanuel
% Version  : 1.0 -- octobre 2026
% ============================================================================

\chapter{De la Gouvernance d'Entreprise à la Gouvernance des SI}
\label{chap:gouv-entreprise}
\chaptermark{Gouvernance d'entreprise et des SI}

\epigraph{%
    « Le gouvernement d'entreprise est le système par lequel les sociétés
    sont dirigées et contrôlées. »%
}{Rapport \textsc{Cadbury}, Royaume-Uni, 1992 (notre traduction)}

\niveauChapitre{Fondamental}{%
    Maîtrise du Chapitre~\ref{chap:prologue} (les trois questions réflexes
    et la distinction entre gouvernance et management).%
}

\begin{boxobjectifs}
\begin{itemize}
    \item Définir la gouvernance d'entreprise et ses mécanismes (conseil,
          comités, contrôle).
    \item Expliquer la théorie de l'agence et la théorie des parties
          prenantes.
    \item Retracer l'émergence de la gouvernance des SI (rapport Cadbury,
          loi Sarbanes-Oxley, OHADA).
    \item Définir la gouvernance des SI et la distinguer du management des
          SI.
    \item Identifier les organes de gouvernance des SI (comité stratégique,
          comité de pilotage).
\end{itemize}
\end{boxobjectifs}

\bigskip

% ============================================================================
\section*{Question directrice : pourquoi le conseil ne savait-il rien~?}
% ============================================================================

Le prologue s'est achevé sur une séance extraordinaire du conseil
d'administration de la BANQUE SANAGA. Un administrateur y a posé la
question qui ouvre ce chapitre~: « Nous avons voté 3,2~milliards FCFA pour
ce projet. Comment avons-nous pu ignorer, jusqu'au jour de la panne, qu'il
était en danger~? »

La question n'est pas technique. Elle porte sur la relation entre ceux qui
possèdent l'entreprise, ceux qui la contrôlent et ceux qui la dirigent~:
qui informe qui, qui surveille qui, et avec quels moyens. Cette relation
est l'objet de la \termecle{gouvernance d'entreprise}, discipline bien plus
ancienne que l'informatique. La gouvernance des systèmes d'information n'en
est pas une invention à part~: elle en est le prolongement, devenu
nécessaire le jour où l'activité de l'entreprise a reposé sur ses systèmes.

Ce chapitre suit ce prolongement en cinq temps~: ce qu'est la gouvernance
d'entreprise (section~\ref{sec:c01-gouv-ent}), les théories qui
l'expliquent (section~\ref{sec:c01-theories}), son histoire et son entrée
dans le champ informatique (section~\ref{sec:c01-histoire}), la définition
de la gouvernance des SI (section~\ref{sec:c01-definition}), et enfin ses
organes (section~\ref{sec:c01-organes}), replacés dans le contexte
camerounais (section~\ref{sec:c01-cameroun}).

\separateur

% ============================================================================
\section{La Gouvernance d'Entreprise : Définition et Mécanismes}
\label{sec:c01-gouv-ent}
% ============================================================================

\subsection{Diriger et contrôler}

La définition la plus citée tient en une phrase, celle de l'épigraphe~: la
gouvernance d'entreprise est le système par lequel les sociétés sont
dirigées et contrôlées~\cite{cadbury1992}. Les deux verbes comptent. Une
entreprise dirigée mais non contrôlée dépend entièrement de la qualité de
ses dirigeants. Une entreprise contrôlée mais mal dirigée se conforme sans
avancer. La gouvernance organise l'équilibre entre les deux.

\begin{boxdef}
\termecle{Gouvernance d'entreprise}\index{gouvernance d'entreprise}
(\textit{corporate governance})~: ensemble des relations, des règles et des
mécanismes par lesquels une société est dirigée et contrôlée. Elle répartit
les droits et les responsabilités entre les actionnaires, le conseil
d'administration, la direction générale et les autres parties prenantes,
et fixe les règles de décision et de reddition de comptes.
\end{boxdef}

Les Principes de gouvernance d'entreprise du G20 et de l'\sigle{OCDE},
publiés pour la première fois en 1999 et révisés en dernier lieu en
2023~\cite{oecd2023principles}, constituent la référence internationale.
Ils traitent des droits des actionnaires, de la transparence de
l'information, du rôle des parties prenantes et des responsabilités du
conseil d'administration.

\subsection{Les mécanismes de la gouvernance}

La gouvernance d'entreprise s'appuie sur deux familles de mécanismes. Les
\textbf{mécanismes internes} sont créés par l'entreprise elle-même~; les
\textbf{mécanismes externes} lui sont imposés ou s'exercent de l'extérieur.

\begin{tableaucours}{L{3.6cm}L{4.4cm}L{4.2cm}}{%
    \tableauHeader{Mécanisme} &
    \tableauHeader{Rôle} &
    \tableauHeader{À la BANQUE SANAGA}
}
\multicolumn{3}{l}{\textit{Mécanismes internes}} \\
Assemblée générale & Désigne les administrateurs, approuve les comptes &
Réunie une fois par an \\
Conseil d'administration & Fixe les orientations, contrôle la direction &
A approuvé le projet HORIZON, puis n'en a plus entendu parler \\
Comités du conseil (audit, risques) & Préparent les décisions du conseil
sur des sujets techniques & Comité d'audit sans point sur le projet \\
Contrôle interne et audit interne & Vérifient que les risques sont maîtrisés
& Alerte d'octobre~2025 restée sans suite \\
\multicolumn{3}{l}{\textit{Mécanismes externes}} \\
Commissaire aux comptes & Certifie les comptes & Interroge la fiabilité des
soldes convertis \\
Régulateur (COBAC) & Contrôle la solidité et la conformité & Demande un
rapport circonstancié \\
Marché et clients & Sanctionnent par la réputation et le départ & Files
d'attente, plaintes, presse \\
\end{tableaucours}
\vspace{-0.8em}
\begin{center}{\small\itshape Tableau~1.1. Mécanismes de la gouvernance d'entreprise}\end{center}

Le tableau montre que la BANQUE SANAGA ne manquait d'aucun mécanisme. Ils
existaient tous, mais aucun ne portait sur le système d'information. Le
conseil votait le budget du projet sans en recevoir le suivi~; le comité
d'audit examinait les comptes sans examiner le risque qui allait les
fausser. C'est la situation typique que la gouvernance des SI vient
corriger.

% ============================================================================
\section{Les Théories qui Expliquent la Gouvernance}
\label{sec:c01-theories}
% ============================================================================

\subsection{La théorie de l'agence}

Pourquoi faut-il contrôler les dirigeants~? La réponse la plus influente
vient de Jensen et Meckling, qui décrivent l'entreprise comme un nœud de
contrats, et la relation entre actionnaires et dirigeants comme une
\termecle{relation d'agence}~\cite{jensen1976agency}.

\begin{boxdef}
\termecle{Relation d'agence}\index{relation d'agence} (\textit{agency
relationship})~: contrat par lequel une personne, le \emph{principal},
engage une autre personne, l'\emph{agent}, pour accomplir en son nom une
tâche qui implique de lui déléguer un pouvoir de décision.

\medskip
\termecle{Coûts d'agence}\index{coûts d'agence} (\textit{agency costs})~:
somme des dépenses de surveillance engagées par le principal, des dépenses
d'engagement (de « dédouanement ») engagées par l'agent pour prouver sa
loyauté, et de la perte résiduelle due aux divergences d'intérêts qui
subsistent malgré tout.
\end{boxdef}

La théorie repose sur deux hypothèses. L'agent a ses propres intérêts, qui
ne coïncident pas toujours avec ceux du principal. Et l'agent en sait plus
que le principal sur ce qu'il fait~: c'est l'\termecle{asymétrie
d'information}\index{asymétrie d'information}. La gouvernance est alors
l'ensemble des dispositifs qui réduisent cette asymétrie (rapports,
audits, comités) et alignent les intérêts (rémunération, responsabilité).

\subsection{La chaîne d'agence des systèmes d'information}

L'informatique allonge la chaîne d'agence et accroît l'asymétrie
d'information à chaque maillon. Le conseil délègue au directeur général,
qui délègue au DSI, qui délègue à un intégrateur. À chaque niveau, l'agent
comprend mieux le système que son principal.

\begin{center}
\begin{tikzpicture}[
    maillon/.style = {
        draw=CouleurBordDef, fill=CouleurDef, rounded corners=4pt,
        text width=2.5cm, minimum height=1.0cm, align=center,
        font=\footnotesize\bfseries\color{CouleurPrimaire}
    },
    fleche/.style = {-{Stealth[length=5pt]}, thick, color=CouleurSecondaire}
]
\node[maillon] (ca)  at (0,0)    {Conseil\\ d'administration};
\node[maillon] (dg)  at (3.9,0)  {Directeur\\ général};
\node[maillon] (dsi) at (7.8,0)  {DSI};
\node[maillon] (int) at (11.7,0) {Intégrateur};
\draw[fleche] (ca) -- node[above, font=\scriptsize\itshape]{délègue} (dg);
\draw[fleche] (dg) -- node[above, font=\scriptsize\itshape]{délègue} (dsi);
\draw[fleche] (dsi) -- node[above, font=\scriptsize\itshape]{délègue} (int);
\draw[fleche, color=CouleurAccent] (0,-0.9) -- node[below, align=center,
    font=\scriptsize\itshape\color{CouleurAccent}]{l'information technique
    décroît en remontant la chaîne} (11.7,-0.9);
\end{tikzpicture}
\end{center}
\vspace{-0.4em}
\begin{center}{\small\itshape Figure~1.1. La chaîne d'agence des systèmes
d'information}\end{center}

\begin{boxscenario}{Opération SANAGA. L'agence}
\textbf{Constatation.} Le directeur général fixe seul la date de bascule
en janvier~2026. La DSI l'alerte par une note le 20~février. Le conseil
n'est informé de rien entre juin~2025 et le 9~mars~2026.

\medskip
\textbf{Analyse} (confiance élevée). Le directeur général, agent du
conseil, avait un intérêt propre au respect de la date~: le lancement
de l'offre mobile, déjà annoncé, était l'objectif sur lequel il serait
jugé. Il disposait d'une information (la note de la DSI) que son principal
n'avait pas. Le conseil n'avait mis en place aucun dispositif pour réduire
cette asymétrie~: ni point d'avancement obligatoire, ni accès direct du
comité d'audit à la DSI. La perte résiduelle s'est élevée à
1,9~milliard FCFA.
\end{boxscenario}

\subsection{La théorie des parties prenantes}

La théorie de l'agence ne connaît que deux acteurs~: le principal et
l'agent. Freeman élargit la perspective~: l'entreprise doit tenir compte
de tous les groupes qui affectent ses objectifs ou sont affectés par
eux~\cite{freeman1984stakeholder}. Pour une banque, ce sont aussi les
déposants, les salariés, le régulateur, l'État et les commerçants qui
dépendent de ses paiements.

\begin{boxdef}
\termecle{Partie prenante}\index{partie prenante} (\textit{stakeholder})~:
tout groupe ou individu qui peut affecter la réalisation des objectifs
d'une organisation ou qui est affecté par elle.
\end{boxdef}

Cette théorie a une conséquence directe pour la gouvernance des SI~: la
valeur d'un système d'information ne se mesure pas seulement pour
l'actionnaire. La panne de la BANQUE SANAGA a touché des clients qui ne
pouvaient plus payer leurs salariés, et des commerçants dont les terminaux
de paiement étaient bloqués. COBIT~2019 retient d'ailleurs, comme premier
principe de tout système de gouvernance, la création de valeur pour les
parties prenantes~\cite{isaca2018cobitintro}.

\begin{boiteRemarque}{Et si les dirigeants étaient loyaux~?}
La théorie de l'agence suppose un agent qui sert d'abord ses propres
intérêts. La théorie de l'intendance (\textit{stewardship theory}) part de
l'hypothèse inverse~: un dirigeant peut trouver sa motivation dans la
réussite de l'organisation elle-même~\cite{davis1997stewardship}. Elle
plaide alors pour la confiance et la délégation plutôt que pour la
surveillance. Les deux lectures ne s'excluent pas. Le directeur général de
la BANQUE SANAGA était sans doute sincèrement convaincu de servir la
banque en tenant sa date. Le problème n'était pas sa loyauté, mais
l'absence d'un mécanisme qui confronte sa conviction aux faits techniques.
Une bonne gouvernance ne présume pas la malhonnêteté~: elle organise la
contradiction.
\end{boiteRemarque}

% ============================================================================
\section{Histoire : de Cadbury à la Gouvernance des SI}
\label{sec:c01-histoire}
% ============================================================================

\subsection{Les scandales fondateurs}

La gouvernance d'entreprise s'est construite par réaction à des scandales.
Au Royaume-Uni, la faillite du groupe Maxwell et celle de la banque BCCI,
en 1991, conduisent à la création du comité présidé par Sir Adrian
Cadbury, dont le rapport de 1992 pose les règles de séparation des
pouvoirs et du rôle des comités~\cite{cadbury1992}. Aux États-Unis, les
faillites d'Enron (2001) et de WorldCom (2002) entraînent l'adoption de la
loi Sarbanes-Oxley~\cite{sox2002}.

\subsection{Sarbanes-Oxley~: l'informatique entre dans la gouvernance}

La loi Sarbanes-Oxley est le moment où l'informatique entre dans la
gouvernance d'entreprise, presque sans l'avoir cherché. Sa section~404
impose à la direction des sociétés cotées aux États-Unis d'évaluer chaque
année l'efficacité de son contrôle interne sur l'information financière,
et à l'auditeur externe d'en attester. Or l'information financière est
produite par des systèmes informatiques. Évaluer le contrôle interne,
c'est donc évaluer aussi les contrôles qui garantissent la fiabilité de
ces systèmes~: qui peut modifier un programme, qui peut accéder aux
données, comment les changements sont autorisés.

\begin{boxdef}
\termecle{Contrôles généraux informatiques}\index{contrôles généraux
informatiques} (\textit{IT general controls}, \sigle{ITGC})~: contrôles qui
s'appliquent à l'ensemble de l'environnement informatique et conditionnent
la fiabilité de toutes les applications~: gestion des accès, gestion des
changements, exploitation, développement et acquisition des systèmes.
\end{boxdef}

Le Chapitre~\ref{chap:audit} reviendra sur ces contrôles, que l'auditeur
teste en priorité. Il suffit ici de noter que les défaillances de la
BANQUE SANAGA en sont deux exemples directs~: un changement non autorisé
et des données dont l'exactitude n'était plus garantie.

\subsection{Les étapes de la gouvernance des SI}

\begin{tableaucours}{L{1.6cm}L{10.6cm}}{%
    \tableauHeader{Année} &
    \tableauHeader{Étape}
}
1992 & Rapport Cadbury (Royaume-Uni)~: séparation des pouvoirs, comités
du conseil \\
1994 & Premier rapport King (Afrique du Sud) \\
1996 & Première édition de COBIT, alors centrée sur l'audit et le contrôle
informatiques \\
1997 & Acte uniforme OHADA relatif au droit des sociétés commerciales et du
groupement d'intérêt économique \\
1998 & Création de l'IT Governance Institute (\sigle{ITGI}) \\
2002 & Loi Sarbanes-Oxley (États-Unis) \\
2004 & Weill et Ross définissent la gouvernance des SI par les droits de
décision~\cite{weill2004itgovernance} \\
2008 & Première édition de la norme ISO/IEC~38500 sur la gouvernance des
technologies de l'information \\
2009 & Rapport King~III~: la gouvernance des technologies de l'information
y devient un chapitre à part entière \\
2014 & Acte uniforme OHADA révisé~: comité d'audit obligatoire dans les
sociétés cotées \\
2016 & Rapport King~IV~: principe~12 sur la gouvernance des technologies et
de l'information~\cite{iod2016king4} \\
2018 & COBIT~2019~: système de gouvernance sur mesure, facteurs de
conception \\
2023 & Révision des Principes de gouvernance d'entreprise du G20 et de
l'OCDE \\
\end{tableaucours}
\vspace{-0.8em}
\begin{center}{\small\itshape Tableau~1.2. Repères historiques de la
gouvernance d'entreprise et de la gouvernance des SI}\end{center}

\subsection{Une voix africaine~: les rapports King}

L'Afrique n'est pas restée en marge de ce mouvement. Les rapports King,
publiés par l'Institute of Directors in Southern Africa, figurent parmi
les codes de gouvernance les plus avancés au monde. Le troisième (2009) a
été l'un des premiers codes nationaux à consacrer un chapitre entier à la
gouvernance des technologies de l'information. Le quatrième (2016) en fait
l'un de ses principes~: l'organe de gouvernance doit gouverner la
technologie et l'information de manière à aider l'organisation à définir
et à atteindre ses objectifs stratégiques~\cite{iod2016king4}. COBIT~2019
cite d'ailleurs le rapport King~IV parmi ses références pour l'objectif de
gouvernance EDM01~\cite{isaca2018cobitobjectifs}.

% ============================================================================
\section{Définir la Gouvernance des SI}
\label{sec:c01-definition}
% ============================================================================

\subsection{Quatre définitions de référence}

Il n'existe pas une définition unique, mais quatre définitions qui se
complètent. Chacune met l'accent sur un aspect différent.

\begin{boxdef}
\textbf{COBIT~4.1 (ITGI, 2007, version française de l'AFAI)}~:
« La gouvernance des SI est de la responsabilité des dirigeants et du
conseil d'administration, et elle est constituée des structures et
processus de commandement et de fonctionnement qui conduisent
l'informatique de l'entreprise à soutenir les stratégies et les objectifs
de l'entreprise, et à lui permettre de les élargir. »~\cite{itgi2007cobit41}

\medskip
\textbf{Weill et Ross (2004)}~: spécifier les droits de décision et le
cadre de responsabilité qui encouragent les comportements souhaitables
dans l'usage des technologies de l'information~\cite{weill2004itgovernance}.

\medskip
\textbf{ISO/IEC~38500}~: système par lequel l'utilisation actuelle et
future de l'informatique est dirigée et contrôlée~\cite{iso38500}.

\medskip
\textbf{COBIT~2019}~: activités par lesquelles l'organe de gouvernance
évalue les besoins des parties prenantes, dirige par la priorisation et la
décision, et surveille la performance et la conformité~\cite{isaca2018cobitintro}.
\end{boxdef}

\begin{tableaucours}{L{2.8cm}L{4.4cm}L{5.0cm}}{%
    \tableauHeader{Source} &
    \tableauHeader{Accent principal} &
    \tableauHeader{Question qu'elle pose}
}
COBIT 4.1 & Responsabilité des dirigeants, structures et processus & Qui
porte la gouvernance, et par quels dispositifs~? \\
Weill et Ross & Droits de décision et comportements & Qui décide de quoi,
et quels comportements cela produit-il~? \\
ISO/IEC 38500 & Direction et contrôle de l'usage de l'informatique &
L'usage de l'informatique est-il dirigé et contrôlé~? \\
COBIT 2019 & Évaluer, diriger, surveiller & Les besoins sont-ils évalués,
les priorités fixées, les résultats suivis~? \\
\end{tableaucours}
\vspace{-0.8em}
\begin{center}{\small\itshape Tableau~1.3. Comparaison des définitions de
la gouvernance des SI}\end{center}

\subsection{Une définition de travail pour ce cours}

De ces quatre définitions, ce cours retient une synthèse, qui servira de
référence dans tous les chapitres.

\begin{boxdef}
\termecle{Gouvernance des SI}\index{gouvernance des SI!définition de
travail}~: ensemble des responsabilités, des droits de décision et des
mécanismes par lesquels l'organe de gouvernance d'une organisation
s'assure que l'usage, présent et futur, de l'information et des
technologies \emph{crée de la valeur} pour les parties prenantes,
\emph{maîtrise les risques} qui en découlent et \emph{utilise
judicieusement les ressources} qui lui sont consacrées.
\end{boxdef}

Les trois finalités soulignées (valeur, risque, ressources) seront
développées au Chapitre~\ref{chap:alignement}. Elles donnent déjà un
critère simple~: une décision relève de la gouvernance des SI lorsqu'elle
engage durablement l'une de ces trois finalités pour l'organisation
entière.

\subsection{Structures, processus et mécanismes relationnels}

De Haes et Van Grembergen proposent de décrire la gouvernance des SI d'une
organisation par trois familles de mécanismes~\cite{dehaes2015egit}. La
grille est utile en diagnostic~: une organisation qui n'a que des
structures, sans processus ni relations, a une gouvernance de façade.

\begin{tableaucours}{L{2.6cm}L{4.6cm}L{5.0cm}}{%
    \tableauHeader{Famille} &
    \tableauHeader{Exemples} &
    \tableauHeader{Signal d'absence}
}
Structures & Comité stratégique des SI, comité de pilotage, rôle de DSI
membre de la direction & Personne n'est désigné pour trancher \\
Processus & Planification stratégique des SI, gestion du portefeuille,
suivi de la valeur, audit & Les décisions se prennent au cas par cas,
sans calendrier ni critères \\
Mécanismes relationnels & Formation des administrateurs, communication
entre métiers et DSI, binômes métier-informatique & Les métiers découvrent
les décisions informatiques après coup \\
\end{tableaucours}
\vspace{-0.8em}
\begin{center}{\small\itshape Tableau~1.4. Les trois familles de
mécanismes de gouvernance des SI}\end{center}

% ============================================================================
\section{Gouvernance et Management des SI : Quatre Confusions}
\label{sec:c01-confusions}
% ============================================================================

Le prologue a posé la distinction entre gouvernance et management
(section~\ref{sec:pro-gouv-management}). Dans la pratique, elle est
brouillée par quatre confusions fréquentes. Chacune conduit à une
gouvernance qui n'en a que le nom.

\begin{enumerate}
    \item \textbf{« La gouvernance des SI, c'est l'affaire du DSI. »} Le DSI
          est un acteur central du management des SI. La gouvernance, elle,
          relève de l'organe qui le contrôle~: le conseil et ses comités.
          Confier la gouvernance au DSI, c'est demander à l'agent de se
          surveiller lui-même.
    \item \textbf{« La gouvernance, c'est la conformité. »} Respecter la
          loi et les normes est une exigence de la gouvernance, mais ne
          suffit pas à la définir. Une organisation peut être conforme et
          investir dans des projets sans valeur.
    \item \textbf{« La gouvernance, c'est la sécurité. »} La sécurité de
          l'information est un domaine de risque parmi d'autres. La
          gouvernance arbitre entre la sécurité, la valeur et le coût~;
          elle ne se réduit à aucun des trois.
    \item \textbf{« La gouvernance, c'est de la bureaucratie. »} Une
          gouvernance qui multiplie les comités sans décisions alourdit
          l'organisation. Une bonne gouvernance est proportionnée~: le
          Chapitre~\ref{chap:cobit-conception} montrera comment COBIT~2019
          l'adapte à la taille et aux risques de chaque organisation.
\end{enumerate}

\begin{boxalert}
\textbf{Violation typique.} Après la panne, le conseil d'administration de
la BANQUE SANAGA envisage de « confier la gouvernance des SI » à un DSI
adjoint nouvellement recruté, chargé de produire des rapports mensuels au
directeur général.

\smallskip
\textbf{Règle~:} on ne corrige pas une défaillance de gouvernance en
ajoutant un acteur de management. Il faut un dispositif qui rende compte
\emph{au conseil}, et non à celui dont les décisions doivent être
contrôlées.
\end{boxalert}

% ============================================================================
\section{Les Organes de la Gouvernance des SI}
\label{sec:c01-organes}
% ============================================================================

\subsection{Deux niveaux de comités}

L'ITGI a popularisé une distinction devenue classique entre deux comités.
Le \termecle{comité stratégique des SI}\index{comité stratégique des SI}
(\textit{IT strategy committee}) se situe au niveau du conseil~: il
l'assiste sur les orientations, la valeur et les risques du numérique. Le
\termecle{comité de pilotage des SI}\index{comité de pilotage des SI}
(\textit{IT steering committee}) se situe au niveau de la direction
générale~: il priorise les projets, arbitre les ressources et suit leur
exécution.

\begin{tableaucours}{L{2.8cm}L{2.2cm}L{3.4cm}L{3.8cm}}{%
    \tableauHeader{Organe} &
    \tableauHeader{Niveau} &
    \tableauHeader{Composition type} &
    \tableauHeader{Décisions}
}
Conseil d'administration & Gouvernance & Administrateurs & Approuve la
stratégie numérique, l'appétence au risque, les grands investissements \\
Comité d'audit, comité des risques & Gouvernance & Administrateurs
indépendants & Surveille le contrôle interne et les risques, dont ceux du
SI \\
Comité stratégique des SI & Gouvernance & Administrateurs, DG, DSI,
experts invités & Évalue l'alignement, la valeur et les risques du
portefeuille numérique \\
Comité de pilotage des SI & Management & DG ou DGA, directeurs métiers,
DSI & Priorise les projets, arbitre les ressources, décide des bascules \\
Comité des changements & Management & DSI, exploitation, sécurité,
métiers & Autorise les changements majeurs (Chapitre~\ref{chap:itil-changements}) \\
\end{tableaucours}
\vspace{-0.8em}
\begin{center}{\small\itshape Tableau~1.5. Les organes de la gouvernance et
du management des SI}\end{center}

\subsection{La gouvernance du système de gouvernance}

COBIT~2019 consacre son premier objectif à la mise en place et au maintien
de ce dispositif lui-même. C'est l'objectif qui manquait le plus
visiblement à la BANQUE SANAGA.

\begin{boxobjet}{EDM01 · Assurer la définition et le maintien du cadre de gouvernance}
\textbf{Finalité~:} doter l'organisation d'une approche cohérente de la
gouvernance de l'information et des technologies, intégrée à la
gouvernance d'entreprise, avec des décisions alignées sur sa stratégie.

\textbf{Pratiques~:} EDM01.01 \emph{évaluer} le système de gouvernance
(identifier les parties prenantes et leurs exigences)~; EDM01.02
\emph{diriger} le système de gouvernance (définir les structures, les
principes, les niveaux d'autorité et l'information nécessaire aux
décisions)~; EDM01.03 \emph{surveiller} le système de gouvernance (vérifier
qu'il fonctionne et assure une supervision effective).

\textbf{Sorties~:} principes de gouvernance, modèle de décision et niveaux
d'autorité, rapports sur l'efficacité du système de gouvernance.

\textbf{RACI (extrait)~:} \emph{A} au conseil d'administration, \emph{R} au
directeur général et au comité stratégique des SI, \emph{C} au DSI et à
l'audit interne.

\textbf{Indicateurs~:} part des décisions stratégiques prises par
l'instance prévue~; délai de traitement des alertes remontées au conseil.
\end{boxobjet}
{\small\itshape Intitulés et pratiques traduits par nos soins à partir de
COBIT~2019~\cite{isaca2018cobitobjectifs}.}

\begin{boxprocedure}
\textbf{Un comité de gouvernance qui fonctionne.} Pour qu'un auditeur
puisse établir que le comité stratégique des SI existe autrement que sur
le papier, il doit trouver~:
\begin{itemize}
    \item une charte approuvée par le conseil (mission, composition,
          pouvoirs, fréquence, quorum)~;
    \item un calendrier de réunions respecté, et des procès-verbaux datés
          et signés~;
    \item des décisions inscrites au procès-verbal, avec leur responsable
          et leur échéance~;
    \item un suivi de ces décisions d'une séance à l'autre~;
    \item un compte rendu régulier au conseil d'administration.
\end{itemize}
\end{boxprocedure}

\begin{boxmetier}{Responsable GRC}
\textbf{Rédiger la charte d'un comité stratégique des SI}
\begin{enumerate}
    \item \textbf{Mission~:} assister le conseil dans l'évaluation de la
          stratégie numérique, de la valeur des investissements et des
          risques liés au SI.
    \item \textbf{Composition~:} au moins deux administrateurs, dont l'un
          préside~; le directeur général et le DSI comme membres~; l'audit
          interne comme invité permanent.
    \item \textbf{Pouvoirs~:} accès direct à la DSI et à l'audit interne~;
          droit de demander toute expertise externe.
    \item \textbf{Fréquence~:} au moins trimestrielle, et sur convocation
          en cas de projet majeur en difficulté.
    \item \textbf{Reddition de comptes~:} rapport écrit au conseil après
          chaque séance.
\end{enumerate}
Le point décisif est le deuxième pouvoir~: un comité qui ne reçoit
l'information que par le directeur général reproduit l'asymétrie qu'il
devait réduire.
\end{boxmetier}

% ============================================================================
\section{Le Contexte Camerounais et Africain}
\label{sec:c01-cameroun}
% ============================================================================

\subsection{Des organisations, des cadres différents}

La gouvernance des SI ne s'exerce pas de la même manière dans toutes les
organisations camerounaises, parce que leur gouvernance d'entreprise
relève de cadres juridiques différents.

\begin{tableaucours}{L{3.2cm}L{4.6cm}L{4.4cm}}{%
    \tableauHeader{Organisation} &
    \tableauHeader{Cadre de gouvernance} &
    \tableauHeader{Cas du cours}
}
Société commerciale & Acte uniforme OHADA sur les sociétés commerciales
(révisé en 2014) & KRIBILOG SA, NOUN PAY \\
Établissement de crédit & Droit OHADA et réglementation de la COBAC &
BANQUE SANAGA \\
Entreprise publique & Loi n°~2017/011 du 12~juillet~2017 & · \\
Établissement public & Loi n°~2017/010 du 12~juillet~2017 & ANPP
(opération MEFOU) \\
\end{tableaucours}
\vspace{-0.8em}
\begin{center}{\small\itshape Tableau~1.6. Cadres de gouvernance des
organisations camerounaises}\end{center}

Les lois de 2017 sur les établissements publics et les entreprises
publiques, précisées par le décret n°~2019/320 du 19~juin~2019, organisent
notamment le conseil d'administration et la tutelle de l'État. Le
Chapitre~\ref{chap:droit} en étudiera les dispositions.

% A VERIFIER au Journal officiel de l'OHADA : libellé exact de l'article
% 829-1 de l'AUSCGIE (source utilisée : brochure de présentation de la
% réforme de 2014, PwC/OHADA).
\begin{boxjuris}
\textbf{Acte uniforme OHADA relatif au droit des sociétés commerciales et
du groupement d'intérêt économique, révisé le 30~janvier~2014,
article~829-1}~\cite{ohada2014auscgie}.

\medskip
Selon la présentation de la réforme publiée par l'OHADA, l'Acte uniforme
révisé rend le comité d'audit obligatoire dans les sociétés cotées. Ce
comité, composé exclusivement d'administrateurs non salariés et n'exerçant
pas de mandat de direction générale, assure notamment le suivi de
l'efficacité des systèmes de contrôle interne et de gestion des risques,
et rend compte régulièrement au conseil d'administration.

\medskip
\textbf{Portée.} Le droit OHADA institue ainsi, pour les sociétés cotées,
une instance du conseil chargée de surveiller les risques. Rien ne
l'oblige à examiner les risques du SI, mais rien ne l'en dispense~: le
système d'information porte aujourd'hui l'essentiel du contrôle interne.
Pour les sociétés non cotées, comme la BANQUE SANAGA dans notre cas, la
création d'un tel comité relève des statuts et de la réglementation
sectorielle.
\end{boxjuris}

\subsection{Quatre défis propres au contexte}

\begin{enumerate}
    \item \textbf{La concentration du pouvoir.} Dans de nombreuses
          organisations, le directeur général concentre la décision, et le
          conseil se réunit peu. La chaîne d'agence y est courte en
          apparence, mais l'asymétrie d'information y est forte.
    \item \textbf{La tutelle.} Dans le secteur public, une partie des
          décisions relève du ministère de tutelle. La gouvernance des SI
          doit préciser ce qui se décide au conseil et ce qui remonte à la
          tutelle.
    \item \textbf{La dépendance aux prestataires.} Les grands systèmes sont
          souvent conçus et maintenus par des intégrateurs étrangers. Le
          dernier maillon de la chaîne d'agence échappe alors au droit et
          au contrôle local, d'où l'importance des contrats
          (Chapitre~\ref{chap:itil-niveaux}).
    \item \textbf{La compétence numérique des conseils.} Peu
          d'administrateurs ont une expérience des systèmes d'information.
          La formation des administrateurs est un mécanisme relationnel de
          première importance.
\end{enumerate}

\begin{boxscenario}{Opération MEFOU. Une agence sans comité stratégique}
L'AGENCE NATIONALE DES PAIEMENTS PUBLICS (ANPP), établissement public
installé à Yaoundé, gère le paiement dématérialisé des salaires et des
fournisseurs de l'État. Son schéma directeur informatique, adopté en 2023
pour 14~milliards FCFA, comprend neuf projets. Le conseil d'administration
se réunit deux fois par an et consacre en moyenne vingt minutes au
numérique. Aucun comité ne suit le schéma directeur. Les projets sont
lancés sur décision du directeur général, M.~ESSOMBA, souvent à la demande
du ministère de tutelle.

\medskip
\textbf{Constatation.} À mi-parcours, trois projets ont consommé 71~\% du
budget et aucun n'est achevé. Le conseil l'apprend par une question d'un
administrateur représentant le ministère des Finances.

\medskip
Ce cas reviendra au Chapitre~\ref{chap:alignement} (arbitrage du
portefeuille) et au Chapitre~\ref{chap:audit} (audit de la Chambre des
comptes).
\end{boxscenario}

% ============================================================================
\separateur
% ============================================================================

\begin{boxresume}
\textbf{Ce qu'il faut retenir de ce chapitre~:}
\begin{itemize}
    \item \textbf{Diriger et contrôler}~: la gouvernance d'entreprise est le
          système par lequel les sociétés sont dirigées et contrôlées. Elle
          repose sur des mécanismes internes (conseil, comités, contrôle
          interne) et externes (commissaire aux comptes, régulateur,
          marché).
    \item \textbf{La théorie de l'agence}~: le dirigeant, agent des
          actionnaires, en sait plus qu'eux. La gouvernance réduit cette
          asymétrie d'information. L'informatique allonge la chaîne
          d'agence et l'accroît.
    \item \textbf{Les parties prenantes}~: la valeur d'un SI se mesure pour
          tous ceux qu'il affecte, et non pour le seul actionnaire.
    \item \textbf{L'informatique est entrée dans la gouvernance par le
          contrôle interne}~: la loi Sarbanes-Oxley a fait des contrôles
          généraux informatiques une condition de la fiabilité financière.
    \item \textbf{La gouvernance des SI} est l'ensemble des
          responsabilités, des droits de décision et des mécanismes par
          lesquels l'organe de gouvernance s'assure que l'usage de
          l'information et des technologies crée de la valeur, maîtrise les
          risques et utilise judicieusement les ressources.
    \item \textbf{Deux niveaux de comités}~: le comité stratégique des SI
          relève du conseil, le comité de pilotage relève de la direction
          générale. Ne pas les confondre.
    \item \textbf{Le contexte africain} a ses repères (rapports King, droit
          OHADA, lois camerounaises de 2017) et ses défis (concentration du
          pouvoir, tutelle, dépendance aux prestataires, compétence des
          conseils).
\end{itemize}
\end{boxresume}

\bigskip

\begin{boxexo}[title={\ding{45}\quad Exercice~1.1. Gouvernance ou management~? \niveaubase}]
Classez chacune des décisions suivantes dans la gouvernance ou dans le
management des SI, et indiquez l'organe qui devrait la prendre
(Tableau~1.5).
\begin{enumerate}[(a)]
    \item Fixer la part du budget de la banque consacrée au numérique pour
          les trois prochaines années.
    \item Choisir entre deux éditeurs pour le nouveau logiciel de gestion
          des crédits.
    \item Décider de reporter la bascule d'un système critique.
    \item Fixer le niveau de risque informatique que la banque accepte de
          supporter.
    \item Autoriser l'installation d'un correctif sur les serveurs de
          messagerie.
\end{enumerate}
\end{boxexo}

\begin{boxexo}[title={\ding{45}\quad Exercice~1.2. La chaîne d'agence de KRIBILOG \niveaumoyen}]
La société KRIBILOG SA confie l'exploitation de son logiciel de suivi des
conteneurs à l'infogérant NETPORT SERVICES. Le conseil d'administration ne
reçoit aucun indicateur sur ce service~; le responsable des services
informatiques, M.~MBARGA, reçoit chaque mois un rapport rédigé par
NETPORT SERVICES elle-même.
\begin{enumerate}[(a)]
    \item Dessinez la chaîne d'agence de ce service, sur le modèle de la
          Figure~1.1.
    \item Identifiez, à chaque maillon, l'asymétrie d'information.
    \item Proposez, pour chaque maillon, un mécanisme qui la réduit, en
          précisant s'il s'agit d'une structure, d'un processus ou d'un
          mécanisme relationnel (Tableau~1.4).
\end{enumerate}
\end{boxexo}

\begin{boxexo}[title={\ding{45}\quad Exercice~1.3. Doter l'ANPP d'une gouvernance des SI \niveauexpert}]
Le nouveau président du conseil d'administration de l'ANPP (opération
MEFOU) vous demande une note de trois pages pour doter l'agence d'une
gouvernance des SI.
\begin{enumerate}[(a)]
    \item Rédigez la constatation de départ, à partir des faits de
          l'opération MEFOU, en la séparant de votre analyse.
    \item Proposez les organes à créer, leur niveau (gouvernance ou
          management) et leur composition, en tenant compte de la tutelle
          ministérielle.
    \item Rédigez les cinq articles essentiels de la charte du comité
          stratégique des SI.
    \item Indiquez les trois pièces qu'un auditeur demandera, dans un an,
          pour vérifier que ce dispositif fonctionne réellement.
\end{enumerate}
\end{boxexo}

\bigskip

\begin{boxinfo}
\textbf{Transition.} Ce chapitre a défini la gouvernance des SI et ses
organes à partir de la gouvernance d'entreprise. Le
Chapitre~\ref{chap:defaillances} met ces notions à l'épreuve des faits~:
il analyse de grandes défaillances réelles et documentées, pour montrer
qu'à chaque fois une décision de gouvernance a manqué derrière l'incident
technique.
\end{boxinfo}

        % ============================================================================
% Fichier  : partie1/P01_C02_grandes_defaillances.tex
% Titre    : Grandes Défaillances et Leçons de Gouvernance
% Auteur   : MINKA MI NGUIDJOÏ Thierry Emmanuel
% Version  : 1.0 -- octobre 2026
% ============================================================================
% Chapitre-cas : une section par affaire (brief, déroulé, révélations,
% synthèse), puis lecture croisée et grille d'erreurs.
% Sources : uniquement des documents publics de première main (décisions
% des régulateurs, analyses de cause publiées par les entreprises).
% ============================================================================

\chapter{Grandes Défaillances et Leçons de Gouvernance}
\label{chap:defaillances}
\chaptermark{Grandes défaillances}

\epigraph{%
    « Chaque grande panne informatique raconte deux histoires~: celle d'un
    composant qui a cédé, et celle d'une décision que personne n'a
    prise. »%
}{Principe d'analyse des défaillances}

\niveauChapitre{Fondamental}{%
    Maîtrise du Chapitre~\ref{chap:gouv-entreprise} (théorie de l'agence,
    gouvernance et management, organes de gouvernance).%
}

\begin{boxobjectifs}
\begin{itemize}
    \item Analyser des défaillances documentées à partir de sources
          publiques.
    \item Isoler, pour chaque cas, la défaillance de gouvernance derrière
          l'incident technique.
    \item Distinguer constatation et analyse dans un rapport de défaillance.
    \item Construire une grille d'erreurs de gouvernance réutilisable.
\end{itemize}
\end{boxobjectifs}

\bigskip

% ============================================================================
\section*{Mode d'emploi du chapitre}
% ============================================================================

Le prologue reposait sur un cas fictif, construit pour l'enseignement. Ce
chapitre passe aux faits. Il étudie trois défaillances réelles, choisies
parce qu'elles sont \textbf{documentées par des sources de première main}~:
la décision d'un régulateur des marchés, celle de deux régulateurs
bancaires, et l'analyse de cause publiée par un éditeur de logiciels. Elles
couvrent trois situations que tout DSI rencontrera~: un déploiement de
code, une migration de plateforme, et la mise à jour d'un logiciel fourni
par un tiers.

Chaque affaire suit le même plan~: le \textbf{brief} (qui, quand, quel
enjeu), le \textbf{déroulé} (la chronologie), les \textbf{révélations}
(ce que les enquêtes ont établi) et la \textbf{synthèse} (la cause de
gouvernance). La section~\ref{sec:c02-croisee} croise les trois cas, et
la section~\ref{sec:c02-grille} en tire une grille d'erreurs que vous
réutiliserez tout au long du cours.

\begin{boxalert}
\textbf{Règle de lecture.} Les faits rapportés dans ce chapitre sont ceux
qu'établissent les documents cités. Les analyses en termes de gouvernance
sont les nôtres~: elles sont signalées comme telles, avec leur niveau de
confiance. Les organisations citées ne sont pas jugées~; leurs défaillances
sont étudiées parce qu'elles ont eu le mérite d'être rendues publiques.
\end{boxalert}

\separateur

% ============================================================================
\section{Une Méthode d'Analyse des Défaillances}
\label{sec:c02-methode}
% ============================================================================

\subsection{Les sources}

Une analyse de défaillance vaut ce que valent ses sources. Les articles de
presse et les billets techniques sont utiles pour s'orienter, mais ils
simplifient et se recopient. Seules les sources de première main
permettent d'établir des constatations.

\begin{boxdef}
\termecle{Source de première main}\index{source de première main}
(\textit{primary source})~: document produit par une partie qui a eu
accès direct aux faits et qui engage sa responsabilité sur leur
exactitude~: décision d'un régulateur ou d'un tribunal, rapport d'enquête
indépendant commandé par le conseil, analyse de cause publiée par
l'entreprise elle-même.
\end{boxdef}

\begin{tableaucours}{L{3.6cm}L{4.2cm}L{4.4cm}}{%
    \tableauHeader{Type de source} &
    \tableauHeader{Force} &
    \tableauHeader{Limite}
}
Décision d'un régulateur & Faits contradictoirement établis, sanctions
motivées & Centrée sur les manquements à la réglementation \\
Rapport d'enquête indépendant & Examen large, accès aux pièces internes &
Commandé par l'entreprise, parfois publié en partie seulement \\
Analyse de cause de l'éditeur & Précision technique & Rédigée par la partie
en cause, prudente sur les responsabilités \\
Presse spécialisée & Rapidité, contexte & Chiffres approximatifs,
sources secondaires \\
\end{tableaucours}
\vspace{-0.8em}
\begin{center}{\small\itshape Tableau~2.1. Hiérarchie des sources d'une
analyse de défaillance}\end{center}

\subsection{La grille d'analyse}

Pour chaque affaire, sept questions structurent l'analyse. Elles reprennent
les réflexes du prologue (Chapitre~\ref{chap:prologue}).

\begin{enumerate}
    \item \textbf{Quel était l'enjeu~?} Activité, taille, dépendance au
          système.
    \item \textbf{Que s'est-il passé, et quand~?} Chronologie établie par
          les sources.
    \item \textbf{Quelle est la cause technique~?} Le composant, le code,
          la donnée.
    \item \textbf{Quelle décision a manqué~?} Qui décide, qui exécute, qui
          rend compte.
    \item \textbf{Quels contrôles auraient dû jouer~?} Avant, pendant et
          après l'incident.
    \item \textbf{Qui a demandé des comptes, et avec quelles
          conséquences~?} Régulateur, clients, marché, dirigeants.
    \item \textbf{Quelle leçon pour une organisation camerounaise~?}
\end{enumerate}

% ============================================================================
\section{Knight Capital, 2012 : 45 Minutes, 460 Millions de Dollars}
\label{sec:c02-knight}
% ============================================================================

\subsection{Brief}

Knight Capital Americas était, en 2012, l'un des principaux teneurs de
marché et intermédiaires sur les actions américaines. Son activité
reposait entièrement sur des systèmes automatisés d'acheminement des
ordres, dont le système SMARS. Les faits qui suivent sont ceux qu'établit
la décision de la Securities and Exchange Commission (\sigle{SEC}) du
16~octobre~2013~\cite{sec2013knight}.

\subsection{Déroulé}

Pour participer à un nouveau programme de la Bourse de New York, Knight
déploie un nouveau code sur les huit serveurs de SMARS dans les jours qui
précèdent le 1\textsuperscript{er}~août~2012. Ce nouveau code réutilise un
indicateur qui servait autrefois à une ancienne fonction, appelée « Power
Peg », abandonnée depuis des années mais jamais retirée du programme. Un
technicien déploie le code sur sept serveurs seulement~; le huitième garde
l'ancien code.

Le 1\textsuperscript{er}~août, à l'ouverture du marché, les ordres qui
portent le nouvel indicateur réveillent l'ancienne fonction sur le
huitième serveur. En 45~minutes, alors qu'il traitait 212~ordres de
clients, SMARS envoie des millions d'ordres sur le marché. Il en résulte
plus de 4~millions d'exécutions sur 154~valeurs, pour plus de
397~millions d'actions. Pour tenter de corriger, les équipes retirent le
nouveau code des sept serveurs où il avait été correctement installé, ce
qui étend le problème à l'ensemble des serveurs.

\subsection{Révélations}

La décision de la SEC établit notamment que~:
\begin{itemize}
    \item Knight ne disposait pas de procédures écrites de déploiement qui
          imposent une revue par un second technicien~;
    \item aucun contrôle ne vérifiait que le code était identique sur tous
          les serveurs~;
    \item avant l'ouverture du marché, le système avait émis des messages
          automatiques signalant une anomalie, que personne n'était chargé
          de traiter~;
    \item l'entreprise n'avait pas de procédure de réponse aux incidents
          adaptée à ce type de défaillance.
\end{itemize}
Knight a perdu plus de 460~millions de dollars sur les positions ainsi
créées. Elle a dû lever des fonds en urgence auprès d'investisseurs, puis
a été absorbée l'année suivante. La SEC lui a infligé une amende de
12~millions de dollars pour manquement à la règle sur l'accès au marché,
qui impose aux intermédiaires des contrôles de gestion des risques sur
leurs ordres automatisés.

\subsection{Synthèse}

\begin{boxscenario}{Knight Capital. Lecture de gouvernance}
\textbf{Cause technique} (établie). Un code obsolète, réactivé par un
indicateur réutilisé, sur un serveur où le déploiement était incomplet.

\medskip
\textbf{Cause de gouvernance} (analyse, confiance élevée). La gestion des
changements reposait sur le savoir-faire individuel, sans procédure
écrite, sans double contrôle et sans vérification après déploiement. La
dette technique (un code mort laissé en place) n'avait jamais été
arbitrée. Et le dispositif de surveillance produisait des alertes sans
qu'un rôle soit chargé de les traiter.

\medskip
\textbf{Les trois questions.} \emph{Qui décide} qu'un déploiement est
terminé~? Personne, formellement. \emph{Qui exécute}~? Un technicien seul.
\emph{Qui rend compte} des alertes avant l'ouverture~? Personne.
\end{boxscenario}

% ============================================================================
\section{TSB, 2018 : la Migration qui a Paralysé une Banque}
\label{sec:c02-tsb}
% ============================================================================

\subsection{Brief}

TSB Bank est une banque de détail britannique, rachetée en 2015 par le
groupe espagnol Sabadell. Elle utilisait jusqu'alors la plateforme
informatique de son ancien propriétaire. Le programme de migration vers
une nouvelle plateforme, fournie par une filiale informatique du groupe
Sabadell, devait lui donner son autonomie technique. Les faits qui
suivent sont ceux qu'établissent les décisions conjointes de la Financial
Conduct Authority (\sigle{FCA}) et de la Prudential Regulation Authority
(\sigle{PRA}) de décembre~2022~\cite{fca2022tsb}.

\subsection{Déroulé}

La migration des données des clients et de l'entreprise vers la nouvelle
plateforme a lieu lors d'un week-end d'avril~2018. Dès la remise en
service, toutes les agences et une part importante des 5,2~millions de
clients de la banque sont touchées~: accès impossibles aux services en
ligne et mobiles, opérations erronées, centres d'appel saturés. Le
rétablissement complet du service prend plusieurs mois. Le directeur
général quitte ses fonctions à l'automne~2018.

\subsection{Révélations}

Les régulateurs ont établi que TSB n'avait pas organisé ni contrôlé de
manière adéquate son programme de migration, et qu'elle n'avait pas géré
de manière adéquate les risques opérationnels nés de son externalisation
auprès d'un fournisseur tiers critique. La banque a versé 32,7~millions de
livres de réparation aux clients lésés. Les deux régulateurs lui ont
infligé une amende totale de 48,65~millions de livres (29,75~millions pour
la FCA, 18,9~millions pour la PRA), après une réduction de 30~\% accordée
pour son règlement amiable. En 2023, la PRA a en outre sanctionné
personnellement l'ancien directeur des systèmes d'information de la banque
pour ne pas avoir pris de mesures raisonnables de supervision du
prestataire~\cite{pra2023abarca}.

\begin{boxinfo}
\textbf{Un précédent pour tous les DSI.} La sanction personnelle de
l'ancien DSI de TSB montre que la redevabilité n'est pas une notion
abstraite. Dans le régime britannique de responsabilité des dirigeants du
secteur financier, le titulaire d'une fonction clé répond personnellement
des mesures raisonnables qu'il a prises, ou omis de prendre, dans son
périmètre. Externaliser une activité ne transfère pas la responsabilité de
la superviser.
\end{boxinfo}

\subsection{Synthèse}

\begin{boxscenario}{TSB. Lecture de gouvernance}
\textbf{Cause technique} (établie dans ses grandes lignes). Une plateforme
et un fournisseur qui n'étaient pas prêts à absorber la charge et les cas
réels de l'ensemble des clients au moment de la bascule.

\medskip
\textbf{Cause de gouvernance} (analyse, confiance élevée). Le programme n'a
pas été piloté avec un contrôle suffisant de son état de préparation, et
la banque s'est fiée à son fournisseur, entité de son propre groupe, sans
superviser elle-même les risques qu'il portait. La décision de basculer a
été prise sans que l'organe de décision dispose d'une évaluation
indépendante de cet état de préparation.

\medskip
\textbf{Les trois questions.} \emph{Qui décide} de la bascule~? La
direction, sur la foi des assurances du programme. \emph{Qui exécute}~?
Un fournisseur dont la banque ne maîtrisait pas les risques. \emph{Qui
rend compte}~? La banque, son directeur général et son DSI, devant les
régulateurs.
\end{boxscenario}

La ressemblance avec l'opération SANAGA n'est pas fortuite~: le cas fictif
du prologue reprend des mécanismes que les décisions des régulateurs
britanniques ont rendus publics. Ils ne sont propres ni à une banque ni à
un pays.

% ============================================================================
\section{CrowdStrike, 2024 : la Mise à Jour qui a Arrêté le Monde}
\label{sec:c02-crowdstrike}
% ============================================================================

\subsection{Brief}

CrowdStrike est un éditeur américain de solutions de cybersécurité. Son
logiciel Falcon protège les postes de travail et les serveurs de très
nombreuses organisations. Pour détecter les attaques, il s'exécute au cœur
du système d'exploitation Windows et reçoit plusieurs fois par jour des
mises à jour de contenu, qui décrivent de nouveaux comportements
malveillants. Les faits techniques qui suivent sont ceux qu'établit
l'analyse de cause publiée par l'éditeur le 6~août~2024~\cite{crowdstrike2024rca}.

\subsection{Déroulé}

Le 19~juillet~2024, CrowdStrike diffuse une mise à jour de contenu pour
son capteur Windows. Elle provoque l'arrêt brutal des machines qui la
reçoivent, puis leur redémarrage en boucle. L'éditeur retire la mise à
jour dans l'heure et demie qui suit, mais les machines déjà touchées ne
peuvent plus la recevoir~: il faut souvent intervenir manuellement, poste
par poste. Microsoft estime à environ 8,5~millions le nombre d'appareils
Windows touchés~\cite{microsoft2024crowdstrike}. Des compagnies aériennes,
des hôpitaux, des banques et des médias sont paralysés pendant plusieurs
jours.

\subsection{Révélations}

L'analyse de cause établit que le capteur attendait 20~champs de données,
alors que la mise à jour en fournissait 21. Ce décalage a provoqué une
lecture hors des limites de la mémoire, et l'arrêt du système. Le défaut
n'avait pas été détecté par les contrôles de validation du contenu.
L'éditeur annonce, parmi ses mesures correctives, davantage de tests, un
déploiement progressif de ces mises à jour, et la possibilité pour les
clients de choisir où et quand elles s'installent.

\subsection{Synthèse}

Ce cas se lit à deux niveaux de gouvernance~: celle de l'éditeur, et celle
de ses clients.

\begin{boxscenario}{CrowdStrike. Lecture de gouvernance}
\textbf{Chez l'éditeur} (analyse, confiance moyenne). Les mises à jour de
contenu n'étaient pas soumises à la même rigueur de gestion des
changements que les versions du logiciel, ni déployées progressivement.
Un type de changement fréquent et jugé peu risqué échappait aux contrôles
les plus stricts.

\medskip
\textbf{Chez les clients} (analyse, confiance élevée). Des milliers
d'organisations avaient accordé à un fournisseur unique le droit de
modifier, sans délai ni contrôle de leur part, un composant qui s'exécute
au cœur de tous leurs postes. Ce choix est rationnel pour la sécurité~:
une protection doit être à jour. Mais il crée une \termecle{dépendance
critique}\index{dépendance critique} qui n'avait, dans la plupart des cas,
jamais été évaluée ni arbitrée par les instances de gouvernance.
\end{boxscenario}

\begin{boxmetier}{RSSI}
\textbf{Évaluer une dépendance critique à un fournisseur de sécurité}
\begin{enumerate}
    \item Recenser les logiciels qui disposent de droits élevés sur
          l'ensemble du parc et reçoivent des mises à jour automatiques.
    \item Pour chacun, vérifier si l'organisation peut retarder ou
          échelonner les mises à jour (groupe pilote, puis déploiement
          général).
    \item Évaluer le scénario « le fournisseur rend tout le parc
          inutilisable »~: délai de restauration, procédure manuelle,
          personnel nécessaire.
    \item Présenter ce risque au comité des risques, avec une
          recommandation~: l'accepter, le réduire ou le transférer.
\end{enumerate}
Ce qui compte n'est pas de supprimer la dépendance, souvent inévitable,
mais qu'elle soit \textbf{connue et décidée} par l'instance compétente.
\end{boxmetier}

% ============================================================================
\section{Lecture Croisée des Trois Affaires}
\label{sec:c02-croisee}
% ============================================================================

\subsection{Comparaison}

\begin{tableaucours}{L{2.9cm}L{3.1cm}L{3.1cm}L{3.1cm}}{%
    \tableauHeader{Question} &
    \tableauHeader{Knight Capital} &
    \tableauHeader{TSB} &
    \tableauHeader{CrowdStrike}
}
Nature du changement & Déploiement de code & Migration de plateforme &
Mise à jour de contenu \\
Durée de l'incident & 45 minutes & Plusieurs mois & Plusieurs jours \\
Contrôle défaillant & Déploiement sans procédure ni double contrôle &
Préparation de la bascule, supervision du fournisseur & Validation et
déploiement progressif \\
Rôle d'un tiers & Aucun & Fournisseur critique du groupe & L'éditeur
lui-même est le tiers des clients \\
Sanction & Amende de la SEC, absorption & Amendes FCA et PRA, sanction du
DSI & Pertes et contentieux des clients \\
\end{tableaucours}
\vspace{-0.8em}
\begin{center}{\small\itshape Tableau~2.2. Comparaison des trois affaires}\end{center}

\subsection{Quatre constantes}

Malgré leurs différences, les trois affaires partagent quatre traits.

\begin{enumerate}
    \item \textbf{Un changement est à l'origine de chaque incident.} Aucune
          de ces pannes ne vient d'une usure ou d'une attaque. Toutes
          viennent d'une modification volontaire du système. La gestion
          des changements (Chapitre~\ref{chap:itil-changements}) est le
          premier contrôle de la gouvernance des SI.
    \item \textbf{Le risque était connu ou connaissable.} Code mort jamais
          retiré, préparation insuffisante, type de mise à jour moins
          contrôlé~: aucun de ces risques n'était imprévisible.
    \item \textbf{La responsabilité était diluée.} Dans chaque cas, il est
          difficile de nommer la personne ou l'instance qui avait
          formellement décidé que le risque était acceptable.
    \item \textbf{Le retour arrière était difficile ou impossible.} Knight a
          aggravé la situation en tentant de revenir en arrière, TSB a mis
          des mois à se rétablir, et les clients de CrowdStrike ont dû
          réparer leurs postes un par un.
\end{enumerate}

\begin{boiteRemarque}{Et en Afrique~?}
Les trois affaires sont américaines ou européennes. Ce n'est pas que
l'Afrique centrale ignore les grandes pannes~: les interruptions de
services bancaires, de paiement mobile ou de télécommunications y sont
fréquentes et parfois longues. C'est que leurs causes sont rarement
rendues publiques. Peu de régulateurs publient leurs décisions motivées,
et peu d'entreprises publient leurs analyses de cause.

Cette opacité est elle-même une question de gouvernance. Une défaillance
dont on ne publie pas l'analyse ne profite à personne d'autre qu'à
l'organisation qui l'a subie, et souvent pas même à elle. Le
Chapitre~\ref{chap:droit} examinera les obligations de déclaration des
incidents prévues par le droit camerounais, et le
Chapitre~\ref{chap:recherche-mesure} la place de cette question dans la
recherche.
\end{boiteRemarque}

% ============================================================================
\section{Une Grille d'Erreurs de Gouvernance}
\label{sec:c02-grille}
% ============================================================================

Les trois affaires réelles et l'opération SANAGA permettent de construire
une grille de dix erreurs. Chacune s'accompagne d'un signal d'alerte
observable et de la question qu'un administrateur, un auditeur ou un DSI
devrait poser. La dernière colonne renvoie à l'objectif COBIT~2019 qui
traite l'erreur~; ces objectifs seront étudiés au
Chapitre~\ref{chap:cobit-objectifs}.

\begin{tableaucours}{L{3.2cm}L{3.4cm}L{3.8cm}L{1.6cm}}{%
    \tableauHeader{Erreur} &
    \tableauHeader{Signal d'alerte} &
    \tableauHeader{Question à poser} &
    \tableauHeader{COBIT}
}
1. Décision sans instance & Date ou budget fixés par une seule personne &
Quelle instance a décidé, et où est le procès-verbal~? & EDM01 \\
2. Risque non arbitré & Alerte écrite restée sans réponse & Qui a accepté
ce risque, par écrit~? & EDM03 \\
3. Changement non autorisé & Mise en production sans comité & Où est
l'autorisation, datée avant la mise en production~? & BAI06 \\
4. Pas de retour arrière & Plan de retour absent ou non testé & Quand le
retour arrière a-t-il été testé, et avec quel résultat~? & BAI07 \\
5. Recette incomplète & Bascule avec des scénarios non testés & Quel
taux de couverture, et qui a accepté le reste~? & BAI07 \\
6. Fournisseur non supervisé & Rapports de suivi rédigés par le
fournisseur & Comment vérifions-nous nous-mêmes ce qu'il affirme~? &
APO10 \\
7. Contrat sans engagement & Pas de niveau de service ni de pénalité &
Que se passe-t-il s'il ne tient pas ses engagements~? & APO09 \\
8. Alertes sans propriétaire & Messages automatiques non traités & Qui est
chargé de lire cette alerte, et dans quel délai~? & DSS02 \\
9. Continuité non testée & Plan de continuité jamais exercé & Quand a eu
lieu le dernier exercice, et qu'a-t-il révélé~? & DSS04 \\
10. Audit sans suite & Recommandations « en cours » depuis des mois &
Qui suit les recommandations, devant quelle instance~? & MEA02 \\
\end{tableaucours}
\vspace{-0.8em}
\begin{center}{\small\itshape Tableau~2.3. Grille des dix erreurs de
gouvernance des SI}\end{center}

\begin{boxprocedure}
\textbf{Le dossier d'analyse d'une défaillance.} Pour qu'une analyse de
défaillance soit opposable à la direction, au conseil ou au régulateur,
elle doit contenir~:
\begin{itemize}
    \item la liste des sources utilisées, classées selon le
          Tableau~2.1~;
    \item une chronologie où chaque entrée renvoie à une pièce~;
    \item des constatations séparées des analyses, chaque analyse portant
          un niveau de confiance (faible, moyen, élevé)~;
    \item la cause technique et la cause de gouvernance, distinguées~;
    \item pour chaque erreur de la grille relevée, la pièce qui l'établit.
\end{itemize}
\end{boxprocedure}

\begin{boxmetier}{Auditeur SI}
\textbf{Utiliser la grille en mission}
\begin{enumerate}
    \item Avant la mission, choisir les erreurs de la grille les plus
          probables au regard de l'activité (une banque qui migre~:
          erreurs 1 à 5~; une entreprise très externalisée~: erreurs 6
          et~7).
    \item Pour chaque erreur retenue, préparer la question et la pièce à
          demander.
    \item En entretien, poser la question, puis demander la pièce~:
          une réponse orale n'est pas une preuve.
    \item Consigner chaque écart sous la forme constatation, critère,
          cause, conséquence, recommandation
          (Chapitre~\ref{chap:audit}).
\end{enumerate}
\end{boxmetier}

\begin{boxmetier}{Responsable des services IT}
\textbf{Trois contrôles à mettre en place dès demain}
\begin{enumerate}
    \item Aucun changement en production sans enregistrement préalable, et
          aucun changement majeur sans autorisation d'un comité.
    \item Une vérification après chaque déploiement~: la version installée
          est la même partout.
    \item Un propriétaire nommé pour chaque alerte automatique, avec un
          délai de traitement.
\end{enumerate}
Ces trois contrôles auraient, à eux seuls, très probablement évité
l'affaire Knight Capital.
\end{boxmetier}

% ============================================================================
\separateur
% ============================================================================

\begin{boxresume}
\textbf{Ce qu'il faut retenir de ce chapitre~:}
\begin{itemize}
    \item \textbf{Les sources de première main d'abord}~: décisions des
          régulateurs, rapports d'enquête, analyses de cause. La presse
          oriente, elle n'établit pas.
    \item \textbf{Knight Capital (2012)}~: un déploiement sans procédure
          écrite ni double contrôle, un code mort réactivé, des alertes
          sans propriétaire~; 460~millions de dollars perdus en
          45~minutes.
    \item \textbf{TSB (2018)}~: une migration mal pilotée, un fournisseur
          critique mal supervisé~; des amendes de 48,65~millions de livres
          et la sanction personnelle de l'ancien DSI.
    \item \textbf{CrowdStrike (2024)}~: une mise à jour de contenu
          insuffisamment validée et déployée en une fois~; une dépendance
          critique que la plupart des clients n'avaient jamais arbitrée.
    \item \textbf{Quatre constantes}~: un changement à l'origine, un risque
          connaissable, une responsabilité diluée, un retour arrière
          difficile.
    \item \textbf{La grille des dix erreurs} relie chaque erreur à un
          signal, une question et un objectif COBIT~: elle servira dans
          tout le cours.
\end{itemize}
\end{boxresume}

\bigskip

\begin{boxexo}[title={\ding{45}\quad Exercice~2.1. Constatation ou analyse~? \niveaubase}]
Pour chacune des phrases suivantes, indiquez s'il s'agit d'une constatation
ou d'une analyse. Pour les analyses, proposez un niveau de confiance et la
pièce qui permettrait de le renforcer.
\begin{enumerate}[(a)]
    \item « Le nouveau code n'a été installé que sur sept des huit
          serveurs. »
    \item « Les équipes de Knight manquaient de formation. »
    \item « TSB a versé 32,7~millions de livres de réparation à ses
          clients. »
    \item « Le conseil de TSB aurait dû refuser la date de bascule. »
    \item « La mise à jour fournissait 21~champs au lieu des 20 attendus. »
\end{enumerate}
\end{boxexo}

\begin{boxexo}[title={\ding{45}\quad Exercice~2.2. La grille appliquée à KRIBILOG \niveaumoyen}]
Chez KRIBILOG SA, le logiciel de suivi des conteneurs s'arrête trois fois
en deux mois. Vous disposez des faits suivants~: l'infogérant NETPORT
SERVICES applique les correctifs le soir, sans prévenir~; son rapport
mensuel ne mentionne aucun incident~; le contrat ne prévoit aucun délai
d'intervention~; le dernier test de continuité date de 2021.
\begin{enumerate}[(a)]
    \item Identifiez les erreurs de la grille (Tableau~2.3) présentes dans
          ce cas.
    \item Pour chacune, rédigez la question à poser et la pièce à
          demander.
    \item Classez-les par ordre de priorité, en justifiant.
\end{enumerate}
\end{boxexo}

\begin{boxexo}[title={\ding{45}\quad Exercice~2.3. Simulation d'examen \niveauexpert}]
\textit{Durée conseillée~: 1~heure.} Choisissez une interruption de service
informatique survenue en Afrique et documentée publiquement (communiqué de
l'organisation, décision d'un régulateur, rapport d'enquête).
\begin{enumerate}[(a)]
    \item Classez les sources disponibles selon le Tableau~2.1, et
          indiquez ce qu'elles permettent d'établir.
    \item Appliquez les sept questions de la grille d'analyse
          (section~\ref{sec:c02-methode}).
    \item Rédigez une synthèse d'une page, en séparant constatations et
          analyses.
    \item Si les sources ne permettent pas d'établir la cause de
          gouvernance, dites-le, et indiquez les pièces qu'il faudrait
          obtenir. Cette conclusion est une réponse valable.
\end{enumerate}
\end{boxexo}

\bigskip

\begin{boxinfo}
\textbf{Transition.} Les défaillances montrent ce que coûte une
gouvernance absente. Le Chapitre~\ref{chap:alignement} montre ce qu'une
gouvernance présente doit rechercher~: l'alignement du SI sur la
stratégie, la création de valeur, et l'arbitrage des décisions
d'investissement.
\end{boxinfo}

        % ============================================================================
% Fichier  : partie1/P01_C03_alignement_valeur_decision.tex
% Titre    : Théories et Modèles : Alignement, Valeur et Décision
% Auteur   : MINKA MI NGUIDJOÏ Thierry Emmanuel
% Version  : 1.0 -- octobre 2026
% ============================================================================
% Chapitre associé au laboratoire GCI-LAB-01 (arbitrage du portefeuille de
% l'ANPP, opération MEFOU).
% ============================================================================

\chapter{Théories et Modèles : Alignement, Valeur et Décision}
\label{chap:alignement}
\chaptermark{Alignement, valeur et décision}

\epigraph{%
    « On voit l'ère de l'ordinateur partout, sauf dans les statistiques
    de productivité. »%
}{Robert \textsc{Solow}, prix Nobel d'économie, 1987 (notre traduction)}

\niveauChapitre{Fondamental}{%
    Maîtrise du Chapitre~\ref{chap:gouv-entreprise} (définition de travail
    de la gouvernance des SI et ses trois finalités).%
}

\begin{boxobjectifs}
\begin{itemize}
    \item Appliquer le modèle d'alignement stratégique de Henderson et
          Venkatraman.
    \item Classer les droits de décision selon les archétypes de Weill et
          Ross.
    \item Expliquer les trois finalités~: réalisation des bénéfices,
          optimisation des risques, optimisation des ressources.
    \item Évaluer un investissement informatique (coût total de possession,
          valeur actuelle nette, score multicritère).
    \item Arbitrer un portefeuille de projets.
\end{itemize}
\end{boxobjectifs}

\bigskip

% ============================================================================
\section*{Question directrice : 14 milliards pour quoi faire~?}
% ============================================================================

À l'ANPP (opération MEFOU, Chapitre~\ref{chap:gouv-entreprise}), le
schéma directeur adopté en 2023 prévoyait neuf projets pour 14~milliards
FCFA. À mi-parcours, trois projets ont consommé 71~\% du budget sans être
achevés. Le nouveau président du conseil d'administration pose trois
questions simples~: « Ces projets servaient-ils la mission de l'agence~?
Qui les a choisis, et selon quels critères~? Combien rapportent-ils, ou
combien de risques évitent-ils~? »

Ces trois questions sont celles de ce chapitre. La première relève de
l'\textbf{alignement stratégique} (section~\ref{sec:c03-alignement}). La
deuxième relève des \textbf{droits de décision}
(section~\ref{sec:c03-weill-ross}). La troisième relève de la
\textbf{valeur}, que le chapitre définit (section~\ref{sec:c03-valeur}),
mesure (section~\ref{sec:c03-evaluation}) et arbitre à l'échelle d'un
portefeuille (section~\ref{sec:c03-portefeuille}). Le laboratoire GCI-LAB-01
de NDAP SUKULU applique la démarche au portefeuille complet de l'ANPP.

\separateur

% ============================================================================
\section{Le Paradoxe de la Productivité}
\label{sec:c03-paradoxe}
% ============================================================================

La phrase de Solow, en épigraphe, a ouvert un long débat. Dans les années
1980, les entreprises américaines investissent massivement dans
l'informatique, sans que la productivité mesurée progresse en proportion.
Les travaux ultérieurs ont largement nuancé ce constat~: les gains
existent, mais ils sont inégaux, différés, et dépendent moins de la
technologie elle-même que de la manière dont l'organisation la met en
œuvre.

Pour la gouvernance, la leçon est décisive~: \textbf{la valeur n'est pas
dans la technologie, elle est dans son usage}. Deux organisations qui
achètent le même logiciel n'en tirent pas la même valeur. Ce qui les
distingue, c'est la manière dont elles choisissent leurs projets,
transforment leurs processus et suivent les bénéfices. C'est précisément
le champ de la gouvernance des SI.

\begin{boxinfo}
\textbf{Les quatre questions de Thorp.} John Thorp a résumé la
gouvernance de la valeur en quatre questions~\cite{thorp1998paradox},
reprises plus tard par le référentiel Val~IT de l'ISACA~:
\emph{fait-on les bonnes choses~?} (la stratégie), \emph{les fait-on de la
bonne manière~?} (l'architecture), \emph{les mène-t-on à bien~?} (la
livraison), \emph{en tire-t-on les bénéfices~?} (la valeur). Les deux
premières relèvent de la gouvernance, les deux dernières du management,
sous le contrôle de la gouvernance.
\end{boxinfo}

% ============================================================================
\section{Le Modèle d'Alignement Stratégique}
\label{sec:c03-alignement}
% ============================================================================

\subsection{Quatre domaines, deux dimensions}

Le modèle d'alignement stratégique (\textit{Strategic Alignment Model},
\sigle{SAM}) de Henderson et Venkatraman~\cite{henderson1993alignment}
reste le cadre de référence de l'alignement. Il distingue quatre domaines,
organisés selon deux dimensions.

\begin{itemize}
    \item La dimension de l'\textbf{ajustement stratégique}
          (\textit{strategic fit}) relie le niveau externe (la stratégie~:
          sur quel marché, avec quelles compétences distinctives, avec
          quelle gouvernance) et le niveau interne (l'infrastructure et les
          processus qui la mettent en œuvre).
    \item La dimension de l'\textbf{intégration fonctionnelle}
          (\textit{functional integration}) relie le domaine de l'activité
          et le domaine des technologies de l'information.
\end{itemize}

\begin{center}
\begin{tikzpicture}[
    dom/.style = {
        draw=CouleurBordDef, fill=CouleurDef, rounded corners=4pt,
        text width=4.2cm, minimum height=1.3cm, align=center,
        font=\footnotesize\bfseries\color{CouleurPrimaire}
    },
    lien/.style = {{Stealth[length=5pt]}-{Stealth[length=5pt]}, thick,
        color=CouleurSecondaire}
]
\node[dom] (bs) at (0,2.6)   {Stratégie de l'organisation};
\node[dom] (ts) at (7.6,2.6) {Stratégie des technologies\\ de l'information};
\node[dom] (oi) at (0,0)     {Infrastructure et processus\\ de l'organisation};
\node[dom] (ii) at (7.6,0)   {Infrastructure et processus\\ des SI};
\draw[lien] (bs) -- node[above, font=\scriptsize\itshape, align=center]{intégration\\
    fonctionnelle} (ts);
\draw[lien] (oi) -- node[below, font=\scriptsize\itshape, align=center]{intégration\\
    fonctionnelle} (ii);
\draw[lien] (bs) -- node[left, font=\scriptsize\itshape, align=center]{ajustement\\
    stratégique} (oi);
\draw[lien] (ts) -- node[right, font=\scriptsize\itshape, align=center]{ajustement\\
    stratégique} (ii);
\node[font=\scriptsize\bfseries\color{CouleurAccent}, above=0.1cm of bs]
    {Activité};
\node[font=\scriptsize\bfseries\color{CouleurAccent}, above=0.1cm of ts]
    {Technologies};
\end{tikzpicture}
\end{center}
\vspace{-0.4em}
\begin{center}{\small\itshape Figure~3.1. Le modèle d'alignement
stratégique (d'après Henderson et Venkatraman, 1993)}\end{center}

\begin{boxdef}
\termecle{Alignement stratégique}\index{alignement stratégique}
(\textit{strategic alignment})~: degré de cohérence entre la stratégie de
l'organisation, la stratégie des technologies de l'information et les
infrastructures et processus qui les mettent en œuvre. Pour Henderson et
Venkatraman, l'alignement n'est pas un état atteint une fois pour toutes,
mais un processus continu d'adaptation.
\end{boxdef}

\subsection{Quatre perspectives d'alignement}

Le modèle décrit quatre manières de parcourir les domaines, selon le point
de départ et le point d'arrivée. Chacune correspond à un rôle différent
de la direction générale et du DSI.

\begin{tableaucours}{L{3.0cm}L{4.4cm}L{4.8cm}}{%
    \tableauHeader{Perspective} &
    \tableauHeader{Parcours} &
    \tableauHeader{Exemple camerounais}
}
Exécution de la stratégie & Stratégie de l'organisation $\rightarrow$
infrastructure de l'organisation $\rightarrow$ infrastructure des SI &
L'ANPP doit payer tous les fournisseurs de l'État sous 30~jours~: elle
réorganise ses circuits, puis outille la chaîne de paiement \\
Transformation technologique & Stratégie de l'organisation $\rightarrow$
stratégie des TI $\rightarrow$ infrastructure des SI & La BANQUE SANAGA
veut conquérir les jeunes urbains~: elle choisit une stratégie mobile,
puis la plateforme qui la porte \\
Potentiel concurrentiel & Stratégie des TI $\rightarrow$ stratégie de
l'organisation $\rightarrow$ infrastructure de l'organisation & NOUN PAY
découvre que la notation par IA permet le microcrédit instantané~: elle
redéfinit son offre \\
Niveau de service & Stratégie des TI $\rightarrow$ infrastructure des SI
$\rightarrow$ infrastructure de l'organisation & KRIBILOG fiabilise son
suivi des conteneurs, ce qui permet aux équipes du port de réorganiser
leurs rotations \\
\end{tableaucours}
\vspace{-0.8em}
\begin{center}{\small\itshape Tableau~3.1. Les quatre perspectives
d'alignement}\end{center}

\begin{boxalert}
\textbf{Violation typique.} Un projet est lancé parce que « la
technologie existe » ou parce qu'un autre organisme l'a adoptée, sans que
personne n'ait identifié l'objectif de l'organisation qu'il sert. Le
projet est livré, mais aucun processus ne change, et aucun bénéfice n'est
mesuré.

\smallskip
\textbf{Règle~:} tout projet doit pouvoir être rattaché, par écrit, à un
objectif de l'organisation et à l'une des quatre perspectives. Un projet
qu'on ne peut rattacher à rien n'est pas aligné~: il doit être justifié
autrement, ou abandonné.
\end{boxalert}

% ============================================================================
\section{Qui Décide de Quoi : les Archétypes de Weill et Ross}
\label{sec:c03-weill-ross}
% ============================================================================

\subsection{Cinq décisions}

Weill et Ross~\cite{weill2004itgovernance} identifient cinq grandes
catégories de décisions relatives aux technologies de l'information.
Elles sont liées~: les principes encadrent les quatre autres.

\begin{enumerate}
    \item \textbf{Les principes}~: le rôle que les TI doivent jouer dans
          l'organisation (réduire les coûts, soutenir la croissance,
          transformer le métier).
    \item \textbf{L'architecture}~: l'organisation logique des données, des
          applications et de l'infrastructure, et les standards retenus.
    \item \textbf{L'infrastructure}~: les services techniques partagés
          (réseaux, centres de données, sécurité) et leur niveau.
    \item \textbf{Les besoins applicatifs}~: les applications à acquérir ou
          à développer pour les métiers.
    \item \textbf{L'investissement et la priorisation}~: combien investir,
          dans quoi, et dans quel ordre.
\end{enumerate}

\subsection{Six archétypes}

Pour chaque décision, Weill et Ross se demandent \emph{qui} détient le
droit de décider, et \emph{qui} est consulté avant. Ils distinguent six
archétypes de gouvernance.

\begin{tableaucours}{L{2.8cm}L{4.6cm}L{4.8cm}}{%
    \tableauHeader{Archétype} &
    \tableauHeader{Qui décide} &
    \tableauHeader{Risque typique}
}
Monarchie métier & La direction générale ou un groupe de dirigeants des
métiers & Décisions déconnectées des contraintes techniques \\
Monarchie informatique & Les responsables informatiques & Décisions
techniquement justes, mais sans adhésion des métiers \\
Féodalité & Chaque direction métier pour son propre périmètre & Doublons,
incompatibilités, coûts non mutualisés \\
Fédéralisme & La direction générale et les métiers ensemble, avec ou sans
l'informatique & Lenteur, compromis difficiles \\
Duopole & L'informatique et un groupe de métiers & Dépend de la qualité
de la relation bilatérale \\
Anarchie & Chaque utilisateur ou petit groupe & Aucune cohérence, risques
non maîtrisés \\
\end{tableaucours}
\vspace{-0.8em}
\begin{center}{\small\itshape Tableau~3.2. Les six archétypes de Weill et
Ross}\end{center}

\subsection{La matrice des droits de décision}

L'outil le plus utile est la \termecle{matrice des droits de
décision}\index{matrice des droits de décision}~: en lignes les
archétypes, en colonnes les cinq décisions, et dans chaque case qui
apporte un avis (\emph{consultation}) et qui tranche (\emph{décision}).
Elle décrit la gouvernance telle qu'elle fonctionne réellement, et non
telle qu'elle est écrite.

\begin{boxscenario}{Opération MEFOU. Les droits de décision réels}
L'entretien avec les responsables de l'ANPP fait apparaître la situation
suivante.

\medskip
\begin{tableaucours}{L{3.6cm}L{4.0cm}L{4.4cm}}{%
    \tableauHeader{Décision} &
    \tableauHeader{Qui décide réellement} &
    \tableauHeader{Archétype}
}
Principes & Personne (aucun document) & Absence de décision \\
Architecture & Chaque prestataire, projet par projet & Anarchie \\
Infrastructure & Cheffe de la division informatique & Monarchie
informatique \\
Besoins applicatifs & Chaque direction, avec son prestataire & Féodalité \\
Investissement & Directeur général, à la demande de la tutelle &
Monarchie métier \\
\end{tableaucours}

\medskip
\textbf{Analyse} (confiance élevée). La décision la plus lourde,
l'investissement, est prise par une seule personne, sans consultation de
la division informatique ni critères écrits. Les besoins applicatifs sont
définis direction par direction, ce qui explique que trois projets aient
absorbé l'essentiel du budget sans cohérence d'ensemble. L'absence de
principes rend tout arbitrage impossible~: il n'existe aucun critère pour
dire non.
\end{boxscenario}

\begin{boiteRemarque}{Existe-t-il un bon archétype~?}
Il n'existe pas d'archétype bon en soi. Weill et Ross observent que les
organisations les plus performantes combinent des archétypes différents
selon les décisions~: par exemple une monarchie informatique pour
l'architecture et l'infrastructure, où l'expertise prime, et un
fédéralisme ou un duopole pour les besoins applicatifs et
l'investissement, où l'adhésion des métiers est indispensable. La seule
configuration à proscrire est celle qui n'a pas été choisie~: la
gouvernance de fait, qui s'installe faute de décision.
\end{boiteRemarque}

% ============================================================================
\section{Les Trois Finalités de la Gouvernance}
\label{sec:c03-valeur}
% ============================================================================

\subsection{Créer de la valeur}

COBIT~2019 définit l'objectif de la gouvernance comme la \textbf{création
de valeur}, qui résulte de l'équilibre entre trois
finalités~\cite{isaca2018cobitintro}. La définition de travail du
Chapitre~\ref{chap:gouv-entreprise} les reprend.

\begin{boxdef}
\termecle{Réalisation des bénéfices}\index{réalisation des bénéfices}
(\textit{benefits realization})~: créer de la valeur pour l'organisation
grâce aux technologies de l'information, en entretenant et en accroissant
la valeur des investissements existants, et en éliminant les initiatives
et les actifs qui n'en créent pas assez.

\medskip
\termecle{Optimisation des risques}\index{optimisation des risques}
(\textit{risk optimization})~: traiter les risques liés à l'usage, à la
propriété, à l'exploitation et à l'adoption des technologies de
l'information, en les maintenant dans les limites de l'appétence au
risque fixée par l'organe de gouvernance.

\medskip
\termecle{Optimisation des ressources}\index{optimisation des
ressources} (\textit{resource optimization})~: s'assurer que les
compétences et les ressources nécessaires (personnes, processus,
technologies, données) sont disponibles, suffisantes et utilisées de
manière efficace.
\end{boxdef}

Le mot clé est \textbf{équilibre}. Maximiser les bénéfices sans se soucier
des risques conduit à l'opération SANAGA. Minimiser les risques sans se
soucier des bénéfices conduit à l'immobilisme. Économiser les ressources
sans se soucier des deux autres conduit à des systèmes vieillissants et
fragiles.

\subsection{Trois objectifs de gouvernance}

À chaque finalité correspond un objectif du domaine de gouvernance de
COBIT~2019. Ils seront étudiés au Chapitre~\ref{chap:cobit-objectifs}.

\begin{tableaucours}{L{3.4cm}L{4.4cm}L{4.4cm}}{%
    \tableauHeader{Finalité} &
    \tableauHeader{Objectif COBIT 2019} &
    \tableauHeader{Question au conseil}
}
Réalisation des bénéfices & EDM02 · Assurer la réalisation des bénéfices
& Les investissements produisent-ils les bénéfices annoncés~? \\
Optimisation des risques & EDM03 · Assurer l'optimisation des risques &
Les risques sont-ils dans les limites que nous avons fixées~? \\
Optimisation des ressources & EDM04 · Assurer l'optimisation des
ressources & Avons-nous les compétences et les moyens de notre
stratégie~? \\
\end{tableaucours}
\vspace{-0.8em}
\begin{center}{\small\itshape Tableau~3.3. Finalités et objectifs de
gouvernance (intitulés traduits par nos soins)}\end{center}

\begin{boxobjet}{EDM02 · Assurer la réalisation des bénéfices}
\textbf{Finalité~:} obtenir la valeur optimale des initiatives, des
services et des actifs informatiques, et une gestion des coûts efficace.

\textbf{Pratiques~:} évaluer l'optimisation de la valeur~; diriger
l'optimisation de la valeur (critères d'investissement, rôles)~;
surveiller l'optimisation de la valeur (indicateurs, écarts).

\textbf{Sorties~:} principes et critères d'investissement, retour sur
le portefeuille, actions correctives.

\textbf{RACI (extrait)~:} \emph{A} au conseil d'administration, \emph{R}
au directeur général et au comité stratégique des SI, \emph{C} au
directeur financier et au DSI.

\textbf{Indicateurs~:} part des investissements dont les bénéfices sont
suivis après la mise en service~; écart entre bénéfices annoncés et
bénéfices constatés.
\end{boxobjet}

% ============================================================================
\section{Évaluer un Investissement Informatique}
\label{sec:c03-evaluation}
% ============================================================================

\subsection{Le coût total de possession}

Le prix d'achat d'un logiciel n'est qu'une partie de son coût. Le
\termecle{coût total de possession}\index{coût total de possession}
(\textit{total cost of ownership}, \sigle{TCO}) additionne tous les coûts
supportés sur la durée de vie de la solution~: acquisition, intégration,
matériel, formation, exploitation, support et sortie.

\begin{tableaucours}{L{8.6cm}r}{%
    \tableauHeader{Poste (sur 5 ans)} &
    \tableauHeader{Millions FCFA}
}
Licences (prix affiché par l'éditeur) & 300 \\
Intégration et reprise des données & 250 \\
Matériel et hébergement & 150 \\
Formation des utilisateurs & 50 \\
Exploitation interne (90 par an) & 450 \\
Support et maintenance de l'éditeur (60 par an) & 300 \\
\textbf{Coût total de possession} & \textbf{1~500} \\
\end{tableaucours}
\vspace{-0.8em}
\begin{center}{\small\itshape Tableau~3.4. Coût total de possession d'un
logiciel de gestion des paiements (exemple)}\end{center}

Le coût total est cinq fois le prix affiché. C'est un ordre de grandeur
fréquent, et l'une des raisons pour lesquelles les budgets des schémas
directeurs dérapent~: ils ont été construits sur les prix d'acquisition.

\subsection{La valeur actuelle nette}

Un franc reçu dans quatre ans vaut moins qu'un franc dépensé aujourd'hui.
La \termecle{valeur actuelle nette}\index{valeur actuelle nette}
(\sigle{VAN}, \textit{net present value}) ramène tous les flux d'un projet
à leur valeur d'aujourd'hui, au moyen d'un taux d'actualisation $r$~:
\[
    \mathrm{VAN} = -I_0 + \sum_{t=1}^{n} \frac{F_t}{(1+r)^t}
\]
où $I_0$ est l'investissement initial et $F_t$ le flux net (bénéfices
moins coûts d'exploitation) de l'année~$t$. Un projet dont la VAN est
positive crée de la valeur au taux retenu.

Lorsque le flux annuel $F$ est constant, la somme se calcule avec le
facteur d'annuité~:
\[
    \sum_{t=1}^{n} \frac{1}{(1+r)^t} = \frac{1-(1+r)^{-n}}{r},
    \qquad \text{soit } 3{,}7908 \text{ pour } r=10~\% \text{ et } n=5.
\]

\begin{boxscenario}{Opération MEFOU. Trois projets à évaluer}
L'ANPP retient un taux d'actualisation de 10~\% et un horizon de cinq ans.
Trois projets sont en concurrence.

\medskip
\begin{tableaucours}{L{4.4cm}rrr}{%
    \tableauHeader{Projet} &
    \tableauHeader{$I_0$} &
    \tableauHeader{$F$ annuel} &
    \tableauHeader{VAN}
}
A. Paiement dématérialisé des fournisseurs & 1~200 & 450 & $+506$ \\
B. Refonte du portail des ressources humaines & 800 & 180 & $-118$ \\
C. Site de secours (continuité) & 600 & 120 & $-145$ \\
\end{tableaucours}
{\small Montants en millions FCFA.}

\medskip
\textbf{Calcul du projet A.} $450 \times 3{,}7908 = 1~705{,}9$~; VAN $=
1~705{,}9 - 1~200 = 505{,}9$, soit environ $+506$~millions FCFA. Le délai
de récupération est de $1~200 / 450 \approx 2{,}7$~ans, et le taux de
rentabilité interne, pour lequel la VAN s'annule, est d'environ 25~\%.

\medskip
\textbf{Le cas du projet C.} Le site de secours ne produit aucune recette.
Son flux annuel est une \emph{perte évitée}~: l'ANPP estime à 5~\% par an
la probabilité d'un sinistre majeur sur son centre de données, pour une
perte de 2~400~millions FCFA (pénalités de retard de paiement, reprise
manuelle), soit une perte annuelle attendue de $0{,}05 \times 2~400 =
120$~millions FCFA.
\end{boxscenario}

\subsection{Les limites de l'évaluation financière}

Si l'ANPP ne retenait que la VAN, elle financerait le seul projet~A. Ce
serait une erreur de gouvernance, pour trois raisons.
\begin{itemize}
    \item \textbf{Les risques se chiffrent mal.} La perte attendue du
          projet~C repose sur une probabilité très incertaine. Une
          variation de quelques points change le signe de la VAN.
    \item \textbf{Certains projets sont obligatoires.} Une obligation
          légale ou réglementaire ne se discute pas en VAN~: elle se
          respecte, au moindre coût.
    \item \textbf{L'alignement ne se chiffre pas.} Un projet peut avoir
          une VAN modeste et être indispensable à la mission de
          l'organisation.
\end{itemize}

\subsection{Le score multicritère}

Pour dépasser ces limites, l'organe de gouvernance fixe à l'avance des
critères pondérés, puis note chaque projet. Les poids traduisent ses
priorités~: les fixer est une décision de gouvernance, qui ne doit pas
être laissée à ceux qui proposent les projets.

\begin{tableaucours}{L{3.8cm}crrr}{%
    \tableauHeader{Critère} &
    \tableauHeader{Poids} &
    \tableauHeader{A} &
    \tableauHeader{B} &
    \tableauHeader{C}
}
Alignement sur la mission & 30~\% & 5 & 3 & 3 \\
Valeur financière (VAN) & 25~\% & 5 & 2 & 1 \\
Réduction des risques & 20~\% & 2 & 1 & 5 \\
Obligation réglementaire & 15~\% & 3 & 1 & 4 \\
Faisabilité & 10~\% & 4 & 5 & 3 \\
\textbf{Score pondéré} & & \textbf{4,00} & \textbf{2,25} & \textbf{3,05} \\
\end{tableaucours}
\vspace{-0.8em}
\begin{center}{\small\itshape Tableau~3.5. Score multicritère des trois
projets de l'ANPP (notes de 1 à 5)}\end{center}

Le calcul du projet~A est le suivant~: $0{,}30 \times 5 + 0{,}25 \times 5 +
0{,}20 \times 2 + 0{,}15 \times 3 + 0{,}10 \times 4 = 4{,}00$. Avec une
enveloppe de 1~900~millions FCFA, l'ANPP finance A et C (1~800~millions)
et écarte B. Le site de secours, rejeté par la VAN seule, est retenu au
titre des risques et de l'obligation de continuité du service public.

\begin{boxprocedure}
\textbf{Une décision d'investissement auditable.} Pour qu'un auditeur, ou
la Chambre des comptes, puisse établir que le portefeuille a été arbitré,
le dossier doit contenir~:
\begin{itemize}
    \item les critères et leurs poids, approuvés par l'organe de
          gouvernance \emph{avant} l'examen des projets~;
    \item pour chaque projet, un dossier d'investissement (objectif servi,
          coût total de possession, bénéfices attendus et leur
          responsable, risques)~;
    \item la grille de notation, signée par ceux qui l'ont établie~;
    \item le procès-verbal de la décision, y compris pour les projets
          écartés~;
    \item le plan de suivi des bénéfices après la mise en service.
\end{itemize}
\end{boxprocedure}

% ============================================================================
\section{Arbitrer un Portefeuille}
\label{sec:c03-portefeuille}
% ============================================================================

\subsection{Projet, programme, portefeuille}

\begin{boxdef}
\termecle{Portefeuille}\index{portefeuille} (\textit{portfolio})~:
ensemble des projets, des programmes et des services informatiques d'une
organisation, géré comme un tout pour maximiser la valeur obtenue sous
contrainte de ressources et de risque. L'unité de décision de la
gouvernance n'est pas le projet isolé, mais le portefeuille.
\end{boxdef}

Raisonner en portefeuille change la question posée. Il ne s'agit plus de
savoir si un projet est bon, mais s'il est \emph{meilleur que les autres
usages possibles des mêmes ressources}. C'est ce qui a manqué à l'ANPP~:
chaque projet a été approuvé seul, sans jamais être comparé aux autres.

\subsection{Équilibrer le portefeuille}

Weill et Broadbent proposent de classer les investissements informatiques
selon leur finalité~\cite{weill1998infrastructure}. Un portefeuille
équilibré combine les quatre catégories, dans des proportions qui
reflètent la stratégie.

\begin{tableaucours}{L{2.8cm}L{4.6cm}L{4.8cm}}{%
    \tableauHeader{Catégorie} &
    \tableauHeader{Finalité} &
    \tableauHeader{Exemple à l'ANPP}
}
Infrastructure & Fournir les services partagés sur lesquels tout repose &
Réseau sécurisé entre l'agence et les trésoreries régionales \\
Transactionnel & Réduire les coûts et automatiser les opérations
répétitives & Paiement dématérialisé des fournisseurs (projet A) \\
Informationnel & Mieux piloter, mieux contrôler & Tableau de bord des
délais de paiement \\
Stratégique & Changer la position ou la mission de l'organisation &
Paiement des usagers par téléphone mobile \\
\end{tableaucours}
\vspace{-0.8em}
\begin{center}{\small\itshape Tableau~3.6. Les quatre catégories
d'investissement (d'après Weill et Broadbent)}\end{center}

\subsection{Le cadre de la gestion publique au Cameroun}

Pour une administration ou un établissement public, l'arbitrage du
portefeuille s'inscrit dans le cadre budgétaire de l'État.

\begin{boxjuris}
\textbf{\loicam{n°~2018/012} du 11~juillet~2018 portant régime financier
de l'État et des autres entités publiques, article~69.}

\medskip
Selon la présentation qu'en fait le ministère des Finances~\cite{minfi2023regimefinancier},
la loi organise le budget par programmes et institue un
\emph{responsable de programme}, désigné par le ministre. Ce responsable
formule le programme dans le \emph{Projet de performance annuel}
(\sigle{PPA}), affecte les ressources en son sein, organise le dialogue
de gestion et élabore le \emph{Rapport annuel de performance}
(\sigle{RAP}). Il veille au respect des dispositifs de contrôle interne et
de contrôle de gestion.

\medskip
\textbf{Portée.} La gestion publique camerounaise est désormais pilotée
par la performance~: objectifs, indicateurs, compte rendu annuel. Un
projet informatique public doit donc pouvoir être rattaché à un programme
et à ses indicateurs de performance. C'est l'équivalent public de
l'alignement stratégique, et un argument fort pour imposer, à l'ANPP
comme ailleurs, des dossiers d'investissement et un suivi des bénéfices.
\end{boxjuris}
% A VERIFIER au Journal officiel : libellé exact de l'article 69 de la loi
% n° 2018/012 (source utilisée : publication du MINFI, collection
% « Les essentiels de la réforme »).

\begin{boxmetier}{DSI}
\textbf{Présenter un portefeuille au comité stratégique}
\begin{enumerate}
    \item Une page par projet~: objectif servi, perspective d'alignement,
          coût total de possession, VAN ou perte évitée, risques, responsable
          des bénéfices.
    \item Une vue d'ensemble~: répartition du budget entre les quatre
          catégories (Tableau~3.6) et comparaison avec la stratégie.
    \item Le résultat de la grille multicritère, avec les poids approuvés
          par le comité.
    \item Une recommandation explicite~: financer, différer, arrêter.
\end{enumerate}
Le dernier point est le plus difficile, et le plus important~: proposer
d'arrêter un projet en cours est souvent la meilleure décision de
gouvernance qu'un DSI puisse provoquer.
\end{boxmetier}

\begin{boxmetier}{Auditeur SI}
\textbf{Vérifier qu'un portefeuille est réellement arbitré}
\begin{enumerate}
    \item Demander les critères d'investissement, et vérifier leur date
          d'approbation au regard des dates de décision des projets.
    \item Prendre un échantillon de projets approuvés \emph{et} de projets
          écartés~: une organisation qui n'a jamais écarté un projet
          n'arbitre pas.
    \item Pour les projets en service, demander le suivi des bénéfices
          annoncés dans le dossier d'investissement.
\end{enumerate}
\end{boxmetier}

% ============================================================================
\separateur
% ============================================================================

\begin{boxresume}
\textbf{Ce qu'il faut retenir de ce chapitre~:}
\begin{itemize}
    \item \textbf{La valeur est dans l'usage}~: le paradoxe de la
          productivité montre que la technologie seule ne crée pas de
          valeur~; la manière de la choisir et de l'exploiter, si.
    \item \textbf{L'alignement stratégique} relie quatre domaines (stratégie
          et infrastructure, de l'activité et des TI). Il se parcourt selon
          quatre perspectives, et c'est un processus continu.
    \item \textbf{Les droits de décision} portent sur cinq décisions
          (principes, architecture, infrastructure, besoins applicatifs,
          investissement)~; six archétypes décrivent qui décide. Le pire
          archétype est celui qu'on n'a pas choisi.
    \item \textbf{Trois finalités en équilibre}~: réalisation des
          bénéfices, optimisation des risques, optimisation des ressources
          (EDM02, EDM03, EDM04).
    \item \textbf{Évaluer}~: coût total de possession, VAN, délai de
          récupération~; mais aussi perte évitée pour les projets de
          risque, et score multicritère aux poids fixés par la gouvernance.
    \item \textbf{Arbitrer un portefeuille}~: comparer les projets entre
          eux, équilibrer les catégories, et savoir arrêter un projet.
\end{itemize}
\end{boxresume}

\bigskip

\begin{boxexo}[title={\ding{45}\quad Exercice~3.1. Calculer une VAN \niveaubase}]
KRIBILOG SA envisage un logiciel de planification des quais~:
investissement de 450~millions FCFA, gain net de 140~millions FCFA par an
pendant cinq ans, taux d'actualisation de 12~\%.
\begin{enumerate}[(a)]
    \item Calculez le facteur d'annuité pour $r=12~\%$ et $n=5$.
    \item Calculez la VAN du projet. Est-il créateur de valeur~?
    \item Calculez son délai de récupération simple.
\end{enumerate}
\end{boxexo}

\begin{boxexo}[title={\ding{45}\quad Exercice~3.2. La matrice de la BANQUE SANAGA \niveaumoyen}]
À partir des faits du prologue (Chapitre~\ref{chap:prologue}), construisez
la matrice des droits de décision de la BANQUE SANAGA pour le projet
HORIZON.
\begin{enumerate}[(a)]
    \item Pour chacune des cinq décisions de Weill et Ross, indiquez qui
          a réellement décidé et l'archétype correspondant.
    \item Proposez la matrice cible, en justifiant l'archétype retenu pour
          chaque décision.
    \item Rédigez la constatation et l'analyse correspondantes, avec leur
          niveau de confiance.
\end{enumerate}
\end{boxexo}

\begin{boxexo}[title={\ding{45}\quad Exercice~3.3. Arbitrer sous contrainte \niveauexpert}]
Le comité stratégique de l'ANPP révise ses poids~: l'obligation
réglementaire passe à 25~\% et la valeur financière à 15~\%, les autres
poids restant inchangés. Un quatrième projet~D (tableau de bord des
délais de paiement, 300~millions FCFA) obtient les notes 4, 2, 2, 5, 5.
\begin{enumerate}[(a)]
    \item Recalculez les scores des projets A, B, C et D.
    \item Avec une enveloppe de 1~900~millions FCFA, quel portefeuille
          recommandez-vous~? Justifiez.
    \item Le directeur général conteste le changement des poids, qui
          « favorise » le projet~D. Qui doit trancher, et comment
          l'auditeur saura-t-il que les poids n'ont pas été fixés après
          coup pour justifier un choix déjà fait~?
\end{enumerate}
\end{boxexo}

\bigskip

\begin{boxinfo}
\textbf{Transition.} Ce chapitre a posé les modèles qui guident les
décisions de gouvernance~: alignement, droits de décision, valeur. Pour
les mettre en œuvre, les organisations s'appuient sur des référentiels.
Le Chapitre~\ref{chap:referentiels} dresse le panorama de ces
référentiels (COBIT, ITIL, ISO, COSO, TOGAF\ldots) et montre comment ils
s'articulent.
\end{boxinfo}

        % ============================================================================
% Fichier  : partie1/P01_C04_panorama_referentiels.tex
% Titre    : Panorama et Articulation des Référentiels
% Auteur   : MINKA MI NGUIDJOÏ Thierry Emmanuel
% Version  : 1.0 -- octobre 2026
% ============================================================================

\chapter{Panorama et Articulation des Référentiels}
\label{chap:referentiels}
\chaptermark{Panorama des référentiels}

\epigraph{%
    « Un référentiel ne pense pas à votre place.\\
    Il vous évite seulement d'oublier ce que d'autres ont appris à leurs
    dépens. »%
}{Adage des praticiens de la gouvernance}

\niveauChapitre{Fondamental}{%
    Maîtrise du Chapitre~\ref{chap:alignement} (alignement, droits de
    décision, trois finalités de la gouvernance).%
}

\begin{boxobjectifs}
\begin{itemize}
    \item Situer COBIT, ITIL, ISO/IEC~38500, ISO/IEC~20000, ISO/IEC~27001,
          COSO, CMMI, TOGAF et PMBOK.
    \item Distinguer cadre, norme certifiable et bonnes pratiques.
    \item Associer chaque référentiel au besoin qu'il couvre.
    \item Construire une carte d'articulation des référentiels pour une
          organisation donnée.
\end{itemize}
\end{boxobjectifs}

\bigskip

% ============================================================================
\section*{Question directrice : faut-il tout appliquer~?}
% ============================================================================

Après la panne, le conseil d'administration de la BANQUE SANAGA reçoit en
trois semaines quatre propositions de cabinets de conseil. La première
propose « une mise en conformité COBIT », la deuxième « un déploiement
ITIL », la troisième « une certification ISO~27001 », la quatrième « un
programme de maturité CMMI ». Chacune promet de régler le problème. Le
président du conseil demande à la nouvelle responsable GRC~: « Faut-il
tout faire~? Dans quel ordre~? Et ces référentiels se contredisent-ils~? »

Ce chapitre répond à ces questions. Il commence par distinguer les natures
de référentiels (section~\ref{sec:c04-natures}), présente ensuite les
principaux d'entre eux, regroupés par besoin
(sections~\ref{sec:c04-gouvernance} à~\ref{sec:c04-projets}), puis montre
comment les articuler en une carte cohérente
(section~\ref{sec:c04-carte}). La réponse courte est~: \textbf{non, il ne
faut pas tout appliquer}, et les référentiels sérieux ne se contredisent
pas~: ils se répartissent le travail.

\separateur

% ============================================================================
\section{Trois Natures de Référentiels}
\label{sec:c04-natures}
% ============================================================================

\subsection{Cadre, norme, bonnes pratiques}

Le mot « référentiel » recouvre des documents de nature très différente.
Les confondre conduit à des promesses impossibles, comme celle d'une
« certification COBIT » d'entreprise, qui n'existe pas.

\begin{boxdef}
\termecle{Cadre de référence}\index{cadre de référence}
(\textit{framework})~: ensemble structuré de principes, de concepts et
d'objectifs qui organise un domaine et guide sa mise en œuvre, sans
imposer d'exigences vérifiables par un tiers. Exemples~: COBIT, COSO.

\medskip
\termecle{Norme certifiable}\index{norme certifiable}
(\textit{certifiable standard})~: norme qui formule des \emph{exigences}
(le verbe « doit »), dont le respect peut être audité par un organisme de
certification accrédité, qui délivre alors un certificat à
l'organisation. Exemples~: ISO/IEC~27001, ISO/IEC~20000-1, ISO~22301.

\medskip
\termecle{Recueil de bonnes pratiques}\index{bonnes pratiques}
(\textit{body of knowledge, guidance})~: ensemble de pratiques reconnues,
présentées comme des recommandations (« il convient de »), que
l'organisation adapte librement. Exemples~: ITIL, PMBOK, ISO/IEC~27002,
ISO~31000.
\end{boxdef}

\subsection{Certifier une organisation, certifier une personne}

La seconde confusion porte sur ce qui est certifié. Une organisation peut
être certifiée conforme à une norme d'exigences. Une personne peut être
certifiée compétente dans un référentiel. Les deux n'ont pas la même
portée.

\begin{tableaucours}{L{3.0cm}L{4.4cm}L{4.8cm}}{%
    \tableauHeader{Référentiel} &
    \tableauHeader{Certification de l'organisation} &
    \tableauHeader{Certification des personnes (exemples)}
}
COBIT & Non (évaluation de capacité possible) & COBIT Foundation, CGEIT,
CISA (ISACA) \\
ITIL & Non (voir ISO/IEC~20000-1) & ITIL Foundation et niveaux supérieurs
(PeopleCert) \\
ISO/IEC 20000-1 & Oui & Auditeur ou responsable de mise en œuvre \\
ISO/IEC 27001 & Oui & Auditeur ou responsable de mise en œuvre \\
ISO/IEC 38500 & Non & Formations spécialisées \\
COSO & Non & Certificats de l'IIA et d'autres organismes \\
CMMI & Évaluation de niveau de maturité (pas une certification ISO) &
Évaluateurs agréés \\
TOGAF & Non & TOGAF (The Open Group) \\
PMBOK & Non & PMP et autres (PMI) \\
\end{tableaucours}
\vspace{-0.8em}
\begin{center}{\small\itshape Tableau~4.1. Ce qui se certifie, et pour qui}\end{center}

\begin{boxalert}
\textbf{Violation typique.} Une organisation obtient la certification
ISO~27001 et en conclut que sa gouvernance des SI est « conforme ». Or le
périmètre du certificat peut ne couvrir qu'un centre de données, et la
norme ne traite que du management de la sécurité de l'information~: elle
ne dit rien de la valeur des investissements ni de la qualité des
services.

\smallskip
\textbf{Règle~:} lire toujours le \emph{périmètre} d'un certificat, et ne
jamais confondre la certification d'un système de management avec une
gouvernance des SI d'ensemble.
\end{boxalert}

\subsection{Le cadre légal camerounais}

Une certification volontaire ne remplace pas les obligations légales.

\begin{boxjuris}
\textbf{\loicam{n°~2010/012} du 21~décembre~2010 relative à la
cybersécurité et la cybercriminalité au Cameroun, article~11} (version
anglaise officielle)~:

\medskip
\textit{« The following activities may be subject to authorization:
\textelp{} the provision of security auditing, security programs editing,
and other authorized security services to the public. »}

\medskip
Soit~: la fourniture au public de services d'audit de sécurité peut être
soumise à autorisation.

\medskip
\textbf{Portée.} Au Cameroun, l'audit de sécurité des systèmes
d'information n'est pas une activité libre au même titre que le conseil
en organisation. Les articles~13 et~32 de la même loi
(Chapitre~\ref{chap:prologue}) rendent par ailleurs cet audit obligatoire
pour certains systèmes. Un certificat ISO/IEC~27001, délivré par un
organisme étranger, ne dispense donc pas des audits prévus par la loi~:
les deux démarches coexistent, et une bonne gouvernance les coordonne. Le
Chapitre~\ref{chap:droit} précisera les textes d'application.
\end{boxjuris}

% ============================================================================
\section{Les Référentiels de la Gouvernance}
\label{sec:c04-gouvernance}
% ============================================================================

\subsection{COBIT : le cadre de gouvernance et de management}

COBIT est publié par l'ISACA, association internationale des
professionnels de l'audit et de la gouvernance des SI. Né en 1996 comme
un référentiel de \emph{contrôle} à l'usage des auditeurs, il est devenu
un cadre de \emph{gouvernance et de management} de l'ensemble de
l'information et des technologies de l'organisation. Sa version
actuelle, COBIT~2019~\cite{isaca2018cobitintro}, compte 40~objectifs
répartis entre un domaine de gouvernance (EDM) et quatre domaines de
management (APO, BAI, DSS, MEA). Elle fait l'objet de la
Partie~\ref{part:cobit}.

\textbf{Sa force}~: c'est le seul référentiel qui couvre de bout en bout
le domaine, de la décision du conseil au contrôle de l'exploitation. Il
sert de « référentiel parapluie » sous lequel les autres viennent se
ranger. \textbf{Sa limite}~: il dit \emph{quoi} maîtriser, rarement
\emph{comment} le faire en détail.

\subsection{ISO/IEC 38500 : les principes pour l'organe de gouvernance}

La norme ISO/IEC~38500 s'adresse directement à l'organe de gouvernance
(conseil d'administration, direction). Sa troisième édition date de
février~2024~\cite{iso38500_2024}. Courte, elle ne décrit pas de
processus, mais des principes et un modèle de gouvernance fondé sur trois
tâches~: \emph{évaluer}, \emph{diriger}, \emph{surveiller}. Elle n'est
pas certifiable. Elle est étudiée au Chapitre~\ref{chap:normes}.

\subsection{COSO : le contrôle interne et la gestion des risques}

Le Committee of Sponsoring Organizations of the Treadway Commission
(\sigle{COSO}) publie deux cadres de référence. Le premier, sur le
\emph{contrôle interne} (1992, mis à jour en 2013), est la référence des
auditeurs et des régulateurs financiers, notamment pour la loi
Sarbanes-Oxley. Le second, sur la \emph{gestion des risques de
l'entreprise} (2004, refondu en 2017 sous le titre \emph{Enterprise Risk
Management: Integrating with Strategy and Performance}), relie le risque à
la stratégie et à la performance~\cite{coso2017erm}. COSO n'est pas
propre à l'informatique, mais la gouvernance des SI s'y inscrit~: les
contrôles généraux informatiques font partie du contrôle interne.

% ============================================================================
\section{Les Référentiels des Services et de la Sécurité}
\label{sec:c04-services}
% ============================================================================

\subsection{ITIL : les bonnes pratiques de la gestion des services}

ITIL est né à la fin des années 1980 au sein de l'administration
britannique, qui cherchait à rationaliser la gestion de ses services
informatiques. Ses versions successives l'ont fait passer d'un catalogue
de processus (version~2) à un cycle de vie des services (version~3, 2007),
puis à un système de valeur des services organisé en pratiques
(ITIL~4, 2019). Il est aujourd'hui la propriété de PeopleCert, qui a
lancé en 2026 \emph{ITIL version~5}~\cite{peoplecert2026itil5}.

\textbf{Sa force}~: une description détaillée et éprouvée de la manière de
gérer les services au quotidien (incidents, changements, niveaux de
service\ldots). \textbf{Sa limite}~: il ne traite la gouvernance qu'à la
marge, et ce n'est pas une norme~: on ne « certifie » pas une organisation
ITIL.

\begin{boiteRemarque}{De ITIL~4 à ITIL version~5}
ITIL version~5 a été annoncée par PeopleCert en janvier~2026, et ses
modules de formation sont publiés progressivement depuis. Selon la
présentation de l'éditeur, elle passe d'une approche centrée sur la
gestion des services à une approche intégrée des \emph{produits et des
services numériques}, couvre leur cycle de vie de bout en bout, et se
veut conçue d'emblée pour un environnement où l'intelligence artificielle
est omniprésente~\cite{peoplecert2026itil5}. Elle conserve les sept
principes directeurs et l'essentiel des pratiques d'ITIL~4. ITIL~4 reste
commercialisé et ses certifications restent reconnues pendant la
transition.

\smallskip
Ce cours s'appuie sur \textbf{ITIL~4} dans la Partie~\ref{part:pratique}~:
c'est la version documentée par des sources publiques, et celle que
pratiquent aujourd'hui les organisations. Le Chapitre~\ref{chap:itil-svs}
situera ITIL version~5 par rapport à elle.
\end{boiteRemarque}

\subsection{ISO/IEC 20000-1 : la norme du management des services}

ISO/IEC~20000-1 formule les exigences d'un \emph{système de management des
services}. Elle est certifiable, et c'est par elle qu'une organisation
peut faire reconnaître par un tiers la qualité de sa gestion des services.
Ses exigences sont très proches des pratiques d'ITIL, dont elle est
historiquement issue~: ITIL dit \emph{comment faire}, ISO/IEC~20000-1 dit
\emph{ce qui doit être démontré}.

\subsection{ISO/IEC 27001 et 27002 : la sécurité de l'information}

ISO/IEC~27001 formule les exigences d'un \emph{système de management de la
sécurité de l'information} (\sigle{SMSI}) et est certifiable. ISO/IEC~27002
détaille les mesures de sécurité qui peuvent y être mises en œuvre. Leur
édition en vigueur date de 2022. Elles s'appuient sur une appréciation des
risques, que précise ISO/IEC~27005. La famille ISO/IEC~27000 sera reprise
au Chapitre~\ref{chap:normes}.

\subsection{ISO 22301 et ISO 31000 : continuité et risques}

ISO~22301 formule les exigences d'un système de management de la
continuité d'activité et est certifiable. ISO~31000 donne des lignes
directrices générales sur le management du risque, applicables à tout
type de risque, et n'est pas certifiable.

% ============================================================================
\section{Architecture, Projets et Maturité}
\label{sec:c04-projets}
% ============================================================================

\subsection{TOGAF : l'architecture d'entreprise}

TOGAF, publié par The Open Group, est le cadre de référence le plus
répandu en \termecle{architecture d'entreprise}\index{architecture
d'entreprise}~: la description cohérente des processus métier, des
données, des applications et des technologies d'une organisation, et de
leur évolution. Son cœur est une méthode de développement de
l'architecture (\textit{Architecture Development Method}, \sigle{ADM}).
TOGAF répond à la deuxième décision de Weill et Ross~: l'architecture
(Chapitre~\ref{chap:alignement}).

\subsection{PMBOK et PRINCE2 : la gestion de projet}

Le \textit{Guide to the Project Management Body of Knowledge}
(\sigle{PMBOK}), publié par le Project Management Institute, et la
méthode PRINCE2, d'origine britannique, sont les deux grandes références
de la gestion de projet. Le premier est un recueil de connaissances et de
principes~; la seconde, une méthode structurée en processus et en rôles.
Ils outillent l'objectif COBIT BAI11 (gérer les projets).

\subsection{CMMI : la maturité des processus}

Le CMMI (\textit{Capability Maturity Model Integration}), né à
l'université Carnegie Mellon pour évaluer les fournisseurs de logiciels du
département de la Défense américain, est aujourd'hui géré par l'ISACA. Il
décrit des niveaux de maturité, de~1 (processus improvisés) à~5
(processus en amélioration continue). COBIT~2019 s'en inspire directement
pour son modèle de capacité des processus (Chapitre~\ref{chap:cobit-capacite}).

% ============================================================================
\section{À Chaque Besoin son Référentiel}
\label{sec:c04-besoins}
% ============================================================================

Le choix d'un référentiel ne part pas du référentiel, mais du besoin.

\begin{tableaucours}{L{4.4cm}L{3.8cm}L{4.0cm}}{%
    \tableauHeader{Besoin de l'organisation} &
    \tableauHeader{Référentiel principal} &
    \tableauHeader{Compléments}
}
Organiser la gouvernance du conseil & ISO/IEC 38500, COBIT (EDM) & King~IV,
principes OCDE \\
Couvrir l'ensemble du management des SI & COBIT 2019 & · \\
Fiabiliser le contrôle interne & COSO (contrôle interne) & COBIT (MEA02) \\
Gérer les risques de l'entreprise & COSO ERM, ISO 31000 & COBIT (EDM03,
APO12) \\
Gérer les services au quotidien & ITIL & ISO/IEC 20000-1 \\
Faire reconnaître la qualité des services & ISO/IEC 20000-1 & ITIL \\
Sécuriser l'information & ISO/IEC 27001 et 27002 & COBIT (APO13, DSS05) \\
Assurer la continuité & ISO 22301 & ITIL, COBIT (DSS04) \\
Cohérence des systèmes & TOGAF & COBIT (APO03) \\
Conduire les projets & PMBOK, PRINCE2 & COBIT (BAI01, BAI11) \\
Améliorer la maturité des processus & CMMI & COBIT (capacité) \\
\end{tableaucours}
\vspace{-0.8em}
\begin{center}{\small\itshape Tableau~4.2. Correspondance entre besoins et
référentiels}\end{center}

\begin{boxobjet}{APO01 · Gérer le cadre de management des SI}
\textbf{Finalité~:} concevoir le système de management de l'information
et des technologies de l'organisation, et en garantir la cohérence avec
ses objectifs.

\textbf{Pratiques (extrait)~:} concevoir le système de management~;
communiquer les objectifs et l'orientation~; définir la structure
organisationnelle~; définir et entretenir les politiques et procédures.

\textbf{Sorties~:} politiques et procédures, définition des rôles,
\emph{choix et articulation des référentiels retenus}.

\textbf{RACI (extrait)~:} \emph{A} au directeur général, \emph{R} au DSI,
\emph{C} au responsable GRC et au RSSI.

\textbf{Indicateurs~:} part des politiques à jour et approuvées~;
couverture des processus par un référentiel identifié.
\end{boxobjet}
{\small\itshape Intitulé traduit par nos soins à partir de
COBIT~2019~\cite{isaca2018cobitobjectifs}.}

% ============================================================================
\section{La Carte d'Articulation}
\label{sec:c04-carte}
% ============================================================================

\subsection{Trois niveaux}

Les référentiels s'articulent naturellement en trois niveaux. Au
sommet, ceux qui s'adressent à l'organe de gouvernance. Au milieu, le
cadre qui organise l'ensemble du management des SI. À la base, les
référentiels spécialisés, chacun sur son domaine.

\begin{center}
\begin{tikzpicture}[
    niv/.style = {
        draw=CouleurPrimaire, fill=CouleurPrimaire, rounded corners=4pt,
        minimum width=12.6cm, minimum height=0.9cm, align=center,
        font=\footnotesize\bfseries\color{white}
    },
    mil/.style = {
        draw=CouleurBordDef, fill=CouleurDef, rounded corners=4pt,
        minimum width=12.6cm, minimum height=0.9cm, align=center,
        font=\footnotesize\bfseries\color{CouleurPrimaire}
    },
    spe/.style = {
        draw=CouleurBordDef, fill=white, rounded corners=3pt,
        text width=1.4cm, minimum height=1.3cm, align=center,
        anchor=north, inner sep=2pt,
        font=\tiny\bfseries\color{CouleurPrimaire}
    }
]
\node[niv] (g) at (0,2.6) {Gouvernance~: ISO/IEC 38500 · COBIT (EDM) ·
    COSO · King~IV};
\node[mil] (m) at (0,1.3) {Management des SI~: COBIT 2019 (APO · BAI ·
    DSS · MEA)};
\node[spe] at (-5.4,0.6) {Services\\ ITIL\\ ISO 20000};
\node[spe] at (-3.6,0.6) {Sécurité\\ ISO 27001\\ ISO 27002};
\node[spe] at (-1.8,0.6) {Continuité\\ ISO 22301};
\node[spe] at (0,0.6)    {Risques\\ ISO 31000\\ COSO ERM};
\node[spe] at (1.8,0.6)  {Architecture\\ TOGAF};
\node[spe] at (3.6,0.6)  {Projets\\ PMBOK\\ PRINCE2};
\node[spe] at (5.4,0.6)  {Maturité\\ CMMI};
\end{tikzpicture}
\end{center}
\vspace{-0.4em}
\begin{center}{\small\itshape Figure~4.1. La carte d'articulation des
référentiels}\end{center}

La règle d'articulation est simple~: \textbf{COBIT fixe les objectifs, les
référentiels spécialisés fournissent les pratiques}. COBIT~2019 le prévoit
lui-même~: pour chacun de ses objectifs, il indique les normes et
référentiels qui le détaillent.

\subsection{Construire la carte d'une organisation}

\begin{boxscenario}{Opération SANAGA. La carte de la banque}
La responsable GRC de la BANQUE SANAGA propose au conseil la carte
suivante, à partir des défaillances du prologue et des exigences de la
COBAC.
\begin{enumerate}
    \item \textbf{Gouvernance}~: principes de l'ISO/IEC~38500 dans la
          charte du comité stratégique des SI, objectifs EDM de COBIT.
    \item \textbf{Management}~: COBIT~2019 comme cadre unique, avec une
          évaluation de capacité sur les objectifs prioritaires
          (Partie~\ref{part:cobit}).
    \item \textbf{Services}~: les pratiques ITIL de gestion des
          changements, des mises en production et des incidents, en
          priorité.
    \item \textbf{Continuité}~: un plan conforme à ISO~22301, sans viser la
          certification dans un premier temps.
    \item \textbf{Sécurité}~: un SMSI conforme à ISO/IEC~27001, avec une
          certification visée à trois ans sur le périmètre des systèmes de
          paiement.
\end{enumerate}
\textbf{Ce qui est écarté, et pourquoi.} Le programme CMMI proposé par un
cabinet n'est pas retenu~: le modèle de capacité de COBIT suffit au
besoin de mesure, et la banque ne développe pas elle-même ses logiciels.
\end{boxscenario}

\begin{boxprocedure}
\textbf{La carte des référentiels, pièce de gouvernance.} Pour qu'un
auditeur puisse établir que le choix des référentiels est maîtrisé, il
doit trouver~:
\begin{itemize}
    \item un document qui liste les référentiels retenus, leur rôle et leur
          édition de référence~;
    \item la décision qui l'approuve, et sa date~;
    \item pour chaque norme certifiable visée, le périmètre et l'échéance~;
    \item la liste des référentiels écartés, avec la raison~;
    \item une date de revue, pour suivre les nouvelles éditions.
\end{itemize}
\end{boxprocedure}

\begin{boxmetier}{Responsable GRC}
\textbf{Choisir les référentiels sans les empiler}
\begin{enumerate}
    \item Partir des besoins et des obligations (lois, régulateur, clients),
          jamais de l'offre des cabinets.
    \item Retenir un cadre unique de gouvernance et de management, puis un
          référentiel spécialisé par domaine.
    \item Ne viser une certification que si un tiers l'exige ou si elle
          apporte un avantage identifié, et toujours sur un périmètre
          défini.
    \item Suivre les éditions~: une politique rédigée sur une édition
          périmée expose à des écarts lors de l'audit.
\end{enumerate}
\end{boxmetier}

\begin{boiteRemarque}{Les référentiels importés sont-ils adaptés à l'Afrique~?}
Les référentiels internationaux ont été conçus pour des organisations
dotées de moyens, de compétences et de marchés de prestataires que
beaucoup d'organisations camerounaises n'ont pas. Ce n'est pas une raison
de les écarter, mais d'en faire un usage \emph{proportionné}~: retenir les
objectifs essentiels, viser des niveaux de capacité réalistes, et
privilégier les pratiques qui traitent les risques locaux les plus
probables (dépendance aux prestataires, coupures d'énergie, rareté des
compétences). COBIT~2019 offre précisément cet outil d'adaptation avec
ses facteurs de conception (Chapitre~\ref{chap:cobit-conception}).
\end{boiteRemarque}

% ============================================================================
\separateur
% ============================================================================

\begin{boxresume}
\textbf{Ce qu'il faut retenir de ce chapitre~:}
\begin{itemize}
    \item \textbf{Trois natures}~: les cadres (COBIT, COSO) organisent, les
          normes certifiables (ISO/IEC~27001, 20000-1, ISO~22301) exigent,
          les recueils de bonnes pratiques (ITIL, PMBOK, ISO/IEC~27002)
          recommandent.
    \item \textbf{Certifier une organisation n'est pas certifier une
          personne}~: il n'existe pas de certification COBIT ou ITIL d'une
          entreprise.
    \item \textbf{Lire le périmètre}~: un certificat ne vaut que pour le
          périmètre qu'il désigne.
    \item \textbf{Trois niveaux}~: gouvernance (ISO/IEC~38500, COBIT EDM,
          COSO), management (COBIT~2019), domaines spécialisés (ITIL,
          ISO~27001, ISO~22301, TOGAF, PMBOK, CMMI\ldots).
    \item \textbf{COBIT fixe les objectifs, les autres fournissent les
          pratiques}~: ils ne se contredisent pas, ils se répartissent le
          travail.
    \item \textbf{Partir du besoin}~: une carte d'articulation se construit à
          partir des risques et des obligations, et se révise à chaque
          nouvelle édition.
    \item \textbf{La loi s'ajoute aux normes}~: au Cameroun, les audits de
          sécurité prévus par la loi n°~2010/012 coexistent avec les
          certifications volontaires.
\end{itemize}
\end{boxresume}

\bigskip

\begin{boxexo}[title={\ding{45}\quad Exercice~4.1. Quelle nature~? \niveaubase}]
Pour chacun des documents suivants, indiquez s'il s'agit d'un cadre de
référence, d'une norme certifiable ou d'un recueil de bonnes pratiques, et
si une organisation peut être certifiée sur sa base.
\begin{enumerate}[(a)]
    \item ISO/IEC~27001.
    \item ITIL.
    \item COSO ERM.
    \item ISO/IEC~20000-1.
    \item ISO~31000.
    \item COBIT~2019.
\end{enumerate}
\end{boxexo}

\begin{boxexo}[title={\ding{45}\quad Exercice~4.2. Lire un certificat \niveaumoyen}]
Un prestataire candidat à l'infogérance de KRIBILOG SA joint à son offre
un certificat ISO/IEC~27001 dont le périmètre est~: « Hébergement des
serveurs du centre de données de Casablanca ». Le service proposé à
KRIBILOG serait exploité depuis Douala.
\begin{enumerate}[(a)]
    \item Que prouve ce certificat, et que ne prouve-t-il pas pour le
          service proposé~?
    \item Quelles pièces complémentaires M.~MBARGA doit-il exiger~?
    \item Quelle clause proposez-vous d'insérer dans le contrat
          (Chapitre~\ref{chap:itil-niveaux})~?
\end{enumerate}
\end{boxexo}

\begin{boxexo}[title={\ding{45}\quad Exercice~4.3. La carte de NOUN PAY \niveauexpert}]
NOUN PAY, jeune entreprise de paiement mobile de Bafoussam, compte
45~salariés, héberge tout son SI chez un fournisseur d'informatique en
nuage étranger et utilise un modèle d'intelligence artificielle pour
noter les demandes de microcrédit. Ses investisseurs exigent « une
gouvernance des SI au niveau des standards internationaux ».
\begin{enumerate}[(a)]
    \item Identifiez ses besoins et ses obligations prioritaires.
    \item Construisez sa carte d'articulation sur le modèle de la
          Figure~4.1, en justifiant chaque référentiel retenu.
    \item Indiquez au moins deux référentiels que vous écartez, et
          pourquoi.
    \item Proposez un calendrier sur trois ans, compatible avec la taille
          de l'entreprise.
\end{enumerate}
\end{boxexo}

\bigskip

\begin{boxinfo}
\textbf{Transition.} La Partie~\ref{part:fondements} a posé les fondements
de la gouvernance des SI~: ses origines, ses défaillances, ses modèles et
ses référentiels. La Partie~\ref{part:cobit} entre dans le référentiel
central de ce cours. Le Chapitre~\ref{chap:cobit-archi} présente les
principes et l'architecture de COBIT~2019.
\end{boxinfo}


    % ======================================================================
    % PARTIE II — COBIT 2019 : LE RÉFÉRENTIEL DE LA GOUVERNANCE
    % « Un système de gouvernance se conçoit, se mesure et s'améliore. »
    % Standards : COBIT 2019 (Introduction and Methodology, Governance and
    % Management Objectives, Design Guide, Implementation Guide), COBIT 4.1.
    % Niveau : fondamental à intermédiaire
    % ======================================================================
    \part{COBIT 2019 : le Référentiel de la Gouvernance}
    \label{part:cobit}

        % ============================================================================
% Fichier  : partie2/P02_C05_cobit_principes_architecture.tex
% Titre    : COBIT 2019 : Principes et Architecture
% Auteur   : MINKA MI NGUIDJOÏ Thierry Emmanuel
% Version  : 1.0 -- octobre 2026
% ============================================================================
% Source principale : COBIT 2019 Framework: Introduction and Methodology
% (ISACA, 2018), paraphrasée et traduite par nos soins (dossier Normes).
% ============================================================================

\chapter{COBIT 2019 : Principes et Architecture}
\label{chap:cobit-archi}
\chaptermark{COBIT 2019 : principes et architecture}

\epigraph{%
    « COBIT ne prend aucune décision à votre place. Il dit quelles
    décisions doivent être prises, comment, et par qui. »%
}{D'après ISACA, \textit{COBIT 2019 Framework: Introduction and
Methodology} (2018, notre traduction)}

\niveauChapitre{Fondamental}{%
    Maîtrise du Chapitre~\ref{chap:referentiels} (place de COBIT parmi les
    référentiels) et du Chapitre~\ref{chap:alignement} (trois finalités).%
}

\begin{boxobjectifs}
\begin{itemize}
    \item Retracer l'évolution de COBIT, de la version~1 à COBIT~2019.
    \item Énoncer les six principes du système de gouvernance et les trois
          principes du cadre.
    \item Décrire les sept composantes d'un système de gouvernance.
    \item Présenter les cinq domaines (EDM, APO, BAI, DSS, MEA) et les
          domaines d'intérêt.
\end{itemize}
\end{boxobjectifs}

\bigskip

% ============================================================================
\section*{Positionnement : de la carte au territoire}
% ============================================================================

Le Chapitre~\ref{chap:referentiels} a placé COBIT au centre de la carte des
référentiels~: il fixe les objectifs, et les autres référentiels
fournissent les pratiques. La Partie~\ref{part:cobit}, qui s'ouvre ici,
entre dans ce référentiel central. Elle suit la logique même de COBIT~:
comprendre son architecture (ce chapitre), connaître ses objectifs
(Chapitre~\ref{chap:cobit-objectifs}), concevoir un système de gouvernance
adapté (Chapitre~\ref{chap:cobit-conception}), puis le mesurer et le
mettre en œuvre (Chapitre~\ref{chap:cobit-capacite}).

La question directrice de ce chapitre est simple~: \emph{de quoi un
système de gouvernance des SI est-il fait~?} La réponse de COBIT~2019 tient
en trois mots~: des \textbf{principes}, des \textbf{composantes} et des
\textbf{objectifs}. Le chapitre les présente dans cet ordre, après un
retour sur l'histoire du référentiel.

\separateur

% ============================================================================
\section{Trente Ans d'Évolution}
\label{sec:c05-histoire}
% ============================================================================

\subsection{Du contrôle à la gouvernance}

COBIT est né en 1996 sous le nom de \textit{Control Objectives for
Information and Related Technology}. Il était alors un référentiel
d'objectifs de contrôle à l'usage des auditeurs des systèmes
d'information. Chaque nouvelle version a élargi son public et son champ,
de l'audit vers le management, puis vers la gouvernance de l'ensemble de
l'information et des technologies.

\begin{tableaucours}{L{2.2cm}L{4.4cm}L{5.6cm}}{%
    \tableauHeader{Version} &
    \tableauHeader{Public principal} &
    \tableauHeader{Apport}
}
COBIT 1 (1996) & Auditeurs & Objectifs de contrôle des technologies de
l'information \\
1998, 2000 & Auditeurs, puis managers & Éditions successives~: guides de
management, modèles de maturité \\
COBIT 4.0 et 4.1 (2005, 2007) & Direction informatique, direction
générale & Gouvernance des SI~; 34~processus en quatre domaines (PO, AI,
DS, ME)~\cite{itgi2007cobit41} \\
COBIT 5 (2012) & Ensemble de l'entreprise & Gouvernance de l'entreprise
numérique~; cinq principes, sept « facilitateurs », 37~processus~;
intégration de Val~IT et Risk~IT \\
COBIT 2019 (2018) & Conseil, direction, auditeurs, praticiens & Système
de gouvernance sur mesure~; facteurs de conception, domaines d'intérêt,
40~objectifs, mesure de la capacité alignée sur CMMI \\
\end{tableaucours}
\vspace{-0.8em}
\begin{center}{\small\itshape Tableau~5.1. Évolution de COBIT}\end{center}

\subsection{Ce que COBIT 2019 apporte}

Par rapport à COBIT~5, dont il conserve l'essentiel, COBIT~2019 met en
avant quatre améliorations~\cite{isaca2018cobitintro}~:
\begin{itemize}
    \item \textbf{la flexibilité}~: les facteurs de conception permettent
          d'adapter le système de gouvernance au contexte de
          l'organisation, et l'architecture ouverte permet d'ajouter des
          domaines d'intérêt sans toucher au modèle central~;
    \item \textbf{l'actualité}~: le modèle se rattache plus facilement aux
          normes et réglementations récentes~;
    \item \textbf{l'application prescriptive}~: une fois adapté, le modèle
          sert de prescription pour le système de gouvernance de
          l'organisation~;
    \item \textbf{la mesure de la performance}~: les notions de capacité et
          de maturité sont intégrées au modèle et alignées sur CMMI.
\end{itemize}

\subsection{Ce que COBIT n'est pas}

COBIT~2019 prend soin de dissiper quatre malentendus. COBIT \emph{n'est
pas} une description complète de l'environnement informatique d'une
organisation. Il \emph{n'est pas} un cadre d'organisation des processus
métier. Il \emph{n'est pas} un référentiel technique de gestion des
technologies. Surtout, il \emph{ne prend ni ne prescrit aucune décision}~:
il ne dit pas quelle est la meilleure stratégie, la meilleure architecture
ou le bon niveau de dépense, mais quelles décisions doivent être prises,
comment, et par qui. C'est le sens de l'épigraphe.

\begin{boxalert}
\textbf{Violation typique.} Un cabinet promet à la BANQUE SANAGA de
« mettre la banque en conformité COBIT » en six mois, avec à la clé une
« certification ». La promesse repose sur deux erreurs~: COBIT n'est pas
une norme d'exigences, et il n'existe pas de certification COBIT d'une
organisation (Chapitre~\ref{chap:referentiels}).

\smallskip
\textbf{Règle~:} on ne se « met pas en conformité » avec COBIT~; on s'en
sert pour concevoir et évaluer un système de gouvernance, dont on mesure
la capacité des processus (Chapitre~\ref{chap:cobit-capacite}).
\end{boxalert}

% ============================================================================
\section{Les Principes de COBIT 2019}
\label{sec:c05-principes}
% ============================================================================

COBIT~2019 repose sur deux séries de principes~: six principes qui
décrivent ce qu'un \emph{système de gouvernance} doit être, et trois
principes qui décrivent ce qu'un \emph{cadre de gouvernance}, c'est-à-dire
un référentiel comme COBIT lui-même, doit être~\cite{isaca2018cobitintro}.

\subsection{Six principes pour un système de gouvernance}

\begin{tableaucours}{L{3.6cm}L{8.6cm}}{%
    \tableauHeader{Principe} &
    \tableauHeader{Contenu}
}
1. Créer de la valeur pour les parties prenantes & Toute organisation a
besoin d'un système de gouvernance pour satisfaire ses parties prenantes
et tirer de la valeur des technologies. La valeur est un équilibre entre
bénéfices, risques et ressources. \\
2. Approche globale & Le système de gouvernance est fait de composantes de
types différents, qui fonctionnent ensemble. \\
3. Système dynamique & Chaque fois qu'un facteur de conception change
(stratégie, technologie\ldots), l'effet de ce changement sur le système de
gouvernance doit être examiné. \\
4. Gouvernance distincte du management & Le système distingue clairement
les activités et les structures de gouvernance de celles de management. \\
5. Adapté aux besoins de l'organisation & Le système est ajusté au moyen
de facteurs de conception, qui servent à adapter et à prioriser ses
composantes. \\
6. Couverture de bout en bout & Le système couvre toute l'organisation,
pas seulement la fonction informatique~: tout traitement de l'information
et toute technologie, où qu'ils se trouvent. \\
\end{tableaucours}
\vspace{-0.8em}
\begin{center}{\small\itshape Tableau~5.2. Les six principes d'un système
de gouvernance (d'après COBIT~2019, notre traduction)}\end{center}

\begin{boxscenario}{Opération SANAGA. Les principes violés}
\textbf{Constatation.} L'application « SANAGA Mobile » a été conçue par la
direction marketing avec une agence extérieure, sans passer par la DSI.
Son lancement a fixé la date de bascule du système central.

\medskip
\textbf{Analyse} (confiance élevée). Trois principes au moins ont été
méconnus. Le principe~6~: une application critique s'est développée en
dehors du périmètre de la gouvernance, parce qu'elle ne relevait pas de la
DSI. Le principe~4~: le directeur général a pris seul une décision de
gouvernance (l'arbitrage entre enjeu commercial et risque). Le
principe~3~: le changement de stratégie (le lancement mobile) n'a donné
lieu à aucun réexamen du système de gouvernance du projet HORIZON.
\end{boxscenario}

\subsection{Trois principes pour un cadre de gouvernance}

Les trois autres principes s'appliquent au référentiel lui-même. Un cadre
de gouvernance doit être \textbf{fondé sur un modèle conceptuel}, qui
identifie les composantes et leurs relations, pour garantir la cohérence
et permettre l'automatisation. Il doit être \textbf{ouvert et flexible},
pour accueillir de nouveaux contenus sans perdre sa cohérence. Il doit
être \textbf{aligné sur les grandes normes}, référentiels et
réglementations. Ces principes expliquent pourquoi COBIT renvoie, pour
chaque objectif, aux normes qui le détaillent (Chapitre~\ref{chap:referentiels}).

% ============================================================================
\section{Les Composantes du Système de Gouvernance}
\label{sec:c05-composantes}
% ============================================================================

\subsection{Sept types de composantes}

Pour atteindre ses objectifs, une organisation construit son système de
gouvernance à partir de composantes. COBIT~5 les appelait
« facilitateurs » (\textit{enablers})~; COBIT~2019 parle de
\termecle{composantes}\index{composantes du système de gouvernance}.

\begin{boxdef}
\termecle{Composante}\index{composante} (\textit{component})~: élément qui
contribue, seul et avec les autres, au bon fonctionnement du système de
gouvernance de l'information et des technologies. Les composantes
interagissent et forment un système global. Elles sont de sept types.
\end{boxdef}

\begin{tableaucours}{L{4.2cm}L{8.0cm}}{%
    \tableauHeader{Composante} &
    \tableauHeader{Rôle}
}
Processus & Ensemble organisé de pratiques et d'activités qui atteint un
objectif et produit des résultats. \\
Structures organisationnelles & Instances de décision de l'organisation
(conseil, comités, fonctions). \\
Principes, politiques et procédures & Traduisent le comportement attendu
en consignes pratiques pour le management quotidien. \\
Information & Toute information produite et utilisée~; COBIT s'intéresse à
celle qu'exige le fonctionnement du système de gouvernance. \\
Culture, éthique et comportement & Facteurs souvent sous-estimés de la
réussite de la gouvernance. \\
Personnes, aptitudes et compétences & Nécessaires aux bonnes décisions et
à la bonne exécution. \\
Services, infrastructure et applications & Moyens techniques qui
supportent le système de gouvernance lui-même. \\
\end{tableaucours}
\vspace{-0.8em}
\begin{center}{\small\itshape Tableau~5.3. Les sept composantes (d'après
COBIT~2019, notre traduction)}\end{center}

\begin{center}
\begin{tikzpicture}[
    comp/.style = {
        draw=CouleurBordDef, fill=CouleurDef, rounded corners=4pt,
        text width=2.5cm, minimum height=1.0cm, align=center,
        font=\scriptsize\bfseries\color{CouleurPrimaire}
    },
    centre/.style = {
        draw=CouleurPrimaire, fill=CouleurPrimaire, circle,
        text width=1.9cm, align=center, font=\scriptsize\bfseries\color{white}
    }
]
\node[centre] (c) at (0,0) {Système de\\ gouvernance};
\foreach \ang/\txt in {90/Processus,
    39/{Structures\\ organisationnelles},
    -13/{Principes,\\ politiques,\\ procédures},
    -64/Information,
    -116/{Culture, éthique,\\ comportement},
    -167/{Personnes,\\ compétences},
    141/{Services, infra-\\ structure, applications}}
{
    \node[comp] at (\ang:3.3cm) {\txt};
    \draw[CouleurSecondaire, thick] (c) -- (\ang:2.35cm);
}
\end{tikzpicture}
\end{center}
\vspace{-0.4em}
\begin{center}{\small\itshape Figure~5.1. Les sept composantes du système
de gouvernance}\end{center}

\subsection{La leçon des composantes}

La leçon principale tient au principe~2 (approche globale)~: \textbf{un
processus seul ne suffit jamais}. La BANQUE SANAGA avait, sur le papier,
une procédure de mise en production. Il lui manquait les autres
composantes~: la structure (un comité des changements), l'information (un
procès-verbal de recette fiable), la culture (le droit de dire non à une
date) et les compétences (des testeurs formés au nouveau système).

\begin{boxprocedure}
\textbf{Une composante qui existe vraiment.} Pour chaque composante d'un
objectif prioritaire, l'auditeur cherche une preuve de son existence
\emph{et} de son fonctionnement~:
\begin{itemize}
    \item \emph{processus}~: une procédure approuvée, et des
          enregistrements de son exécution~;
    \item \emph{structure}~: une charte, et des procès-verbaux~;
    \item \emph{politique}~: un document approuvé, daté, diffusé~;
    \item \emph{information}~: le document attendu, à jour et utilisé~;
    \item \emph{culture}~: des signaux observables (alertes remontées,
          décisions de report acceptées)~;
    \item \emph{compétences}~: des fiches de poste et un plan de formation
          suivi~;
    \item \emph{services et applications}~: l'outil qui supporte le
          processus, et ses journaux.
\end{itemize}
\end{boxprocedure}

\subsection{Composantes génériques et variantes}

Les composantes décrites dans le modèle central de COBIT sont
\emph{génériques}~: elles valent en principe pour toute organisation, mais
doivent être adaptées avant d'être mises en œuvre. Une \emph{variante} est
une composante générique adaptée à un usage particulier~: la sécurité de
l'information, DevOps, une réglementation donnée.

% ============================================================================
\section{Objectifs et Domaines}
\label{sec:c05-domaines}
% ============================================================================

\subsection{Objectifs de gouvernance et de management}

\begin{boxdef}
\termecle{Objectif de gouvernance ou de management}\index{objectif de
gouvernance}\index{objectif de management} (\textit{governance or
management objective})~: ce qu'il faut atteindre pour que l'information
et les technologies contribuent aux objectifs de l'organisation. Chaque
objectif se rapporte à un processus de même nom (ou d'un nom proche) et à
un ensemble de composantes des autres types qui aident à l'atteindre.
\end{boxdef}

Un objectif de gouvernance se rapporte à un processus de gouvernance, dont
le conseil et la direction générale répondent. Un objectif de management
se rapporte à un processus de management, qui relève de l'encadrement
supérieur et intermédiaire.

\subsection{Cinq domaines, quarante objectifs}

Les 40~objectifs du modèle central sont regroupés en cinq domaines, dont
les noms sont des verbes~\cite{isaca2018cobitintro}.

\begin{tableaucours}{L{1.4cm}L{3.6cm}cL{5.8cm}}{%
    \tableauHeader{Code} &
    \tableauHeader{Domaine} &
    \tableauHeader{Nb} &
    \tableauHeader{Champ}
}
\multicolumn{4}{l}{\textit{Gouvernance}} \\
EDM & Évaluer, diriger et surveiller & 5 & L'organe de gouvernance évalue
les options stratégiques, oriente la direction sur les options retenues et
surveille la réalisation de la stratégie \\
\multicolumn{4}{l}{\textit{Management}} \\
APO & Aligner, planifier et organiser & 14 & Organisation d'ensemble,
stratégie et activités de soutien \\
BAI & Bâtir, acquérir et implémenter & 11 & Définition, acquisition et mise
en œuvre des solutions, et leur intégration dans les processus métier \\
DSS & Délivrer, servir et soutenir & 6 & Exploitation et soutien des
services, y compris la sécurité \\
MEA & Surveiller, évaluer et apprécier & 4 & Suivi de la performance et de
la conformité aux objectifs internes, au contrôle interne et aux
exigences externes \\
\textbf{Total} & & \textbf{40} & \\
\end{tableaucours}
\vspace{-0.8em}
\begin{center}{\small\itshape Tableau~5.4. Les cinq domaines de COBIT~2019
(intitulés traduits par nos soins)}\end{center}

On retrouve dans ce découpage la distinction des Chapitres~\ref{chap:prologue}
et~\ref{chap:gouv-entreprise}~: le domaine EDM correspond aux trois verbes
de la gouvernance, et les quatre domaines de management au cycle
planifier, bâtir, exploiter, surveiller. Le Chapitre~\ref{chap:cobit-objectifs}
parcourt les 40~objectifs un par un.

\begin{boxobjet}{EDM05 · Assurer l'engagement des parties prenantes}
\textbf{Finalité~:} s'assurer que les parties prenantes sont identifiées et
associées au système de gouvernance, et que la mesure et le compte rendu
de la performance et de la conformité des SI leur sont transparents.

\textbf{Pratiques~:} évaluer les exigences de compte rendu des parties
prenantes~; diriger la communication et le compte rendu~; surveiller la
communication.

\textbf{Sorties~:} principes de compte rendu, plan de communication,
rapports aux parties prenantes.

\textbf{RACI (extrait)~:} \emph{A} au conseil d'administration, \emph{R} au
directeur général, \emph{C} au DSI et au responsable GRC.

\textbf{Indicateurs~:} satisfaction des parties prenantes à l'égard des
comptes rendus~; délai de transmission des rapports exigés par le
régulateur.
\end{boxobjet}
{\small\itshape Intitulé et pratiques résumés et traduits par nos soins à
partir de COBIT~2019~\cite{isaca2018cobitobjectifs}.}

\begin{boxmetier}{Responsable GRC}
\textbf{Les parties prenantes de la BANQUE SANAGA, objectif EDM05}
\begin{enumerate}
    \item Dresser la liste des parties prenantes internes (conseil,
          comités, direction, métiers, audit interne) et externes (COBAC,
          commissaire aux comptes, ANTIC, clients, intégrateurs).
    \item Pour chacune, noter ce qu'elle attend comme compte rendu, à quelle
          fréquence, et sous quelle forme.
    \item Construire le calendrier des comptes rendus de l'année.
    \item Faire approuver ce plan par le comité stratégique des SI.
\end{enumerate}
Après la panne, c'est le rapport à la COBAC qui a révélé l'absence de ce
plan~: personne ne savait qui devait le rédiger, ni avec quelles données.
\end{boxmetier}

% ============================================================================
\section{Domaines d'Intérêt et Facteurs de Conception}
\label{sec:c05-focus}
% ============================================================================

\subsection{Les domaines d'intérêt}

\begin{boxdef}
\termecle{Domaine d'intérêt}\index{domaine d'intérêt} (\textit{focus
area})~: thème, domaine ou question de gouvernance traité par un ensemble
d'objectifs de gouvernance et de management et de leurs composantes.
Exemples donnés par COBIT~2019~: petites et moyennes entreprises,
cybersécurité, transformation numérique, informatique en nuage, protection
de la vie privée, DevOps.
\end{boxdef}

Le nombre de domaines d'intérêt n'est pas limité~: c'est ce qui rend le
modèle ouvert. Un domaine d'intérêt peut combiner des composantes
génériques et des variantes. Pour NOUN PAY, jeune entreprise hébergée dans
le nuage, les domaines d'intérêt « petites et moyennes entreprises » et
« informatique en nuage » sont naturellement pertinents.

\subsection{Les facteurs de conception}

\begin{boxdef}
\termecle{Facteur de conception}\index{facteur de conception}
(\textit{design factor})~: facteur susceptible d'influencer la conception
du système de gouvernance d'une organisation. COBIT~2019 en retient onze~:
la stratégie, les objectifs de l'organisation, le profil de risque, les
problèmes liés aux SI, le paysage des menaces, les exigences de
conformité, le rôle des SI, le mode d'approvisionnement des SI, les
méthodes de mise en œuvre, la stratégie d'adoption des technologies et la
taille de l'organisation.
\end{boxdef}

Les facteurs de conception sont l'outil qui rend COBIT proportionné. Une
banque soumise à la COBAC et une jeune entreprise de 45~salariés n'auront
ni les mêmes objectifs prioritaires, ni les mêmes niveaux de capacité
cibles. Le Chapitre~\ref{chap:cobit-conception} montre comment les
appliquer, et le laboratoire GCI-LAB-02 en fait le calcul.

\begin{boiteRemarque}{Un modèle conçu ailleurs, utilisable ici~?}
Les facteurs de conception répondent en partie à l'objection du
Chapitre~\ref{chap:referentiels}~: les référentiels internationaux
seraient taillés pour d'autres contextes. En faisant de la taille, du
mode d'approvisionnement ou des exigences de conformité des paramètres
explicites, COBIT~2019 permet d'intégrer les réalités camerounaises
(dépendance aux intégrateurs, régulation de la COBAC et de l'ANTIC,
ressources limitées) dans la conception même du système. Encore faut-il
renseigner ces facteurs avec lucidité~: un DSI qui se décrit comme une
organisation « innovante » pour justifier un budget fausse le résultat.
\end{boiteRemarque}

\begin{boxmetier}{Auditeur SI}
\textbf{Vérifier qu'une organisation utilise réellement COBIT}
\begin{enumerate}
    \item Demander le document qui identifie les objectifs prioritaires
          retenus, et la manière dont ils ont été choisis (facteurs de
          conception).
    \item Pour deux ou trois objectifs, vérifier l'existence des sept
          composantes (exigence de contrôle ci-dessus).
    \item Demander la dernière évaluation de capacité et ses suites.
\end{enumerate}
Une organisation qui « applique COBIT » sans pouvoir produire ces trois
pièces en a adopté le vocabulaire, pas la méthode.
\end{boxmetier}

% ============================================================================
\separateur
% ============================================================================

\begin{boxresume}
\textbf{Ce qu'il faut retenir de ce chapitre~:}
\begin{itemize}
    \item \textbf{Du contrôle à la gouvernance}~: né en 1996 pour les
          auditeurs, COBIT est devenu avec la version 2019 un cadre de
          gouvernance et de management sur mesure.
    \item \textbf{COBIT ne décide pas}~: il dit quelles décisions doivent
          être prises, comment et par qui. Il n'est ni une norme
          certifiable, ni un référentiel technique.
    \item \textbf{Six principes pour le système}~: valeur pour les parties
          prenantes, approche globale, système dynamique, gouvernance
          distincte du management, adaptation, couverture de bout en bout.
    \item \textbf{Trois principes pour le cadre}~: modèle conceptuel,
          ouverture, alignement sur les normes.
    \item \textbf{Sept composantes}~: processus, structures, politiques,
          information, culture, compétences, services et applications. Un
          processus seul ne suffit jamais.
    \item \textbf{Cinq domaines, quarante objectifs}~: EDM (5) pour la
          gouvernance~; APO (14), BAI (11), DSS (6) et MEA (4) pour le
          management.
    \item \textbf{Domaines d'intérêt et facteurs de conception} rendent le
          modèle ouvert et proportionné.
\end{itemize}
\end{boxresume}

\bigskip

\begin{boxexo}[title={\ding{45}\quad Exercice~5.1. Classer par domaine \niveaubase}]
Rattachez chacune des activités suivantes au domaine COBIT qui convient
(EDM, APO, BAI, DSS, MEA), et justifiez en une phrase.
\begin{enumerate}[(a)]
    \item Le conseil fixe le niveau de risque informatique acceptable.
    \item La DSI élabore le plan stratégique des SI à trois ans.
    \item Une équipe recette le nouveau logiciel de crédit.
    \item Le centre de services traite une panne de messagerie.
    \item L'audit interne évalue les contrôles d'accès.
    \item La direction des achats négocie le contrat d'un infogérant.
\end{enumerate}
\end{boxexo}

\begin{boxexo}[title={\ding{45}\quad Exercice~5.2. Les composantes chez KRIBILOG \niveaumoyen}]
M.~MBARGA, responsable des services informatiques de KRIBILOG SA, veut
améliorer la gestion des incidents du logiciel de suivi des conteneurs.
\begin{enumerate}[(a)]
    \item Pour chacune des sept composantes, décrivez ce qui devrait
          exister chez KRIBILOG.
    \item Pour chacune, indiquez la pièce qui prouverait son existence et
          son fonctionnement.
    \item Lesquelles relèvent de KRIBILOG, et lesquelles de l'infogérant
          NETPORT SERVICES~? Qu'en déduisez-vous pour le contrat~?
\end{enumerate}
\end{boxexo}

\begin{boxexo}[title={\ding{45}\quad Exercice~5.3. De COBIT 4.1 à COBIT 2019 à l'ANPP \niveauexpert}]
Les procédures informatiques de l'ANPP (opération MEFOU) ont été rédigées
en 2010 à partir de COBIT~4.1 et de ses quatre domaines (PO, AI, DS, ME).
\begin{enumerate}[(a)]
    \item Établissez la correspondance entre les quatre domaines de
          COBIT~4.1 et les cinq domaines de COBIT~2019. Lequel n'a pas
          d'équivalent direct, et pourquoi est-ce important pour l'ANPP~?
    \item À l'aide des documents du dossier \texttt{Normes}, identifiez
          deux objectifs de COBIT~2019 sans équivalent clair dans
          COBIT~4.1, et expliquez leur intérêt pour l'agence.
    \item Rédigez une note de deux pages au directeur général proposant
          une démarche de mise à jour, en appliquant les principes~3 et~5.
\end{enumerate}
\end{boxexo}

\bigskip

\begin{boxinfo}
\textbf{Transition.} Ce chapitre a présenté l'architecture de COBIT~2019~:
principes, composantes, domaines. Le Chapitre~\ref{chap:cobit-objectifs}
entre dans le détail des 40~objectifs, apprend à lire leur structure et
montre comment ils s'enchaînent depuis les objectifs de l'organisation
par la cascade des objectifs.
\end{boxinfo}

        % ============================================================================
% Fichier  : partie2/P02_C06_cobit_objectifs.tex
% Titre    : Les 40 Objectifs de Gouvernance et de Management
% Auteur   : MINKA MI NGUIDJOÏ Thierry Emmanuel
% Version  : 1.0 -- octobre 2026
% ============================================================================
% Sources : COBIT 2019 Framework: Governance and Management Objectives et
% Introduction and Methodology (ISACA, 2018), dossier Normes. Intitulés
% traduits par nos soins ; liens objectifs/buts relevés dans les fiches
% officielles (EDM03, APO12, BAI06, BAI07, DSS02, DSS04).
% ============================================================================

\chapter{Les 40 Objectifs de Gouvernance et de Management}
\label{chap:cobit-objectifs}
\chaptermark{Les 40 objectifs de COBIT 2019}

\epigraph{%
    « Quarante objectifs ne font pas quarante chantiers.\\
    Ils font une liste dans laquelle on choisit. »%
}{Principe de mise en œuvre de COBIT}

\niveauChapitre{Fondamental}{%
    Maîtrise du Chapitre~\ref{chap:cobit-archi} (principes, composantes et
    domaines de COBIT~2019).%
}

\begin{boxobjectifs}
\begin{itemize}
    \item Parcourir les objectifs de chaque domaine et en lire la
          structure.
    \item Rédiger la fiche processus d'un objectif (finalité, pratiques,
          entrées et sorties, RACI, indicateurs).
    \item Relier objectifs d'entreprise, objectifs d'alignement et
          objectifs COBIT (cascade des objectifs).
    \item Identifier les objectifs prioritaires d'une banque et d'une
          administration.
\end{itemize}
\end{boxobjectifs}

\bigskip

% ============================================================================
\section*{Question directrice : par où commencer~?}
% ============================================================================

Mme~EKAMBI, directrice de l'audit interne de la BANQUE SANAGA, a remis au
conseil un tableau des 40~objectifs de COBIT~2019. Devant chacun, une case
« conforme » ou « non conforme ». Trente-deux cases sont rouges. Le
président s'inquiète~: « Faut-il lancer trente-deux chantiers~? »

La réponse est non, et ce chapitre explique pourquoi. Il apprend d'abord à
\textbf{lire} un objectif (section~\ref{sec:c06-lire}), puis parcourt les
\textbf{quarante objectifs}, domaine par domaine
(sections~\ref{sec:c06-edm} à~\ref{sec:c06-dss-mea}). Il montre enfin
comment la \textbf{cascade des objectifs} relie les besoins de
l'organisation aux objectifs COBIT, et permet de choisir ceux qui comptent
(sections~\ref{sec:c06-cascade} et~\ref{sec:c06-priorites}).

\separateur

% ============================================================================
\section{Lire un Objectif COBIT}
\label{sec:c06-lire}
% ============================================================================

\subsection{L'anatomie d'une fiche}

Le guide \textit{COBIT~2019 Framework: Governance and Management
Objectives}~\cite{isaca2018cobitobjectifs} présente chaque objectif sous
une forme constante. Savoir la lire est la première compétence du
praticien.

\begin{tableaucours}{L{3.8cm}L{8.4cm}}{%
    \tableauHeader{Rubrique} &
    \tableauHeader{Contenu}
}
En-tête & Domaine, code et intitulé de l'objectif, domaine d'intérêt
(le modèle central, ou un domaine particulier) \\
Description et finalité & Ce que l'objectif couvre, et pourquoi il existe \\
Objectifs liés & Objectifs d'entreprise et objectifs d'alignement
auxquels l'objectif contribue \emph{à titre principal}, avec des exemples
d'indicateurs \\
A. Processus & Les pratiques, leurs activités (chacune rattachée à un
niveau de capacité) et des indicateurs \\
B. Structures organisationnelles & Une matrice RACI par pratique \\
C. Flux d'information & Entrées et sorties de chaque pratique, avec leur
provenance et leur destination \\
D. Personnes et compétences & Compétences requises, avec renvoi aux
référentiels de compétences \\
E. Politiques et procédures & Politiques pertinentes et leur contenu \\
F. Culture et comportement & Éléments culturels clés \\
G. Services et applications & Outils et services qui supportent
l'objectif \\
\end{tableaucours}
\vspace{-0.8em}
\begin{center}{\small\itshape Tableau~6.1. Structure d'une fiche d'objectif
COBIT~2019 (d'après ISACA, notre traduction)}\end{center}

On retrouve dans les rubriques~A à~G les sept composantes du
Chapitre~\ref{chap:cobit-archi}. Une fiche COBIT n'est donc pas la
description d'un processus~: c'est la description d'un \emph{morceau de
système de gouvernance}, avec toutes ses composantes.

\subsection{Pratiques et activités}

\begin{boxdef}
\termecle{Pratique}\index{pratique (COBIT)} (\textit{governance or
management practice})~: action clé qu'il faut accomplir pour atteindre un
objectif. Elle porte un code (par exemple DSS02.04) et un intitulé à
l'infinitif.

\medskip
\termecle{Activité}\index{activité (COBIT)} (\textit{activity})~: action
concrète qui met en œuvre une pratique. Chaque activité est rattachée à
un niveau de capacité (de~2 à~5)~: c'est ce qui permet de mesurer
la capacité d'un processus (Chapitre~\ref{chap:cobit-capacite}).
\end{boxdef}

\subsection{Rédiger une fiche processus}

Le cours présente les objectifs sous une forme condensée, la
\emph{fiche processus}, qui retient l'essentiel des rubriques officielles.
Voici la fiche complète d'un objectif central de l'exploitation.

\begin{boxobjet}{DSS02 · Gérer les demandes de service et les incidents}
\textbf{Description~:} répondre de manière rapide et efficace aux demandes
des utilisateurs et résoudre tous les types d'incidents~; rétablir le
service normal, enregistrer et satisfaire les demandes, enregistrer,
analyser, diagnostiquer, faire remonter et résoudre les incidents.

\textbf{Finalité~:} accroître la productivité et réduire les
perturbations par une résolution rapide des demandes et des incidents.

\textbf{Objectifs liés~:} EG01 (portefeuille de produits et services
compétitifs), EG08 (optimisation des processus internes)~; AG05 (services
conformes aux exigences des métiers).

\textbf{Pratiques~:} DSS02.01 définir les schémas de classification des
incidents et des demandes~; DSS02.02 enregistrer, classer et prioriser~;
DSS02.03 vérifier, approuver et satisfaire les demandes~; DSS02.04
analyser, diagnostiquer et affecter les incidents~; DSS02.05 résoudre les
incidents et rétablir le service~; DSS02.06 clore les demandes et les
incidents~; DSS02.07 suivre l'état et produire des rapports.

\textbf{Entrées et sorties (exemples)~:} entrées~: accords de niveau de
service (APO09), base des erreurs connues (DSS03)~; sorties~: tickets
classés et priorisés, rapports d'incidents, demandes de changement
(vers BAI06).

\textbf{RACI (extrait, adapté à la BANQUE SANAGA)~:} DSS02.05~: \emph{A}
au responsable de l'exploitation, \emph{R} au centre de services et aux
équipes de support, \emph{C} aux métiers concernés, \emph{I} au DSI pour
les incidents majeurs.

\textbf{Indicateurs~:} délai moyen de résolution, taux de résolution au
premier contact, nombre d'incidents réouverts, satisfaction des
utilisateurs.
\end{boxobjet}
{\small\itshape Description, finalité, objectifs liés et pratiques résumés
et traduits par nos soins à partir de COBIT~2019~\cite{isaca2018cobitobjectifs}.}

\begin{boxprocedure}
\textbf{Une fiche processus utilisable.} Pour qu'une fiche processus serve
réellement, et puisse être auditée, elle doit~:
\begin{itemize}
    \item citer l'objectif COBIT et l'édition de référence~;
    \item nommer le propriétaire du processus, au niveau de l'organisation
          (et non du référentiel)~;
    \item remplacer les rôles génériques de COBIT par les fonctions réelles
          de l'organisation dans la matrice RACI~;
    \item préciser pour chaque indicateur sa cible, sa source et sa
          fréquence~;
    \item être approuvée et datée, et revue au moins une fois par an.
\end{itemize}
\end{boxprocedure}

% ============================================================================
\section{Le Domaine de Gouvernance : EDM}
\label{sec:c06-edm}
% ============================================================================

Les cinq objectifs du domaine EDM relèvent de l'organe de gouvernance. Ils
reprennent, objectif par objectif, la définition de travail du
Chapitre~\ref{chap:gouv-entreprise}~: un cadre de gouvernance, trois
finalités et l'engagement des parties prenantes.

\begin{tableaucours}{L{1.4cm}L{4.6cm}L{6.2cm}}{%
    \tableauHeader{Code} &
    \tableauHeader{Intitulé} &
    \tableauHeader{Question pour le conseil}
}
EDM01 & Assurer la définition et le maintien du cadre de gouvernance &
Qui décide de quoi, et le dispositif fonctionne-t-il~? \\
EDM02 & Assurer la réalisation des bénéfices & Nos investissements
produisent-ils la valeur annoncée~? \\
EDM03 & Assurer l'optimisation des risques & Les risques liés aux SI
restent-ils dans notre appétence~? \\
EDM04 & Assurer l'optimisation des ressources & Avons-nous les
compétences et les moyens de notre stratégie~? \\
EDM05 & Assurer l'engagement des parties prenantes & Rendons-nous compte
à ceux qui doivent l'être~? \\
\end{tableaucours}
\vspace{-0.8em}
\begin{center}{\small\itshape Tableau~6.2. Les objectifs du domaine EDM}\end{center}

Les fiches EDM01, EDM02 et EDM05 ont été présentées aux
Chapitres~\ref{chap:gouv-entreprise}, \ref{chap:alignement}
et~\ref{chap:cobit-archi}. EDM03 a pour finalité de s'assurer que les
risques d'entreprise liés aux SI ne dépassent pas l'appétence et la
tolérance au risque de l'organisation~; il sera repris avec APO12.

% ============================================================================
\section{Aligner, Planifier et Organiser : APO}
\label{sec:c06-apo}
% ============================================================================

Le domaine APO est le plus vaste. Il couvre l'organisation d'ensemble des
SI, leur stratégie et les activités de soutien.

\begin{tableaucours}{L{1.4cm}L{4.6cm}L{6.2cm}}{%
    \tableauHeader{Code} &
    \tableauHeader{Intitulé} &
    \tableauHeader{Ce qu'il couvre}
}
APO01 & Gérer le cadre de management des SI & Politiques, rôles, choix des
référentiels (Chapitre~\ref{chap:referentiels}) \\
APO02 & Gérer la stratégie & Plan stratégique des SI, feuille de route \\
APO03 & Gérer l'architecture d'entreprise & Cohérence des processus,
données, applications et technologies \\
APO04 & Gérer l'innovation & Veille, expérimentation des technologies
émergentes \\
APO05 & Gérer le portefeuille & Arbitrage des investissements
(Chapitre~\ref{chap:alignement}) \\
APO06 & Gérer le budget et les coûts & Budget des SI, refacturation,
suivi des coûts \\
APO07 & Gérer les ressources humaines & Effectifs, compétences,
formation, dépendance à des personnes clés \\
APO08 & Gérer les relations & Relation entre la DSI et les métiers \\
APO09 & Gérer les accords de service & Catalogue de services, accords de
niveau de service \\
APO10 & Gérer les fournisseurs & Sélection, contrats, supervision des
prestataires \\
APO11 & Gérer la qualité & Système de management de la qualité des SI \\
APO12 & Gérer le risque & Identification, évaluation et traitement des
risques liés aux SI \\
APO13 & Gérer la sécurité & Système de management de la sécurité de
l'information \\
APO14 & Gérer les données & Gouvernance et qualité des données
(Chapitre~\ref{chap:donnees-ia}) \\
\end{tableaucours}
\vspace{-0.8em}
\begin{center}{\small\itshape Tableau~6.3. Les objectifs du domaine APO
(intitulés traduits par nos soins)}\end{center}

\begin{boxalert}
\textbf{Violation typique.} L'ANPP (opération MEFOU) confie l'objectif
APO10 à la direction des achats, qui le réduit à la passation des
marchés. Une fois le contrat signé, plus personne ne supervise les
prestataires.

\smallskip
\textbf{Règle~:} gérer un fournisseur ne s'arrête pas à la signature.
APO10 couvre toute la relation, jusqu'à la supervision de la
performance et des risques du prestataire, que l'ancien DSI de TSB a été
sanctionné pour avoir négligée (Chapitre~\ref{chap:defaillances}).
\end{boxalert}

% ============================================================================
\section{Bâtir, Acquérir et Implémenter : BAI}
\label{sec:c06-bai}
% ============================================================================

Le domaine BAI couvre tout ce qui fait \emph{changer} le système
d'information~: projets, acquisitions, mises en production. C'est
le domaine où se sont produites toutes les défaillances du
Chapitre~\ref{chap:defaillances}.

\begin{tableaucours}{L{1.4cm}L{4.6cm}L{6.2cm}}{%
    \tableauHeader{Code} &
    \tableauHeader{Intitulé} &
    \tableauHeader{Ce qu'il couvre}
}
BAI01 & Gérer les programmes & Pilotage des programmes et de leurs
bénéfices \\
BAI02 & Gérer la définition des exigences & Recueil et validation des
besoins des métiers \\
BAI03 & Gérer l'identification et la construction des solutions &
Conception, développement, acquisition \\
BAI04 & Gérer la disponibilité et la capacité & Dimensionnement et
performance (Chapitre~\ref{chap:itil-niveaux}) \\
BAI05 & Gérer le changement organisationnel & Accompagnement des
utilisateurs et des métiers \\
BAI06 & Gérer les changements informatiques & Évaluation, autorisation et
suivi des changements (Chapitre~\ref{chap:itil-changements}) \\
BAI07 & Gérer l'acceptation et la transition des changements & Recette,
mise en production, retour arrière \\
BAI08 & Gérer la connaissance & Capitalisation et partage des
connaissances \\
BAI09 & Gérer les actifs & Inventaire et cycle de vie des actifs \\
BAI10 & Gérer la configuration & Base de configuration et ses relations \\
BAI11 & Gérer les projets & Conduite des projets \\
\end{tableaucours}
\vspace{-0.8em}
\begin{center}{\small\itshape Tableau~6.4. Les objectifs du domaine BAI
(intitulés traduits par nos soins)}\end{center}

% ============================================================================
\section{Délivrer, Servir, Soutenir et Surveiller : DSS et MEA}
\label{sec:c06-dss-mea}
% ============================================================================

Le domaine DSS couvre l'exploitation quotidienne des services, sécurité
comprise. Le domaine MEA ferme la boucle~: il vérifie la performance et
la conformité, et fournit l'assurance.

\begin{tableaucours}{L{1.4cm}L{4.6cm}L{6.2cm}}{%
    \tableauHeader{Code} &
    \tableauHeader{Intitulé} &
    \tableauHeader{Ce qu'il couvre}
}
DSS01 & Gérer l'exploitation & Exploitation courante, ordonnancement,
environnement physique \\
DSS02 & Gérer les demandes de service et les incidents & Centre de
services, incidents (Chapitre~\ref{chap:itil-incidents}) \\
DSS03 & Gérer les problèmes & Causes profondes, erreurs connues \\
DSS04 & Gérer la continuité & Plans de continuité, tests, reprise \\
DSS05 & Gérer les services de sécurité & Protection contre les logiciels
malveillants, accès, surveillance \\
DSS06 & Gérer les contrôles des processus métier & Contrôles applicatifs
intégrés aux processus \\
MEA01 & Gérer la surveillance de la performance et de la conformité &
Indicateurs, tableaux de bord \\
MEA02 & Gérer le système de contrôle interne & Autoévaluation, suivi des
déficiences \\
MEA03 & Gérer la conformité aux exigences externes & Lois, règlements,
contrats \\
MEA04 & Gérer l'assurance & Missions d'assurance et d'audit
(Chapitre~\ref{chap:audit}) \\
\end{tableaucours}
\vspace{-0.8em}
\begin{center}{\small\itshape Tableau~6.5. Les objectifs des domaines DSS et
MEA (intitulés traduits par nos soins)}\end{center}

\begin{boxmetier}{Responsable des services IT}
\textbf{Les six objectifs du quotidien}
\begin{enumerate}
    \item DSS01, DSS02 et DSS03 forment le noyau de l'exploitation~:
          faire tourner, réparer, éviter que cela recommence.
    \item BAI06 et BAI07 encadrent tout ce qui modifie la production.
    \item APO09 fixe les engagements de service sur lesquels l'exploitation
          est jugée.
\end{enumerate}
Ces six objectifs correspondent aux pratiques d'ITIL~4 étudiées dans la
Partie~\ref{part:pratique}~: COBIT dit ce qui doit être maîtrisé, ITIL
comment le faire.
\end{boxmetier}

% ============================================================================
\section{La Cascade des Objectifs}
\label{sec:c06-cascade}
% ============================================================================

\subsection{Des besoins aux objectifs COBIT}

Comment choisir, parmi les quarante objectifs, ceux qui comptent pour une
organisation donnée~? COBIT~2019 propose un premier outil~: la
\termecle{cascade des objectifs}\index{cascade des objectifs}
(\textit{goals cascade})~\cite{isaca2018cobitintro}.

\noindent\begin{minipage}{\linewidth}
\begin{center}
\begin{tikzpicture}[
    etape/.style = {
        draw=CouleurBordDef, fill=CouleurDef, rounded corners=4pt,
        text width=2.5cm, minimum height=1.2cm, align=center,
        font=\footnotesize\bfseries\color{CouleurPrimaire}
    },
    fleche/.style = {-{Stealth[length=5pt]}, thick, color=CouleurSecondaire}
]
\node[etape] (a) at (0,0)    {Besoins des\\ parties prenantes};
\node[etape] (b) at (3.4,0)  {13 objectifs\\ d'entreprise};
\node[etape] (c) at (6.8,0)  {13 objectifs\\ d'alignement};
\node[etape] (d) at (10.2,0) {40 objectifs\\ COBIT};
\draw[fleche] (a) -- (b);
\draw[fleche] (b) -- (c);
\draw[fleche] (c) -- (d);
\end{tikzpicture}
\end{center}
\vspace{-0.4em}
\begin{center}{\small\itshape Figure~6.1. La cascade des objectifs de
COBIT~2019}\end{center}
\end{minipage}

\begin{boxdef}
\termecle{Objectif d'entreprise}\index{objectif d'entreprise}
(\textit{enterprise goal}, EG)~: but de l'organisation dans son ensemble,
classé selon les quatre dimensions du tableau de bord prospectif
(finances, clients, processus internes, apprentissage et croissance).

\medskip
\termecle{Objectif d'alignement}\index{objectif d'alignement}
(\textit{alignment goal}, AG)~: but que l'informatique et les technologies
doivent atteindre pour soutenir les objectifs d'entreprise. COBIT~5 les
appelait « objectifs liés aux SI »~; le nouveau terme souligne qu'ils ne
sont pas des objectifs internes de la DSI.
\end{boxdef}

\subsection{Les vingt-six objectifs intermédiaires}

\begin{tableaucours}{L{1.0cm}L{5.0cm}L{1.0cm}L{5.0cm}}{%
    \tableauHeader{Code} &
    \tableauHeader{Objectif d'entreprise} &
    \tableauHeader{Code} &
    \tableauHeader{Objectif d'alignement}
}
EG01 & Portefeuille de produits et services compétitifs & AG01 &
Conformité des SI et soutien à la conformité aux lois \\
EG02 & Risque métier maîtrisé & AG02 & Risque lié aux SI maîtrisé \\
EG03 & Conformité aux lois et règlements & AG03 & Bénéfices réalisés des
investissements et services \\
EG04 & Qualité de l'information financière & AG04 & Qualité de
l'information financière liée aux technologies \\
EG05 & Culture de service orientée client & AG05 & Services conformes aux
exigences des métiers \\
EG06 & Continuité et disponibilité des services métier & AG06 & Agilité
pour transformer les exigences en solutions \\
EG07 & Qualité de l'information de gestion & AG07 & Sécurité de
l'information, des infrastructures et des applications, et vie privée \\
EG08 & Optimisation des fonctionnalités des processus internes & AG08 &
Soutien aux processus métier par l'intégration des applications \\
EG09 & Optimisation des coûts des processus & AG09 & Programmes livrés dans
les délais, les budgets et la qualité \\
EG10 & Compétences, motivation et productivité du personnel & AG10 &
Qualité de l'information de gestion des SI \\
EG11 & Conformité aux politiques internes & AG11 & Conformité des SI aux
politiques internes \\
EG12 & Programmes de transformation numérique maîtrisés & AG12 &
Personnel compétent, compréhension mutuelle des technologies et des
métiers \\
EG13 & Innovation des produits et des processus & AG13 & Connaissances et
initiatives pour l'innovation \\
\end{tableaucours}
\vspace{-0.8em}
\begin{center}{\small\itshape Tableau~6.6. Objectifs d'entreprise et
objectifs d'alignement de COBIT~2019 (intitulés abrégés et traduits par
nos soins)}\end{center}

\subsection{Lire la cascade dans les fiches}

Chaque fiche d'objectif indique les objectifs d'entreprise et
d'alignement auxquels il contribue \emph{à titre principal}. En les
relevant, on peut parcourir la cascade dans les deux sens. Le relevé
suivant est tiré des fiches officielles~\cite{isaca2018cobitobjectifs}.

\begin{tableaucours}{L{2.6cm}L{4.4cm}L{5.2cm}}{%
    \tableauHeader{Objectif} &
    \tableauHeader{Objectifs d'entreprise} &
    \tableauHeader{Objectifs d'alignement}
}
EDM03 & EG02, EG06 & AG02, AG07 \\
APO12 & EG02, EG06 & AG02, AG07 \\
DSS04 & EG01, EG02, EG06, EG08 & AG05, AG07 \\
DSS02 & EG01, EG08 & AG05 \\
BAI06, BAI07 & EG01 & AG06 \\
\end{tableaucours}
\vspace{-0.8em}
\begin{center}{\small\itshape Tableau~6.7. Liens principaux de quelques
objectifs (relevés dans COBIT~2019)}\end{center}

\begin{boiteRemarque}{Pourquoi la gestion des changements n'est-elle pas reliée au risque~?}
Le Tableau~6.7 réserve une surprise. On s'attendrait à ce que BAI06 et
BAI07, au cœur de toutes les défaillances du Chapitre~\ref{chap:defaillances},
soient rattachés en priorité au risque (EG02, AG02). Or leurs liens
principaux sont l'agilité (AG06) et la compétitivité de l'offre (EG01)~:
COBIT voit d'abord dans la gestion des changements le moyen de changer
\emph{vite et bien}.

La leçon est double. D'abord, la cascade ne relève que les liens
\emph{principaux}~; elle ne suffit pas, à elle seule, à prioriser. Ensuite,
une organisation qui a connu une défaillance due à un changement doit
remonter BAI06 et BAI07 par un autre chemin~: le profil de risque et les
problèmes liés aux SI, qui sont des facteurs de conception
(Chapitre~\ref{chap:cobit-conception}). C'est ce que fera la BANQUE
SANAGA.
\end{boiteRemarque}

% ============================================================================
\section{Objectifs Prioritaires : une Banque, une Administration}
\label{sec:c06-priorites}
% ============================================================================

\subsection{La BANQUE SANAGA}

\begin{boxscenario}{Opération SANAGA. Choisir dans la liste}
La responsable GRC part des objectifs d'entreprise que le conseil a
déclarés prioritaires après la panne~: EG06 (continuité et disponibilité),
EG02 (risque métier maîtrisé) et EG03 (conformité, au regard de la COBAC).

\medskip
\textbf{Par la cascade.} EG06 et EG02 conduisent, par les fiches
officielles, à EDM03, APO12 et DSS04 (Tableau~6.7). EG03 conduit à
l'objectif d'alignement AG01, et de là aux objectifs de conformité, dont
MEA03.

\medskip
\textbf{Par l'histoire de la banque.} La panne impose d'ajouter BAI06 et
BAI07 (changements), APO10 (fournisseurs), MEA02 (suivi des
recommandations) et EDM01 (cadre de gouvernance).

\medskip
\textbf{Résultat.} Dix objectifs prioritaires, et non trente-deux~: EDM01,
EDM03, APO10, APO12, BAI06, BAI07, DSS02, DSS04, MEA02, MEA03. Les autres
objectifs ne sont pas abandonnés~: ils seront visés à un niveau de
capacité plus modeste (Chapitre~\ref{chap:cobit-capacite}).
\end{boxscenario}

\subsection{L'ANPP}

Pour l'ANPP (opération MEFOU), le point de départ est différent. Les
objectifs d'entreprise prioritaires sont EG12 (programmes de
transformation numérique maîtrisés), EG09 (optimisation des coûts) et
EG03 (conformité, au regard des règles de la gestion publique et de la
Chambre des comptes). L'histoire de l'agence (Chapitres~\ref{chap:gouv-entreprise}
et~\ref{chap:alignement}) désigne d'elle-même les objectifs de
gouvernance du portefeuille (EDM02, APO05), de conduite des programmes et
des projets (BAI01, BAI11) et de gestion des fournisseurs (APO10). Le
choix définitif se fera avec les facteurs de conception, au
Chapitre~\ref{chap:cobit-conception}, puis avec les tables de
correspondance complètes de COBIT~2019 dans le laboratoire GCI-LAB-03.

\begin{boxmetier}{Auditeur SI}
\textbf{Répondre au « trente-deux cases rouges »}
\begin{enumerate}
    \item Une grille « conforme / non conforme » sur les 40~objectifs n'a
          pas de sens~: COBIT ne fixe pas d'exigences, mais des objectifs à
          atteindre à un niveau de capacité choisi.
    \item Demander d'abord quels objectifs ont été retenus comme
          prioritaires, et selon quelle démarche.
    \item Évaluer ensuite la capacité des seuls objectifs prioritaires, au
          regard de la cible fixée.
    \item Présenter au conseil l'écart entre capacité constatée et capacité
          cible, objectif par objectif.
\end{enumerate}
Mme~EKAMBI reprendra son tableau sur ces bases au Chapitre~\ref{chap:cobit-capacite}.
\end{boxmetier}

% ============================================================================
\separateur
% ============================================================================

\begin{boxresume}
\textbf{Ce qu'il faut retenir de ce chapitre~:}
\begin{itemize}
    \item \textbf{Une fiche COBIT décrit un morceau de système de
          gouvernance}~: description, finalité, objectifs liés, puis les
          sept composantes (processus, RACI, flux d'information,
          compétences, politiques, culture, outils).
    \item \textbf{Pratiques et activités}~: chaque activité est rattachée à
          un niveau de capacité, ce qui permet la mesure.
    \item \textbf{Quarante objectifs en cinq domaines}~: EDM pour le
          conseil~; APO pour l'organisation et la stratégie~; BAI pour tout
          ce qui change~; DSS pour l'exploitation~; MEA pour la surveillance
          et l'assurance.
    \item \textbf{La cascade des objectifs} relie les besoins des parties
          prenantes aux 13~objectifs d'entreprise, aux 13~objectifs
          d'alignement, puis aux 40~objectifs COBIT.
    \item \textbf{La cascade ne suffit pas}~: elle ne relève que les liens
          principaux~; l'histoire et le profil de risque de l'organisation
          complètent le choix.
    \item \textbf{Choisir, c'est retenir une dizaine d'objectifs
          prioritaires}, et viser pour les autres une capacité plus
          modeste.
\end{itemize}
\end{boxresume}

\bigskip

\begin{boxexo}[title={\ding{45}\quad Exercice~6.1. Trouver l'objectif \niveaubase}]
Pour chacune des situations suivantes, indiquez l'objectif COBIT~2019
concerné (code et intitulé).
\begin{enumerate}[(a)]
    \item Personne ne sait combien de licences de logiciels la banque
          possède.
    \item Le même incident de messagerie se reproduit chaque lundi.
    \item Le budget informatique est dépassé de 30~\% sans que la direction
          financière en soit informée.
    \item Les recommandations de l'audit interne restent sans suite.
    \item Les métiers se plaignent de ne jamais être consultés sur les
          projets informatiques.
\end{enumerate}
\end{boxexo}

\begin{boxexo}[title={\ding{45}\quad Exercice~6.2. La fiche processus de DSS04 \niveaumoyen}]
À l'aide du guide \textit{COBIT~2019 Governance and Management Objectives}
(dossier \texttt{Normes}), rédigez la fiche processus de l'objectif DSS04
(gérer la continuité) pour la BANQUE SANAGA, sur le modèle de la fiche
DSS02 de la section~\ref{sec:c06-lire}.
\begin{enumerate}[(a)]
    \item Résumez sa description et sa finalité, sans recopier le texte.
    \item Relevez ses objectifs d'entreprise et d'alignement liés.
    \item Listez ses pratiques et proposez une matrice RACI avec les
          fonctions réelles de la banque.
    \item Proposez trois indicateurs, avec leur cible et leur source.
\end{enumerate}
\end{boxexo}

\begin{boxexo}[title={\ding{45}\quad Exercice~6.3. Les priorités de NOUN PAY \niveauexpert}]
NOUN PAY (Bafoussam) fait face à deux enjeux~: l'enquête de l'autorité des
données personnelles sur sa notation de crédit par IA, et sa dépendance à
un fournisseur unique d'informatique en nuage.
\begin{enumerate}[(a)]
    \item Choisissez trois objectifs d'entreprise prioritaires, et
          justifiez.
    \item Descendez la cascade jusqu'aux objectifs d'alignement, puis aux
          objectifs COBIT, en vous appuyant sur les fiches officielles.
    \item Ajoutez les objectifs que la cascade ne fait pas apparaître mais
          que la situation impose, et expliquez pourquoi.
    \item Limitez votre liste finale à huit objectifs au plus, et
          justifiez chaque exclusion importante.
\end{enumerate}
\end{boxexo}

\bigskip

\begin{boxinfo}
\textbf{Transition.} La cascade des objectifs n'est qu'un des facteurs qui
orientent le choix des objectifs prioritaires. Le
Chapitre~\ref{chap:cobit-conception} présente les onze facteurs de
conception et la démarche complète de COBIT~2019 pour concevoir un
système de gouvernance sur mesure, avec les niveaux de capacité cibles.
\end{boxinfo}

        % ============================================================================
% Fichier  : partie2/P02_C07_cobit_conception.tex
% Titre    : Concevoir un Système de Gouvernance sur Mesure
% Auteur   : MINKA MI NGUIDJOÏ Thierry Emmanuel
% Version  : 1.0 -- octobre 2026
% ============================================================================
% Source : COBIT 2019 Framework: Introduction and Methodology, chapitres 4.5
% et 7 (dossier Normes). Le COBIT 2019 Design Guide (outil de calcul
% officiel) n'est pas dans le dossier : la méthode de calcul de la
% section 4 est une méthode pédagogique simplifiée, signalée comme telle.
% Laboratoires associés : GCI-LAB-02 (facteurs de conception) et GCI-LAB-03
% (cascade des objectifs).
% ============================================================================

\chapter{Concevoir un Système de Gouvernance sur Mesure}
\label{chap:cobit-conception}
\chaptermark{Concevoir un système sur mesure}

\epigraph{%
    « Il n'existe pas de formule magique. La conception finale est une
    décision au cas par cas. »%
}{D'après ISACA, \textit{COBIT 2019 Framework: Introduction and
Methodology} (2018, notre traduction)}

\niveauChapitre{Intermédiaire}{%
    Maîtrise du Chapitre~\ref{chap:cobit-objectifs} (les 40~objectifs et la
    cascade des objectifs).%
}

\begin{boxobjectifs}
\begin{itemize}
    \item Appliquer les onze facteurs de conception.
    \item Dérouler le processus de conception en quatre étapes.
    \item Calculer un profil d'objectifs prioritaires et en fixer les
          niveaux de capacité cibles.
    \item Justifier les arbitrages devant une direction générale.
\end{itemize}
\end{boxobjectifs}

\bigskip

% ============================================================================
\section*{Question directrice : dix objectifs, mais jusqu'où~?}
% ============================================================================

Au Chapitre~\ref{chap:cobit-objectifs}, la BANQUE SANAGA a retenu dix
objectifs prioritaires. Le directeur général, M.~FOTSO, pose aussitôt la
question suivante~: « Faut-il porter les dix au même niveau d'exigence~?
Notre budget de mise en œuvre ne le permettra pas. » À Bafoussam, la
fondatrice de NOUN PAY, Mme~KAMGA, pose la question inverse~: « Avec
quarante-cinq salariés, avons-nous seulement besoin de COBIT~? »

Les deux questions ont la même réponse~: COBIT~2019 ne s'applique pas
tel quel, il se \textbf{conçoit}. Les quarante objectifs sont, dans le
modèle, d'importance égale~; ce sont les caractéristiques de chaque
organisation, les \emph{facteurs de conception}, qui les hiérarchisent.
Ce chapitre présente leurs effets (section~\ref{sec:c07-effets}), les
onze facteurs (section~\ref{sec:c07-facteurs}), la démarche en quatre
étapes (section~\ref{sec:c07-etapes}), une méthode de calcul
(section~\ref{sec:c07-calcul}) et l'art de l'arbitrage
(section~\ref{sec:c07-arbitrage}). Les laboratoires GCI-LAB-02 et GCI-LAB-03 du
NDAP SUKULU en sont l'application.

\separateur

% ============================================================================
\section{Les Trois Effets des Facteurs de Conception}
\label{sec:c07-effets}
% ============================================================================

Le modèle central de COBIT~2019 traite ses quarante objectifs comme
intrinsèquement équivalents~: il n'existe entre eux aucun ordre de priorité
naturel. Les facteurs de conception rompent cette équivalence de trois
manières~\cite{isaca2018cobitintro}.

\begin{enumerate}
    \item \textbf{La priorité des objectifs et leur capacité cible.} Certains
          objectifs deviennent plus importants que d'autres, et quelques-uns
          peuvent même devenir négligeables. En pratique, l'importance se
          traduit par un \emph{niveau de capacité cible} plus élevé. Une
          organisation très averse au risque, par exemple, donnera plus de
          poids à EDM03, APO12, APO13 et DSS05.
    \item \textbf{La variation des composantes.} Certains facteurs changent
          l'importance d'une composante, ou en exigent une variante. Une
          petite organisation n'a pas besoin de toutes les structures du
          modèle~; une organisation très réglementée accorde plus
          d'importance aux procédures écrites et à la fonction conformité.
    \item \textbf{Le recours à un domaine d'intérêt.} Certains facteurs
          (menaces, méthodes de développement, taille) appellent un contenu
          adapté, décrit dans un domaine d'intérêt~: DevOps, petites et
          moyennes entreprises, cybersécurité.
\end{enumerate}

\begin{boxdef}
\termecle{Niveau de capacité cible}\index{niveau de capacité cible}
(\textit{target capability level})~: niveau de capacité, de~0 à~5, que
l'organisation décide d'atteindre pour un processus donné. Il traduit
l'importance de l'objectif pour elle. Le Chapitre~\ref{chap:cobit-capacite}
définit chaque niveau et la manière de mesurer la capacité constatée.
\end{boxdef}

% ============================================================================
\section{Les Onze Facteurs de Conception}
\label{sec:c07-facteurs}
% ============================================================================

\subsection{Le tableau des facteurs}

COBIT~2019 retient onze facteurs, dont les huit derniers sont parfois
appelés \emph{facteurs contextuels}. Pour chacun, il propose des valeurs
types~\cite{isaca2018cobitintro}. Le tableau suivant les applique à deux
organisations du cours.

\begin{tableaucours}{L{2.7cm}L{3.5cm}L{3.0cm}L{3.0cm}}{%
    \tableauHeader{Facteur} &
    \tableauHeader{Valeurs types} &
    \tableauHeader{BANQUE SANAGA} &
    \tableauHeader{ANPP}
}
1. Stratégie & Croissance, innovation, maîtrise des coûts, service et
stabilité & Service et stabilité (après la panne) & Service et stabilité \\
2. Objectifs d'entreprise & EG01 à EG13 & EG06, EG02, EG03 & EG12, EG09,
EG03 \\
3. Profil de risque & 19~catégories de risques liés aux SI & Incidents
d'exploitation, fournisseurs, non-conformité & Projets, décisions
d'investissement, fournisseurs \\
4. Problèmes liés aux SI & 20~problèmes types & Incidents majeurs,
défaillance du prestataire, constats d'audit répétés & Projets en retard,
doublons, dépenses hors budget \\
5. Paysage des menaces & Normal, élevé & Élevé (secteur bancaire) &
Normal \\
6. Exigences de conformité & Faibles, normales, élevées & Élevées (COBAC,
loi n°~2010/012) & Élevées (gestion publique) \\
7. Rôle des SI & Soutien, usine, redressement, stratégique & Stratégique
& Usine \\
8. Approvisionnement & Externalisé, nuage, interne, hybride & Hybride &
Externalisé \\
9. Méthodes de mise en œuvre & Agile, DevOps, traditionnelle, hybride &
Traditionnelle & Traditionnelle \\
10. Adoption des technologies & Pionnier, suiveur, tardif & Suiveur &
Tardif \\
11. Taille & Grande (plus de 250), PME (50 à 250) & Grande & Grande \\
\end{tableaucours}
\vspace{-0.8em}
\begin{center}{\small\itshape Tableau~7.1. Les onze facteurs de conception
appliqués à deux organisations (valeurs types d'après COBIT~2019, notre
traduction)}\end{center}

\subsection{Quatre facteurs à bien comprendre}

\paragraph{Le rôle des SI.} COBIT reprend une typologie classique en
quatre rôles. Les SI ont un rôle de \emph{soutien} quand ils ne sont
cruciaux ni pour le fonctionnement ni pour l'innovation~; d'\emph{usine}
quand leur panne arrête immédiatement l'activité, sans qu'ils soient un
moteur d'innovation~; de \emph{redressement} quand ils sont vus comme un
moteur d'innovation sans que l'activité courante en dépende encore~;
\emph{stratégiques} quand ils sont critiques à la fois pour fonctionner et
pour innover. L'ANPP est une usine~: une panne arrête les paiements, mais
l'agence n'innove pas par le numérique. La BANQUE SANAGA est stratégique.

\paragraph{Les problèmes liés aux SI.} Le profil de risque regarde vers
l'avenir~; les problèmes regardent ce qui s'est déjà produit. COBIT en
propose une liste type~: frustrations entre métiers et DSI, incidents
majeurs, problèmes de service des prestataires, constats d'audit répétés,
dépenses informatiques cachées, projets en retard, solutions développées
par les métiers sans la DSI, méconnaissance des règles de protection de la
vie privée. C'est par ce facteur que la panne de la BANQUE SANAGA remonte
BAI06 et BAI07 dans ses priorités (Chapitre~\ref{chap:cobit-objectifs}).

\paragraph{Les exigences de conformité.} Elles mesurent le poids des
obligations légales, réglementaires et contractuelles. Au Cameroun, elles
sont élevées pour les banques, les opérateurs de communications
électroniques et, depuis 2026, pour toute organisation qui traite des
données personnelles à grande échelle.

\paragraph{La taille.} COBIT distingue les grandes organisations (plus de
250~salariés à temps plein) et les petites et moyennes (de 50 à~250).
Les très petites organisations, de moins de 50~salariés, ne sont pas
couvertes par le modèle.

\subsection{L'ancrage juridique camerounais}

\begin{boxjuris}
\textbf{\loicam{n°~2010/012} du 21~décembre~2010 relative à la
cybersécurité et la cybercriminalité au Cameroun, article~31, alinéa~1}
(version anglaise officielle)~:

\medskip
\textit{« Electronic communication networks and information systems
content providers shall be bound to ensure the availability of material,
as well as the data stored in their installations. »}

\medskip
Soit~: les fournisseurs de contenus des réseaux de communications
électroniques et des systèmes d'information sont tenus d'assurer la
disponibilité des contenus et des données stockées dans leurs
installations.

\medskip
\textbf{\loicam{n°~2024/017} du 23~décembre~2024 relative à la protection
des données à caractère personnel.} D'après les présentations publiées
lors de son adoption~\cite{cio2026donnees}, cette loi s'applique
pleinement depuis le 23~juin~2026, au terme d'un délai de mise en
conformité de dix-huit mois. Elle institue une Autorité de protection des
données à caractère personnel et encadre strictement le transfert des
données hors du territoire national.

\medskip
\textbf{Portée pour la conception.} L'obligation de disponibilité de
l'article~31, l'obligation d'audit des articles~13 et~32
(Chapitre~\ref{chap:prologue}) et, depuis 2026, les obligations relatives
aux données personnelles justifient de classer les \emph{exigences de
conformité} en « élevées » pour la BANQUE SANAGA et pour NOUN PAY. Ce
classement relève en priorité DSS04 (continuité), MEA03 (conformité aux
exigences externes) et APO14 (données). Le Chapitre~\ref{chap:droit}
étudiera ces textes en détail.
\end{boxjuris}
% A VERIFIER au Journal officiel : articles de la loi n° 2024/017 (date
% d'application, autorité, transferts) ; sources utilisées : presse
% spécialisée (CIO Mag, Droit Médias Finance).

% ============================================================================
\section{La Démarche en Quatre Étapes}
\label{sec:c07-etapes}
% ============================================================================

COBIT~2019 propose une démarche de conception en quatre étapes, chacune
décomposée en sous-étapes~\cite{isaca2018cobitintro}.

\noindent\begin{minipage}{\linewidth}
\begin{center}
\begin{tikzpicture}[
    etape/.style = {
        draw=CouleurPrimaire, fill=CouleurPrimaire, rounded corners=4pt,
        text width=2.7cm, minimum height=1.0cm, align=center,
        font=\footnotesize\bfseries\color{white}
    },
    detail/.style = {
        draw=CouleurBordDef, fill=CouleurDef, rounded corners=3pt,
        text width=2.7cm, align=left, anchor=north,
        font=\scriptsize\color{CouleurPrimaire}
    },
    fleche/.style = {-{Stealth[length=5pt]}, thick, color=CouleurSecondaire}
]
\node[etape] (e1) at (0,0)    {1. Comprendre le contexte et la stratégie};
\node[etape] (e2) at (3.4,0)  {2. Déterminer le périmètre initial};
\node[etape] (e3) at (6.8,0)  {3. Affiner le périmètre};
\node[etape] (e4) at (10.2,0) {4. Conclure la conception};
\draw[fleche] (e1) -- (e2);
\draw[fleche] (e2) -- (e3);
\draw[fleche] (e3) -- (e4);
\node[detail] at (0,-0.75)    {Stratégie, objectifs d'entreprise, profil de
    risque, problèmes liés aux SI};
\node[detail] at (3.4,-0.75)  {Appliquer ces quatre facteurs, dont la
    cascade des objectifs};
\node[detail] at (6.8,-0.75)  {Menaces, conformité, rôle des SI,
    approvisionnement, méthodes, adoption, taille};
\node[detail] at (10.2,-0.75) {Résoudre les conflits de priorités,
    arrêter la conception};
\end{tikzpicture}
\end{center}
\vspace{-0.4em}
\begin{center}{\small\itshape Figure~7.1. La démarche de conception de
COBIT~2019 (d'après ISACA, notre traduction)}\end{center}
\end{minipage}

\medskip

Les deux premières étapes travaillent sur les quatre premiers facteurs~:
la première les \emph{comprend}, la deuxième les \emph{applique} pour
dégager un premier périmètre. La troisième étape affine ce périmètre avec
les sept facteurs contextuels. La quatrième résout les contradictions et
arrête la conception.

COBIT recommande de reporter toutes les indications obtenues sur un
\termecle{canevas de conception}\index{canevas de conception}
(\textit{design canvas}), puis de résoudre, autant que possible, les
conflits qui y apparaissent. Ces conflits sont inévitables~: avec onze
facteurs et des correspondances génériques, certaines indications se
contredisent.

% ============================================================================
\section{Calculer un Profil d'Objectifs}
\label{sec:c07-calcul}
% ============================================================================

\subsection{Une méthode pédagogique}

Le guide de conception de COBIT~2019 (\textit{Design Guide}) et son outil
de calcul attribuent, pour chaque facteur et chaque valeur, des
coefficients d'importance à chacun des quarante objectifs. Ces tables ne
sont pas reproduites ici. Ce cours utilise une \textbf{méthode
pédagogique simplifiée}, qui suit la même logique~: chaque facteur
contribue à l'importance de chaque objectif, et l'importance se traduit en
niveau de capacité cible. Le laboratoire GCI-LAB-02 applique cette méthode~;
en entreprise, on utilisera l'outil officiel.

\begin{boxalert}
\textbf{À ne pas confondre.} Les poids, les notes et les seuils de la
méthode ci-dessous sont des conventions du cours, fixées pour
l'apprentissage. Ce ne sont pas les coefficients officiels du
\textit{Design Guide}. Un rapport professionnel qui les présenterait comme
tels commettrait une erreur de source.
\end{boxalert}

\subsection{Les quatre temps du calcul}

\begin{enumerate}
    \item \textbf{Retenir les facteurs déterminants} et leur donner un poids
          (la somme des poids vaut 100~\%). Le choix des poids est une
          décision de gouvernance.
    \item \textbf{Noter, pour chaque objectif, la contribution de chaque
          facteur}, de~0 (aucune) à~3 (très forte), en justifiant chaque
          note.
    \item \textbf{Calculer le score d'importance} de l'objectif~:
          \[
              S = \frac{\sum_{f} w_f \, n_f}{3} \times 100,
          \]
          où $w_f$ est le poids du facteur et $n_f$ la note. Le score va
          de~0 à~100.
    \item \textbf{Traduire le score en capacité cible} selon la convention
          du cours~: 75 et plus, niveau~4~; de 50 à~74, niveau~3~; de 25
          à~49, niveau~2~; moins de~25, niveau~1.
\end{enumerate}

\subsection{Application à la BANQUE SANAGA}

\begin{boxscenario}{Opération SANAGA. Le profil de la banque}
Le comité stratégique des SI retient cinq facteurs et leurs poids~:
cascade des objectifs (25~\%), profil de risque (25~\%), problèmes liés aux
SI (20~\%), exigences de conformité (15~\%), approvisionnement (15~\%).

\medskip
\begin{tableaucours}{L{1.5cm}cccccrc}{%
    \tableauHeader{Objectif} &
    \tableauHeader{Casc.} &
    \tableauHeader{Risq.} &
    \tableauHeader{Prob.} &
    \tableauHeader{Conf.} &
    \tableauHeader{Appr.} &
    \tableauHeader{Score} &
    \tableauHeader{Cible}
}
DSS04 & 3 & 3 & 3 & 3 & 1 & 90,0 & 4 \\
EDM03 & 3 & 3 & 2 & 3 & 1 & 83,3 & 4 \\
APO12 & 3 & 3 & 2 & 2 & 1 & 78,3 & 4 \\
BAI06 & 1 & 3 & 3 & 2 & 2 & 73,3 & 3 \\
APO10 & 1 & 2 & 3 & 2 & 3 & 70,0 & 3 \\
BAI07 & 1 & 3 & 3 & 1 & 2 & 68,3 & 3 \\
MEA02 & 1 & 2 & 3 & 3 & 1 & 65,0 & 3 \\
APO04 & 1 & 0 & 0 & 0 & 0 & 8,3 & 1 \\
\end{tableaucours}

\medskip
\textbf{Calcul de DSS04.} $0{,}25 \times 3 + 0{,}25 \times 3 + 0{,}20
\times 3 + 0{,}15 \times 3 + 0{,}15 \times 1 = 2{,}70$~; $S = 2{,}70 / 3
\times 100 = 90{,}0$, d'où une cible de niveau~4.

\medskip
\textbf{Lecture.} Les notes de la colonne « cascade » reprennent les liens
du Tableau~6.7~: DSS04, EDM03 et APO12 sont liés à EG06 ou EG02, BAI06 et
BAI07 ne le sont pas, d'où leur note de~1. Ils remontent par le risque et
les problèmes. APO04 (innovation) tombe au niveau~1~: la banque n'en fait
pas une priorité pendant le redressement, ce qui ne veut pas dire qu'elle
renonce à innover.
\end{boxscenario}

\begin{boxprocedure}
\textbf{Le dossier de conception.} Pour qu'un auditeur, ou le régulateur,
puisse comprendre pourquoi un objectif a reçu telle capacité cible, le
dossier de conception doit contenir~:
\begin{itemize}
    \item la valeur retenue pour chacun des onze facteurs, avec sa
          justification~;
    \item les poids des facteurs et la décision qui les a approuvés~;
    \item les notes et les scores, avec la méthode de calcul et sa source
          (outil officiel ou méthode interne)~;
    \item le canevas de conception, avec les conflits relevés et la manière
          dont ils ont été résolus~;
    \item la liste finale des objectifs et de leurs capacités cibles,
          approuvée par le comité stratégique des SI.
\end{itemize}
\end{boxprocedure}

% ============================================================================
\section{Arbitrer et Justifier}
\label{sec:c07-arbitrage}
% ============================================================================

\subsection{Résoudre les conflits}

Le calcul propose~; l'organe de gouvernance décide. La quatrième étape de
la démarche consiste à confronter le résultat aux contraintes et aux
intentions de l'organisation.

\begin{boxscenario}{Opération SANAGA. Le conflit sur BAI06}
\textbf{Le conflit.} Le calcul donne une cible de niveau~3 pour BAI06. Le
président du comité d'audit demande le niveau~4~: « Après ce qui s'est
passé, la gestion des changements doit être exemplaire. »

\medskip
\textbf{Les arguments.} Le niveau~4 suppose un processus piloté par des
mesures quantitatives (Chapitre~\ref{chap:cobit-capacite}). La DSI estime
qu'il faut d'abord atteindre le niveau~3, c'est-à-dire un processus
défini et appliqué partout, ce qui n'est pas le cas aujourd'hui. Viser
directement le niveau~4 disperserait l'effort.

\medskip
\textbf{La décision} (procès-verbal du comité stratégique). Cible de
niveau~3 à dix-huit mois, puis réexamen pour le niveau~4. En contrepartie,
un indicateur trimestriel (part des changements majeurs dotés d'un
retour arrière testé) est présenté au comité d'audit dès le premier
trimestre.
\end{boxscenario}

\subsection{Trois types de conflits fréquents}

\begin{tableaucours}{L{3.4cm}L{4.4cm}L{4.4cm}}{%
    \tableauHeader{Conflit} &
    \tableauHeader{Exemple} &
    \tableauHeader{Piste de résolution}
}
Contrôle contre agilité & Une jeune entreprise en DevOps soumise à des
exigences de conformité élevées & Automatiser les contrôles dans la chaîne
de déploiement plutôt que d'ajouter des validations manuelles \\
Ambition contre moyens & Dix objectifs visés au niveau~4 avec une équipe
de cinq personnes & Échelonner les cibles dans le temps, concentrer
l'effort sur trois objectifs \\
Facteurs contradictoires & Stratégie d'innovation, mais profil de risque
très défavorable & Dissocier les périmètres~: innover sur un périmètre
isolé, stabiliser le cœur \\
\end{tableaucours}
\vspace{-0.8em}
\begin{center}{\small\itshape Tableau~7.2. Conflits de conception et
pistes de résolution}\end{center}

\begin{boiteRemarque}{NOUN PAY a-t-elle besoin de COBIT~?}
Avec 45~salariés, NOUN PAY se situe en dessous du seuil de 50~salariés à
partir duquel COBIT~2019 considère une organisation comme une PME~: les
très petites organisations ne sont pas couvertes par le modèle. Faut-il
pour autant renoncer~? Non, mais il faut l'utiliser comme une liste de
contrôle plutôt que comme un système complet~: retenir cinq ou six
objectifs, viser des capacités de niveau~2 ou~3, s'appuyer sur le domaine
d'intérêt consacré aux PME pour réduire les structures, et réserver
l'effort aux risques qui la menacent vraiment (dépendance au fournisseur
de nuage, données personnelles, notation par IA). C'est l'application
directe du principe~5~: adapter le système aux besoins de l'organisation.
\end{boiteRemarque}

\begin{boxmetier}{DSI}
\textbf{Présenter un profil de conception à la direction générale}
\begin{enumerate}
    \item Commencer par les facteurs, en langage métier~: « nous sommes une
          banque dont les SI sont stratégiques, sous forte contrainte
          réglementaire, qui sort d'une panne majeure ».
    \item Présenter la liste des objectifs prioritaires et leur capacité
          cible, sur une seule page.
    \item Montrer ce qui \emph{n'est pas} retenu, et pourquoi~: c'est ce qui
          rend le profil crédible.
    \item Chiffrer l'effort de mise en œuvre de chaque cible, et proposer un
          échelonnement.
    \item Demander une décision formelle, et la date de réexamen.
\end{enumerate}
\end{boxmetier}

% ============================================================================
\separateur
% ============================================================================

\begin{boxresume}
\textbf{Ce qu'il faut retenir de ce chapitre~:}
\begin{itemize}
    \item \textbf{COBIT se conçoit, il ne s'applique pas tel quel}~: les
          40~objectifs sont équivalents dans le modèle, les facteurs de
          conception les hiérarchisent.
    \item \textbf{Trois effets}~: priorité et capacité cible des objectifs,
          variation des composantes, recours à des domaines d'intérêt.
    \item \textbf{Onze facteurs}~: stratégie, objectifs d'entreprise, profil
          de risque, problèmes liés aux SI, menaces, conformité, rôle des
          SI, approvisionnement, méthodes, adoption des technologies,
          taille.
    \item \textbf{Quatre étapes}~: comprendre le contexte, déterminer le
          périmètre initial, l'affiner, conclure en résolvant les conflits
          sur un canevas de conception.
    \item \textbf{Le calcul propose, la gouvernance décide}~: la méthode de
          score aide à fixer les capacités cibles, mais les arbitrages
          relèvent du comité stratégique et se consignent au procès-verbal.
    \item \textbf{Au Cameroun}, les obligations de disponibilité et d'audit
          de la loi n°~2010/012 et la loi n°~2024/017 sur les données
          personnelles placent les exigences de conformité au niveau
          « élevé » pour les banques et les acteurs du paiement.
\end{itemize}
\end{boxresume}

\bigskip

\begin{boxexo}[title={\ding{45}\quad Exercice~7.1. Les facteurs de KRIBILOG \niveaubase}]
Renseignez, pour KRIBILOG SA (opérateur logistique portuaire, 1~200
salariés, exploitation confiée à l'infogérant NETPORT SERVICES), la valeur
de chacun des onze facteurs de conception, sur le modèle du Tableau~7.1.
Justifiez en une phrase les facteurs 4 (problèmes liés aux SI), 7 (rôle
des SI) et 8 (approvisionnement).
\end{boxexo}

\begin{boxexo}[title={\ding{45}\quad Exercice~7.2. Le profil de l'ANPP \niveaumoyen}]
Le comité stratégique de l'ANPP retient les poids suivants~: cascade
(30~\%), problèmes liés aux SI (30~\%), exigences de conformité (20~\%),
approvisionnement (20~\%). Les notes proposées sont~:
EDM02~(3,~3,~2,~1), APO05~(3,~3,~1,~1), APO10~(1,~3,~2,~3),
BAI11~(2,~3,~1,~2), DSS04~(1,~1,~2,~1).
\begin{enumerate}[(a)]
    \item Calculez le score d'importance de chaque objectif.
    \item Déduisez-en la capacité cible selon la convention du cours.
    \item Le directeur général demande que DSS04 soit porté au niveau~3.
          Quels arguments pourriez-vous avancer pour et contre~?
\end{enumerate}
\end{boxexo}

\begin{boxexo}[title={\ding{45}\quad Exercice~7.3. Concevoir le système de NOUN PAY \niveauexpert}]
NOUN PAY développe ses applications en DevOps, héberge tout son SI chez un
fournisseur de nuage étranger et doit se conformer, depuis juin~2026, à la
loi n°~2024/017 sur les données personnelles.
\begin{enumerate}[(a)]
    \item Renseignez les onze facteurs de conception, en signalant ceux pour
          lesquels COBIT n'est pas pleinement adapté à la taille de
          l'entreprise.
    \item Identifiez au moins deux conflits sur le canevas de conception et
          proposez leur résolution.
    \item Proposez une liste de six objectifs au plus, avec leur capacité
          cible et leur échéance.
    \item Rédigez la page de synthèse que Mme~KAMGA présentera à ses
          investisseurs.
\end{enumerate}
\end{boxexo}

\bigskip

\begin{boxinfo}
\textbf{Transition.} La conception a fixé des capacités cibles. Il reste à
savoir ce qu'elles signifient, comment mesurer la capacité réelle d'un
processus, et comment combler l'écart. C'est l'objet du
Chapitre~\ref{chap:cobit-capacite}~: le modèle de capacité, la maturité,
et la démarche de mise en œuvre en sept phases.
\end{boxinfo}

        % ============================================================================
% Fichier  : partie2/P02_C08_cobit_capacite_mise_en_oeuvre.tex
% Titre    : Mesurer et Mettre en Œuvre : Capacité, Maturité et Feuille de Route
% Auteur   : MINKA MI NGUIDJOÏ Thierry Emmanuel
% Version  : 1.0 -- octobre 2026
% ============================================================================
% Source : COBIT 2019 Framework: Introduction and Methodology, chapitres 6
% et 8 (dossier Normes), paraphrasés. Laboratoire associé : GCI-LAB-04
% (évaluation de capacité de DSS02 et BAI06).
% ============================================================================

\chapter{Mesurer et Mettre en Œuvre : Capacité, Maturité et Feuille de Route}
\label{chap:cobit-capacite}
\chaptermark{Capacité, maturité, mise en œuvre}

\epigraph{%
    « Ce qui ne se mesure pas ne s'améliore pas.\\
    Ce qui se mesure mal s'améliore dans la mauvaise direction. »%
}{Adage des praticiens de la qualité}

\niveauChapitre{Intermédiaire}{%
    Maîtrise du Chapitre~\ref{chap:cobit-conception} (facteurs de
    conception et niveaux de capacité cibles).%
}

\begin{boxobjectifs}
\begin{itemize}
    \item Appliquer le modèle de capacité des processus (niveaux 0 à 5).
    \item Distinguer capacité d'un processus et maturité d'un domaine.
    \item Évaluer un niveau de capacité à partir de pièces justificatives.
    \item Dérouler le cycle de mise en œuvre en sept phases.
    \item Construire une feuille de route d'amélioration.
\end{itemize}
\end{boxobjectifs}

\bigskip

% ============================================================================
\section*{Question directrice : où en sommes-nous vraiment~?}
% ============================================================================

Au Chapitre~\ref{chap:cobit-objectifs}, Mme~EKAMBI, directrice de l'audit
interne de la BANQUE SANAGA, avait présenté au conseil un tableau de
quarante cases « conforme / non conforme ». Au Chapitre~\ref{chap:cobit-conception},
le comité stratégique a fixé, pour dix objectifs prioritaires, une
capacité cible. Il reste à répondre à deux questions~: \emph{quelle est la
capacité réelle de chaque processus aujourd'hui~?} et \emph{comment combler
l'écart~?}

La première question relève de la \textbf{mesure de la performance}
(sections~\ref{sec:c08-principes} à~\ref{sec:c08-evaluer}). La seconde
relève de la \textbf{mise en œuvre} (sections~\ref{sec:c08-phases}
et~\ref{sec:c08-feuille}). Le laboratoire GCI-LAB-04 de NDAP SUKULU fait
évaluer la capacité de deux processus de la banque à partir d'un dossier
de pièces.

\separateur

% ============================================================================
\section{Mesurer la Performance dans COBIT}
\label{sec:c08-principes}
% ============================================================================

\subsection{Cinq principes}

COBIT~2019 désigne par \emph{gestion de la performance} l'ensemble des
méthodes qui disent dans quelle mesure le système de gouvernance et ses
composantes fonctionnent, et comment les améliorer. Il fixe cinq
principes à cette gestion~\cite{isaca2018cobitintro}~: elle doit être
\textbf{simple} à comprendre et à utiliser~; \textbf{cohérente} avec le
modèle conceptuel, et applicable à tous les types de composantes~;
\textbf{fiable}, c'est-à-dire produire des résultats reproductibles et
pertinents~; \textbf{flexible}, pour s'adapter aux priorités de chaque
organisation~; et capable de soutenir \textbf{différents types
d'évaluation}, de l'autoévaluation à l'audit formel.

\subsection{Un modèle aligné sur CMMI}

Le modèle de COBIT~2019 s'aligne largement sur la version~2.0 du CMMI et
l'étend. Il repose sur deux notions~:
\begin{itemize}
    \item la \textbf{capacité}, qui se mesure \emph{processus par
          processus}~: chaque activité d'un processus est rattachée à un
          niveau de capacité~;
    \item la \textbf{maturité}, qui se mesure \emph{domaine d'intérêt par
          domaine d'intérêt}~: elle est atteinte quand tous les processus du
          domaine atteignent le niveau de capacité requis.
\end{itemize}

\begin{boiteRemarque}{Et le modèle de COBIT~5~?}
COBIT~5 mesurait la capacité selon la norme ISO/IEC~15504, devenue
ISO/IEC~33000, dont les niveaux ont un sens assez différent. COBIT~2019
lui substitue un modèle inspiré du CMMI. Les organisations qui veulent
conserver l'approche ISO/IEC~33000 le peuvent~: COBIT~2019 précise que ses
pratiques tiennent lieu de « pratiques de base » et ses flux d'information
de « produits de travail ». Il faut seulement ne pas mélanger les deux
échelles dans un même rapport.
\end{boiteRemarque}

% ============================================================================
\section{Les Niveaux de Capacité}
\label{sec:c08-niveaux}
% ============================================================================

\subsection{Six niveaux}

\begin{boxdef}
\termecle{Capacité d'un processus}\index{capacité d'un processus}
(\textit{process capability})~: mesure de la manière dont un processus est
mis en œuvre et fonctionne. COBIT~2019 la mesure sur une échelle de~0
à~5.
\end{boxdef}

\begin{tableaucours}{cL{3.2cm}L{8.0cm}}{%
    \tableauHeader{Niv.} &
    \tableauHeader{Caractère} &
    \tableauHeader{Description}
}
0 & Absence & Aucune capacité de base~; approche incomplète~; les pratiques
ne sont pas, ou pas vraiment, mises en œuvre. \\
1 & Initial, intuitif & Le processus atteint plus ou moins sa finalité, par
un ensemble incomplet d'activités, peu organisé. \\
2 & Réalisé & Le processus atteint sa finalité par un ensemble d'activités
élémentaire, mais complet. \\
3 & Défini & Le processus atteint sa finalité de manière beaucoup plus
organisée, en s'appuyant sur des ressources de l'organisation~; il est
généralement bien défini. \\
4 & Mesuré & Le processus atteint sa finalité, il est bien défini, et sa
performance est mesurée de manière quantitative. \\
5 & Optimisé & Le processus est bien défini, sa performance est mesurée
pour l'améliorer, et l'amélioration continue est recherchée. \\
\end{tableaucours}
\vspace{-0.8em}
\begin{center}{\small\itshape Tableau~8.1. Les niveaux de capacité des
processus (d'après COBIT~2019, notre traduction)}\end{center}

\subsection{Des activités rattachées à des niveaux}

Le guide des objectifs rattache chaque activité à un niveau de capacité.
Dans la fiche DSS02, par exemple, les activités qui consistent à
enregistrer, classer et prioriser les demandes et les incidents
(pratique DSS02.02) relèvent du niveau~2~; celles qui définissent les
schémas de classification et les règles d'escalade (DSS02.01) relèvent du
niveau~3~; l'analyse des tendances et son utilisation pour l'amélioration
continue (DSS02.07) relèvent des niveaux~4 et~5~\cite{isaca2018cobitobjectifs}.

La progression est logique~: on commence par \emph{faire} (niveau~2),
puis on \emph{standardise} (niveau~3), puis on \emph{mesure} (niveau~4),
enfin on \emph{améliore} en continu (niveau~5).

\begin{boxalert}
\textbf{Violation typique.} Une organisation se déclare « au niveau~4 »
parce qu'elle produit un tableau de bord mensuel sur ses incidents. Mais
les incidents ne sont pas tous enregistrés, et chaque équipe les classe à
sa façon~: le tableau de bord mesure un processus qui n'est ni complet
(niveau~2) ni défini (niveau~3).

\smallskip
\textbf{Règle~:} un niveau ne s'acquiert que si les niveaux inférieurs
sont atteints. Mesurer un processus qui n'est pas défini produit des
chiffres, pas une capacité.
\end{boxalert}

% ============================================================================
\section{Capacité et Maturité}
\label{sec:c08-maturite}
% ============================================================================

\begin{boxdef}
\termecle{Maturité}\index{maturité} (\textit{maturity})~: mesure de
performance au niveau d'un domaine d'intérêt (un ensemble d'objectifs et
de leurs composantes). Un domaine d'intérêt atteint un niveau de maturité
donné lorsque \emph{tous} les processus qu'il contient atteignent ce
niveau de capacité.
\end{boxdef}

\begin{tableaucours}{cL{3.0cm}L{8.2cm}}{%
    \tableauHeader{Niv.} &
    \tableauHeader{Maturité} &
    \tableauHeader{Description}
}
0 & Incomplète & Le travail n'est pas forcément accompli. \\
1 & Initiale & Le travail est accompli, mais l'objectif du domaine n'est pas
pleinement atteint. \\
2 & Gérée & Planification et mesure de la performance existent, mais pas
encore de manière standardisée. \\
3 & Définie & Des standards communs à toute l'organisation guident le
travail. \\
4 & Quantitative & L'organisation est pilotée par les données, avec une
amélioration quantitative de la performance. \\
5 & Optimisée & L'organisation est tournée vers l'amélioration continue. \\
\end{tableaucours}
\vspace{-0.8em}
\begin{center}{\small\itshape Tableau~8.2. Les niveaux de maturité des
domaines d'intérêt (d'après COBIT~2019, notre traduction)}\end{center}

La règle d'agrégation est exigeante~: la maturité d'un domaine est celle
de son processus le plus faible. Une banque dont neuf processus de
sécurité sont au niveau~3 et le dixième au niveau~1 a une maturité de
niveau~1 dans le domaine de la sécurité. C'est voulu~: une chaîne a la
solidité de son maillon le plus faible.

\subsection{Mesurer les autres composantes}

COBIT~2019 reconnaît qu'il n'existe pas de méthode généralement admise
pour évaluer les composantes autres que les processus, mais il propose des
critères. Pour une structure organisationnelle, par exemple~: elle est
formellement créée~; son mandat est clair et documenté~; ses règles de
fonctionnement sont écrites~; elle se réunit comme prévu~; ses
procès-verbaux existent et sont utiles~; ses droits de décision sont
définis et respectés~; ses procédures d'escalade sont appliquées. Ces
critères rejoignent l'exigence de contrôle « un comité qui fonctionne » du
Chapitre~\ref{chap:gouv-entreprise}.

% ============================================================================
\section{Évaluer une Capacité}
\label{sec:c08-evaluer}
% ============================================================================

\subsection{L'échelle de notation}

Un niveau de capacité peut être atteint plus ou moins complètement.
COBIT~2019 propose, pour les évaluations orientées vers l'amélioration,
l'échelle suivante~\cite{isaca2018cobitintro}.

\begin{tableaucours}{L{3.6cm}L{3.6cm}L{5.0cm}}{%
    \tableauHeader{Note} &
    \tableauHeader{Taux de réalisation} &
    \tableauHeader{Lecture}
}
Complètement (\textit{Fully}, F) & Plus de 85~\% & Le niveau est acquis \\
Largement (\textit{Largely}, L) & De 50 à 85~\% & Acquis, avec des
faiblesses \\
Partiellement (\textit{Partially}, P) & De 15 à 50~\% & Non acquis \\
Pas (\textit{Not}, N) & Moins de 15~\% & Non acquis \\
\end{tableaucours}
\vspace{-0.8em}
\begin{center}{\small\itshape Tableau~8.3. L'échelle de notation des
niveaux de capacité (d'après COBIT~2019)}\end{center}

COBIT~2019 souligne que la note reste un jugement, qui doit être étayé par
l'examen des composantes~: activités réalisées, pièces produites,
structures en fonctionnement. Les méthodes formelles, qui débouchent sur
une certification indépendante, peuvent préférer une notation binaire
(atteint ou non).

\subsection{La règle d'attribution du niveau}

Le cours retient la règle usuelle des évaluations de processus, issue de
la famille ISO/IEC~33000~: \textbf{un processus atteint le niveau~$n$ si
ce niveau est noté au moins « largement » et si tous les niveaux
inférieurs sont notés « complètement »}. Un niveau noté « largement » est
donc acquis, mais il bloque l'accès au niveau suivant tant qu'il n'est pas
porté à « complètement ».

\subsection{Une évaluation pas à pas}

\begin{boxscenario}{Opération SANAGA. La capacité de BAI06}
L'auditrice interne, Mme~EKAMBI, évalue le processus BAI06 (gérer les
changements informatiques), dont la cible est le niveau~3. Elle examine,
pour chaque activité, les pièces de vingt changements tirés au hasard sur
six mois. \textit{(Le nombre d'activités par niveau est simplifié pour
l'exercice~; la fiche officielle BAI06 en donne la liste exacte.)}

\medskip
\begin{tableaucours}{cccccc}{%
    \tableauHeader{Niveau} &
    \tableauHeader{Activités} &
    \tableauHeader{Réalisées} &
    \tableauHeader{Taux} &
    \tableauHeader{Note} &
    \tableauHeader{Acquis~?}
}
2 & 10 & 6 & 60~\% & L & Oui \\
3 & 8 & 2 & 25~\% & P & Non \\
4 & 4 & 0 & 0~\% & N & Non \\
\end{tableaucours}

\medskip
\textbf{Constatation.} Six activités de niveau~2 sur dix sont réalisées~:
les demandes de changement sont enregistrées et évaluées, mais les
changements d'urgence ne font l'objet d'aucune procédure, et
l'autorisation n'est pas tracée pour les changements jugés mineurs. Deux
activités de niveau~3 sur huit sont réalisées.

\medskip
\textbf{Analyse} (confiance élevée). Capacité constatée~: \textbf{niveau~2,
largement atteint}. L'écart avec la cible est d'un niveau. Pour progresser,
il faut d'abord porter le niveau~2 à « complètement » (procédure
d'urgence, traçabilité de toutes les autorisations), puis réaliser les
activités de niveau~3.
\end{boxscenario}

\begin{boxprocedure}
\textbf{Une évaluation de capacité opposable.} Pour qu'une évaluation de
capacité puisse être présentée au conseil ou au régulateur, le dossier
doit contenir~:
\begin{itemize}
    \item l'objectif évalué, l'édition de COBIT et l'échelle de notation
          utilisées~;
    \item le périmètre (entités, systèmes, période) et la méthode
          d'échantillonnage~;
    \item pour chaque activité, la pièce examinée et le constat
          (réalisée ou non)~;
    \item le taux et la note de chaque niveau, et la règle d'attribution
          appliquée~;
    \item la capacité constatée, la capacité cible et l'écart~;
    \item la date, le nom de l'évaluateur et son indépendance par rapport
          au processus évalué.
\end{itemize}
\end{boxprocedure}

\begin{boxmetier}{Auditeur SI}
\textbf{Conduire une évaluation de capacité}
\begin{enumerate}
    \item Partir de la fiche officielle de l'objectif~: lister ses
          activités, niveau par niveau.
    \item Pour chaque activité, définir la pièce qui prouve sa réalisation
          (procès-verbal, ticket, journal, rapport).
    \item Échantillonner~: une activité réalisée une fois n'est pas une
          activité réalisée.
    \item Noter chaque niveau sur l'échelle F, L, P, N, puis appliquer la
          règle d'attribution.
    \item Restituer à l'audité avant de conclure~: il peut détenir une pièce
          que vous n'avez pas vue.
\end{enumerate}
\end{boxmetier}

% ============================================================================
\section{La Mise en Œuvre en Sept Phases}
\label{sec:c08-phases}
% ============================================================================

\subsection{Un programme, pas un projet}

Le guide de mise en œuvre de COBIT~2019 part d'un constat~: beaucoup de
démarches de gouvernance échouent parce qu'elles ne sont pas lancées et
pilotées comme de véritables \emph{programmes}, parrainés par la direction
générale, au périmètre maîtrisé et aux objectifs atteignables. Le but final
n'est pas de terminer un projet, mais d'installer une pratique durable~:
le programme se clôt quand l'amélioration de la gouvernance produit des
bénéfices mesurables et fait partie du fonctionnement normal.

\subsection{Sept phases, trois anneaux}

La démarche comporte sept phases, chacune formulée comme une question.
Chaque phase se lit sur trois anneaux~: le \emph{management du programme}
(anneau extérieur), l'\emph{accompagnement du changement} (anneau
intermédiaire) et le \emph{cycle d'amélioration continue} (anneau
intérieur)~\cite{isaca2018cobitintro}.

\begin{tableaucours}{L{3.3cm}L{2.9cm}L{2.9cm}L{3.1cm}}{%
    \tableauHeader{Phase} &
    \tableauHeader{Programme} &
    \tableauHeader{Changement} &
    \tableauHeader{Amélioration}
}
1. Quels sont les moteurs~? & Lancer le programme & Créer le désir de
changer & Reconnaître la nécessité d'agir \\
2. Où en sommes-nous~? & Définir problèmes et opportunités & Former
l'équipe de mise en œuvre & Évaluer l'état actuel \\
3. Où voulons-nous aller~? & Définir la feuille de route & Communiquer le
résultat visé & Définir l'état cible \\
4. Que faut-il faire~? & Planifier le programme & Identifier les acteurs &
Construire les améliorations \\
5. Comment y arriver~? & Exécuter le plan & Mettre en service et utiliser
& Mettre en œuvre les améliorations \\
6. Y sommes-nous arrivés~? & Réaliser les bénéfices & Ancrer les nouvelles
pratiques & Exploiter et mesurer \\
7. Comment garder l'élan~? & Examiner l'efficacité & Pérenniser &
Surveiller et évaluer \\
\end{tableaucours}
\vspace{-0.8em}
\begin{center}{\small\itshape Tableau~8.4. La feuille de route de mise en
œuvre de COBIT~2019 (d'après ISACA, notre traduction)}\end{center}

\subsection{Le lien avec la conception}

Les deux démarches s'emboîtent. La conception du système de gouvernance
(Chapitre~\ref{chap:cobit-conception}) nourrit les trois premières phases
de la mise en œuvre~: la compréhension du contexte alimente la phase~1,
la détermination et l'affinage du périmètre alimentent la phase~2, et la
conclusion de la conception fixe l'état cible de la phase~3. Les
évaluations de capacité, enfin, mesurent l'état actuel en phase~2 et les
progrès en phase~6.

% ============================================================================
\section{Construire une Feuille de Route}
\label{sec:c08-feuille}
% ============================================================================

\subsection{De l'écart au plan}

La phase~3 compare l'état cible et l'état actuel, puis cherche des
solutions. COBIT recommande de distinguer les \emph{gains rapides} des
chantiers de long terme, de donner la priorité aux actions faciles et à
fort bénéfice, et de découper les chantiers longs en étapes maîtrisables.

\begin{boxscenario}{Opération SANAGA. La feuille de route à dix-huit mois}
Mme~EKAMBI reprend son tableau sur ces bases. Il ne compte plus quarante
cases rouges, mais une ligne par objectif prioritaire. Voici les sept
objectifs du profil de conception (Chapitre~\ref{chap:cobit-conception})~;
EDM01, DSS02 et MEA03 suivent le même format.

\medskip
\begin{tableaucours}{L{1.5cm}cccL{6.2cm}}{%
    \tableauHeader{Objectif} &
    \tableauHeader{Constatée} &
    \tableauHeader{Cible} &
    \tableauHeader{Écart} &
    \tableauHeader{Première action}
}
DSS04 & 1 & 4 & 3 & Tester le plan de continuité du système central
(gain rapide~: 3~mois) \\
EDM03 & 1 & 4 & 3 & Faire approuver l'appétence au risque par le conseil \\
APO12 & 2 & 4 & 2 & Registre des risques SI relié au registre de la banque \\
BAI06 & 2 & 3 & 1 & Procédure des changements d'urgence (gain rapide) \\
APO10 & 1 & 3 & 2 & Renégocier le contrat de DIGICORE AFRIQUE \\
BAI07 & 1 & 3 & 2 & Recette et retour arrière obligatoires avant bascule \\
MEA02 & 2 & 3 & 1 & Suivi trimestriel des recommandations au comité
d'audit \\
\end{tableaucours}

\medskip
\textbf{Échelonnement.} Mois 1 à 3~: gains rapides (DSS04, BAI06, MEA02).
Mois 4 à 12~: chantiers de fond (EDM03, APO12, APO10, BAI07). Mois 13 à
18~: mesure des progrès (phase~6) et nouvelle évaluation de capacité.
\end{boxscenario}

\begin{boxmetier}{Responsable GRC}
\textbf{Piloter la feuille de route}
\begin{enumerate}
    \item Donner à chaque ligne un propriétaire, une échéance et un
          indicateur de progrès.
    \item Rendre compte chaque trimestre au comité stratégique des SI, et
          chaque semestre au conseil.
    \item Réévaluer la capacité des objectifs traités à mi-parcours, par une
          personne indépendante du processus.
    \item Mettre à jour la conception (Chapitre~\ref{chap:cobit-conception})
          si un facteur change~: c'est le principe du système dynamique.
\end{enumerate}
\end{boxmetier}

\begin{boiteRemarque}{Viser le niveau~5 partout~?}
Rien ne serait plus coûteux, ni plus inutile. Le niveau~5 suppose une
amélioration continue fondée sur des mesures quantitatives~: il se
justifie pour quelques processus réellement critiques, rarement plus. Pour
la plupart des organisations camerounaises, atteindre durablement le
niveau~3 sur les objectifs prioritaires serait déjà un progrès décisif~:
des processus définis, appliqués partout, qui ne dépendent plus de la
bonne volonté d'une personne. La gouvernance, ici encore, est affaire de
proportion.
\end{boiteRemarque}

% ============================================================================
\separateur
% ============================================================================

\begin{boxresume}
\textbf{Ce qu'il faut retenir de ce chapitre~:}
\begin{itemize}
    \item \textbf{Capacité et maturité}~: la capacité se mesure processus
          par processus (0 à 5)~; la maturité se mesure par domaine
          d'intérêt, et vaut celle du processus le plus faible.
    \item \textbf{Six niveaux}~: absence, initial, réalisé, défini, mesuré,
          optimisé. On fait, on standardise, on mesure, on améliore.
    \item \textbf{Des activités rattachées à des niveaux}~: c'est ce qui rend
          la mesure objective, à condition de l'étayer par des pièces.
    \item \textbf{Quatre notes}~: complètement (plus de 85~\%), largement
          (50 à 85~\%), partiellement (15 à 50~\%), pas (moins de 15~\%).
          Un niveau ne s'acquiert que si les niveaux inférieurs sont
          complètement atteints.
    \item \textbf{Sept phases, trois anneaux}~: programme, changement,
          amélioration continue. La gouvernance s'installe par un
          programme, pas par un projet.
    \item \textbf{La feuille de route} part des écarts entre capacité
          constatée et cible, commence par des gains rapides, et se pilote
          devant le comité stratégique.
\end{itemize}
\end{boxresume}

\bigskip

\begin{boxexo}[title={\ding{45}\quad Exercice~8.1. Noter un niveau \niveaubase}]
Pour chacun des cas suivants, donnez la note (F, L, P ou N) du niveau
concerné, puis la capacité du processus.
\begin{enumerate}[(a)]
    \item Niveau~2~: 9 activités sur 10 réalisées~; niveau~3~: 3 sur 6.
    \item Niveau~2~: 7 activités sur 10 réalisées~; niveau~3~: 6 sur 6.
    \item Niveau~2~: 1 activité sur 10 réalisée.
    \item Niveau~2~: 12 sur 12~; niveau~3~: 8 sur 9~; niveau~4~: 1 sur 5.
\end{enumerate}
\end{boxexo}

\begin{boxexo}[title={\ding{45}\quad Exercice~8.2. La maturité de la continuité \niveaumoyen}]
Le domaine d'intérêt « continuité » de KRIBILOG SA regroupe quatre
processus. Leurs capacités constatées sont~: DSS04 niveau~2, BAI04
niveau~3, DSS01 niveau~3, APO10 niveau~1.
\begin{enumerate}[(a)]
    \item Quelle est la maturité du domaine~? Justifiez.
    \item Le directeur général affirme que « la continuité est au niveau~3
          en moyenne ». Que lui répondez-vous~?
    \item Quelle action aurait le plus d'effet sur la maturité du domaine~?
\end{enumerate}
\end{boxexo}

\begin{boxexo}[title={\ding{45}\quad Exercice~8.3. La feuille de route de l'ANPP \niveauexpert}]
À partir des capacités cibles calculées à l'Exercice~7.2, et des capacités
constatées suivantes~: EDM02 niveau~0, APO05 niveau~1, APO10 niveau~1,
BAI11 niveau~2, DSS04 niveau~1.
\begin{enumerate}[(a)]
    \item Construisez le tableau des écarts.
    \item Identifiez deux gains rapides et deux chantiers de long terme.
    \item Organisez la feuille de route selon les sept phases, en précisant
          pour chaque phase une action sur chacun des trois anneaux.
    \item Indiquez comment la feuille de route s'inscrit dans le budget
          programme de l'agence (Chapitre~\ref{chap:alignement}).
\end{enumerate}
\end{boxexo}

\bigskip

\begin{boxinfo}
\textbf{Transition.} La Partie~\ref{part:cobit} a construit, avec COBIT~2019,
un système de gouvernance mesuré et une feuille de route. La
Partie~\ref{part:normatif} replace ce système dans son cadre normatif et
juridique. Le Chapitre~\ref{chap:normes} présente les normes
internationales de la gouvernance et des services (ISO/IEC~38500,
ISO/IEC~20000-1, ISO/IEC~27001, ISO~31000, COSO).
\end{boxinfo}


    % ======================================================================
    % PARTIE III — CADRE NORMATIF, JURIDIQUE ET DÉONTOLOGIQUE
    % « Ce qui n'est pas prouvé n'est pas maîtrisé. »
    % Standards : ISO/IEC 38500, ISO/IEC 20000-1, ISO/IEC 27001 et 27002,
    % ISO 31000, COSO ERM, normes IIA, INTOSAI GUID 5100 ; lois camerounaises
    % n° 2010/012 et 2010/013, loi de 2024 sur les données personnelles,
    % OHADA, COBAC, Convention de Malabo.
    % Niveau : fondamental
    % ======================================================================
    \part{Cadre Normatif, Juridique et Déontologique}
    \label{part:normatif}

        % ============================================================================
% Fichier  : partie3/P03_C09_normes_internationales.tex
% Titre    : Les Normes Internationales de la Gouvernance et des Services
% Auteur   : MINKA MI NGUIDJOÏ Thierry Emmanuel
% Version  : 1.0 -- octobre 2026
% ============================================================================
% Sources : normes du dossier Normes (ISO/IEC 27000:2018, 27001:2013,
% 27002:2013, 27006 ; résumé COSO ERM 2004 en français) ; éditions en
% vigueur citées d'après les catalogues ISO. Principes de l'ISO/IEC
% 38500:2024 repris d'ISO 37000 : liste à vérifier dans le texte de la
% norme.
% ============================================================================

\chapter{Les Normes Internationales de la Gouvernance et des Services}
\label{chap:normes}
\chaptermark{Normes internationales}

\epigraph{%
    « Une norme ne dit pas ce qu'il faut penser.\\
    Elle dit ce qu'il faudra pouvoir démontrer. »%
}{Adage des auditeurs de systèmes de management}

\niveauChapitre{Fondamental}{%
    Maîtrise du Chapitre~\ref{chap:referentiels} (natures des référentiels
    et carte d'articulation).%
}

\begin{boxobjectifs}
\begin{itemize}
    \item Présenter les principes et le modèle « évaluer, diriger,
          surveiller » de l'ISO/IEC~38500.
    \item Décrire les exigences d'un système de management des services
          (ISO/IEC~20000-1).
    \item Articuler la gouvernance des SI avec l'ISO/IEC~27001, l'ISO~31000
          et COSO ERM.
    \item Distinguer certification, conformité et alignement volontaire.
\end{itemize}
\end{boxobjectifs}

\bigskip

% ============================================================================
\section*{Question directrice : faut-il un certificat~?}
% ============================================================================

KRIBILOG SA vient de recevoir l'appel d'offres d'un grand armateur pour
la gestion de ses conteneurs au port de Kribi. Une clause l'arrête~: « Le
prestataire devra justifier, dans un délai de dix-huit mois, d'une
certification ISO/IEC~20000-1 sur le périmètre du service. » Au même
moment, le conseil de la BANQUE SANAGA demande s'il ne faudrait pas
« se faire certifier ISO~27001 pour rassurer la COBAC ».

La Partie~\ref{part:normatif}, qui s'ouvre ici, place la gouvernance des SI
dans son cadre normatif, puis juridique (Chapitre~\ref{chap:droit}), puis
dans celui de l'audit (Chapitre~\ref{chap:audit}). Ce chapitre présente les
normes internationales qui comptent pour la gouvernance des SI~: leur
architecture commune (section~\ref{sec:c09-structure}), la norme de la
gouvernance (section~\ref{sec:c09-38500}), celles des services et de la
sécurité (sections~\ref{sec:c09-20000} et~\ref{sec:c09-27001}), celles du
risque (section~\ref{sec:c09-risque}), et enfin la question du certificat
(section~\ref{sec:c09-certification}).

\separateur

% ============================================================================
\section{L'Architecture des Systèmes de Management}
\label{sec:c09-structure}
% ============================================================================

\subsection{Qu'est-ce qu'un système de management~?}

\begin{boxdef}
\termecle{Système de management}\index{système de management}
(\textit{management system})~: ensemble d'éléments corrélés ou en
interaction d'une organisation (politiques, objectifs, rôles, processus)
qui sert à établir des politiques et des objectifs, et à atteindre ces
objectifs. Un système de management de la qualité, de la sécurité de
l'information ou des services en sont des exemples.
\end{boxdef}

\subsection{Une structure harmonisée}

Les normes ISO de systèmes de management certifiables partagent une même
structure, dite \termecle{structure harmonisée}\index{structure
harmonisée} (anciennement « structure de haut niveau »). Elle facilite la
mise en place de systèmes intégrés~: une organisation certifiée à la fois
ISO/IEC~27001 et ISO/IEC~20000-1 n'a qu'un seul cycle de revue de
direction, d'audit interne et d'amélioration.

\begin{tableaucours}{cL{3.6cm}L{8.0cm}}{%
    \tableauHeader{Art.} &
    \tableauHeader{Titre} &
    \tableauHeader{Contenu}
}
4 & Contexte de l'organisation & Enjeux internes et externes, parties
intéressées, périmètre du système \\
5 & Leadership & Engagement de la direction, politique, rôles et
responsabilités \\
6 & Planification & Risques et opportunités, objectifs et plans pour les
atteindre \\
7 & Support & Ressources, compétences, sensibilisation, communication,
informations documentées \\
8 & Fonctionnement & Planification et maîtrise opérationnelles~: c'est là
que chaque norme place ses exigences propres \\
9 & Évaluation des performances & Surveillance, mesure, audit interne,
revue de direction \\
10 & Amélioration & Non-conformités, actions correctives, amélioration
continue \\
\end{tableaucours}
\vspace{-0.8em}
\begin{center}{\small\itshape Tableau~9.1. La structure harmonisée des
normes ISO de systèmes de management (articles~4 à~10)}\end{center}

Cette structure suit le cycle \emph{planifier, réaliser, vérifier, agir}~:
les articles~4 à~7 planifient, l'article~8 réalise, l'article~9 vérifie,
l'article~10 améliore. On y retrouve, à une autre échelle, la logique de
COBIT~: le domaine EDM correspond aux articles~5 et~9 (direction et
surveillance), les domaines de management aux articles~6 à~8 et~10.

% ============================================================================
\section{ISO/IEC 38500 : la Norme de la Gouvernance}
\label{sec:c09-38500}
% ============================================================================

\subsection{Une norme pour l'organe de gouvernance}

L'ISO/IEC~38500 s'adresse à l'organe de gouvernance~: conseil
d'administration, conseil de surveillance, direction d'un établissement
public. Elle ne décrit ni processus ni contrôles, mais des principes et un
modèle. Elle n'est pas certifiable. Sa troisième édition a été publiée en
février~2024~\cite{iso38500_2024}.

\subsection{Le modèle évaluer, diriger, surveiller}

\noindent\begin{minipage}{\linewidth}
\begin{center}
\begin{tikzpicture}[
    tache/.style = {
        draw=CouleurPrimaire, fill=CouleurPrimaire, rounded corners=4pt,
        text width=2.4cm, minimum height=0.9cm, align=center,
        font=\footnotesize\bfseries\color{white}
    },
    bloc/.style = {
        draw=CouleurBordDef, fill=CouleurDef, rounded corners=4pt,
        text width=3.0cm, minimum height=0.9cm, align=center,
        font=\footnotesize\bfseries\color{CouleurPrimaire}
    },
    fleche/.style = {-{Stealth[length=5pt]}, thick, color=CouleurSecondaire}
]
\node[tache] (ev) at (0,2.2)   {Évaluer};
\node[tache] (di) at (3.6,2.2) {Diriger};
\node[tache] (su) at (7.2,2.2) {Surveiller};
\node[bloc] (pr) at (1.8,0)    {Propositions,\\ plans};
\node[bloc] (op) at (5.4,0)    {Activités,\\ performance};
\draw[fleche] (ev) -- (di);
\draw[fleche] (di) -- node[left, font=\scriptsize\itshape, align=right]{politiques,\\
    plans} (op);
\draw[fleche] (op) -- node[right, font=\scriptsize\itshape]{mesures} (su);
\draw[fleche] (pr) -- (ev);
\draw[fleche] (su.north) to[out=120, in=60] (ev.north);
\node[font=\scriptsize\bfseries\color{CouleurAccent}] at (3.6,4.0)
    {Organe de gouvernance};
\node[font=\scriptsize\bfseries\color{CouleurAccent}] at (3.6,-0.9)
    {Management};
\end{tikzpicture}
\end{center}
\vspace{-0.4em}
\begin{center}{\small\itshape Figure~9.1. Le modèle évaluer, diriger,
surveiller (représentation simplifiée)}\end{center}
\end{minipage}

\medskip

L'organe de gouvernance \textbf{évalue} l'usage actuel et futur de
l'informatique, ainsi que les propositions que lui soumet le management~;
il \textbf{dirige} la préparation et la mise en œuvre des plans et des
politiques~; il \textbf{surveille} la performance et la conformité. Ce sont
les trois verbes du domaine EDM de COBIT, qui s'inspire directement de la
norme.

\subsection{Des principes de 2015 à ceux de 2024}

L'édition de 2015~\cite{iso38500} reposait sur six principes~:
\emph{responsabilité}, \emph{stratégie}, \emph{acquisition},
\emph{performance}, \emph{conformité} et \emph{comportement humain}.
L'édition de 2024 les remplace par les onze principes de la norme
ISO~37000 sur la gouvernance des organisations, et précise comment ils
s'appliquent à la gouvernance des SI. Elle fait de la gouvernance des SI
une partie intégrante de la gouvernance de l'organisation, et l'oriente
vers trois résultats~: une performance effective, une gestion
responsable et un comportement éthique.

\begin{tableaucours}{L{3.2cm}L{9.0cm}}{%
    \tableauHeader{Catégorie} &
    \tableauHeader{Principes (d'après ISO~37000)}
}
Principe premier & Raison d'être (\textit{purpose}) \\
Principes fondateurs & Création de valeur~; stratégie~; supervision~;
redevabilité \\
Principes facilitateurs & Engagement des parties prenantes~; leadership~;
données et décisions~; gouvernance des risques~; responsabilité
sociétale~; viabilité et performance dans la durée \\
\end{tableaucours}
\vspace{-0.8em}
\begin{center}{\small\itshape Tableau~9.2. Les onze principes de
gouvernance repris par l'ISO/IEC~38500:2024 (notre traduction)}\end{center}
% A VERIFIER dans le texte de l'ISO/IEC 38500:2024 : intitulés exacts et
% ordre des onze principes (source : présentations publiques d'ISO 37000).

\begin{boxscenario}{Opération MEFOU. Une charte inspirée de l'ISO/IEC 38500}
Le nouveau président du conseil d'administration de l'ANPP demande que la
charte du comité stratégique des SI (Chapitre~\ref{chap:gouv-entreprise})
s'appuie sur l'ISO/IEC~38500. La division informatique propose trois
articles~:
\begin{enumerate}
    \item \emph{Raison d'être}~: tout projet informatique doit être rattaché
          à la mission de l'agence (payer l'État à bonne date) et à un
          programme budgétaire.
    \item \emph{Supervision}~: le comité \emph{évalue} chaque année le
          portefeuille, \emph{dirige} par la décision d'investissement et
          \emph{surveille} chaque trimestre les projets en cours.
    \item \emph{Données et décisions}~: le conseil reçoit, avant toute
          décision d'investissement, un dossier normalisé (coût total,
          bénéfices, risques).
\end{enumerate}
La charte ne cite pas la norme pour la forme~: elle en traduit les
principes en règles vérifiables.
\end{boxscenario}

% ============================================================================
\section{ISO/IEC 20000-1 : le Management des Services}
\label{sec:c09-20000}
% ============================================================================

\subsection{Les exigences}

L'ISO/IEC~20000-1, dont l'édition en vigueur date de 2018, formule les
exigences d'un \termecle{système de management des services}\index{système
de management des services} (\sigle{SMS}). Elle suit la structure
harmonisée, et place à l'article~8 ses exigences propres au
fonctionnement des services.

\begin{tableaucours}{cL{3.6cm}L{8.0cm}}{%
    \tableauHeader{Art.} &
    \tableauHeader{Thème} &
    \tableauHeader{Exigences principales}
}
8.2 & Portefeuille de services & Planifier les services, gérer le catalogue,
les actifs et la configuration \\
8.3 & Relations et accords & Relations avec les clients, gestion des
niveaux de service, gestion des fournisseurs \\
8.4 & Offre et demande & Budgétisation des services, demande et capacité \\
8.5 & Conception, construction et transition & Gestion des changements,
conception et transition des services, mises en production et
déploiement \\
8.6 & Résolution et réalisation & Gestion des incidents, des demandes de
service et des problèmes \\
8.7 & Assurance des services & Disponibilité, continuité, sécurité de
l'information \\
\end{tableaucours}
\vspace{-0.8em}
\begin{center}{\small\itshape Tableau~9.3. Les exigences de l'article~8 de
l'ISO/IEC~20000-1:2018 (synthèse, notre traduction)}\end{center}

\subsection{ITIL, COBIT et ISO/IEC 20000-1}

Les trois référentiels se recouvrent largement sur la gestion des services,
et c'est précisément pourquoi il faut savoir ce que chacun apporte.
\textbf{ITIL} décrit \emph{comment} mener chaque pratique~;
\textbf{ISO/IEC~20000-1} dit \emph{ce qui doit être démontré} pour obtenir
un certificat~; \textbf{COBIT} relie ces pratiques aux objectifs de
l'organisation et à sa gouvernance. Une organisation qui vise la
certification ISO/IEC~20000-1 s'appuie le plus souvent sur ITIL pour
construire ses pratiques. Le Chapitre~\ref{chap:processus} établira la
correspondance détaillée entre les trois.

\begin{boxmetier}{Responsable des services IT}
\textbf{Préparer KRIBILOG à l'ISO/IEC 20000-1}
\begin{enumerate}
    \item Définir le périmètre~: le service de suivi des conteneurs, depuis
          Douala et Kribi, et non toute l'informatique de l'entreprise.
    \item Établir que KRIBILOG reste \emph{responsable} du service même s'il
          est exploité par NETPORT SERVICES~: la norme exige la maîtrise
          des processus confiés à des tiers.
    \item Mettre d'abord en place les pratiques de l'article~8.6
          (incidents, demandes, problèmes) et de l'article~8.5
          (changements), qui sont aujourd'hui les plus défaillantes.
    \item Conduire un audit interne et une revue de direction avant de
          solliciter l'organisme de certification.
\end{enumerate}
Dix-huit mois sont un délai serré, mais réaliste si le périmètre reste
limité au service demandé par l'armateur.
\end{boxmetier}

% ============================================================================
\section{ISO/IEC 27001 et 27002 : la Sécurité de l'Information}
\label{sec:c09-27001}
% ============================================================================

\subsection{Le système de management de la sécurité}

L'ISO/IEC~27001 formule les exigences d'un \termecle{système de management
de la sécurité de l'information}\index{système de management de la
sécurité de l'information} (\sigle{SMSI}). Le vocabulaire de la famille
est fixé par l'ISO/IEC~27000~\cite{iso27000_2018}. La démarche est fondée
sur le risque~: l'organisation apprécie ses risques de sécurité (selon une
méthode que précise l'ISO/IEC~27005), choisit les mesures qui les
traitent, et justifie ses choix.

\begin{boxdef}
\termecle{Déclaration d'applicabilité}\index{déclaration d'applicabilité}
(\textit{statement of applicability}, \sigle{SoA})~: document qui liste les
mesures de sécurité de l'annexe~A de l'ISO/IEC~27001, indique pour
chacune si elle est retenue ou exclue, et justifie ce choix. C'est la
pièce centrale d'un audit de certification ISO/IEC~27001.
\end{boxdef}

\subsection{Deux éditions à ne pas confondre}

\begin{tableaucours}{L{3.4cm}L{4.4cm}L{4.4cm}}{%
    \tableauHeader{Élément} &
    \tableauHeader{Édition 2013} &
    \tableauHeader{Édition 2022}
}
Mesures de l'annexe~A & 114~mesures en 14~domaines & 93~mesures en
4~thèmes \\
Thèmes & · & Organisationnelles (37), humaines (8), physiques (14),
technologiques (34) \\
Nouvelles mesures & · & Notamment renseignement sur les menaces,
sécurité de l'informatique en nuage, préparation des TIC à la continuité \\
Statut & Remplacée & En vigueur \\
\end{tableaucours}
\vspace{-0.8em}
\begin{center}{\small\itshape Tableau~9.4. ISO/IEC~27001 et 27002~: éditions
2013 et 2022}\end{center}

\begin{boxalert}
\textbf{Violation typique.} Une organisation rédige en 2026 sa politique de
sécurité et sa déclaration d'applicabilité à partir des éditions 2013 des
normes, trouvées dans un dossier partagé. L'auditeur de certification
relève aussitôt l'écart~: la référence n'est plus celle en vigueur, et les
mesures relatives à l'informatique en nuage manquent.

\smallskip
\textbf{Règle~:} toujours vérifier l'édition en vigueur d'une norme dans le
catalogue de l'ISO avant de s'en servir de référence, et indiquer l'édition
dans chaque document.
\end{boxalert}

\subsection{Le lien avec la gouvernance}

La sécurité de l'information n'est pas la gouvernance des SI
(Chapitre~\ref{chap:gouv-entreprise}), mais elle en relève par trois
voies. Le conseil fixe l'appétence au risque de sécurité (EDM03). Le
management conduit le SMSI (APO13, dont la finalité recouvre celle de
l'ISO/IEC~27001). L'exploitation met en œuvre les mesures (DSS05). Un
SMSI certifié est un excellent moyen d'atteindre APO13~; il ne dit rien
de la valeur des investissements ni de la qualité des services.

\begin{boxmetier}{RSSI}
\textbf{Répondre au conseil de la BANQUE SANAGA}
\begin{enumerate}
    \item Une certification ISO/IEC~27001 ne « rassure » le régulateur que
          si son périmètre couvre les systèmes qui l'intéressent~: les
          systèmes de paiement et le système central.
    \item Elle ne remplace ni les audits prévus par la loi n°~2010/012
          (Chapitre~\ref{chap:referentiels}), ni les exigences de la COBAC
          sur le contrôle interne.
    \item Elle n'aurait pas évité la panne de mars~2026, qui ne relevait
          pas de la sécurité mais de la gestion des changements.
    \item Recommandation~: viser la certification à trois ans sur le
          périmètre des paiements, comme le prévoit la carte des
          référentiels (Chapitre~\ref{chap:referentiels}), et traiter
          d'abord la feuille de route COBIT (Chapitre~\ref{chap:cobit-capacite}).
\end{enumerate}
\end{boxmetier}

% ============================================================================
\section{Le Risque : ISO 31000 et COSO}
\label{sec:c09-risque}
% ============================================================================

\subsection{ISO 31000 : des lignes directrices pour tout risque}

L'ISO~31000 (édition 2018) donne des lignes directrices sur le management
du risque, applicables à tout type de risque et à toute organisation. Elle
n'est pas certifiable. Elle s'articule en trois parties~: des
\emph{principes} (un management du risque intégré, structuré, adapté,
inclusif, dynamique, fondé sur la meilleure information disponible,
attentif aux facteurs humains et culturels, et en amélioration continue),
un \emph{cadre organisationnel}, et un \emph{processus}~: établir le
contexte et les critères, apprécier les risques (les identifier, les
analyser, les évaluer), les traiter, puis surveiller, communiquer et
rendre compte.

\subsection{COSO : contrôle interne et risques de l'entreprise}

Les deux cadres du COSO (Chapitre~\ref{chap:referentiels}) sont la
référence des auditeurs et des régulateurs financiers.

\begin{tableaucours}{L{3.6cm}L{8.6cm}}{%
    \tableauHeader{Cadre} &
    \tableauHeader{Composantes}
}
Contrôle interne (2013) & Environnement de contrôle~; évaluation des
risques~; activités de contrôle~; information et communication~;
activités de pilotage. Cinq composantes et 17~principes. \\
Management des risques (2004) & Huit éléments~: environnement interne~;
fixation des objectifs~; identification des événements~; évaluation des
risques~; traitement des risques~; activités de contrôle~; information et
communication~; pilotage~\cite{coso2004ermfr}. \\
Management des risques (2017) & Cinq composantes~: gouvernance et
culture~; stratégie et définition des objectifs~; performance~; revue et
révision~; information, communication et reporting. Vingt
principes~\cite{coso2017erm}. \\
\end{tableaucours}
\vspace{-0.8em}
\begin{center}{\small\itshape Tableau~9.5. Les cadres du COSO}\end{center}

Le résumé en français disponible dans le dossier \texttt{Normes} est celui
de l'édition de 2004. L'édition de 2017 en déplace le centre de gravité~:
elle ne traite plus le risque comme un sujet de contrôle, mais comme une
composante de la stratégie et de la performance. C'est la lecture que
retient aussi COBIT avec la finalité d'optimisation des risques
(Chapitre~\ref{chap:alignement}).

\subsection{Articuler les risques}

La règle d'articulation est la suivante~: \textbf{le risque informatique
est un risque de l'entreprise}. Le registre des risques liés aux SI
(APO12) doit alimenter le registre des risques de l'organisation, selon
la même échelle et les mêmes critères, et l'appétence au risque fixée par
le conseil (EDM03) s'applique aux deux. Une organisation qui tient deux
registres séparés, avec deux échelles, ne peut pas comparer un risque
informatique à un risque de crédit~: elle ne peut donc pas arbitrer.

% ============================================================================
\section{Certification, Conformité, Alignement}
\label{sec:c09-certification}
% ============================================================================

\subsection{Trois démarches}

\begin{boxdef}
\termecle{Certification}\index{certification}~: attestation, par un
organisme tiers accrédité, qu'un système de management est conforme aux
exigences d'une norme, sur un périmètre donné et pour une durée donnée.

\medskip
\termecle{Conformité}\index{conformité}~: respect d'exigences obligatoires
(lois, règlements, contrats). Elle ne se choisit pas~; elle se démontre.

\medskip
\termecle{Alignement volontaire}\index{alignement volontaire}~: choix
d'une organisation de s'inspirer d'un référentiel ou d'une norme sans
viser de certificat. C'est le cas de tout usage de COBIT, d'ITIL ou de
l'ISO~31000.
\end{boxdef}

\subsection{Le cycle de certification}

Une certification de système de management suit un cycle de trois ans~:
un audit initial en deux étapes (revue de la documentation et de l'état
de préparation, puis audit de la mise en œuvre), la délivrance du
certificat, des audits de surveillance au moins annuels, puis un audit de
renouvellement. Les organismes de certification doivent eux-mêmes
satisfaire aux exigences de l'ISO/IEC~17021-1, complétées pour la sécurité
de l'information par l'ISO/IEC~27006, et être accrédités par un organisme
d'accréditation. Un certificat délivré par un organisme non accrédité n'a
qu'une valeur commerciale.

\begin{boxprocedure}
\textbf{Préparer un audit de certification.} Avant de solliciter
l'organisme de certification, l'organisation doit pouvoir présenter~:
\begin{itemize}
    \item le périmètre du système, écrit et justifié~;
    \item la politique et les objectifs, approuvés par la direction~;
    \item l'appréciation des risques et, pour l'ISO/IEC~27001, la
          déclaration d'applicabilité~;
    \item les procédures exigées et les enregistrements de leur
          fonctionnement sur plusieurs mois~;
    \item au moins un cycle complet d'audit interne et de revue de
          direction~;
    \item le suivi des non-conformités relevées et des actions correctives.
\end{itemize}
\end{boxprocedure}

\begin{boxscenario}{Opération KRIBI. Le calendrier de la certification}
\textbf{Contrainte.} Certification ISO/IEC~20000-1 exigée dans dix-huit
mois sur le service de suivi des conteneurs.

\medskip
\textbf{Plan retenu par M.~MBARGA.} Mois 1 à 2~: périmètre, politique et
objectifs, analyse d'écart. Mois 3 à 9~: mise en place des pratiques
(incidents, demandes, problèmes, changements, niveaux de service,
fournisseurs), avec un avenant au contrat de NETPORT SERVICES. Mois 10 à
13~: fonctionnement et enregistrements. Mois 14~: audit interne. Mois 15~:
revue de direction et actions correctives. Mois 16 à 17~: audit de
certification en deux étapes. Mois 18~: certificat.

\medskip
\textbf{Analyse} (confiance moyenne). Le calendrier tient si l'avenant avec
l'infogérant est signé dans les trois premiers mois~: sans lui, KRIBILOG
ne pourra démontrer la maîtrise d'un service exploité par un tiers.
\end{boxscenario}

\begin{boxmetier}{Auditeur SI}
\textbf{Lire un certificat en cinq questions}
\begin{enumerate}
    \item Quelle norme, et quelle édition~?
    \item Quel périmètre exact (sites, services, systèmes)~?
    \item Quel organisme l'a délivré, et est-il accrédité~?
    \item Quelle date de validité, et le dernier audit de surveillance
          a-t-il eu lieu~?
    \item Le périmètre couvre-t-il le service qui vous intéresse~?
\end{enumerate}
\end{boxmetier}

% ============================================================================
\separateur
% ============================================================================

\begin{boxresume}
\textbf{Ce qu'il faut retenir de ce chapitre~:}
\begin{itemize}
    \item \textbf{Une structure harmonisée}~: les normes de systèmes de
          management partagent les articles~4 à~10 (contexte, leadership,
          planification, support, fonctionnement, évaluation,
          amélioration), ce qui permet des systèmes intégrés.
    \item \textbf{ISO/IEC~38500} s'adresse à l'organe de gouvernance~:
          évaluer, diriger, surveiller. Depuis 2024, elle reprend les onze
          principes de l'ISO~37000. Elle n'est pas certifiable.
    \item \textbf{ISO/IEC~20000-1} exige un système de management des
          services~; ITIL dit comment le construire, COBIT le relie à la
          gouvernance.
    \item \textbf{ISO/IEC~27001 et 27002}~: un SMSI fondé sur le risque, une
          déclaration d'applicabilité, 93~mesures depuis 2022. Vérifier
          toujours l'édition en vigueur.
    \item \textbf{ISO~31000 et COSO}~: le risque informatique est un risque
          de l'entreprise, à inscrire dans un registre et une échelle
          communs.
    \item \textbf{Certification, conformité, alignement}~: la certification
          atteste un périmètre, la conformité répond à une obligation,
          l'alignement est un choix.
\end{itemize}
\end{boxresume}

\bigskip

\begin{boxexo}[title={\ding{45}\quad Exercice~9.1. Quel article~? \niveaubase}]
Indiquez l'article de la structure harmonisée (4 à 10) auquel se rattache
chacun des éléments suivants.
\begin{enumerate}[(a)]
    \item La revue annuelle du système par la direction générale.
    \item La liste des parties intéressées et de leurs attentes.
    \item Le plan de formation des techniciens du centre de services.
    \item Le traitement d'une non-conformité relevée en audit interne.
    \item La procédure de gestion des incidents.
\end{enumerate}
\end{boxexo}

\begin{boxexo}[title={\ding{45}\quad Exercice~9.2. Deux registres, une échelle \niveaumoyen}]
La BANQUE SANAGA tient un registre des risques de crédit et de marché
(échelle de 1 à 4) et, à la DSI, un registre des risques informatiques
(échelle de 1 à 5, définie différemment).
\begin{enumerate}[(a)]
    \item Expliquez pourquoi cette situation empêche le conseil
          d'arbitrer.
    \item Proposez une échelle commune, avec des critères d'impact
          exprimés en FCFA et en durée d'indisponibilité.
    \item Reclassez sur cette échelle le risque « échec d'une migration du
          système central ».
\end{enumerate}
\end{boxexo}

\begin{boxexo}[title={\ding{45}\quad Exercice~9.3. Une stratégie normative pour NOUN PAY \niveauexpert}]
Les investisseurs de NOUN PAY exigent « des certifications
internationales ». Les moyens de l'entreprise ne permettent d'en viser
qu'une dans les deux ans.
\begin{enumerate}[(a)]
    \item Comparez l'intérêt d'une certification ISO/IEC~27001 et d'une
          certification ISO/IEC~20000-1 pour NOUN PAY.
    \item Proposez un périmètre réaliste pour la certification retenue.
    \item Indiquez les référentiels que l'entreprise suivra par simple
          alignement volontaire, et pourquoi.
    \item Rédigez la réponse de Mme~KAMGA aux investisseurs, en distinguant
          certification, conformité légale (loi n°~2024/017) et
          alignement.
\end{enumerate}
\end{boxexo}

\bigskip

\begin{boxinfo}
\textbf{Transition.} Les normes sont des engagements volontaires~; la loi,
elle, s'impose. Le Chapitre~\ref{chap:droit} présente le cadre juridique
et réglementaire camerounais et africain de la gouvernance des SI~:
cybersécurité, communications électroniques, données personnelles,
réglementation bancaire et droit des sociétés.
\end{boxinfo}

        % ============================================================================
% Fichier  : partie3/P03_C10_cadre_juridique_cameroun_afrique.tex
% Titre    : Cadre Juridique et Réglementaire Camerounais et Africain
% Auteur   : MINKA MI NGUIDJOÏ Thierry Emmanuel
% Version  : 1.0 -- octobre 2026
% ============================================================================
% Chapitre juridique. Sources et degré de vérification :
%  - Loi n° 2010/012 : texte officiel anglais consulté (sections citées
%    textuellement).
%  - Règlement COBAC R-2016/04 : articles 1 et 2 relevés sur une base
%    réglementaire en ligne ; à collationner avec la version officielle.
%  - Loi n° 2024/017 : présentations publiées (presse spécialisée) ; aucun
%    article cité textuellement. A VERIFIER au JO.
%  - AUSCGIE (OHADA) : numéros d'articles d'après la doctrine et la
%    brochure OHADA de 2014 ; aucun article cité textuellement.
%  - Convention de Malabo : entrée en vigueur vérifiée ; statut du Cameroun
%    non établi.
% ============================================================================

\chapter{Cadre Juridique et Réglementaire Camerounais et Africain}
\label{chap:droit}
\chaptermark{Cadre juridique camerounais et africain}

\epigraph{%
    « Nul n'est censé ignorer la loi, pas même le DSI. »%
}{Adage revisité}

\niveauChapitre{Fondamental}{%
    Maîtrise du Chapitre~\ref{chap:normes} (normes, certification,
    conformité).%
}

\begin{boxobjectifs}
\begin{itemize}
    \item Identifier les obligations des lois n°~2010/012 et 2010/013 qui
          touchent la gouvernance des SI.
    \item Exposer le rôle de l'ANTIC et des autres régulateurs (ART,
          COBAC).
    \item Appliquer les exigences de la loi de 2024 sur les données
          personnelles à un système d'information.
    \item Situer les responsabilités des dirigeants dans le droit OHADA.
    \item Identifier les obligations de contrôle interne et de risque
          opérationnel des banques en zone CEMAC.
    \item Situer la Convention de Malabo.
\end{itemize}
\end{boxobjectifs}

\bigskip

% ============================================================================
\section*{Question directrice : quelles obligations ont été méconnues~?}
% ============================================================================

Trois semaines après la panne, le conseil d'administration de la BANQUE
SANAGA reçoit une note de son avocat. Elle commence ainsi~: « La banque
risque d'avoir à répondre de ses manquements devant trois autorités au
moins~: la COBAC, au titre du contrôle interne~; l'ANTIC, au titre de la
sécurité de ses systèmes d'information~; et, le cas échéant, les
juridictions, au titre de la responsabilité de ses dirigeants. » À
Bafoussam, NOUN PAY reçoit au même moment une demande d'information de
l'Autorité de protection des données à caractère personnel sur son
modèle de notation.

Ce chapitre dresse la carte de ces obligations. Il suit l'ordre du
modèle~: le fait, la règle, le texte, l'analyse. Il présente d'abord le
panorama des textes (section~\ref{sec:c10-panorama}), puis la loi sur la
cybersécurité (section~\ref{sec:c10-2010-012}), les régulateurs
(section~\ref{sec:c10-regulateurs}), les données personnelles
(section~\ref{sec:c10-donnees}), la réglementation bancaire
(section~\ref{sec:c10-cobac}), le droit des sociétés
(section~\ref{sec:c10-ohada}) et le cadre africain
(section~\ref{sec:c10-afrique}). Il se termine par les erreurs de
qualification les plus fréquentes (section~\ref{sec:c10-erreurs}).

\begin{boxalert}
\textbf{Règle de lecture.} Ce chapitre cite textuellement les seuls
articles dont le texte a pu être vérifié. Pour les autres, il en décrit la
substance d'après des sources secondaires identifiées, sans guillemets.
Avant tout usage professionnel (avis juridique, rapport au régulateur),
reportez-vous au texte publié au Journal officiel.
\end{boxalert}

\separateur

% ============================================================================
\section{Panorama des Textes}
\label{sec:c10-panorama}
% ============================================================================

\begin{tableaucours}{L{3.6cm}L{3.8cm}L{2.4cm}L{2.4cm}}{%
    \tableauHeader{Texte} &
    \tableauHeader{Objet} &
    \tableauHeader{Autorité} &
    \tableauHeader{Cas du cours}
}
Loi n°~2010/012 du 21~décembre~2010 & Cybersécurité et cybercriminalité
& ANTIC & Tous \\
Loi n°~2010/013 du 21~décembre~2010, modifiée en 2015 & Communications
électroniques & ART, MINPOSTEL & NOUN PAY \\
Loi n°~2024/017 du 23~décembre~2024 & Protection des données à caractère
personnel & Autorité de protection des données & NOUN PAY, SANAGA \\
Règlement COBAC R-2016/04 & Contrôle interne des établissements de crédit
& COBAC & SANAGA \\
Acte uniforme OHADA (révisé en 2014) & Sociétés commerciales,
responsabilité des dirigeants & Juridictions & SANAGA, KRIBILOG \\
Lois n°~2017/010 et 2017/011 & Établissements et entreprises publics &
Tutelle & ANPP \\
Loi n°~2018/012 & Régime financier de l'État & MINFI, Chambre des comptes &
ANPP \\
Convention de Malabo (2014) & Cybersécurité et données personnelles en
Afrique & Union africaine & · \\
\end{tableaucours}
\vspace{-0.8em}
\begin{center}{\small\itshape Tableau~10.1. Les principaux textes
applicables à la gouvernance des SI au Cameroun}\end{center}

% ============================================================================
\section{La Loi sur la Cybersécurité}
\label{sec:c10-2010-012}
% ============================================================================

\subsection{L'ANTIC, régulateur de la sécurité}

La loi n°~2010/012 confie la régulation de la sécurité électronique à
l'Agence nationale des technologies de l'information et de la
communication (\sigle{ANTIC}), instituée par la loi sur les
communications électroniques~\cite{loi2010012}.

\begin{boxjuris}
\textbf{\loicam{n°~2010/012}, article~7, alinéa~2} (version anglaise
officielle)~:

\medskip
\textit{« The Agency \textelp{} shall be responsible for the regulation,
control and monitoring of activities related to the security of
electronic communication networks, information systems, and electronic
certification on behalf of the State. »}

\medskip
Soit~: l'Agence est chargée, pour le compte de l'État, de la régulation,
du contrôle et du suivi des activités liées à la sécurité des réseaux de
communications électroniques, des systèmes d'information et de la
certification électronique. Parmi ses missions, la loi énumère notamment
le contrôle des activités de sécurité des systèmes d'information,
la régularité et l'efficacité des audits de sécurité, et la veille sur
les risques informatiques et la cybercriminalité.
\end{boxjuris}

\subsection{Les obligations des exploitants de systèmes d'information}

La loi impose aux exploitants de systèmes d'information
(\textit{operators of information systems}) un ensemble d'obligations qui
recoupent directement plusieurs objectifs de COBIT.

\begin{boxjuris}
\textbf{\loicam{n°~2010/012}, article~26, alinéa~1}~:

\medskip
\textit{« Operators of information systems shall take every technical and
administrative measure to ensure the security of services offered. To
this end, they shall have standardized systems enabling them to at all
times identify, assess, process or manage any risk relating to the
security of the information systems of the services provided directly or
indirectly. »}

\medskip
\textbf{Article~29, alinéa~1}~: \textit{« Operators of information systems
shall be bound to conserve the connection and traffic data of their
information systems for a period of 10 (ten) years. »}

\medskip
\textbf{Article~30, alinéa~1}~: \textit{« Operators of information systems
shall assess and revise their security systems and, where necessary, make
the appropriate modifications to their security practices, measures and
techniques according to technological change. »}

\medskip
\textbf{Analyse.} L'article~26 impose plus qu'une sécurité technique~: il
exige des \emph{dispositifs normalisés} permettant d'identifier, d'évaluer
et de traiter en permanence les risques de sécurité. C'est, en termes de
gouvernance, l'exigence d'un processus de gestion des risques (APO12) et
d'un système de management de la sécurité (APO13), dont l'ISO/IEC~27001
fournit précisément le modèle. L'article~30 consacre l'amélioration
continue. L'article~29 impose une durée de conservation des données de
connexion et de trafic de dix ans, qui pèse sur l'architecture et le
budget de stockage.
\end{boxjuris}

\begin{tableaucours}{L{2.4cm}L{5.4cm}L{4.4cm}}{%
    \tableauHeader{Article} &
    \tableauHeader{Obligation} &
    \tableauHeader{Objectifs COBIT concernés}
}
13 et 32 & Audit de sécurité obligatoire et périodique & MEA03, MEA04 \\
26 & Mesures de sécurité, dispositifs normalisés de gestion des risques &
APO12, APO13, DSS05 \\
29 & Conservation des données de connexion et de trafic pendant dix ans &
APO14, DSS01, DSS05 \\
30 & Évaluation et révision des dispositifs de sécurité & APO13, MEA01 \\
31 & Disponibilité des contenus et des données & DSS04, BAI04 \\
42 & Confidentialité des informations acheminées & DSS05, APO14 \\
\end{tableaucours}
\vspace{-0.8em}
\begin{center}{\small\itshape Tableau~10.2. Obligations de la loi
n°~2010/012 et objectifs COBIT correspondants}\end{center}

\subsection{Les sanctions et la responsabilité des personnes morales}

\begin{boxjuris}
\textbf{\loicam{n°~2010/012}, article~64, alinéas~1 à~3}~:

\medskip
\textit{« (1) Corporate bodies shall be criminally liable for offences
committed on their account by their management structures. (2) The
criminal liability of corporate bodies shall not preclude that of natural
persons who commit such offences or are accomplices. (3) The penalties to
be meted out on defaulting corporate bodies shall be fines of from
5~000~000 (five million) to 50~000~000 (fifty million) CFA francs. »}

\medskip
\textbf{Analyse.} La personne morale répond pénalement des infractions
commises pour son compte par ses organes de direction, sans que cela
exonère les personnes physiques. L'article prévoit en outre des peines
complémentaires (dissolution, interdiction d'activité, fermeture
d'établissement) dans les cas les plus graves. L'article~63 punit par
ailleurs d'emprisonnement et d'amende le dirigeant de droit ou de fait de
certains prestataires (fournisseurs d'accès et d'hébergement) qui ne
conserve pas les données exigées par les articles~25 et~29. La
gouvernance des SI n'est donc pas qu'une question de performance~: elle
engage la responsabilité pénale de l'organisation et, dans certains cas,
celle de ses dirigeants.
\end{boxjuris}

\begin{boxmetier}{DSI}
\textbf{La conservation des données pendant dix ans}
\begin{enumerate}
    \item Établir si l'organisation est « exploitant de système
          d'information » au sens de la loi~: une banque, un opérateur, une
          administration le sont presque toujours.
    \item Recenser les données de connexion et de trafic produites
          (journaux d'accès, journaux réseau, traces des applications).
    \item Définir l'architecture de conservation~: durée, intégrité (les
          journaux ne doivent pas pouvoir être modifiés), confidentialité,
          coût sur dix ans.
    \item Concilier cette obligation avec le principe de limitation de la
          conservation des données personnelles (section~\ref{sec:c10-donnees})~:
          les deux textes doivent être lus ensemble, avec un avis
          juridique.
\end{enumerate}
\end{boxmetier}

% ============================================================================
\section{Les Régulateurs}
\label{sec:c10-regulateurs}
% ============================================================================

La gouvernance des SI d'une organisation camerounaise s'exerce sous le
regard de plusieurs autorités, qu'il faut savoir distinguer.

\begin{tableaucours}{L{2.6cm}L{4.6cm}L{5.0cm}}{%
    \tableauHeader{Autorité} &
    \tableauHeader{Fondement} &
    \tableauHeader{Ce qui intéresse la gouvernance des SI}
}
MINPOSTEL & Ministère chargé des télécommunications & Politique du
secteur~; destinataire des rapports d'audit de sécurité (loi
n°~2010/012, art.~32) \\
ANTIC & Lois n°~2010/012 et 2010/013 & Sécurité des systèmes
d'information, audits, certification électronique, veille \\
ART & Loi n°~2010/013 modifiée & Régulation des communications
électroniques et des opérateurs \\
Autorité de protection des données & Loi n°~2024/017 & Traitements de
données personnelles, plaintes, sanctions \\
COBAC & Réglementation bancaire de la CEMAC & Contrôle interne et risques
des établissements de crédit, dont le risque informatique \\
Chambre des comptes & Contrôle des finances publiques & Audit des
établissements et entreprises publics, y compris de leurs SI \\
\end{tableaucours}
\vspace{-0.8em}
\begin{center}{\small\itshape Tableau~10.3. Les autorités de la gouvernance
des SI au Cameroun}\end{center}

La loi n°~2010/013 régit les communications électroniques~; elle a été
modifiée et complétée par la loi n°~2015/006 du 20~avril~2015. Elle
institue l'Agence de régulation des télécommunications (\sigle{ART}) et
l'ANTIC. Ses dispositions intéressent surtout les opérateurs et les
fournisseurs de services de communications électroniques~: une
entreprise de paiement mobile comme NOUN PAY, qui s'appuie sur les
réseaux des opérateurs, doit en connaître les grandes règles. Le recueil
des textes publié par l'ART en donne la version consolidée.

% ============================================================================
\section{La Protection des Données Personnelles}
\label{sec:c10-donnees}
% ============================================================================

\subsection{De la loi de 2010 à la loi de 2024}

La loi n°~2010/012 posait déjà un principe général.

\begin{boxjuris}
\textbf{\loicam{n°~2010/012}, articles~41 et~42}~:

\medskip
\textit{« 41. Every individual shall have the right to the protection of
their privacy. \textelp{} 42. The confidentiality of information
channelled through electronic communication and information systems
networks, including traffic data, shall be ensured by operators of
electronic communication and networks information systems. »}

\medskip
\textbf{Portée.} Ces articles affirment le droit au respect de la vie
privée et l'obligation de confidentialité des exploitants. Ils ne
constituaient pas, à eux seuls, un régime complet de protection des
données personnelles~: c'est l'objet de la loi de 2024.
\end{boxjuris}

\subsection{La loi n° 2024/017}

La loi n°~2024/017 du 23~décembre~2024 relative à la protection des
données à caractère personnel a doté le Cameroun d'un régime complet.
D'après les présentations publiées~\cite{cio2026donnees}, elle a accordé
aux organisations un délai de dix-huit mois pour se mettre en conformité,
et s'applique pleinement depuis le 23~juin~2026. Elle institue une
Autorité de protection des données à caractère personnel, organisme
public indépendant, encadre strictement le transfert des données hors du
territoire national, consacre une part importante de ses dispositions aux
obligations des responsables du traitement et de leurs sous-traitants,
et prévoit des sanctions pénales.

\begin{boxalert}
\textbf{À vérifier avant usage.} Les numéros d'articles, les définitions
exactes (responsable du traitement, sous-traitant, données sensibles),
les conditions du transfert des données hors du Cameroun et l'éventuelle
obligation de désigner un délégué à la protection des données doivent
être lus dans le texte publié et ses décrets d'application. Les
présentations disponibles divergent sur certains de ces points.
\end{boxalert}

\subsection{Ce que la loi change pour la gouvernance des SI}

Quelle que soit la lettre exacte de ses articles, une loi de protection
des données de cette génération a des conséquences directes sur la
gouvernance des SI~:
\begin{itemize}
    \item \textbf{un inventaire des traitements}~: savoir quelles données
          personnelles sont traitées, pourquoi, où et par qui (APO14)~;
    \item \textbf{la sécurité des traitements}~: mesures techniques et
          organisationnelles proportionnées aux risques (APO13, DSS05)~;
    \item \textbf{la maîtrise des sous-traitants}~: un hébergeur ou un
          fournisseur de nuage traite les données pour le compte de
          l'organisation, qui en reste responsable (APO10)~;
    \item \textbf{la localisation des données}~: tout transfert hors du
          territoire doit être examiné au regard de la loi (BAI03, APO10)~;
    \item \textbf{la conformité démontrable}~: registres, analyses
          d'impact, preuves (MEA03).
\end{itemize}

\begin{boxscenario}{Opération NOUN. La demande de l'Autorité}
\textbf{Faits.} L'Autorité de protection des données demande à NOUN PAY~:
la liste des données utilisées par son modèle de notation de crédit, le
lieu où elles sont hébergées, la base juridique du traitement, et la
manière dont un client refusé peut contester la décision.

\medskip
\textbf{Constatation.} NOUN PAY ne dispose d'aucun inventaire de ses
traitements. Ses données sont hébergées dans un centre de données situé
hors d'Afrique, en vertu d'un contrat type du fournisseur. Aucune
procédure ne permet à un client de contester un refus.

\medskip
\textbf{Analyse} (confiance moyenne, sous réserve de la lecture du texte
de la loi). Trois manquements probables~: l'absence d'inventaire, un
transfert hors du territoire non examiné au regard de la loi, et
l'absence de recours contre une décision automatisée. Le
Chapitre~\ref{chap:donnees-ia} reprendra le cas sous l'angle de la
gouvernance de l'intelligence artificielle.
\end{boxscenario}

% ============================================================================
\section{La Réglementation Bancaire de la CEMAC}
\label{sec:c10-cobac}
% ============================================================================

\subsection{Le règlement COBAC R-2016/04}

Les établissements de crédit de la zone CEMAC sont soumis à la
réglementation de la Commission bancaire de l'Afrique centrale. Le
règlement COBAC R-2016/04, relatif au contrôle interne dans les
établissements de crédit et les holdings financières, est le texte de
référence pour la gouvernance des SI d'une banque~\cite{cobac2016r04}.

\begin{boxjuris}
\textbf{Règlement COBAC R-2016/04, articles~1 et~2} (texte relevé sur une
base réglementaire en ligne)~:

\medskip
Article~1~: \textit{« Le présent règlement fixe les exigences minimales
relatives au système de contrôle interne que les établissements assujettis
doivent mettre en œuvre. »}

\medskip
Article~2, définition du risque opérationnel~: \textit{« risque de pertes
résultant de l'inadéquation ou de la défaillance des processus internes,
des personnes et des systèmes, ou d'évènements externes. »} Le même
article précise que ce risque inclut le risque juridique, et exclut les
risques stratégiques et de réputation.

\medskip
\textbf{Analyse.} La défaillance d'un système informatique est, par
définition, un risque opérationnel au sens du règlement. La panne de la
BANQUE SANAGA relève donc du dispositif de contrôle interne exigé par la
COBAC, et non d'un simple incident technique. Selon les présentations du
règlement, celui-ci impose en outre un dispositif d'identification, de
mesure et de surveillance de chaque risque, appuyé sur une cartographie
des risques, et compte un système de traitement de l'information parmi
les composantes du dispositif de contrôle interne.
\end{boxjuris}
% A COLLATIONNER avec la version officielle publiée par le SG-COBAC
% (sgcobac.org) : texte exact des articles 1 et 2, contenu des articles 3,
% 10 et 12.

\subsection{Ce que cela implique}

\begin{boxmetier}{Responsable GRC}
\textbf{Rédiger la partie SI du rapport à la COBAC}
\begin{enumerate}
    \item Qualifier la panne~: incident de risque opérationnel, avec sa
          chronologie et son coût (Chapitre~\ref{chap:prologue}).
    \item Situer la défaillance dans le dispositif de contrôle interne~:
          quels contrôles existaient, lesquels n'ont pas fonctionné.
    \item Relier le risque informatique à la cartographie des risques de la
          banque, avec l'échelle commune (Chapitre~\ref{chap:normes}).
    \item Présenter le plan d'action, objectif par objectif, à partir de la
          feuille de route COBIT (Chapitre~\ref{chap:cobit-capacite}).
\end{enumerate}
\end{boxmetier}

% ============================================================================
\section{Le Droit des Sociétés OHADA}
\label{sec:c10-ohada}
% ============================================================================

\subsection{La responsabilité des dirigeants}

L'Acte uniforme OHADA relatif au droit des sociétés commerciales et du
groupement d'intérêt économique, révisé le 30~janvier~2014~\cite{ohada2014auscgie},
organise la responsabilité civile des dirigeants sociaux. Selon la
doctrine, ses articles~161 et suivants posent le régime général~: le
dirigeant répond des fautes qu'il commet dans l'exercice de ses
fonctions, individuellement ou solidairement selon sa participation, et
peut être poursuivi par la société (action sociale) ou par les tiers
lésés. Pour la société anonyme, l'article~740 rend les administrateurs
responsables, individuellement ou solidairement, envers la société ou
envers les tiers.

\begin{boiteRemarque}{Un défaut de gouvernance des SI peut-il être une faute de gestion~?}
La question n'a pas, à notre connaissance, été tranchée par une
jurisprudence publiée en zone OHADA. Mais le raisonnement est
accessible~: si une banque soumise à des obligations précises (contrôle
interne, sécurité des systèmes) subit une perte majeure parce que ses
dirigeants ont maintenu une bascule malgré une alerte écrite, sans
décision formelle d'acceptation du risque, la qualification de faute de
gestion pourra être discutée. La trace écrite des décisions
(Chapitre~\ref{chap:prologue}) n'est donc pas une formalité~: elle est la
première défense d'un dirigeant diligent, et la première pièce d'un
dossier contre un dirigeant négligent.
\end{boiteRemarque}

\subsection{Le comité d'audit et le contrôle interne}

Le Chapitre~\ref{chap:gouv-entreprise} a présenté l'article~829-1 de l'Acte
uniforme révisé, qui rend le comité d'audit obligatoire dans les sociétés
cotées et le charge du suivi de l'efficacité des systèmes de contrôle
interne et de gestion des risques. Pour les entreprises et
établissements publics, les lois n°~2017/010 et 2017/011 et leur décret
d'application organisent le conseil d'administration et la tutelle~: la
gouvernance des SI de l'ANPP doit s'inscrire dans ce cadre, et dans celui
du régime financier de l'État (Chapitre~\ref{chap:alignement}).

% ============================================================================
\section{Le Cadre Africain}
\label{sec:c10-afrique}
% ============================================================================

\subsection{La Convention de Malabo}

La Convention de l'Union africaine sur la cybersécurité et la protection
des données à caractère personnel, adoptée à Malabo le 27~juin~2014, est
entrée en vigueur le 8~juin~2023, trente jours après le dépôt du
quinzième instrument de ratification. Elle couvre trois domaines~: les
transactions électroniques, la protection des données personnelles, et
la cybersécurité et la lutte contre la cybercriminalité. Elle oblige les
États parties à se doter d'une législation et d'autorités dans ces
domaines.

\begin{boxalert}
\textbf{À vérifier.} Nous n'avons pas trouvé de source établissant que le
Cameroun a ratifié la Convention de Malabo. Avant de l'invoquer comme
source de droit applicable au Cameroun, consultez la liste officielle des
signatures et ratifications tenue par l'Union africaine. Elle reste en
tout état de cause une référence de politique publique pour le
continent.
\end{boxalert}

\subsection{Une dynamique continentale}

Le Cameroun n'est pas seul. Plusieurs États africains se sont dotés de
lois de protection des données personnelles depuis la fin des années
2000, par exemple le Sénégal (loi n°~2008-12 du 25~janvier~2008) ou la
Côte d'Ivoire (loi n°~2013-450 du 19~juin~2013). Pour une organisation qui
opère dans plusieurs pays de la sous-région, comme une banque de la zone
CEMAC ou une entreprise de paiement mobile, la gouvernance des SI doit
tenir un registre des obligations pays par pays.

% ============================================================================
\section{Les Erreurs de Qualification}
\label{sec:c10-erreurs}
% ============================================================================

\begin{tableaucours}{L{4.6cm}L{7.6cm}}{%
    \tableauHeader{Erreur} &
    \tableauHeader{Pourquoi c'est une erreur}
}
« La panne est un incident technique, pas un sujet réglementaire. » & Pour
une banque, la défaillance d'un système est un risque opérationnel au sens
du règlement COBAC R-2016/04. \\
« Nous sommes certifiés ISO~27001, l'audit de l'ANTIC est inutile. » & La
certification est volontaire~; l'audit prévu par la loi n°~2010/012 est
obligatoire. \\
« Nos données sont chez un prestataire, c'est lui le responsable. » &
L'externalisation ne transfère pas la responsabilité du traitement ni
celle de la sécurité (loi n°~2010/012, art.~26~; loi n°~2024/017). \\
« Nos serveurs sont à l'étranger, la loi camerounaise ne s'applique
pas. » & Le transfert hors du territoire est précisément encadré par la
loi n°~2024/017. \\
« Seule la société est poursuivie, pas ses dirigeants. » & La
responsabilité de la personne morale n'exclut pas celle des personnes
physiques (loi n°~2010/012, art.~64~; AUSCGIE). \\
« La Convention de Malabo s'applique directement. » & Son applicabilité
au Cameroun suppose une ratification qui reste à établir. \\
\end{tableaucours}
\vspace{-0.8em}
\begin{center}{\small\itshape Tableau~10.4. Erreurs de qualification
fréquentes}\end{center}

\begin{boxprocedure}
\textbf{Le registre des obligations légales.} Pour démontrer la
conformité (objectif MEA03), l'organisation tient un registre qui, pour
chaque obligation, indique~:
\begin{itemize}
    \item le texte et l'article, avec la date de la version consultée~;
    \item l'obligation, reformulée en termes opérationnels~;
    \item le processus et le propriétaire qui la mettent en œuvre~;
    \item la preuve de conformité (document, enregistrement, rapport
          d'audit)~;
    \item la date de la dernière vérification et son résultat.
\end{itemize}
Le registre est revu au moins une fois par an, et à chaque nouveau texte.
\end{boxprocedure}

\begin{boxmetier}{Auditeur SI}
\textbf{Auditer la conformité légale d'un SI}
\begin{enumerate}
    \item Demander le registre des obligations, et vérifier qu'il couvre au
          moins les textes du Tableau~10.1 applicables à l'organisation.
    \item Pour un échantillon d'obligations, demander la preuve de
          conformité.
    \item Vérifier la conservation des journaux (durée, intégrité) au
          regard de l'article~29 de la loi n°~2010/012.
    \item Vérifier que les contrats des sous-traitants traitent de la
          sécurité, de la localisation des données et de l'audit.
\end{enumerate}
\end{boxmetier}

% ============================================================================
\separateur
% ============================================================================

\begin{boxresume}
\textbf{Ce qu'il faut retenir de ce chapitre~:}
\begin{itemize}
    \item \textbf{La loi n°~2010/012} confie la sécurité électronique à
          l'ANTIC et impose aux exploitants de systèmes d'information des
          mesures de sécurité et des dispositifs normalisés de gestion des
          risques (art.~26), la conservation des données de connexion
          pendant dix ans (art.~29), la révision des dispositifs (art.~30)
          et des audits obligatoires (art.~13 et~32).
    \item \textbf{La responsabilité est pénale aussi}~: la personne morale
          répond des infractions commises par ses organes, sans exonérer
          les personnes physiques (art.~64).
    \item \textbf{La loi n°~2024/017}, applicable depuis juin~2026, organise
          la protection des données personnelles et crée une autorité
          dédiée~: inventaire, sécurité, sous-traitants, localisation,
          preuve.
    \item \textbf{Pour les banques}, la défaillance d'un système est un
          risque opérationnel au sens du règlement COBAC R-2016/04.
    \item \textbf{Le droit OHADA} organise la responsabilité civile des
          dirigeants et le comité d'audit des sociétés cotées~: la trace
          écrite des décisions est la première défense.
    \item \textbf{La Convention de Malabo} est en vigueur depuis 2023~;
          son applicabilité au Cameroun est à vérifier.
    \item \textbf{Ne citer que ce qu'on a vérifié}~: un rapport au
          régulateur se fonde sur le texte publié, pas sur un résumé.
\end{itemize}
\end{boxresume}

\bigskip

\begin{boxexo}[title={\ding{45}\quad Exercice~10.1. Qui est compétent~? \niveaubase}]
Pour chacune des situations suivantes, indiquez l'autorité principalement
compétente (Tableau~10.3) et le texte en cause.
\begin{enumerate}[(a)]
    \item Une banque n'a pas fait réaliser l'audit de sécurité de ses
          systèmes depuis trois ans.
    \item Un client de NOUN PAY conteste l'usage de ses données par le
          modèle de notation.
    \item Le dispositif de contrôle interne d'une banque n'a pas détecté
          une migration non autorisée.
    \item Un établissement public a acquis un logiciel hors de tout
          programme budgétaire.
\end{enumerate}
\end{boxexo}

\begin{boxexo}[title={\ding{45}\quad Exercice~10.2. Le registre de KRIBILOG \niveaumoyen}]
KRIBILOG SA exploite un système de suivi des conteneurs qui enregistre les
données des chauffeurs (nom, numéro de permis, téléphone) et les journaux
d'accès des clients au portail web. L'exploitation est confiée à NETPORT
SERVICES, qui héberge une partie des données à l'étranger.
\begin{enumerate}[(a)]
    \item Établissez les cinq premières lignes du registre des obligations
          légales de KRIBILOG, sur le modèle de l'exigence de contrôle de
          la section~\ref{sec:c10-erreurs}.
    \item Pour chacune, indiquez si le texte a pu être vérifié, et sinon,
          où le vérifier.
    \item Quelles clauses faut-il ajouter au contrat de NETPORT SERVICES~?
\end{enumerate}
\end{boxexo}

\begin{boxexo}[title={\ding{45}\quad Exercice~10.3. La note au conseil de la BANQUE SANAGA \niveauexpert}]
Rédigez, à l'attention du conseil d'administration, la note juridique
évoquée en ouverture du chapitre.
\begin{enumerate}[(a)]
    \item Qualifiez la panne au regard du règlement COBAC R-2016/04 et de la
          loi n°~2010/012, en citant uniquement les articles vérifiés.
    \item Identifiez les obligations probablement méconnues, en séparant
          constatations et analyses, et en donnant un niveau de confiance.
    \item Discutez l'exposition éventuelle des dirigeants au regard du
          droit OHADA, en signalant ce qui reste incertain.
    \item Proposez trois mesures immédiates pour réduire le risque
          juridique de la banque.
\end{enumerate}
\end{boxexo}

\bigskip

\begin{boxinfo}
\textbf{Transition.} Les obligations ne valent que si quelqu'un vérifie
qu'elles sont tenues. Le Chapitre~\ref{chap:audit} présente l'audit, le
contrôle et l'assurance de la gouvernance des SI~: le modèle des trois
lignes, les normes de l'audit interne, l'audit public selon l'INTOSAI, et
le test des contrôles généraux informatiques.
\end{boxinfo}

        % ============================================================================
% Fichier  : partie3/P03_C11_audit_controle_assurance.tex
% Titre    : Audit, Contrôle et Assurance de la Gouvernance des SI
% Auteur   : MINKA MI NGUIDJOÏ Thierry Emmanuel
% Version  : 1.0 -- octobre 2026
% ============================================================================
% Sources du dossier Normes : IIA, Cadre de référence international des
% pratiques professionnelles (édition 2017) et Code de déontologie ; INTOSAI
% GUID 5100 (2019) ; IDI, Manuel de mise en œuvre des ISSAI d'audit de
% conformité (2021). Les Normes mondiales de l'IIA (2024) sont présentées
% d'après leur structure publique. Laboratoire associé : GCI-LAB-05 (test des
% contrôles généraux informatiques à l'ANPP).
% ============================================================================

\chapter{Audit, Contrôle et Assurance de la Gouvernance des SI}
\label{chap:audit}
\chaptermark{Audit, contrôle et assurance}

\epigraph{%
    « Une réponse orale n'est pas une preuve. »%
}{Règle des auditeurs}

\niveauChapitre{Intermédiaire}{%
    Maîtrise des Chapitres~\ref{chap:cobit-capacite} (évaluation de
    capacité) et~\ref{chap:droit} (obligations légales).%
}

\begin{boxobjectifs}
\begin{itemize}
    \item Appliquer le modèle des trois lignes.
    \item Situer les Normes mondiales d'audit interne de l'IIA (2024) et le
          cadre de 2017.
    \item Mener un audit des SI selon l'INTOSAI GUID~5100 dans le secteur
          public.
    \item Tester des contrôles généraux informatiques (accès, changements,
          exploitation).
    \item Rédiger un constat d'audit qui sépare constatation, critère,
          cause, conséquence et recommandation.
    \item Identifier les exigences déontologiques de l'auditeur et du DSI.
\end{itemize}
\end{boxobjectifs}

\bigskip

% ============================================================================
\section*{Question directrice : qui vérifie que tout cela fonctionne~?}
% ============================================================================

Les chapitres précédents ont construit un système de gouvernance, fixé des
capacités cibles et dressé la liste des obligations légales. Il reste une
question~: \emph{qui vérifie, et avec quelle indépendance, que ce système
fonctionne réellement~?} À Yaoundé, la Chambre des comptes vient d'annoncer
à l'ANPP (opération MEFOU) un audit de ses systèmes d'information. À
Douala, Mme~EKAMBI, directrice de l'audit interne de la BANQUE SANAGA, doit
expliquer au comité d'audit pourquoi sa recommandation d'octobre~2025 est
restée sans suite.

Ce chapitre présente l'organisation de l'assurance (le modèle des trois
lignes, section~\ref{sec:c11-trois-lignes}), les normes de l'audit interne
(section~\ref{sec:c11-iia}) et de l'audit public
(section~\ref{sec:c11-intosai}), la pratique du test des contrôles
généraux informatiques (section~\ref{sec:c11-itgc}), la rédaction des
constats (section~\ref{sec:c11-constat}) et la déontologie
(section~\ref{sec:c11-deontologie}). Le laboratoire GCI-LAB-05 fait tester des
contrôles à partir des pièces de l'ANPP.

\separateur

% ============================================================================
\section{Le Modèle des Trois Lignes}
\label{sec:c11-trois-lignes}
% ============================================================================

\subsection{Qui fait quoi dans la maîtrise des risques}

L'Institut des auditeurs internes (\sigle{IIA}) a publié en 2020 le
\termecle{modèle des trois lignes}\index{modèle des trois lignes}
(\textit{Three Lines Model}), qui a succédé au modèle des « trois lignes de
défense ». Il répartit les rôles entre l'organe de gouvernance, le
management et l'audit interne~\cite{iia2020troislignes}.

\begin{tableaucours}{L{2.8cm}L{4.2cm}L{5.2cm}}{%
    \tableauHeader{Acteur} &
    \tableauHeader{Rôle} &
    \tableauHeader{À la BANQUE SANAGA}
}
Organe de gouvernance & Rend compte aux parties prenantes, fixe la
direction, supervise & Conseil d'administration, comité d'audit, comité
stratégique des SI \\
Première ligne & Le management opérationnel, qui porte et maîtrise les
risques dans ses activités & DSI, exploitation, chefs de projet, métiers \\
Deuxième ligne & Les fonctions d'expertise, de soutien et de surveillance
des risques & Gestion des risques, conformité, RSSI, responsable GRC \\
Troisième ligne & L'audit interne, qui fournit une assurance indépendante
et objective & Direction de l'audit interne (Mme~EKAMBI) \\
Prestataires d'assurance externes & Assurance complémentaire,
indépendante de l'organisation & Commissaire aux comptes, COBAC, audit de
sécurité prévu par la loi n°~2010/012 \\
\end{tableaucours}
\vspace{-0.8em}
\begin{center}{\small\itshape Tableau~11.1. Le modèle des trois lignes
appliqué à la BANQUE SANAGA}\end{center}

\subsection{Trois erreurs fréquentes}

\begin{enumerate}
    \item \textbf{Confier la deuxième ligne à la première.} Un RSSI
          rattaché au DSI, et noté par lui, ne peut pas surveiller librement
          les risques que le DSI prend. Le modèle ne prescrit pas un
          rattachement, mais il exige que la deuxième ligne puisse remonter
          ses alertes sans filtre.
    \item \textbf{Faire faire le travail de la première ligne par la
          troisième.} Un audit interne qui rédige les procédures qu'il
          auditera ensuite perd son indépendance.
    \item \textbf{Croire que l'audit est responsable du risque.} La
          troisième ligne donne une assurance~; elle ne maîtrise pas le
          risque. La recommandation non suivie de Mme~EKAMBI n'est pas une
          faute de l'audit~: c'est une défaillance du suivi par la première
          ligne et par le comité d'audit (MEA02).
\end{enumerate}

% ============================================================================
\section{Les Normes de l'Audit Interne}
\label{sec:c11-iia}
% ============================================================================

\subsection{Du cadre de 2017 aux Normes mondiales de 2024}

L'IIA définit l'audit interne comme une activité indépendante et objective
qui donne à une organisation une assurance sur le degré de maîtrise de ses
opérations, lui apporte des conseils pour les améliorer et contribue à
créer de la valeur, en évaluant de manière systématique ses processus de
gouvernance, de management des risques et de contrôle~\cite{iia2017cripp}.

Jusqu'en 2024, la profession se référait au \emph{Cadre de référence
international des pratiques professionnelles} (\sigle{CRIPP}), dans son
édition de 2017, disponible en français dans le dossier \texttt{Normes}.
Il comprenait notamment la définition de l'audit interne, le Code de
déontologie et des normes réparties entre normes de qualification (série
1000) et normes de fonctionnement (série 2000).

En janvier~2024, l'IIA a publié les \textbf{Normes mondiales d'audit
interne} (\textit{Global Internal Audit Standards}), applicables depuis
janvier~2025~\cite{iia2024gias}. Elles remplacent le CRIPP et
s'organisent en cinq domaines~:

\begin{tableaucours}{cL{4.0cm}L{8.0cm}}{%
    \tableauHeader{Dom.} &
    \tableauHeader{Intitulé} &
    \tableauHeader{Contenu}
}
I & Raison d'être de l'audit interne & Ce que l'audit interne apporte à
l'organisation \\
II & Éthique et professionnalisme & Intégrité, objectivité, compétence,
conscience professionnelle, confidentialité \\
III & Gouvernance de la fonction d'audit interne & Relation avec le
conseil, mandat, indépendance, supervision \\
IV & Management de la fonction d'audit interne & Stratégie, ressources,
méthodologie, qualité \\
V & Réalisation des missions & Planification, conduite, communication des
résultats, suivi \\
\end{tableaucours}
\vspace{-0.8em}
\begin{center}{\small\itshape Tableau~11.2. Les cinq domaines des Normes
mondiales d'audit interne (2024, notre traduction)}\end{center}

Le domaine~III est le plus directement lié à la gouvernance~: il place
expressément sous la responsabilité du conseil la supervision de la
fonction d'audit interne, son mandat et les moyens de son indépendance.

\begin{boxalert}
\textbf{Violation typique.} La charte d'audit interne de la BANQUE SANAGA
date de 2015 et renvoie aux « normes de l'IIA » sans préciser l'édition.
Le plan d'audit est approuvé par le directeur général, et non par le
comité d'audit.

\smallskip
\textbf{Règle~:} la charte d'audit doit se référer aux Normes mondiales de
2024, et le plan d'audit comme les moyens de la fonction doivent être
approuvés par l'organe de gouvernance, non par la direction auditée.
\end{boxalert}

% ============================================================================
\section{L'Audit Public des SI : l'INTOSAI GUID 5100}
\label{sec:c11-intosai}
% ============================================================================

\subsection{Les institutions supérieures de contrôle}

Dans le secteur public, l'audit externe relève des institutions
supérieures de contrôle des finances publiques (\sigle{ISC}), qui
appliquent les normes de l'INTOSAI (les \sigle{ISSAI}). Au Cameroun, cette
fonction est assurée par la Chambre des comptes de la Cour suprême. Le
GUID~5100 de l'INTOSAI (2019) fournit le cadre général de la vérification
des systèmes d'information~\cite{intosai2019guid5100}.

\begin{boxdef}
\termecle{Vérification des systèmes d'information}\index{vérification des
systèmes d'information} (\textit{information systems audit})~: selon le
GUID~5100, examen des contrôles liés aux systèmes d'information pour
déceler les écarts par rapport à des critères, définis selon le type de
mission (audit financier, de conformité ou de performance).

\medskip
\termecle{Contrôles généraux}\index{contrôles généraux}~: procédures
manuelles ou automatisées qui visent à assurer la confidentialité,
l'intégrité et la disponibilité de l'information dans l'environnement où
les systèmes sont développés, gérés et exploités.

\medskip
\termecle{Contrôles d'application}\index{contrôles d'application}~:
procédures, dépendantes des TI, intégrées à un système et qui agissent sur
le traitement des transactions~: validation des données d'entrée,
exactitude du traitement, contrôle des sorties, intégrité des données de
référence.
\end{boxdef}

\subsection{La démarche du GUID 5100}

Le GUID~5100 suit les quatre étapes de tout audit public~: planification,
réalisation, rapport et suivi. Trois de ses apports sont particulièrement
utiles.

\paragraph{Les trois risques d'audit.} L'auditeur apprécie le
\emph{risque inhérent} (la probabilité que les caractéristiques mêmes d'un
système nuisent à la fonction qu'il remplit), le \emph{risque de contrôle}
(la probabilité que les contrôles mis en place ne suffisent pas) et le
\emph{risque de détection} (la probabilité que l'auditeur lui-même ne
décèle pas une défaillance). Le portail de paiement de l'ANPP, ouvert à
tous les fournisseurs de l'État, porte un risque inhérent élevé de
saturation.

\paragraph{Treize domaines.} Le GUID~5100 propose de délimiter la
vérification parmi treize domaines, dont la plupart sont adaptés de
l'ISO/IEC~27001~: politique de l'organisation en matière de TI,
structures de gouvernance, contrôles généraux, gestion des actifs,
développement et maintenance, exploitation, environnement physique,
ressources humaines, communications, sécurité de l'information,
conformité réglementaire, continuité et reprise, contrôles
d'application.

\paragraph{Conception, mise en œuvre, efficacité.} Le GUID~5100 distingue
trois niveaux d'examen d'un contrôle. Pour la \emph{conception}, un
entretien ou l'examen des règles documentées peut suffire. Pour la
\emph{mise en œuvre}, l'entretien ne suffit plus~: il faut une revue
structurée ou une analyse de données. Pour l'\emph{efficacité
opérationnelle}, il faut tester un échantillon d'opérations sur toute la
période.

\begin{boxprocedure}
\textbf{Recevoir une extraction de données.} Le GUID~5100 recommande que
toute extraction de données remise par l'entité auditée soit accompagnée
d'une lettre qui précise~:
\begin{itemize}
    \item la source des données et leur horodatage, avec l'empreinte
          numérique (valeur de hachage) du fichier, pour en garantir
          l'intégrité~;
    \item les paramètres d'extraction (requêtes utilisées, états
          exécutés)~;
    \item l'environnement d'origine (production ou autre).
\end{itemize}
À défaut de lettre, l'auditeur documente lui-même ces informations.
\end{boxprocedure}

\begin{boxscenario}{Opération MEFOU. La lettre de mission de la Chambre des comptes}
\textbf{Mandat.} Audit de conformité et de performance des systèmes
d'information de l'ANPP, période 2023--2026.

\medskip
\textbf{Délimitation retenue} (parmi les treize domaines)~: structures de
gouvernance des TI, développement et acquisition (le schéma directeur de
14~milliards FCFA), contrôles généraux (accès et changements) et
conformité réglementaire.

\medskip
\textbf{Critères.} Le régime financier de l'État (loi n°~2018/012), la loi
n°~2010/012, les procédures internes de l'agence et, comme référentiel de
bonnes pratiques, COBIT~2019.

\medskip
\textbf{Équipe.} Le magistrat vérificateur, M.~NJOYA, est assisté d'un
auditeur des SI certifié, recruté sur contrat, conformément à la
possibilité qu'ouvre le GUID~5100 de recourir à des ressources externes
supervisées.
\end{boxscenario}

% ============================================================================
\section{Tester les Contrôles Généraux Informatiques}
\label{sec:c11-itgc}
% ============================================================================

\subsection{Trois familles de contrôles}

Les contrôles généraux informatiques (Chapitre~\ref{chap:gouv-entreprise})
conditionnent la fiabilité de toutes les applications. L'auditeur les
teste en priorité, parce que leur défaillance invalide la confiance que
l'on peut accorder aux contrôles automatisés des applications~: si
n'importe qui peut modifier un programme, le contrôle que ce programme
réalise ne vaut plus rien.

\begin{tableaucours}{L{2.6cm}L{4.6cm}L{5.0cm}}{%
    \tableauHeader{Famille} &
    \tableauHeader{Contrôles typiques} &
    \tableauHeader{Test typique}
}
Accès & Création et suppression des comptes, revue périodique des droits,
comptes à privilèges & Rapprocher la liste des comptes actifs de la liste
des agents présents \\
Changements & Autorisation, test, séparation entre développement et
production, mise en production tracée & Vérifier, sur un échantillon de
changements, l'autorisation préalable et le test \\
Exploitation & Sauvegardes et tests de restauration, ordonnancement des
traitements, gestion des incidents & Vérifier les journaux de sauvegarde
et la preuve d'une restauration réussie \\
\end{tableaucours}
\vspace{-0.8em}
\begin{center}{\small\itshape Tableau~11.3. Les trois familles de contrôles
généraux informatiques}\end{center}

\subsection{L'échantillonnage}

On ne teste pas tous les comptes ni tous les changements~: on
échantillonne. La taille de l'échantillon dépend de la fréquence du
contrôle et du risque. Les cabinets d'audit utilisent couramment des
ordres de grandeur comme ceux du tableau suivant~; ce ne sont pas des
exigences normatives, et chaque organisation d'audit fixe sa propre
méthode.

\begin{tableaucours}{L{4.6cm}L{4.6cm}}{%
    \tableauHeader{Fréquence du contrôle} &
    \tableauHeader{Échantillon indicatif}
}
Annuelle & 1 \\
Trimestrielle & 2 \\
Mensuelle & 2 à 5 \\
Hebdomadaire & 5 à 15 \\
Quotidienne & 20 à 40 \\
Plusieurs fois par jour & 25 à 60 \\
\end{tableaucours}
\vspace{-0.8em}
\begin{center}{\small\itshape Tableau~11.4. Ordres de grandeur
d'échantillonnage pratiqués en audit (indicatifs)}\end{center}

\begin{boxscenario}{Opération MEFOU. Le test des accès}
\textbf{Test.} L'auditeur rapproche la liste des 312~comptes actifs du
système de paiement de l'ANPP et la liste du personnel au 30~juin~2026
fournie par les ressources humaines. Il examine en outre un échantillon de
25~comptes pour vérifier que les droits correspondent aux fonctions.

\medskip
\textbf{Résultats.} 14~comptes appartiennent à des agents partis de
l'agence, dont 3 depuis plus d'un an~; l'un de ces comptes a été utilisé
en mai~2026. Sur les 25~comptes examinés, 6 cumulent les droits de saisie
et de validation d'un paiement. Aucune revue des droits n'a été
documentée depuis 2023.

\medskip
\textbf{Qualification.} Le contrôle « revue périodique des accès » n'est
pas mis en œuvre~; le contrôle « séparation des tâches » est défaillant.
Ces écarts relèvent des objectifs DSS05 et DSS06 de COBIT et de
l'obligation de dispositifs de sécurité de l'article~26 de la loi
n°~2010/012 (Chapitre~\ref{chap:droit}).
\end{boxscenario}

\begin{boxmetier}{Auditeur SI}
\textbf{Tester un contrôle en cinq temps}
\begin{enumerate}
    \item Décrire le contrôle tel qu'il est censé fonctionner (qui, quoi,
          quand, avec quelle preuve).
    \item Vérifier sa conception~: le contrôle, s'il fonctionne, traite-t-il
          le risque~?
    \item Vérifier sa mise en œuvre sur un cas (test de cheminement).
    \item Tester son efficacité sur un échantillon couvrant la période.
    \item Documenter chaque pièce examinée, et conclure~: efficace,
          partiellement efficace, inefficace.
\end{enumerate}
\end{boxmetier}

% ============================================================================
\section{Rédiger un Constat d'Audit}
\label{sec:c11-constat}
% ============================================================================

\subsection{Les cinq éléments}

Un constat d'audit n'est pas une opinion~: c'est un raisonnement en cinq
éléments, que l'on retrouve dans les normes de l'IIA comme dans celles de
l'INTOSAI.

\begin{tableaucours}{L{2.8cm}L{4.6cm}L{4.8cm}}{%
    \tableauHeader{Élément} &
    \tableauHeader{Question} &
    \tableauHeader{Exemple (ANPP)}
}
Constatation & Que se passe-t-il réellement~? & 14 comptes d'agents partis
restent actifs~; l'un a été utilisé en mai~2026 \\
Critère & Que devrait-il se passer~? & Procédure interne~: désactivation
sous 48~heures~; loi n°~2010/012, art.~26 \\
Cause & Pourquoi l'écart existe-t-il~? & Les départs ne sont pas notifiés
par les RH à la division informatique \\
Conséquence & Quel risque ou quel effet~? & Paiement frauduleux possible
au nom d'un agent parti~; exposition juridique \\
Recommandation & Que faut-il faire~? & Notification automatique des
départs, revue trimestrielle des accès \\
\end{tableaucours}
\vspace{-0.8em}
\begin{center}{\small\itshape Tableau~11.5. Les cinq éléments d'un constat
d'audit}\end{center}

\begin{boxalert}
\textbf{Violation typique.} « La sécurité des accès est insuffisante et
doit être renforcée. » Cette phrase ne contient ni constatation chiffrée,
ni critère, ni cause~: elle ne peut être ni vérifiée ni mise en œuvre, et
l'audité peut la contester sans peine.

\smallskip
\textbf{Règle~:} chaque constat chiffre la constatation, cite son critère,
identifie la cause, mesure la conséquence, et formule une recommandation
qui traite la \emph{cause}, et non le symptôme.
\end{boxalert}

\subsection{Le suivi des recommandations}

Une recommandation non suivie est une assurance perdue. Le suivi relève de
la première ligne (qui met en œuvre), de l'audit interne (qui vérifie) et
de l'organe de gouvernance (qui exige). C'est l'objectif MEA02 de COBIT.

\begin{boxprocedure}
\textbf{Le tableau de suivi des recommandations.} Pour chaque
recommandation, le tableau présenté au comité d'audit indique~:
\begin{itemize}
    \item la date du rapport et la priorité de la recommandation~;
    \item le responsable de la mise en œuvre et l'échéance convenue~;
    \item l'état (non commencée, en cours, mise en œuvre, vérifiée)~;
    \item pour l'état « mise en œuvre », la pièce qui le prouve~;
    \item les reports d'échéance, avec leur motif et leur approbation.
\end{itemize}
La mention « en cours » sans échéance ni pièce, qui figurait dans le
rapport de la BANQUE SANAGA d'octobre~2025, est précisément ce qu'il faut
proscrire.
\end{boxprocedure}

% ============================================================================
\section{La Déontologie}
\label{sec:c11-deontologie}
% ============================================================================

\subsection{L'auditeur}

Le Code de déontologie de l'IIA (édition 2017) pose quatre principes, que
les Normes mondiales de 2024 reprennent dans leur domaine~II en y ajoutant
la conscience professionnelle~\cite{iia2017cripp}.

\begin{tableaucours}{L{2.8cm}L{9.4cm}}{%
    \tableauHeader{Principe} &
    \tableauHeader{Exigence (résumé)}
}
Intégrité & Elle fonde la confiance accordée au jugement de l'auditeur~:
honnêteté, diligence, responsabilité. \\
Objectivité & L'auditeur évalue équitablement tous les éléments
pertinents, sans se laisser influencer par ses intérêts ou par autrui. \\
Confidentialité & Il respecte la valeur et la propriété des informations
reçues, et ne les divulgue qu'avec autorisation, sauf obligation légale
ou professionnelle. \\
Compétence & Il met en œuvre les connaissances, les savoir-faire et
l'expérience que requiert sa mission. \\
\end{tableaucours}
\vspace{-0.8em}
\begin{center}{\small\itshape Tableau~11.6. Les principes du Code de
déontologie de l'IIA (résumé d'après l'édition française de 2017)}\end{center}

\subsection{Le DSI}

Le DSI n'est pas soumis à un code de déontologie professionnel
obligatoire, mais sa position l'expose à des situations délicates~: il
choisit des fournisseurs et dispose d'accès étendus aux données. Trois
règles de conduite en découlent.
\begin{itemize}
    \item \textbf{Déclarer ses conflits d'intérêts}, notamment dans le choix
          des prestataires, et s'abstenir dans les décisions concernées.
    \item \textbf{Rendre compte loyalement}, y compris des mauvaises
          nouvelles~: la note de la DSI de la BANQUE SANAGA du
          20~février~2026 en est l'exemple.
    \item \textbf{Ne pas faire obstacle à l'audit}~: donner accès aux pièces
          et aux systèmes, et ne pas « préparer » les données remises.
\end{itemize}

\begin{boiteRemarque}{L'auditeur peut-il conseiller~?}
Oui, et les normes de l'IIA le prévoient expressément~: l'audit interne
donne une assurance \emph{et} des conseils. Mais il y a une limite. Un
auditeur qui conçoit un dispositif ne peut pas, ensuite, en donner une
assurance objective. Dans une petite organisation, où l'audit interne
compte une ou deux personnes, cette tension est permanente. La solution
consiste à déclarer la situation au comité d'audit, et à faire réaliser
l'évaluation du dispositif par un tiers.
\end{boiteRemarque}

% ============================================================================
\separateur
% ============================================================================

\begin{boxresume}
\textbf{Ce qu'il faut retenir de ce chapitre~:}
\begin{itemize}
    \item \textbf{Trois lignes}~: le management porte le risque, les
          fonctions de deuxième ligne le surveillent, l'audit interne donne
          une assurance indépendante, sous la supervision de l'organe de
          gouvernance.
    \item \textbf{Les Normes mondiales de l'IIA (2024)} remplacent le cadre
          de 2017~: cinq domaines, dont la gouvernance de la fonction d'audit
          par le conseil.
    \item \textbf{Le GUID~5100 de l'INTOSAI} guide l'audit public des SI~:
          trois risques d'audit, treize domaines, et la distinction entre
          conception, mise en œuvre et efficacité des contrôles.
    \item \textbf{Les contrôles généraux informatiques} (accès,
          changements, exploitation) se testent en priorité, par
          échantillonnage.
    \item \textbf{Un constat a cinq éléments}~: constatation, critère,
          cause, conséquence, recommandation. La recommandation traite la
          cause.
    \item \textbf{Le suivi est une responsabilité de gouvernance}~: une
          recommandation « en cours » sans échéance ni pièce est une
          assurance perdue.
    \item \textbf{La déontologie} repose sur l'intégrité, l'objectivité, la
          confidentialité et la compétence~; elle concerne aussi le DSI.
\end{itemize}
\end{boxresume}

\bigskip

\begin{boxexo}[title={\ding{45}\quad Exercice~11.1. Quelle ligne~? \niveaubase}]
Rattachez chacun des acteurs suivants à la ligne qui convient (première,
deuxième, troisième, ou prestataire d'assurance externe), et justifiez.
\begin{enumerate}[(a)]
    \item Le responsable de l'exploitation de la BANQUE SANAGA.
    \item Le RSSI, rattaché au directeur des risques.
    \item Le commissaire aux comptes.
    \item La directrice de l'audit interne.
    \item Le cabinet agréé qui réalise l'audit de sécurité prévu par la loi
          n°~2010/012.
\end{enumerate}
\end{boxexo}

\begin{boxexo}[title={\ding{45}\quad Exercice~11.2. Rédiger un constat \niveaumoyen}]
À partir des résultats du test des accès de l'ANPP (section~\ref{sec:c11-itgc}),
rédigez le constat relatif à la séparation des tâches (6~comptes sur
25 cumulant saisie et validation d'un paiement).
\begin{enumerate}[(a)]
    \item Rédigez les cinq éléments du constat.
    \item Proposez une priorité (élevée, moyenne, faible) et justifiez.
    \item Proposez l'indicateur qui permettra au comité d'audit de suivre la
          mise en œuvre.
\end{enumerate}
\end{boxexo}

\begin{boxexo}[title={\ding{45}\quad Exercice~11.3. Planifier l'audit des changements de la BANQUE SANAGA \niveauexpert}]
Le comité d'audit demande à Mme~EKAMBI un audit de la gestion des
changements après la panne.
\begin{enumerate}[(a)]
    \item Rédigez l'objectif, le périmètre et les critères de la mission.
    \item Listez les contrôles à tester, en distinguant conception, mise en
          œuvre et efficacité, et indiquez pour chacun la pièce attendue.
    \item Déterminez la taille de l'échantillon de changements à examiner,
          en justifiant au regard de la fréquence des changements.
    \item Mme~EKAMBI avait elle-même proposé, en 2025, la procédure de
          gestion des changements. Quelle difficulté déontologique cela
          pose-t-il, et comment la traiter~?
\end{enumerate}
\end{boxexo}

\bigskip

\begin{boxinfo}
\textbf{Transition.} La Partie~\ref{part:normatif} a placé la gouvernance
des SI sous le regard des normes, de la loi et de l'audit. La
Partie~\ref{part:pratique} descend au niveau de la pratique quotidienne.
Le Chapitre~\ref{chap:itil-svs} ouvre la gestion des services avec ITIL~4
et son système de valeur des services.
\end{boxinfo}


    % ======================================================================
    % PARTIE IV — ITIL 4 ET PRATIQUE OPÉRATIONNELLE
    % « La valeur d'un service se mesure chez celui qui le reçoit. »
    % Standards : ITIL 4 (sources publiques), avec une note sur ITIL
    % version 5 (PeopleCert, 2026), ISO/IEC 20000-1, COBIT 2019 (BAI, DSS).
    % Niveau : intermédiaire à expert
    % ======================================================================
    \part{ITIL 4 et Pratique Opérationnelle}
    \label{part:pratique}

        % ============================================================================
% Fichier  : partie4/P04_C12_itil_systeme_valeur.tex
% Titre    : ITIL 4 : le Système de Valeur des Services
% Auteur   : MINKA MI NGUIDJOÏ Thierry Emmanuel
% Version  : 1.0 -- octobre 2026
% ============================================================================
% Décision de l'auteur (29/10/2026) : la Partie IV s'appuie sur ITIL 4,
% avec une note sur ITIL version 5 (PeopleCert, 2026). Sources publiques
% uniquement : les définitions sont reformulées, aucune n'est citée
% textuellement.
% ============================================================================

\chapter{ITIL 4 : le Système de Valeur des Services}
\label{chap:itil-svs}
\chaptermark{ITIL 4 : le système de valeur}

\epigraph{%
    « La valeur ne se livre pas.\\
    Elle se crée avec celui qui reçoit le service. »%
}{Idée centrale d'ITIL~4}

\niveauChapitre{Intermédiaire}{%
    Maîtrise du Chapitre~\ref{chap:cobit-archi} (architecture de COBIT~2019)
    et du Chapitre~\ref{chap:referentiels} (place d'ITIL parmi les
    référentiels).%
}

\begin{boxobjectifs}
\begin{itemize}
    \item Retracer l'évolution d'ITIL (v2, v3 et son cycle de vie, ITIL~4).
    \item Définir service, valeur, produit et relation de service.
    \item Décrire le système de valeur des services et sa chaîne de valeur.
    \item Appliquer les sept principes directeurs et les quatre dimensions.
    \item Situer ITIL version~5 (2026) par rapport à ITIL~4.
\end{itemize}
\end{boxobjectifs}

\bigskip

% ============================================================================
\section*{Question directrice : que veut vraiment le client~?}
% ============================================================================

Chez KRIBILOG SA, le centre de services traite 4~800~tickets par mois, et
30~\% d'entre eux sont rouverts. Les indicateurs de l'infogérant NETPORT
SERVICES sont pourtant au vert~: les serveurs sont disponibles à
99,8~\%. Le directeur des opérations portuaires résume le malentendu~:
« Vos serveurs fonctionnent, mais mes conteneurs restent bloqués. »

Ce malentendu est l'objet d'ITIL~4. La Partie~\ref{part:pratique}, qui
s'ouvre ici, passe de la gouvernance à la pratique quotidienne des
services. Ce chapitre en pose les fondations~: l'histoire d'ITIL
(section~\ref{sec:c12-histoire}), ses concepts clés
(section~\ref{sec:c12-concepts}), le système de valeur des services et sa
chaîne de valeur (sections~\ref{sec:c12-svs} et~\ref{sec:c12-chaine}), les
quatre dimensions et les sept principes (sections~\ref{sec:c12-dimensions}
et~\ref{sec:c12-principes}), et enfin la place d'ITIL version~5
(section~\ref{sec:c12-itil5}).

\separateur

% ============================================================================
\section{Trente-Cinq Ans d'ITIL}
\label{sec:c12-histoire}
% ============================================================================

\begin{tableaucours}{L{2.4cm}L{4.0cm}L{5.8cm}}{%
    \tableauHeader{Version} &
    \tableauHeader{Logique} &
    \tableauHeader{Apport}
}
Fin des années 1980 & Recueil de pratiques & L'administration britannique
rassemble les bonnes pratiques de gestion de ses services informatiques \\
ITIL v2 (2000--2001) & Processus & Deux ouvrages centraux~: le soutien et la
fourniture des services~; diffusion mondiale \\
ITIL v3 (2007, révisé en 2011) & Cycle de vie & Cinq étapes~: stratégie,
conception, transition, exploitation et amélioration continue des
services \\
ITIL 4 (2019) & Système de valeur & Co-création de la valeur, 34~pratiques,
intégration de l'agilité, de DevOps et du \textit{lean} \\
ITIL version 5 (2026) & Produits et services & Voir la
section~\ref{sec:c12-itil5} \\
\end{tableaucours}
\vspace{-0.8em}
\begin{center}{\small\itshape Tableau~12.1. Évolution d'ITIL}\end{center}

Le passage d'ITIL~v3 à ITIL~4 est le plus profond. ITIL~v3 décrivait
\emph{comment le fournisseur gère ses services} au long de leur cycle de
vie. ITIL~4 décrit \emph{comment la valeur se crée entre le fournisseur et
ses clients}, et intègre les méthodes agiles et DevOps qui s'étaient
développées en dehors d'ITIL. Le mot « processus » cède la place au mot
« pratique », plus large~: une pratique comprend des processus, mais aussi
des personnes, des informations et des outils, comme les composantes de
COBIT (Chapitre~\ref{chap:cobit-archi}).

% ============================================================================
\section{Les Concepts Clés}
\label{sec:c12-concepts}
% ============================================================================

\subsection{Service et valeur}

\begin{boxdef}
\termecle{Service}\index{service (ITIL)}~: moyen de permettre la
co-création de valeur, en facilitant l'obtention des résultats que le
client veut atteindre, sans qu'il ait à gérer lui-même certains coûts et
certains risques.

\medskip
\termecle{Valeur}\index{valeur (ITIL)}~: bénéfices, utilité et importance
perçus de quelque chose. La valeur est \emph{perçue} par celui qui la
reçoit~: elle ne se décrète pas chez le fournisseur.

\medskip
\termecle{Co-création de valeur}\index{co-création de valeur}~: idée selon
laquelle la valeur naît de la collaboration entre le fournisseur et le
consommateur du service, et non de la seule action du fournisseur.
\end{boxdef}
{\small\itshape Définitions reformulées d'après ITIL~4~\cite{axelos2019itil4}.}

\subsection{Résultats, produits et relations}

ITIL~4 distingue soigneusement ce que le fournisseur \emph{produit} de ce
que le client \emph{obtient}.

\begin{tableaucours}{L{3.2cm}L{5.0cm}L{4.0cm}}{%
    \tableauHeader{Notion} &
    \tableauHeader{Sens} &
    \tableauHeader{Exemple KRIBILOG}
}
Livrable (\textit{output}) & Ce qu'une activité produit & Un écran de suivi
affiché, un ticket clos \\
Résultat (\textit{outcome}) & Ce que le client obtient grâce aux livrables
& Un conteneur localisé et chargé à l'heure \\
Produit & Configuration de ressources conçue pour offrir de la valeur &
Le logiciel de suivi, l'infrastructure, l'équipe de support \\
Offre de service & Description d'un ou plusieurs services, pour un groupe
de clients & « Suivi des conteneurs 24~h/24, support en 30~minutes » \\
Relation de service & Coopération entre fournisseur et consommateur &
Fourniture, consommation et gestion du service \\
\end{tableaucours}
\vspace{-0.8em}
\begin{center}{\small\itshape Tableau~12.2. Les notions clés d'ITIL~4
(reformulées)}\end{center}

Le malentendu de KRIBILOG s'éclaire~: NETPORT SERVICES mesure des
\emph{livrables} (serveurs disponibles, tickets clos), le directeur des
opérations attend un \emph{résultat} (conteneurs chargés). Un ticket clos
qui est rouvert le lendemain est un livrable sans résultat.

\subsection{Utilité et garantie}

\begin{boxdef}
\termecle{Utilité}\index{utilité} (\textit{utility})~: fonctionnalité
offerte par un service pour répondre à un besoin~: ce que fait le
service.

\medskip
\termecle{Garantie}\index{garantie} (\textit{warranty})~: assurance que le
service répondra aux exigences convenues~: disponibilité, capacité,
sécurité, continuité. C'est \emph{comment} le service le fait.
\end{boxdef}

Un service n'a de valeur que s'il a les deux~: une utilité sans garantie
(un logiciel parfait, mais indisponible un jour sur deux) n'en a pas
davantage qu'une garantie sans utilité.

% ============================================================================
\section{Le Système de Valeur des Services}
\label{sec:c12-svs}
% ============================================================================

\subsection{Cinq composantes}

Le cœur d'ITIL~4 est le \termecle{système de valeur des services}\index{système
de valeur des services} (\textit{service value system}, \sigle{SVS}). Il
décrit comment les composantes et les activités d'une organisation
fonctionnent ensemble pour transformer une \emph{opportunité} ou une
\emph{demande} en \emph{valeur}.

\noindent\begin{minipage}{\linewidth}
\begin{center}
\begin{tikzpicture}[
    entree/.style = {
        draw=CouleurAccent, fill=CouleurAccent!15, rounded corners=4pt,
        text width=2.2cm, minimum height=1.0cm, align=center,
        font=\footnotesize\bfseries\color{CouleurPrimaire}
    },
    comp/.style = {
        draw=CouleurBordDef, fill=CouleurDef, rounded corners=3pt,
        text width=3.6cm, minimum height=0.6cm, align=center,
        font=\scriptsize\bfseries\color{CouleurPrimaire}
    },
    cadre/.style = {
        draw=CouleurPrimaire, thick, rounded corners=6pt
    },
    fleche/.style = {-{Stealth[length=5pt]}, thick, color=CouleurSecondaire}
]
\node[entree] (op) at (0,0)   {Opportunité\\ et demande};
\node[comp] (p) at (4.2,1.3)  {Principes directeurs};
\node[comp] (g) at (4.2,0.6)  {Gouvernance};
\node[comp, fill=CouleurPrimaire, text=white] (c) at (4.2,-0.1)
    {Chaîne de valeur des services};
\node[comp] (pr) at (4.2,-0.8) {Pratiques};
\node[comp] (am) at (4.2,-1.5) {Amélioration continue};
\node[cadre, fit=(p)(am), inner sep=6pt] (svs) {};
\node[font=\scriptsize\bfseries\color{CouleurAccent}, above=0.05cm of svs]
    {Système de valeur des services};
\node[entree] (va) at (8.4,0) {Valeur};
\draw[fleche] (op) -- (svs.west |- op);
\draw[fleche] (svs.east |- va) -- (va);
\end{tikzpicture}
\end{center}
\vspace{-0.4em}
\begin{center}{\small\itshape Figure~12.1. Le système de valeur des
services d'ITIL~4 (représentation simplifiée)}\end{center}
\end{minipage}

\medskip

\begin{tableaucours}{L{3.6cm}L{8.6cm}}{%
    \tableauHeader{Composante} &
    \tableauHeader{Rôle}
}
Principes directeurs & Recommandations qui guident les décisions en toutes
circonstances (section~\ref{sec:c12-principes}) \\
Gouvernance & Moyens par lesquels l'organisation est dirigée et
contrôlée~: évaluer, diriger, surveiller \\
Chaîne de valeur des services & Modèle des six activités qui répondent à
la demande et créent la valeur (section~\ref{sec:c12-chaine}) \\
Pratiques & Ensembles de ressources organisationnelles conçus pour
accomplir un travail ou atteindre un objectif (34~pratiques) \\
Amélioration continue & Activité récurrente, à tous les niveaux, qui
assure que la performance répond aux attentes \\
\end{tableaucours}
\vspace{-0.8em}
\begin{center}{\small\itshape Tableau~12.3. Les composantes du système de
valeur des services}\end{center}

\subsection{ITIL et la gouvernance}

La deuxième composante est la \emph{gouvernance}. ITIL~4 y reprend le
modèle de l'ISO/IEC~38500~: évaluer, diriger, surveiller
(Chapitre~\ref{chap:normes}). C'est le point d'ancrage d'ITIL dans la
gouvernance des SI~: ITIL ne prétend pas organiser la gouvernance, mais il
s'y soumet. Le système de valeur des services d'une organisation doit être
dirigé et contrôlé par son organe de gouvernance, celui que COBIT organise
avec le domaine EDM.

% ============================================================================
\section{La Chaîne de Valeur des Services}
\label{sec:c12-chaine}
% ============================================================================

\subsection{Six activités}

\begin{tableaucours}{L{3.0cm}L{9.2cm}}{%
    \tableauHeader{Activité} &
    \tableauHeader{Finalité}
}
Planifier & Partager une vision commune de l'état actuel, des orientations
et des améliorations, pour tous les produits et services \\
Améliorer & Améliorer en continu les produits, les services et les
pratiques \\
S'engager & Comprendre les besoins des parties prenantes, entretenir la
transparence et de bonnes relations \\
Concevoir et faire la transition & Faire en sorte que les produits et
services répondent aux attentes de qualité, de coût et de délai \\
Obtenir ou construire & Disposer des composants de service nécessaires,
quand et où il le faut \\
Fournir et soutenir & Livrer et soutenir les services selon les
spécifications et les attentes convenues \\
\end{tableaucours}
\vspace{-0.8em}
\begin{center}{\small\itshape Tableau~12.4. Les six activités de la chaîne
de valeur des services (finalités reformulées)}\end{center}

\subsection{Des chaînes de valeur, pas un ordre fixe}

Les six activités ne forment pas un cycle à parcourir dans l'ordre, comme
les étapes d'ITIL~v3. Elles se combinent, dans un ordre propre à chaque
situation, en \termecle{flux de valeur}\index{flux de valeur}
(\textit{value streams})~: la suite d'étapes qu'une organisation enchaîne
pour créer un résultat donné pour un client. Une même pratique peut
intervenir dans plusieurs activités.

\begin{boxscenario}{Opération KRIBI. Conteneur bloqué}
\textbf{Déclencheur.} Un agent du port signale qu'un conteneur n'apparaît
plus dans le logiciel de suivi.

\medskip
\textbf{Flux de valeur.}
\begin{enumerate}
    \item \emph{S'engager}~: le centre de services enregistre l'appel et
          comprend l'impact (un navire attend).
    \item \emph{Fournir et soutenir}~: le support diagnostique, contourne
          le problème et relocalise le conteneur.
    \item \emph{Obtenir ou construire}~: l'éditeur corrige l'anomalie
          logicielle.
    \item \emph{Concevoir et faire la transition}~: la correction est
          testée, autorisée, puis déployée.
    \item \emph{Améliorer}~: l'analyse de cause alimente le registre des
          problèmes.
\end{enumerate}

\textbf{Constatation.} Chez KRIBILOG, le flux s'arrête aujourd'hui à
l'étape~2~: le ticket est clos dès le contournement, sans correction de
la cause. Il est rouvert lors du blocage suivant.

\medskip
\textbf{Analyse} (confiance élevée). Le taux de réouverture de 30~\%
s'explique par l'absence des étapes~3 à~5 dans le flux réel~; les
pratiques de gestion des problèmes et d'habilitation des changements
(Chapitres~\ref{chap:itil-incidents} et~\ref{chap:itil-changements}) sont
les premières à mettre en place.
\end{boxscenario}

% ============================================================================
\section{Les Quatre Dimensions}
\label{sec:c12-dimensions}
% ============================================================================

ITIL~4 rappelle qu'un service ne se réduit pas à sa technique. Toute
démarche doit tenir compte de quatre dimensions, sous l'influence de
facteurs externes (politiques, économiques, sociaux, technologiques,
juridiques et environnementaux).

\begin{tableaucours}{L{3.6cm}L{4.4cm}L{4.2cm}}{%
    \tableauHeader{Dimension} &
    \tableauHeader{Contenu} &
    \tableauHeader{Chez KRIBILOG}
}
Organisations et personnes & Structures, rôles, compétences, culture &
Centre de services sous-dimensionné, pas de responsable des problèmes \\
Information et technologie & Données, applications, infrastructures,
outils & Logiciel de suivi, outil de tickets, base de connaissances
absente \\
Partenaires et fournisseurs & Relations avec les autres organisations qui
contribuent au service & NETPORT SERVICES, l'éditeur du logiciel \\
Flux de valeur et processus & Enchaînement des activités et des
processus & Flux « conteneur bloqué » interrompu \\
\end{tableaucours}
\vspace{-0.8em}
\begin{center}{\small\itshape Tableau~12.5. Les quatre dimensions de la
gestion des services}\end{center}

\begin{boxalert}
\textbf{Violation typique.} Un projet ITIL est réduit à l'achat d'un outil
de gestion des tickets. L'outil est installé, mais les rôles ne sont pas
définis, l'infogérant n'est pas tenu par le contrat d'y saisir ses
interventions, et aucun flux de valeur n'est décrit.

\smallskip
\textbf{Règle~:} une démarche ITIL traite toujours les quatre dimensions.
Un outil ne couvre qu'une partie de la deuxième.
\end{boxalert}

% ============================================================================
\section{Les Sept Principes Directeurs}
\label{sec:c12-principes}
% ============================================================================

\begin{tableaucours}{L{3.8cm}L{8.4cm}}{%
    \tableauHeader{Principe} &
    \tableauHeader{Application}
}
Privilégier la valeur & Tout ce que l'on fait doit, directement ou
indirectement, créer de la valeur pour les parties prenantes \\
Commencer là où l'on est & Partir de l'existant, l'observer et le
mesurer, plutôt que de tout reconstruire \\
Progresser par itérations, avec retour d'information & Avancer par
étapes courtes, en recueillant les retours à chaque étape \\
Collaborer et promouvoir la visibilité & Associer les bonnes personnes,
rendre le travail et ses résultats visibles \\
Penser et travailler de manière globale & Traiter le service comme un
tout, avec ses quatre dimensions \\
Rester simple et pratique & Supprimer ce qui ne crée pas de valeur,
préférer le nombre minimal d'étapes \\
Optimiser et automatiser & Optimiser d'abord, automatiser ensuite~; ne pas
automatiser un processus défaillant \\
\end{tableaucours}
\vspace{-0.8em}
\begin{center}{\small\itshape Tableau~12.6. Les sept principes directeurs
d'ITIL~4 (notre traduction)}\end{center}

\begin{boxmetier}{Responsable des services IT}
\textbf{Appliquer les principes au projet ITIL de KRIBILOG}
\begin{enumerate}
    \item \emph{Privilégier la valeur}~: mesurer les conteneurs bloqués et
          les délais de chargement, pas seulement la disponibilité des
          serveurs.
    \item \emph{Commencer là où l'on est}~: analyser les 4~800~tickets du
          dernier mois avant de choisir un nouvel outil.
    \item \emph{Progresser par itérations}~: traiter d'abord la gestion des
          problèmes sur le logiciel de suivi, puis étendre.
    \item \emph{Collaborer}~: associer les agents du port et NETPORT
          SERVICES à la définition des priorités.
    \item \emph{Optimiser et automatiser}~: n'automatiser l'affectation des
          tickets qu'une fois la classification stabilisée.
\end{enumerate}
\end{boxmetier}

\subsection{Les 34 pratiques}

ITIL~4 décrit 34~pratiques, réparties en trois familles~\cite{axelos2019itil4}.

\begin{tableaucours}{L{3.2cm}cL{8.0cm}}{%
    \tableauHeader{Famille} &
    \tableauHeader{Nb} &
    \tableauHeader{Exemples}
}
Management général & 14 & Amélioration continue, gestion des risques,
gestion de la sécurité de l'information, gestion des fournisseurs, gestion
des projets, gestion du portefeuille, gestion des relations \\
Gestion des services & 17 & Centre de services, gestion des incidents,
des problèmes, des demandes de service, habilitation des changements,
gestion des mises en production, des niveaux de service, de la
disponibilité, de la continuité, de la configuration \\
Gestion technique & 3 & Gestion des déploiements, gestion des
infrastructures et des plateformes, développement et gestion des
logiciels \\
\end{tableaucours}
\vspace{-0.8em}
\begin{center}{\small\itshape Tableau~12.7. Les trois familles de pratiques
d'ITIL~4}\end{center}

Les Chapitres~\ref{chap:itil-incidents} à~\ref{chap:itil-niveaux}
présentent les pratiques de gestion des services les plus importantes pour
la gouvernance~; le Chapitre~\ref{chap:processus} les met en
correspondance avec COBIT et l'ISO/IEC~20000-1.

\begin{boxobjet}{Amélioration continue (pratique ITIL 4)}
\textbf{Finalité~:} aligner en permanence les pratiques et les services
sur les besoins de l'organisation, par l'amélioration des produits, des
services et des pratiques.

\textbf{Modèle en sept étapes~:} quelle est la vision~? Où en
sommes-nous~? Où voulons-nous aller~? Comment y arriver~? Agir. Y
sommes-nous arrivés~? Comment garder l'élan~?

\textbf{Sorties~:} registre des améliorations, plans d'action, mesures
avant et après.

\textbf{RACI (extrait, KRIBILOG)~:} \emph{A} au responsable des services
informatiques, \emph{R} aux propriétaires des pratiques, \emph{C} aux
métiers et à l'infogérant.

\textbf{Indicateurs~:} nombre d'améliorations menées à terme, gain mesuré
par amélioration.

\textbf{Correspondance COBIT~:} objectifs APO11 (qualité) et MEA01
(surveillance de la performance).
\end{boxobjet}

On remarquera que le modèle d'amélioration continue d'ITIL reprend les
sept questions de la feuille de route de COBIT~2019
(Chapitre~\ref{chap:cobit-capacite})~: les deux référentiels partagent la
même logique d'amélioration.

% ============================================================================
\section{Et ITIL Version 5~?}
\label{sec:c12-itil5}
% ============================================================================

ITIL version~5 a été annoncée par PeopleCert en janvier~2026, et ses
modules de formation sont publiés progressivement depuis~\cite{peoplecert2026itil5}.
Les sources publiques disponibles permettent d'établir les points
suivants.

\begin{tableaucours}{L{3.4cm}L{4.2cm}L{4.6cm}}{%
    \tableauHeader{Aspect} &
    \tableauHeader{ITIL 4 (2019)} &
    \tableauHeader{ITIL version 5 (2026)}
}
Objet & Gestion des services & Produits et services numériques, gérés
ensemble \\
Modèle central & Système et chaîne de valeur des services & Cycle de vie
des produits et des services, de bout en bout \\
Principes directeurs & Sept principes & Les mêmes sept principes, selon
les présentations disponibles \\
Pratiques & 34~pratiques en trois familles & 34~pratiques, regroupées
autrement selon les présentations disponibles \\
Intelligence artificielle & Peu traitée & Présentée comme intégrée dès la
conception, avec un module consacré à la gouvernance de l'IA \\
Certifications & Toujours reconnues & Nouveau schéma, introduit
progressivement \\
\end{tableaucours}
\vspace{-0.8em}
\begin{center}{\small\itshape Tableau~12.8. ITIL~4 et ITIL version~5
(d'après les sources publiques disponibles en octobre~2026)}\end{center}

\begin{boiteRemarque}{Pourquoi ce cours reste-t-il sur ITIL~4~?}
Trois raisons. D'abord, ITIL~4 est la version que pratiquent aujourd'hui
les organisations où exerceront les étudiants. Ensuite, ses concepts sont
largement repris par la version~5~: principes, pratiques, co-création de
la valeur. Enfin, le modèle central de la version~5 n'est pas encore
documenté en détail dans les sources publiques~: l'enseigner exigerait de
s'appuyer sur la publication officielle. Les étudiants qui passeront une
certification devront vérifier la version exigée par l'examen.
\end{boiteRemarque}

% ============================================================================
\separateur
% ============================================================================

\begin{boxresume}
\textbf{Ce qu'il faut retenir de ce chapitre~:}
\begin{itemize}
    \item \textbf{De processus en système de valeur}~: ITIL est passé d'un
          recueil de processus (v2) à un cycle de vie (v3), puis à un
          système de valeur des services (ITIL~4, 2019).
    \item \textbf{La valeur est co-créée}~: elle est perçue par le client.
          Ne pas confondre livrable (ce que l'on produit) et résultat (ce
          que le client obtient).
    \item \textbf{Utilité et garantie}~: ce que fait le service, et
          l'assurance qu'il le fera comme convenu.
    \item \textbf{Cinq composantes}~: principes directeurs, gouvernance,
          chaîne de valeur, pratiques, amélioration continue. La
          gouvernance reprend le modèle évaluer, diriger, surveiller.
    \item \textbf{Six activités} qui se combinent en flux de valeur, sans
          ordre imposé.
    \item \textbf{Quatre dimensions et sept principes}~: un projet ITIL n'est
          jamais un simple projet d'outil.
    \item \textbf{ITIL version~5} (2026) recentre le référentiel sur les
          produits et services numériques~; ce cours reste sur ITIL~4.
\end{itemize}
\end{boxresume}

\bigskip

\begin{boxexo}[title={\ding{45}\quad Exercice~12.1. Livrable ou résultat~? \niveaubase}]
Pour chacun des éléments suivants, indiquez s'il s'agit d'un livrable ou
d'un résultat, et précisez pour quel client.
\begin{enumerate}[(a)]
    \item Un ticket clos dans le délai.
    \item Un fournisseur de l'État payé à bonne date par l'ANPP.
    \item Un serveur disponible à 99,8~\%.
    \item Un client de la BANQUE SANAGA qui retire de l'argent un dimanche.
    \item Un rapport mensuel de NETPORT SERVICES.
\end{enumerate}
\end{boxexo}

\begin{boxexo}[title={\ding{45}\quad Exercice~12.2. Un flux de valeur à la BANQUE SANAGA \niveaumoyen}]
Décrivez le flux de valeur « ouverture d'un compte par l'application
mobile », de la demande du client jusqu'au premier virement.
\begin{enumerate}[(a)]
    \item Identifiez les étapes et rattachez chacune à une activité de la
          chaîne de valeur.
    \item Pour chaque étape, indiquez la pratique ITIL concernée.
    \item Identifiez, pour chacune des quatre dimensions, un risque qui
          pourrait interrompre le flux.
\end{enumerate}
\end{boxexo}

\begin{boxexo}[title={\ding{45}\quad Exercice~12.3. Le contrat de KRIBILOG au prisme d'ITIL~4 \niveauexpert}]
Le contrat entre KRIBILOG et NETPORT SERVICES fixe un seul engagement~: la
disponibilité des serveurs à 99,5~\%.
\begin{enumerate}[(a)]
    \item Expliquez, avec les notions d'utilité, de garantie, de livrable et
          de résultat, pourquoi cet engagement ne protège pas KRIBILOG.
    \item Proposez trois engagements orientés vers le résultat, mesurables.
    \item Montrez comment ces engagements s'inscrivent dans la gouvernance
          de KRIBILOG (qui les fixe, qui les surveille, qui en rend compte).
    \item Ce travail sera repris au Chapitre~\ref{chap:itil-niveaux}~:
          indiquez les pratiques ITIL qu'il mobilisera.
\end{enumerate}
\end{boxexo}

\bigskip

\begin{boxinfo}
\textbf{Transition.} Le système de valeur fixe le cadre. Le
Chapitre~\ref{chap:itil-incidents} entre dans les pratiques du quotidien~:
le centre de services, la gestion des incidents, des problèmes et des
demandes de service, là où se joue la perception du service par ses
utilisateurs.
\end{boxinfo}

        % ============================================================================
% Fichier  : partie4/P04_C13_incidents_problemes_demandes.tex
% Titre    : Centre de Services, Incidents, Problèmes et Demandes
% Auteur   : MINKA MI NGUIDJOÏ Thierry Emmanuel
% Version  : 1.0 -- octobre 2026
% ============================================================================
% Sources : ITIL 4 (sources publiques, définitions reformulées) ; COBIT 2019,
% fiches DSS02 et DSS03 (dossier Normes). Laboratoire associé : GCI-LAB-06
% (priorisation de tickets, MTTR, taux de réouverture, chez KRIBILOG).
% ============================================================================

\chapter{Centre de Services, Incidents, Problèmes et Demandes}
\label{chap:itil-incidents}
\chaptermark{Incidents, problèmes et demandes}

\epigraph{%
    « Réparer vite, c'est bien.\\
    Ne plus avoir à réparer, c'est mieux. »%
}{Adage des équipes de support}

\niveauChapitre{Intermédiaire}{%
    Maîtrise du Chapitre~\ref{chap:itil-svs} (concepts et système de valeur
    d'ITIL~4).%
}

\begin{boxobjectifs}
\begin{itemize}
    \item Organiser un centre de services.
    \item Prioriser un incident avec une matrice impact $\times$ urgence.
    \item Distinguer incident, problème, erreur connue et demande de
          service.
    \item Calculer les indicateurs (MTTR, taux de résolution au premier
          contact).
    \item Conduire une analyse de cause racine.
\end{itemize}
\end{boxobjectifs}

\bigskip

% ============================================================================
\section*{Question directrice : pourquoi les mêmes pannes reviennent-elles~?}
% ============================================================================

Le Chapitre~\ref{chap:itil-svs} a montré que le flux de valeur
« conteneur bloqué » de KRIBILOG SA s'arrête au contournement~: le
ticket est clos, la cause demeure, et le blocage revient. Avec
4~800~tickets par mois et 30~\% de réouvertures, le centre de services
est saturé, et M.~MBARGA, responsable des services informatiques, ne sait
plus par quoi commencer.

Ce chapitre présente les quatre pratiques qui forment le cœur de
l'exploitation quotidienne~: le centre de services
(section~\ref{sec:c13-desk}), la gestion des incidents
(section~\ref{sec:c13-incidents}), celle des demandes de service
(section~\ref{sec:c13-demandes}) et celle des problèmes
(section~\ref{sec:c13-problemes}). Il apprend à mesurer leur performance
(section~\ref{sec:c13-mesure}) et à les relier à la gouvernance
(section~\ref{sec:c13-gouvernance}). Le laboratoire GCI-LAB-06 fait
prioriser et analyser vingt tickets réels de KRIBILOG.

\separateur

% ============================================================================
\section{Quatre Notions à ne pas Confondre}
\label{sec:c13-notions}
% ============================================================================

\begin{boxdef}
\termecle{Incident}\index{incident}~: interruption non planifiée d'un
service, ou dégradation de sa qualité.

\medskip
\termecle{Problème}\index{problème}~: cause, réelle ou potentielle, d'un ou
de plusieurs incidents.

\medskip
\termecle{Erreur connue}\index{erreur connue} (\textit{known error})~:
problème qui a été analysé, dont la cause est identifiée, mais qui n'est
pas encore résolu.

\medskip
\termecle{Solution de contournement}\index{solution de contournement}
(\textit{workaround})~: moyen de réduire ou d'éliminer l'effet d'un
incident ou d'un problème, sans en traiter la cause.

\medskip
\termecle{Demande de service}\index{demande de service} (\textit{service
request})~: demande d'un utilisateur portant sur une prestation déjà
prévue dans le service (un accès, une information, un équipement).
\end{boxdef}
{\small\itshape Définitions reformulées d'après ITIL~4~\cite{axelos2019itil4}.}

\begin{tableaucours}{L{3.6cm}L{4.6cm}L{4.0cm}}{%
    \tableauHeader{Situation} &
    \tableauHeader{Qualification} &
    \tableauHeader{Pratique}
}
Le logiciel de suivi ne répond plus & Incident & Gestion des incidents \\
Le logiciel perd la position des conteneurs chaque fois que la balise
change de réseau & Problème, puis erreur connue une fois la cause
établie & Gestion des problèmes \\
On relance la synchronisation à la main & Solution de contournement &
Incidents et problèmes \\
Un nouvel agent demande un accès au portail & Demande de service & Gestion
des demandes de service \\
Un agent demande une nouvelle fonction du logiciel & Ni incident ni demande~:
demande de changement & Habilitation des changements
(Chapitre~\ref{chap:itil-changements}) \\
\end{tableaucours}
\vspace{-0.8em}
\begin{center}{\small\itshape Tableau~13.1. Qualifier une situation}\end{center}

La confusion la plus coûteuse est celle de l'incident et du problème. Le
centre de services de KRIBILOG traite chaque blocage comme un incident
isolé~: il rétablit le service, et le blocage revient. Ce qui manque n'est
pas la réactivité, c'est la \emph{gestion des problèmes}.

% ============================================================================
\section{Le Centre de Services}
\label{sec:c13-desk}
% ============================================================================

\subsection{Le point d'entrée unique}

Le \termecle{centre de services}\index{centre de services} (\textit{service
desk}) est le point de contact entre le fournisseur et les utilisateurs
du service. Il reçoit les signalements d'incidents et les demandes, les
enregistre, les qualifie, les résout lorsqu'il le peut, et les oriente
sinon. ITIL~4 insiste sur sa dimension humaine~: il est la vitrine du
service, et c'est par lui que l'utilisateur perçoit la valeur.

\subsection{Organiser le centre de services}

\begin{tableaucours}{L{3.4cm}L{8.8cm}}{%
    \tableauHeader{Choix} &
    \tableauHeader{Options et critères}
}
Canaux & Téléphone, courriel, portail en libre-service, messagerie
instantanée~; au port, le téléphone et la messagerie mobile restent
indispensables \\
Localisation & Centralisé (un seul site), local (près des utilisateurs),
virtuel (équipes réparties, un outil commun) \\
Horaires & Aligner sur l'activité~: le port de Kribi travaille 24~h/24,
le centre de services doit le faire aussi pour les incidents prioritaires \\
Niveaux & Niveau 1 (centre de services), niveau 2 (support spécialisé),
niveau 3 (éditeur, infogérant) \\
Compétences & Écoute, connaissance du métier des utilisateurs, maîtrise
de l'outil et de la base de connaissances \\
\end{tableaucours}
\vspace{-0.8em}
\begin{center}{\small\itshape Tableau~13.2. Les choix d'organisation d'un
centre de services}\end{center}

\begin{boxmetier}{Responsable des services IT}
\textbf{Désengorger le centre de services de KRIBILOG}
\begin{enumerate}
    \item Analyser la répartition des tickets~: 20~\% sont des
          réinitialisations de mot de passe, automatisables par un portail
          en libre-service.
    \item Créer une base de connaissances pour les vingt incidents les plus
          fréquents, avec leur contournement.
    \item Définir des modèles de demandes (accès, équipements) avec leur
          circuit d'approbation.
    \item Exiger, par avenant, que NETPORT SERVICES enregistre toutes ses
          interventions dans l'outil de KRIBILOG.
    \item Nommer un responsable de la gestion des problèmes, distinct du
          responsable du centre de services.
\end{enumerate}
\end{boxmetier}

% ============================================================================
\section{La Gestion des Incidents}
\label{sec:c13-incidents}
% ============================================================================

\subsection{Le cycle d'un incident}

La gestion des incidents a pour but de \emph{rétablir le service le plus
vite possible} et de limiter l'effet de l'incident sur l'activité. Son
cycle comprend l'enregistrement, la classification, la priorisation, le
diagnostic, l'escalade éventuelle, la résolution, puis la clôture après
vérification auprès de l'utilisateur. Ce sont les pratiques DSS02.02 à
DSS02.06 de COBIT (Chapitre~\ref{chap:cobit-objectifs}).

\subsection{La matrice de priorité}

La priorité d'un incident combine deux critères~:
l'\termecle{impact}\index{impact} (l'étendue des effets sur l'activité~:
combien d'utilisateurs, quels processus) et
l'\termecle{urgence}\index{urgence} (la vitesse à laquelle l'effet
s'aggrave). La priorité fixe l'ordre de traitement et le délai cible de
résolution.

\begin{tableaucours}{L{2.6cm}ccc}{%
    \tableauHeader{Impact} &
    \tableauHeader{Urgence élevée} &
    \tableauHeader{Urgence moyenne} &
    \tableauHeader{Urgence faible}
}
Élevé & P1 & P2 & P3 \\
Moyen & P2 & P3 & P4 \\
Faible & P3 & P4 & P5 \\
\end{tableaucours}
\vspace{-0.8em}
\begin{center}{\small\itshape Tableau~13.3. Matrice de priorité impact
$\times$ urgence (exemple)}\end{center}

\begin{tableaucours}{cL{6.0cm}L{4.0cm}}{%
    \tableauHeader{Priorité} &
    \tableauHeader{Exemple à KRIBILOG} &
    \tableauHeader{Délai cible de résolution}
}
P1 & Logiciel de suivi indisponible pour tout le port, navire à quai & 4
heures \\
P2 & Suivi indisponible sur un terminal, chargement en cours & 8 heures \\
P3 & Lenteur générale du logiciel & 24 heures \\
P4 & Imprimante d'un bureau en panne & 72 heures \\
P5 & Erreur d'affichage sans effet sur le travail & 5 jours ouvrés \\
\end{tableaucours}
\vspace{-0.8em}
\begin{center}{\small\itshape Tableau~13.4. Priorités et délais cibles
(convention retenue par KRIBILOG)}\end{center}

Les délais cibles ne sont pas fixés par ITIL~: ils relèvent des
engagements de service convenus avec les métiers
(Chapitre~\ref{chap:itil-niveaux}). Ce que la gouvernance exige, c'est
qu'ils soient \emph{écrits, connus et mesurés}.

\subsection{L'incident majeur}

Un \termecle{incident majeur}\index{incident majeur} est un incident dont
l'impact sur l'activité exige une réponse dépassant le traitement
ordinaire~: une coordination dédiée, une communication spécifique, une
escalade immédiate. La panne de la BANQUE SANAGA
(Chapitre~\ref{chap:prologue}) en est l'exemple type.

\begin{boxprocedure}
\textbf{Le dossier d'incident majeur.} Pour qu'un incident majeur puisse
être analysé, présenté au comité de direction ou au régulateur, le
dossier doit contenir~:
\begin{itemize}
    \item l'heure de détection, l'heure de déclaration en incident majeur
          et le nom de la personne qui l'a déclaré~;
    \item la chronologie des actions, avec leurs auteurs~;
    \item les communications émises (destinataires, heures, contenu)~;
    \item l'heure du rétablissement et la vérification auprès des
          utilisateurs~;
    \item la revue post-incident, avec l'ouverture d'un problème et les
          actions décidées.
\end{itemize}
\end{boxprocedure}

\begin{boxmetier}{DSI}
\textbf{Les trois premières heures d'un incident majeur}
\begin{enumerate}
    \item Désigner immédiatement un responsable de l'incident, seul habilité
          à coordonner, distinct de ceux qui réparent.
    \item Informer la direction générale dans l'heure, avec un point de
          situation régulier, même sans nouvelle.
    \item Décider, à une heure fixée à l'avance, du recours au plan de
          continuité si le service n'est pas rétabli
          (Chapitre~\ref{chap:itil-niveaux}).
    \item Consigner chaque décision~: c'est la chronologie dont le
          régulateur demandera la copie.
\end{enumerate}
\end{boxmetier}

% ============================================================================
\section{La Gestion des Demandes de Service}
\label{sec:c13-demandes}
% ============================================================================

Les demandes de service sont prévisibles~: elles se standardisent. ITIL~4
recommande de définir pour chaque type de demande un modèle (contenu,
circuit d'approbation, délai) et d'offrir autant que possible un
libre-service. La pratique DSS02.03 de COBIT exige de vérifier le droit du
demandeur, d'obtenir les approbations nécessaires et de satisfaire la
demande selon la procédure prévue.

\begin{boxalert}
\textbf{Violation typique.} À l'ANPP, les demandes d'accès au système de
paiement sont traitées par courriel, sans vérification du droit du
demandeur ni approbation de son responsable. Le Chapitre~\ref{chap:audit}
a montré la conséquence~: des comptes ouverts sans justification, et
jamais fermés.

\smallskip
\textbf{Règle~:} toute demande d'accès suit un modèle qui impose
l'approbation du responsable du demandeur et du propriétaire de
l'application, et laisse une trace dans l'outil.
\end{boxalert}

% ============================================================================
\section{La Gestion des Problèmes}
\label{sec:c13-problemes}
% ============================================================================

\subsection{Trois phases}

La gestion des problèmes vise à réduire la probabilité et l'effet des
incidents, en identifiant leurs causes et en gérant les contournements et
les erreurs connues. ITIL~4 la décrit en trois phases~:
\begin{itemize}
    \item l'\textbf{identification des problèmes}~: analyse des tendances
          des incidents, des incidents majeurs, des informations des
          fournisseurs~;
    \item le \textbf{contrôle des problèmes}~: analyse de la cause,
          documentation des contournements et des erreurs connues~;
    \item le \textbf{contrôle des erreurs}~: suivi des erreurs connues
          jusqu'à leur correction, souvent par un changement.
\end{itemize}

\begin{boxobjet}{DSS03 · Gérer les problèmes}
\textbf{Description~:} identifier et classer les problèmes et leurs causes
racines~; les résoudre à temps pour éviter la répétition des incidents~;
formuler des recommandations d'amélioration.

\textbf{Finalité~:} accroître la disponibilité, améliorer les niveaux de
service, réduire les coûts et accroître la satisfaction des clients, en
réduisant le nombre de problèmes d'exploitation.

\textbf{Objectifs liés~:} EG01, EG08~; AG05.

\textbf{Pratiques~:} DSS03.01 identifier et classer les problèmes~;
DSS03.02 analyser et diagnostiquer~; DSS03.03 enregistrer les erreurs
connues~; DSS03.04 résoudre et clore les problèmes~; DSS03.05 conduire une
gestion proactive des problèmes.

\textbf{RACI (extrait, KRIBILOG)~:} DSS03.02~: \emph{A} au responsable des
services informatiques, \emph{R} au responsable des problèmes, \emph{C} à
l'éditeur et à NETPORT SERVICES.

\textbf{Indicateurs~:} nombre de problèmes ouverts et leur âge, part des
incidents liés à une erreur connue, nombre d'incidents évités.
\end{boxobjet}
{\small\itshape Description, finalité, objectifs liés et pratiques résumés
et traduits par nos soins à partir de COBIT~2019~\cite{isaca2018cobitobjectifs}.}

\subsection{L'analyse de cause racine}

\begin{tableaucours}{L{3.2cm}L{4.8cm}L{4.2cm}}{%
    \tableauHeader{Technique} &
    \tableauHeader{Principe} &
    \tableauHeader{Usage}
}
Les cinq pourquoi & Demander « pourquoi~? » jusqu'à la cause (prologue,
section~\ref{sec:pro-causes}) & Causes simples, enchaînées \\
Diagramme d'Ishikawa & Classer les causes possibles par familles
(méthodes, matériel, personnes, environnement\ldots) & Causes multiples \\
Analyse de Pareto & Classer les catégories d'incidents par fréquence & Choisir
les problèmes à traiter en premier \\
Analyse chronologique & Reconstituer la suite des événements & Incidents
majeurs \\
\end{tableaucours}
\vspace{-0.8em}
\begin{center}{\small\itshape Tableau~13.5. Techniques d'analyse de cause}\end{center}

\begin{boxscenario}{Opération KRIBI. Le Pareto des tickets}
\textbf{Faits.} Sur les 4~800~tickets du mois, la répartition est la
suivante~: suivi des conteneurs 1~920 (40~\%), mots de passe 960 (20~\%),
réseau 720 (15~\%), imprimantes 480 (10~\%), divers 720 (15~\%).

\medskip
\textbf{Analyse de Pareto.} Deux catégories, suivi des conteneurs et mots
de passe, représentent 60~\% des tickets. La seconde se traite par le
libre-service. La première justifie l'ouverture d'un problème.

\medskip
\textbf{Cinq pourquoi sur le suivi des conteneurs.} Pourquoi le conteneur
disparaît-il~? Parce que sa position n'est plus mise à jour. Pourquoi~?
Parce que la balise perd la connexion en changeant de réseau mobile entre
le port et l'entrepôt. Pourquoi le logiciel ne la retrouve-t-il pas~?
Parce que la reconnexion n'est pas prévue dans la version installée.
Pourquoi n'a-t-elle pas été corrigée~? Parce que l'éditeur a livré le
correctif il y a six mois, et qu'aucune demande de changement n'a été
ouverte pour l'installer.

\medskip
\textbf{Conclusion.} Erreur connue, contournement documenté (relance
manuelle), correction disponible~: le problème se résout par un
changement (Chapitre~\ref{chap:itil-changements}). Cause de gouvernance~:
personne n'était chargé de suivre les correctifs de l'éditeur.
\end{boxscenario}

% ============================================================================
\section{Mesurer la Performance}
\label{sec:c13-mesure}
% ============================================================================

\subsection{Les indicateurs}

\begin{tableaucours}{L{3.6cm}L{5.0cm}L{3.6cm}}{%
    \tableauHeader{Indicateur} &
    \tableauHeader{Calcul} &
    \tableauHeader{Ce qu'il révèle}
}
Délai moyen de résolution (\sigle{MTTR}) & Somme des durées de résolution
/ nombre d'incidents résolus & Rapidité du rétablissement \\
Résolution au premier contact & Incidents résolus par le niveau~1 sans
escalade / incidents reçus & Efficacité du centre de services \\
Taux de réouverture & Incidents rouverts / incidents clos & Qualité de la
résolution \\
Respect des délais cibles & Incidents résolus dans le délai de leur
priorité / incidents résolus & Tenue des engagements \\
Âge des problèmes ouverts & Durée depuis l'ouverture & Capacité à traiter
les causes \\
\end{tableaucours}
\vspace{-0.8em}
\begin{center}{\small\itshape Tableau~13.6. Indicateurs de la gestion des
incidents et des problèmes}\end{center}

\subsection{Un calcul pas à pas}

\begin{boxscenario}{Opération KRIBI. Huit tickets}
\begin{tableaucours}{ccccc}{%
    \tableauHeader{Ticket} &
    \tableauHeader{Priorité} &
    \tableauHeader{Durée (h)} &
    \tableauHeader{Niveau 1~?} &
    \tableauHeader{Rouvert~?}
}
T1 & P3 & 2 & Oui & Non \\
T2 & P2 & 5 & Non & Oui \\
T3 & P3 & 26 & Non & Non \\
T4 & P4 & 1 & Oui & Non \\
T5 & P1 & 49 & Non & Oui \\
T6 & P2 & 3 & Non & Non \\
T7 & P3 & 8 & Non & Oui \\
T8 & P4 & 30 & Non & Non \\
\end{tableaucours}

\medskip
\textbf{Calculs.} MTTR $= (2+5+26+1+49+3+8+30)/8 = 124/8 = 15{,}5$~heures.
Résolution au premier contact~: $2/8 = 25~\%$. Taux de réouverture~:
$3/8 = 37{,}5~\%$. Respect des délais cibles (Tableau~13.4)~: T3 (26~h
pour P3, cible 24~h) et T5 (49~h pour P1, cible 4~h) sont hors délai, soit
$6/8 = 75~\%$ dans les délais.

\medskip
\textbf{Lecture.} Le MTTR moyen masque l'essentiel~: l'incident P1, le plus
grave, a duré douze fois son délai cible. Un indicateur moyen ne remplace
jamais l'examen des incidents prioritaires.
\end{boxscenario}

\begin{boxalert}
\textbf{Violation typique.} Pour améliorer son taux de résolution,
l'infogérant clôt les tickets dès qu'une solution de contournement est
appliquée, sans vérifier auprès de l'utilisateur. Les indicateurs
s'améliorent, le service se dégrade.

\smallskip
\textbf{Règle~:} tout indicateur doit être accompagné de son indicateur
de contrôle (le taux de réouverture pour le taux de résolution), et la
clôture d'un incident exige la confirmation de l'utilisateur
(DSS02.06).
\end{boxalert}

% ============================================================================
\section{Le Lien avec la Gouvernance}
\label{sec:c13-gouvernance}
% ============================================================================

Les incidents et les problèmes relèvent du management, mais ils
intéressent la gouvernance par trois voies.
\begin{itemize}
    \item \textbf{L'information du conseil}~: les incidents majeurs et les
          tendances doivent remonter au comité des risques ou au comité
          stratégique des SI. Le conseil de la BANQUE SANAGA n'avait reçu
          aucun rapport d'incident en 2025.
    \item \textbf{Le risque}~: un problème non résolu est un risque connu
          et accepté de fait~; il doit figurer au registre des risques
          (APO12) et être accepté explicitement s'il n'est pas traité.
    \item \textbf{Les fournisseurs}~: les incidents imputables à un
          prestataire alimentent sa supervision (APO10) et les pénalités
          prévues au contrat (Chapitre~\ref{chap:itil-niveaux}).
\end{itemize}

\begin{boxmetier}{Auditeur SI}
\textbf{Auditer la gestion des incidents et des problèmes}
\begin{enumerate}
    \item Extraire de l'outil tous les tickets de la période, avec la
          lettre de transmission prévue par le GUID~5100
          (Chapitre~\ref{chap:audit}).
    \item Vérifier sur un échantillon la cohérence de la priorité avec la
          matrice, et le respect des délais cibles.
    \item Rapprocher les tickets rouverts des problèmes ouverts~: une
          catégorie très rouverte sans problème ouvert est un constat.
    \item Vérifier que les incidents majeurs ont donné lieu à une revue
          post-incident.
\end{enumerate}
\end{boxmetier}

% ============================================================================
\separateur
% ============================================================================

\begin{boxresume}
\textbf{Ce qu'il faut retenir de ce chapitre~:}
\begin{itemize}
    \item \textbf{Quatre notions}~: l'incident est une interruption, le
          problème sa cause, l'erreur connue une cause identifiée non
          corrigée, la demande de service une prestation prévue.
    \item \textbf{Le centre de services} est le point d'entrée unique et la
          vitrine du service~; il se désengorge par le libre-service, la
          base de connaissances et les modèles de demandes.
    \item \textbf{La priorité} combine impact et urgence~; les délais cibles
          relèvent des engagements de service et doivent être écrits.
    \item \textbf{L'incident majeur} exige un responsable dédié, une
          communication réglée et un dossier complet.
    \item \textbf{La gestion des problèmes} (DSS03) traite les causes~:
          identification, contrôle des problèmes, contrôle des erreurs.
          Sans elle, les mêmes incidents reviennent.
    \item \textbf{Mesurer}~: MTTR, résolution au premier contact, taux de
          réouverture, respect des délais. Une moyenne ne remplace pas
          l'examen des incidents prioritaires.
\end{itemize}
\end{boxresume}

\bigskip

\begin{boxexo}[title={\ding{45}\quad Exercice~13.1. Qualifier et prioriser \niveaubase}]
Qualifiez chacune des situations suivantes (incident, problème, demande de
service, demande de changement), puis donnez la priorité des incidents
selon la matrice du Tableau~13.3.
\begin{enumerate}[(a)]
    \item Le portail de paiement de l'ANPP est inaccessible la veille de la
          paie des fonctionnaires.
    \item Un agent de la BANQUE SANAGA demande l'installation d'un logiciel
          de bureautique.
    \item Les distributeurs d'une agence de Garoua tombent en panne chaque
          soir vers 22~heures.
    \item Le directeur financier de KRIBILOG ne peut pas imprimer un
          rapport non urgent.
\end{enumerate}
\end{boxexo}

\begin{boxexo}[title={\ding{45}\quad Exercice~13.2. Calculer les indicateurs \niveaumoyen}]
Le centre de services de la BANQUE SANAGA a traité dix incidents dans la
journée, avec les durées de résolution suivantes (en heures)~: 1, 3, 2,
12, 4, 1, 6, 30, 2, 9. Les incidents de 1~heure et le premier incident de
2~heures ont été résolus au premier contact. Deux incidents ont été
rouverts.
\begin{enumerate}[(a)]
    \item Calculez le MTTR, le taux de résolution au premier contact et le
          taux de réouverture.
    \item L'incident de 30~heures était un P1 (cible~: 4~heures). Que
          concluez-vous sur la pertinence du MTTR comme indicateur unique~?
    \item Proposez un tableau de bord de quatre indicateurs pour le comité
          de direction.
\end{enumerate}
\end{boxexo}

\begin{boxexo}[title={\ding{45}\quad Exercice~13.3. La revue post-incident de la BANQUE SANAGA \niveauexpert}]
Conduisez la revue post-incident de la panne de mars~2026
(Chapitre~\ref{chap:prologue}).
\begin{enumerate}[(a)]
    \item Reconstituez la chronologie de l'incident majeur, et indiquez les
          pièces qui manquent au dossier (exigence de contrôle de la
          section~\ref{sec:c13-incidents}).
    \item Ouvrez le problème correspondant et conduisez une analyse
          d'Ishikawa des causes.
    \item Distinguez les actions qui relèvent de la gestion des problèmes,
          de la gestion des changements et de la gouvernance.
    \item Rédigez la partie « incidents » du rapport trimestriel au comité
          des risques.
\end{enumerate}
\end{boxexo}

\bigskip

\begin{boxinfo}
\textbf{Transition.} À KRIBILOG comme à la BANQUE SANAGA, la résolution des
problèmes passe par un changement~: l'installation d'un correctif, la
bascule d'un système. Le Chapitre~\ref{chap:itil-changements} présente
l'habilitation des changements, la gestion des mises en production et la
gestion de la configuration.
\end{boxinfo}

        % ============================================================================
% Fichier  : partie4/P04_C14_changements_configuration.tex
% Titre    : Changements, Mises en Production et Configuration
% Auteur   : MINKA MI NGUIDJOÏ Thierry Emmanuel
% Version  : 1.0 -- octobre 2026
% ============================================================================
% Sources : ITIL 4 (sources publiques, notions reformulées) ; COBIT 2019,
% fiches BAI06, BAI07 et BAI10 (dossier Normes). Laboratoire associé :
% GCI-LAB-07 (évaluation d'un changement, analyse d'impact, décision).
% ============================================================================

\chapter{Changements, Mises en Production et Configuration}
\label{chap:itil-changements}
\chaptermark{Changements et configuration}

\epigraph{%
    « Toutes les grandes pannes de ce cours ont commencé par un
    changement. »%
}{Leçon des grandes défaillances}

\niveauChapitre{Intermédiaire}{%
    Maîtrise du Chapitre~\ref{chap:itil-incidents} (incidents et
    problèmes) et du Chapitre~\ref{chap:defaillances} (grandes
    défaillances).%
}

\begin{boxobjectifs}
\begin{itemize}
    \item Distinguer changement standard, normal et urgent.
    \item Évaluer le risque d'un changement et organiser son autorisation.
    \item Planifier une mise en production et son retour arrière.
    \item Exploiter une CMDB pour une analyse d'impact.
    \item Relier ces pratiques à l'objectif COBIT BAI06.
\end{itemize}
\end{boxobjectifs}

\bigskip

% ============================================================================
\section*{Question directrice : comment changer sans casser~?}
% ============================================================================

Knight Capital, TSB, CrowdStrike, la BANQUE SANAGA~: le
Chapitre~\ref{chap:defaillances} a montré que chaque grande panne avait
commencé par un changement. Pourtant, une organisation qui ne change
rien meurt aussi sûrement~: le correctif que KRIBILOG n'a jamais installé
(Chapitre~\ref{chap:itil-incidents}) est la cause de 40~\% de ses tickets.
Et la BANQUE SANAGA doit relancer, en 2027, la bascule de son système
central.

Le dilemme est celui de COBIT, qui rattache la gestion des changements à
l'\emph{agilité} (Chapitre~\ref{chap:cobit-objectifs})~: il faut changer
\emph{vite} et changer \emph{bien}. Ce chapitre présente les trois
pratiques qui y répondent~: l'habilitation des changements
(sections~\ref{sec:c14-types} et~\ref{sec:c14-risque}), la gestion des
mises en production et des déploiements (section~\ref{sec:c14-mep}), et la
gestion de la configuration (section~\ref{sec:c14-config}). Le laboratoire
GCI-LAB-07 fait évaluer et décider un changement de la banque.

\separateur

% ============================================================================
\section{Les Changements et leurs Types}
\label{sec:c14-types}
% ============================================================================

\subsection{Qu'est-ce qu'un changement~?}

\begin{boxdef}
\termecle{Changement}\index{changement} (\textit{change})~: ajout,
modification ou retrait de tout élément susceptible d'avoir un effet,
direct ou indirect, sur un service.

\medskip
\termecle{Habilitation des changements}\index{habilitation des
changements} (\textit{change enablement})~: pratique qui vise à maximiser
le nombre de changements réussis, en veillant à ce que les risques soient
correctement évalués, à ce que les changements soient autorisés et à ce
que leur calendrier soit maîtrisé.
\end{boxdef}
{\small\itshape Définitions reformulées d'après ITIL~4~\cite{axelos2019itil4}.
La pratique s'appelait \emph{change control} dans la première édition
d'ITIL~4.}

Le mot \emph{habilitation} n'est pas anodin. ITIL~4 ne cherche pas à
freiner les changements, mais à en permettre le plus grand nombre possible
\emph{avec succès}. Une organisation où chaque changement attend trois
semaines l'avis d'un comité contourne tôt ou tard la procédure~: c'est
l'erreur symétrique de celle de Knight Capital.

\subsection{Trois types de changements}

\begin{tableaucours}{L{2.4cm}L{4.8cm}L{5.0cm}}{%
    \tableauHeader{Type} &
    \tableauHeader{Caractéristiques} &
    \tableauHeader{Exemple}
}
Standard & Risque faible, bien connu, procédure documentée~; autorisé
d'avance & Créer un compte utilisateur, ajouter de la mémoire à un serveur
virtuel \\
Normal & Doit être évalué et autorisé selon son risque, par l'autorité
compétente & Installer le correctif de l'éditeur chez KRIBILOG~; basculer
le système central de la BANQUE SANAGA \\
Urgent & Doit être mis en œuvre au plus vite, par exemple pour résoudre un
incident ou combler une faille~; procédure accélérée & Corriger une faille
de sécurité activement exploitée \\
\end{tableaucours}
\vspace{-0.8em}
\begin{center}{\small\itshape Tableau~14.1. Les trois types de changements}\end{center}

\begin{boxalert}
\textbf{Violation typique.} Pour échapper au comité des changements, une
équipe qualifie systématiquement ses changements d'« urgents ». À la BANQUE
SANAGA, 35~\% des changements de 2025 étaient enregistrés comme urgents~;
aucun n'avait fait l'objet de la revue a posteriori prévue.

\smallskip
\textbf{Règle~:} un changement urgent n'est pas un changement non
contrôlé. Il est autorisé par une autorité désignée à l'avance (souvent
restreinte), documenté, et revu après coup (BAI06.02). Un taux de
changements urgents élevé est un indicateur d'alerte.
\end{boxalert}

\subsection{L'autorité de changement}

ITIL~4 parle d'\termecle{autorité de changement}\index{autorité de
changement} plutôt que de comité~: la personne ou le groupe qui autorise
un changement. Elle n'est pas unique~: l'autorisation est
\emph{déléguée} selon le risque. Un changement standard est autorisé
d'avance~; un changement normal à faible risque peut l'être par le
responsable de l'équipe~; un changement à risque élevé exige un
\termecle{comité des changements}\index{comité des changements}
(\textit{change advisory board}, \sigle{CAB}), voire la direction.

Cette délégation est une question de \textbf{droits de décision}
(Chapitre~\ref{chap:alignement})~: qui peut décider de quel changement.
Elle doit être fixée par écrit et approuvée par la gouvernance.

% ============================================================================
\section{Évaluer et Autoriser}
\label{sec:c14-risque}
% ============================================================================

\subsection{Une grille d'évaluation du risque}

Chaque organisation définit sa grille. Celle que retient la BANQUE SANAGA
après la panne note six critères de~1 (favorable) à~3 (défavorable)~; le
total, de 6 à~18, détermine l'autorité compétente.

\begin{tableaucours}{L{3.6cm}L{2.8cm}L{2.8cm}L{2.8cm}}{%
    \tableauHeader{Critère} &
    \tableauHeader{1} &
    \tableauHeader{2} &
    \tableauHeader{3}
}
Criticité du service & Soutien & Important & Critique (paiements,
système central) \\
Utilisateurs touchés & Moins de 50 & 50 à 500 & Plus de 500 ou clients \\
Complexité & Un composant & Plusieurs composants & Plusieurs systèmes,
données converties \\
Réversibilité & Retour arrière testé & Retour arrière prévu non testé &
Retour arrière impossible \\
Couverture des tests & Plus de 95~\% & 80 à 95~\% & Moins de 80~\% \\
Expérience & Changement déjà réalisé & Changement analogue & Première fois \\
\end{tableaucours}
\vspace{-0.8em}
\begin{center}{\small\itshape Tableau~14.2. Grille d'évaluation du risque
d'un changement (convention de la BANQUE SANAGA)}\end{center}

\begin{tableaucours}{L{2.4cm}L{2.6cm}L{7.2cm}}{%
    \tableauHeader{Score} &
    \tableauHeader{Risque} &
    \tableauHeader{Autorité de changement}
}
6 à 9 & Faible & Responsable de l'équipe concernée \\
10 à 13 & Moyen & Responsable des changements, après avis du RSSI \\
14 à 18 & Élevé & Comité des changements~; et, si la réversibilité vaut~3,
décision formelle du comité de pilotage avec acceptation écrite du risque \\
\end{tableaucours}
\vspace{-0.8em}
\begin{center}{\small\itshape Tableau~14.3. Score de risque et autorité
compétente}\end{center}

\begin{boxscenario}{Opération SANAGA. La bascule notée}
\textbf{Notation a posteriori.} Criticité~3, utilisateurs~3, complexité~3,
réversibilité~3 (conversion irréversible), tests~3 (63~\%), expérience~3
(première migration). Score~: \textbf{18 sur 18}.

\medskip
\textbf{Analyse} (confiance élevée). Avec la grille, la bascule aurait
relevé du comité des changements \emph{et} d'une décision formelle du
comité de pilotage, avec acceptation écrite du risque. Aucune des deux
instances ne s'est prononcée. La grille n'aurait pas empêché la décision
de basculer, mais elle aurait obligé quelqu'un à la prendre, par écrit,
en connaissance de cause.
\end{boxscenario}

\subsection{Le calendrier des changements}

Le \termecle{calendrier des changements}\index{calendrier des changements}
(\textit{change schedule}) planifie les changements autorisés et en
informe les parties concernées. Il permet d'éviter les collisions (deux
changements sur le même système le même soir) et d'appliquer des
\emph{périodes de gel} pendant les moments sensibles~: la clôture des
comptes pour une banque, la paie des fonctionnaires pour l'ANPP, l'arrivée
d'un grand navire pour KRIBILOG.

\begin{boxobjet}{BAI06 · Gérer les changements informatiques}
\textbf{Finalité~:} permettre la livraison rapide et fiable des
changements, en limitant le risque d'affecter la stabilité ou l'intégrité
de l'environnement modifié.

\textbf{Objectifs liés~:} EG01 (portefeuille compétitif)~; AG06 (agilité).

\textbf{Pratiques~:} BAI06.01 évaluer, prioriser et autoriser les demandes
de changement~; BAI06.02 gérer les changements d'urgence~; BAI06.03 suivre
l'état des changements et en rendre compte~; BAI06.04 clore et documenter
les changements.

\textbf{RACI (extrait, BANQUE SANAGA)~:} BAI06.01 pour un risque élevé~:
\emph{A} au président du comité des changements, \emph{R} au responsable
des changements, \emph{C} au RSSI, aux métiers et à l'exploitation,
\emph{I} au centre de services.

\textbf{Indicateurs~:} taux de changements en échec~; taux de changements
urgents~; nombre de changements non autorisés détectés~; part des
changements à risque élevé dotés d'un retour arrière testé.
\end{boxobjet}
{\small\itshape Finalité, objectifs liés et pratiques résumés et traduits
par nos soins à partir de COBIT~2019~\cite{isaca2018cobitobjectifs}.}

\begin{boxmetier}{Responsable des services IT}
\textbf{La charte du comité des changements}
\begin{enumerate}
    \item \textbf{Mandat}~: autoriser les changements à risque élevé,
          arbitrer les conflits de calendrier, examiner les changements en
          échec.
    \item \textbf{Composition}~: responsable des changements (président),
          exploitation, sécurité, représentants des métiers concernés,
          fournisseur si besoin.
    \item \textbf{Fréquence}~: hebdomadaire, et sur convocation pour un
          changement urgent à risque élevé.
    \item \textbf{Pièces exigées}~: évaluation du risque, plan de mise en
          production, plan de retour arrière et preuve de son test,
          résultats de recette.
    \item \textbf{Trace}~: décision motivée, inscrite au registre des
          changements avant toute mise en production.
\end{enumerate}
\end{boxmetier}

% ============================================================================
\section{Mises en Production et Déploiements}
\label{sec:c14-mep}
% ============================================================================

\subsection{Deux pratiques distinctes}

ITIL~4 distingue la \termecle{gestion des mises en production}\index{mise en
production} (\textit{release management}), qui rend disponibles des
services ou des fonctions nouveaux ou modifiés, et la \termecle{gestion
des déploiements}\index{déploiement} (\textit{deployment management}),
qui transfère les composants vers les environnements cibles. On peut
déployer un logiciel sans le rendre disponible aux utilisateurs, puis
l'activer plus tard~: c'est le principe des fonctions activables à la
demande.

\subsection{Stratégies de mise en production}

\begin{tableaucours}{L{3.0cm}L{4.8cm}L{4.4cm}}{%
    \tableauHeader{Stratégie} &
    \tableauHeader{Principe} &
    \tableauHeader{Risque}
}
Bascule unique (\textit{big bang}) & Tout le monde passe au nouveau système
au même moment & Élevé~: tout échec touche tout le monde \\
Par étapes & Agence par agence, région par région & Modéré~: l'échec est
circonscrit \\
Pilote & Un site ou un groupe d'utilisateurs d'abord & Faible pour le
pilote, mais long \\
Fonctionnement en parallèle & Ancien et nouveau systèmes tournent
ensemble pendant une période & Faible, mais coûteux \\
Déploiement progressif (\textit{canary}) & Une petite part du trafic d'abord,
puis le reste & Faible~; suppose une architecture adaptée \\
\end{tableaucours}
\vspace{-0.8em}
\begin{center}{\small\itshape Tableau~14.4. Stratégies de mise en
production}\end{center}

La BANQUE SANAGA avait choisi une bascule unique, avec une conversion des
données irréversible, pour ses 38~agences. CrowdStrike avait déployé sa
mise à jour en une fois sur tout son parc (Chapitre~\ref{chap:defaillances}).
Dans les deux cas, la stratégie a transformé une erreur en catastrophe.

\subsection{Le retour arrière}

\begin{boxdef}
\termecle{Plan de retour arrière}\index{retour arrière} (\textit{back-out
plan, rollback})~: ensemble des actions qui permettent de revenir à l'état
antérieur au changement si celui-ci échoue, avec les critères qui
déclenchent ce retour et le délai dans lequel il doit être décidé.
\end{boxdef}

Un retour arrière n'existe que s'il a été \emph{testé}. Quand il est
impossible (conversion irréversible des données, mise à jour d'un
schéma de base de données), le risque doit être traité autrement~:
fonctionnement en parallèle, sauvegarde complète vérifiée, ou acceptation
formelle par l'instance compétente.

\subsection{Les pratiques de COBIT}

L'objectif BAI07 (gérer l'acceptation et la transition des changements)
détaille ce qu'il faut maîtriser~: établir un plan de mise en œuvre
(BAI07.01), planifier la conversion des processus, des systèmes et des
données (BAI07.02), planifier et réaliser les tests d'acceptation
(BAI07.03 à BAI07.05), passer en production et gérer les versions
(BAI07.06), assurer un soutien renforcé au démarrage (BAI07.07) et
conduire une revue après la mise en œuvre (BAI07.08)~\cite{isaca2018cobitobjectifs}.
La pratique BAI07.02, la conversion des données, est exactement celle qui
a fait défaut à la BANQUE SANAGA.

\begin{boxprocedure}
\textbf{Les critères « go / no-go ».} Avant toute mise en production à
risque élevé, la décision de passer est prise sur une liste de critères
fixée à l'avance, et consignée~:
\begin{itemize}
    \item recette réalisée au taux convenu, anomalies bloquantes closes~;
    \item retour arrière testé, avec son délai de mise en œuvre~;
    \item données converties contrôlées sur un échantillon, avec un
          rapprochement des soldes~;
    \item équipes de soutien et fournisseur mobilisés pour la période de
          démarrage~;
    \item communication aux utilisateurs prête~;
    \item heure limite au-delà de laquelle le retour arrière est
          automatiquement déclenché.
\end{itemize}
Un seul critère non rempli suffit à décider « no-go »~; le passer outre
exige une acceptation écrite du risque.
\end{boxprocedure}

\begin{boxscenario}{Opération SANAGA. La nouvelle bascule}
\textbf{Plan de 2027.} Bascule par étapes~: une agence pilote à Douala, puis
quatre vagues de huit à dix agences, à un mois d'intervalle. Fonctionnement
en parallèle de l'ancien système pendant deux semaines pour chaque vague,
avec rapprochement quotidien des soldes. Gel des changements pendant la
clôture de fin d'année.

\medskip
\textbf{Évaluation avec la grille.} Criticité~3, utilisateurs~2 (une vague),
complexité~3, réversibilité~1 (fonctionnement en parallèle), tests~1,
expérience~2 (après le pilote). Score~: 12, soit un risque moyen par
vague~; mais le comité de pilotage décide de garder le comité des
changements comme autorité pour chaque vague, compte tenu de l'histoire
de la banque.
\end{boxscenario}

% ============================================================================
\section{La Gestion de la Configuration}
\label{sec:c14-config}
% ============================================================================

\subsection{Éléments de configuration et CMDB}

\begin{boxdef}
\termecle{Élément de configuration}\index{élément de configuration}
(\textit{configuration item}, \sigle{CI})~: tout composant qui doit être
géré pour fournir un service~: un serveur, une application, une base de
données, un contrat, une documentation.

\medskip
\termecle{Base de données de gestion de la configuration}\index{CMDB}
(\textit{configuration management database}, \sigle{CMDB})~: base qui
enregistre les éléments de configuration, leurs attributs et surtout
leurs \emph{relations} (dépend de, est hébergé sur, utilise).
\end{boxdef}

La gestion de la configuration se distingue de la gestion des actifs
(BAI09)~: l'actif s'intéresse à la valeur et au cycle de vie financier
d'un bien (achat, amortissement, mise au rebut)~; la configuration
s'intéresse à son rôle dans le service et à ses relations avec les autres
composants.

\subsection{L'analyse d'impact}

La valeur de la CMDB est de répondre, avant un changement, à la question
« qu'est-ce qui sera touché~? ».

\begin{boxscenario}{Opération KRIBI. Installer le correctif}
\textbf{Extrait de la CMDB pour le CI « Logiciel de suivi des conteneurs,
version~3.2 ».}

\medskip
\begin{tableaucours}{L{4.4cm}L{2.6cm}L{5.0cm}}{%
    \tableauHeader{Élément lié} &
    \tableauHeader{Relation} &
    \tableauHeader{Conséquence pour le changement}
}
Serveur d'application SRV-KRI-02 & Hébergé sur & Redémarrage nécessaire \\
Base de données des mouvements & Utilise & Modification du schéma~:
sauvegarde obligatoire \\
Passerelle des balises GPS & Reçoit de & Test de reconnexion requis \\
Portail clients de l'armateur & Fournit à & Interruption à annoncer au
client \\
Contrat NETPORT SERVICES & Exploité par & Intervention de l'infogérant à
planifier \\
\end{tableaucours}

\medskip
\textbf{Évaluation.} Criticité~3, utilisateurs~3, complexité~2,
réversibilité~2 (modification du schéma de la base), tests~2, expérience~2.
Score~14~: risque élevé, autorité~: comité des changements. Fenêtre
retenue~: dimanche de 2~h à 6~h, hors arrivée de navire.

\medskip
\textbf{Ce que révèle l'exercice.} Sans la CMDB, la modification du schéma
de la base et l'effet sur le portail de l'armateur seraient passés
inaperçus. C'est précisément ce que KRIBILOG ignorait lorsqu'elle a
renoncé, il y a six mois, à installer le correctif.
\end{boxscenario}

\subsection{L'objectif BAI10}

COBIT consacre à la configuration l'objectif BAI10 (gérer la
configuration), dont les pratiques couvrent la définition d'un modèle de
configuration, la tenue d'un référentiel et de ses versions de référence,
la maîtrise des éléments de configuration, les états de configuration, et
la vérification périodique de l'intégrité du référentiel. Ses liens
principaux sont le risque métier (EG02), la continuité (EG06) et la
sécurité (AG07)~\cite{isaca2018cobitobjectifs}~: une CMDB fausse est un
risque, car elle conduit à des analyses d'impact fausses.

\begin{boxalert}
\textbf{Violation typique.} Une organisation achète un outil de CMDB, le
remplit une fois par un inventaire automatique, puis ne le met plus à jour.
Deux ans plus tard, la moitié des relations sont fausses, et plus personne
ne s'y fie.

\smallskip
\textbf{Règle~:} la CMDB se met à jour par le processus de changement~:
tout changement clos met à jour les éléments de configuration concernés
(BAI06.04, BAI10.03), et un contrôle périodique vérifie son intégrité
(BAI10.05).
\end{boxalert}

% ============================================================================
\section{Changements et Gouvernance}
\label{sec:c14-gouvernance}
% ============================================================================

\begin{boxmetier}{Auditeur SI}
\textbf{Détecter les changements non autorisés}
\begin{enumerate}
    \item Extraire les journaux de la production (déploiements, modifications
          de configuration) sur la période.
    \item Les rapprocher du registre des changements~: tout déploiement
          sans changement autorisé correspondant est un écart.
    \item Sur un échantillon de changements autorisés, vérifier que
          l'autorisation est antérieure à la mise en production, que les
          tests et le retour arrière sont documentés.
    \item Vérifier la séparation des tâches~: celui qui développe ne
          déploie pas seul en production.
\end{enumerate}
Ce rapprochement est le test le plus puissant des contrôles généraux sur
les changements~: il part de ce qui s'est réellement passé, et non de ce
qui a été déclaré.
\end{boxmetier}

\begin{boiteRemarque}{DevOps tue-t-il la gestion des changements~?}
Les organisations qui déploient plusieurs fois par jour ne peuvent pas
réunir un comité avant chaque mise en production. Elles ne renoncent pas
pour autant au contrôle~: elles le \emph{déplacent} dans la chaîne de
déploiement automatisée (tests obligatoires, revue de code par un pair,
déploiement progressif, retour arrière automatique) et traitent la plupart
de leurs changements comme des changements standard. Le
Chapitre~\ref{chap:agile-cloud} y reviendra. La gouvernance n'exige pas un
comité~; elle exige que le risque soit évalué et que l'autorisation soit
traçable.
\end{boiteRemarque}

% ============================================================================
\separateur
% ============================================================================

\begin{boxresume}
\textbf{Ce qu'il faut retenir de ce chapitre~:}
\begin{itemize}
    \item \textbf{Changer vite et bien}~: l'habilitation des changements
          vise à maximiser le nombre de changements réussis, pas à les
          freiner.
    \item \textbf{Trois types}~: standard (autorisé d'avance), normal
          (évalué et autorisé selon son risque), urgent (accéléré, mais
          contrôlé et revu après coup).
    \item \textbf{L'autorité de changement est déléguée selon le risque}~:
          c'est une question de droits de décision, fixée par la
          gouvernance.
    \item \textbf{La stratégie de mise en production} détermine l'ampleur
          d'un échec~; un retour arrière n'existe que s'il a été testé.
    \item \textbf{Les critères go / no-go} sont fixés à l'avance~; les
          passer outre exige une acceptation écrite du risque.
    \item \textbf{La CMDB} sert l'analyse d'impact~; elle se tient à jour
          par le processus de changement.
    \item \textbf{BAI06, BAI07, BAI10} sont les objectifs COBIT
          correspondants~; le rapprochement entre journaux de production et
          registre des changements en est le meilleur test d'audit.
\end{itemize}
\end{boxresume}

\bigskip

\begin{boxexo}[title={\ding{45}\quad Exercice~14.1. Quel type de changement~? \niveaubase}]
Qualifiez chacun des changements suivants (standard, normal, urgent) et
justifiez.
\begin{enumerate}[(a)]
    \item Ajouter un utilisateur au groupe « consultation » du logiciel de
          suivi.
    \item Corriger, un samedi soir, une faille activement exploitée sur le
          portail de paiement de l'ANPP.
    \item Migrer la messagerie de la BANQUE SANAGA vers un service en
          nuage.
    \item Remplacer un disque défectueux dans une baie redondante.
\end{enumerate}
\end{boxexo}

\begin{boxexo}[title={\ding{45}\quad Exercice~14.2. Évaluer un changement \niveaumoyen}]
L'ANPP veut mettre en production une nouvelle version de son module de
paie, trois jours avant la paie des fonctionnaires. La recette couvre
85~\% des scénarios, le retour arrière est prévu mais non testé, le
changement touche 400~utilisateurs internes et tous les agents de l'État.
\begin{enumerate}[(a)]
    \item Notez le changement avec la grille du Tableau~14.2 et déduisez
          l'autorité compétente.
    \item Quels critères go / no-go ne sont pas remplis~?
    \item Que recommandez-vous, et qui doit le décider~?
\end{enumerate}
\end{boxexo}

\begin{boxexo}[title={\ding{45}\quad Exercice~14.3. Refonder la gestion des changements de la BANQUE SANAGA \niveauexpert}]
Le comité stratégique des SI demande un dispositif de gestion des
changements qui atteigne le niveau de capacité~3 en dix-huit mois
(Chapitre~\ref{chap:cobit-conception}).
\begin{enumerate}[(a)]
    \item Définissez les trois types de changements de la banque, avec des
          exemples, et la liste initiale des changements standard.
    \item Fixez la délégation de l'autorité de changement, et indiquez qui
          doit l'approuver.
    \item Proposez cinq indicateurs pour le comité d'audit, avec leur cible.
    \item Indiquez comment l'auditrice interne vérifiera, dans un an, que le
          dispositif fonctionne.
\end{enumerate}
\end{boxexo}

\bigskip

\begin{boxinfo}
\textbf{Transition.} Les changements modifient les services~; encore
faut-il savoir ce que l'on a promis à leurs utilisateurs. Le
Chapitre~\ref{chap:itil-niveaux} traite des engagements de service, de la
disponibilité, de la continuité et de la gestion des fournisseurs.
\end{boxinfo}

        % ============================================================================
% Fichier  : partie4/P04_C15_niveaux_service_continuite.tex
% Titre    : Niveaux de Service, Disponibilité, Continuité et Fournisseurs
% Auteur   : MINKA MI NGUIDJOÏ Thierry Emmanuel
% Version  : 1.0 -- octobre 2026
% ============================================================================
% Sources : ITIL 4 (sources publiques, notions reformulées) ; COBIT 2019,
% fiches APO09, APO10 et DSS04 (dossier Normes) ; loi n° 2010/012, art. 26
% (texte anglais officiel). Laboratoire associé : GCI-LAB-08 (disponibilité
% mensuelle, pénalités en FCFA, RTO et RPO chez KRIBILOG).
% ============================================================================

\chapter{Niveaux de Service, Disponibilité, Continuité et Fournisseurs}
\label{chap:itil-niveaux}
\chaptermark{Niveaux de service et fournisseurs}

\epigraph{%
    « Un engagement que l'on ne mesure pas n'engage personne. »%
}{Adage des acheteurs de services}

\niveauChapitre{Intermédiaire}{%
    Maîtrise du Chapitre~\ref{chap:itil-incidents} (incidents, priorités et
    délais cibles) et du Chapitre~\ref{chap:itil-svs} (utilité et
    garantie).%
}

\begin{boxobjectifs}
\begin{itemize}
    \item Rédiger un SLA, un OLA et un contrat de sous-traitance cohérents.
    \item Calculer une disponibilité et les pénalités contractuelles en
          FCFA.
    \item Conduire une analyse d'impact sur l'activité et fixer RTO et RPO.
    \item Piloter un fournisseur de services informatiques.
\end{itemize}
\end{boxobjectifs}

\bigskip

% ============================================================================
\section*{Question directrice : que vaut une promesse de 99,5~\%~?}
% ============================================================================

Le contrat entre KRIBILOG SA et son infogérant NETPORT SERVICES tient en
une ligne d'engagement~: « disponibilité des serveurs~: 99,5~\% ».
L'exercice~12.3 a montré que cette promesse ne protège pas KRIBILOG.
Mais combien vaut-elle exactement~? Combien d'heures de panne
autorise-t-elle chaque mois~? Que rembourse-t-elle quand un navire attend
à quai, à 12~millions FCFA de surestaries par jour~? Et que se passe-t-il si
le centre de données de NETPORT SERVICES brûle~?

Ce chapitre répond à ces questions. Il présente la gestion des niveaux de
service (section~\ref{sec:c15-sla}), le calcul de la disponibilité
(section~\ref{sec:c15-dispo}) et des pénalités (section~\ref{sec:c15-penalites}),
la continuité de service (section~\ref{sec:c15-continuite}) et la gestion
des fournisseurs (section~\ref{sec:c15-fournisseurs}). Le laboratoire
GCI-LAB-08 en fait réaliser les calculs sur les données de KRIBILOG.

\separateur

% ============================================================================
\section{Les Engagements de Service}
\label{sec:c15-sla}
% ============================================================================

\subsection{Trois niveaux d'accord}

\begin{boxdef}
\termecle{Accord de niveau de service}\index{accord de niveau de service}
(\textit{service level agreement}, \sigle{SLA})~: accord documenté entre un
fournisseur de services et son client, qui précise le service fourni et
les niveaux de service attendus.

\medskip
\termecle{Accord de niveau opérationnel}\index{accord de niveau
opérationnel} (\textit{operational level agreement}, \sigle{OLA})~: accord
entre le fournisseur et une autre équipe interne de la même organisation,
qui soutient la fourniture du service.

\medskip
\termecle{Contrat de sous-traitance}\index{contrat de sous-traitance}
(\textit{underpinning contract})~: contrat entre le fournisseur et un
tiers extérieur dont il dépend pour fournir le service.
\end{boxdef}
{\small\itshape Notions reformulées d'après ITIL~4~\cite{axelos2019itil4}.}

\noindent\begin{minipage}{\linewidth}
\begin{center}
\begin{tikzpicture}[
    acteur/.style = {
        draw=CouleurBordDef, fill=CouleurDef, rounded corners=4pt,
        text width=2.8cm, minimum height=1.0cm, align=center,
        font=\footnotesize\bfseries\color{CouleurPrimaire}
    },
    accord/.style = {font=\scriptsize\bfseries\color{CouleurAccent}},
    lien/.style = {{Stealth[length=5pt]}-{Stealth[length=5pt]}, thick,
        color=CouleurSecondaire}
]
\node[acteur] (cl) at (0,0)    {Armateur\\ (client)};
\node[acteur] (kr) at (4.6,0)  {KRIBILOG\\ (fournisseur)};
\node[acteur] (in) at (4.6,-2.2) {Équipe réseau\\ interne};
\node[acteur] (np) at (10.0,0) {NETPORT SERVICES\\ (infogérant)};
\draw[lien] (cl) -- node[above, accord]{SLA} (kr);
\draw[lien] (kr) -- node[right, accord]{OLA} (in);
\draw[lien] (kr) -- node[below, accord, align=center]{contrat de\\
    sous-traitance} (np);
\end{tikzpicture}
\end{center}
\vspace{-0.4em}
\begin{center}{\small\itshape Figure~15.1. SLA, OLA et contrat de
sous-traitance du service de suivi des conteneurs}\end{center}
\end{minipage}

\medskip

La règle de cohérence est simple~: \textbf{on ne peut pas promettre au
client plus que ce que garantissent les accords qui soutiennent le
service}. Si KRIBILOG promet à l'armateur un rétablissement en 4~heures,
mais que le contrat de NETPORT SERVICES ne fixe aucun délai
d'intervention, la promesse ne repose sur rien.

\subsection{Rédiger un SLA utile}

\begin{tableaucours}{L{3.2cm}L{4.6cm}L{4.4cm}}{%
    \tableauHeader{Rubrique} &
    \tableauHeader{Contenu} &
    \tableauHeader{Exemple KRIBILOG}
}
Description du service & Ce que le service permet au client
d'obtenir (le résultat) & Localiser en temps réel les conteneurs de
l'armateur \\
Horaires de service & Périodes pendant lesquelles les engagements
s'appliquent & 24~h/24, 7~j/7 \\
Disponibilité & Part du temps de service où le service est utilisable &
99,5~\% par mois \\
Délais de rétablissement & Par priorité (Chapitre~\ref{chap:itil-incidents})
& P1 en 4~heures, P2 en 8~heures \\
Indicateur de résultat & Mesure ce que le client vit & Part des conteneurs
localisés en moins de 5~minutes~: 99~\% \\
Mesure et compte rendu & Source des données, fréquence & Rapport mensuel
produit par KRIBILOG à partir de son outil \\
Pénalités & Conséquences du non-respect & Section~\ref{sec:c15-penalites} \\
Revue & Fréquence de révision de l'accord & Comité de service
trimestriel \\
\end{tableaucours}
\vspace{-0.8em}
\begin{center}{\small\itshape Tableau~15.1. Les rubriques d'un accord de
niveau de service}\end{center}

\begin{boxalert}
\textbf{Violation typique.} Tous les indicateurs du rapport mensuel sont
au vert, mais le client est mécontent. ITIL~4 appelle ce phénomène
l'« effet pastèque »~: vert à l'extérieur, rouge à l'intérieur. Il naît
d'engagements qui mesurent des livrables (serveurs disponibles) au lieu
de résultats (conteneurs localisés).

\smallskip
\textbf{Règle~:} chaque SLA comporte au moins un indicateur de résultat,
défini avec le client, et la satisfaction du client est mesurée à côté
des indicateurs techniques.
\end{boxalert}

\begin{boxobjet}{APO09 · Gérer les accords de service}
\textbf{Finalité~:} s'assurer que les services et les niveaux de service
répondent aux besoins actuels et futurs de l'organisation.

\textbf{Objectifs liés~:} EG01, EG08~; AG05.

\textbf{Pratiques~:} APO09.01 identifier les services~; APO09.02
cataloguer les services~; APO09.03 définir et préparer les accords de
service~; APO09.04 surveiller les niveaux de service et en rendre
compte~; APO09.05 revoir les accords de service et les contrats.

\textbf{RACI (extrait, KRIBILOG)~:} APO09.03~: \emph{A} au directeur des
opérations, \emph{R} au responsable des services informatiques, \emph{C}
à l'armateur et à NETPORT SERVICES.

\textbf{Indicateurs~:} part des services couverts par un SLA, part des
engagements tenus, satisfaction des clients.
\end{boxobjet}
{\small\itshape Pratiques et objectifs liés traduits par nos soins à partir
de COBIT~2019~\cite{isaca2018cobitobjectifs}~; finalité reformulée.}

% ============================================================================
\section{Calculer une Disponibilité}
\label{sec:c15-dispo}
% ============================================================================

\subsection{La formule}

\begin{boxdef}
\termecle{Disponibilité}\index{disponibilité}~: aptitude d'un service à
remplir sa fonction quand on en a besoin. Elle se mesure par
\[
    A = \frac{T_{\text{convenu}} - T_{\text{indisponible}}}{T_{\text{convenu}}}
    \times 100,
\]
où $T_{\text{convenu}}$ est le temps de service convenu sur la période et
$T_{\text{indisponible}}$ la durée cumulée des interruptions pendant ce
temps.
\end{boxdef}

Sur un mois de 30~jours en service 24~h/24, le temps convenu est de
720~heures. Une disponibilité de 99,5~\% autorise donc $0{,}005 \times 720
= 3{,}6$~heures d'interruption par mois~; 99,9~\% en autorisent
43~minutes~; 99,99~\%, un peu plus de 4~minutes.

\begin{tableaucours}{L{3.0cm}L{4.6cm}L{4.6cm}}{%
    \tableauHeader{Disponibilité} &
    \tableauHeader{Interruption par mois (720 h)} &
    \tableauHeader{Interruption par an (8~760 h)}
}
99~\% & 7,2 heures & 87,6 heures \\
99,5~\% & 3,6 heures & 43,8 heures \\
99,9~\% & 43 minutes & 8,8 heures \\
99,99~\% & 4,3 minutes & 53 minutes \\
\end{tableaucours}
\vspace{-0.8em}
\begin{center}{\small\itshape Tableau~15.2. Ce que signifie un taux de
disponibilité}\end{center}

\subsection{Composants en série et en parallèle}

Un service dépend de plusieurs composants. S'ils sont \textbf{en série}
(il les faut tous), la disponibilité du service est le produit de leurs
disponibilités~: $A = A_1 \times A_2 \times \cdots \times A_n$. S'ils sont
\textbf{en parallèle} (un seul suffit, les autres sont des secours), la
disponibilité est $A = 1 - (1-A_1)(1-A_2)\cdots(1-A_n)$.

\begin{boxscenario}{Opération KRIBI. La disponibilité réelle du suivi}
\textbf{Composants en série.} Serveur d'application (99,9~\%), base de
données (99,8~\%), passerelle des balises (99,5~\%), réseau de
l'opérateur mobile (99,0~\%).

\medskip
\textbf{Calcul.} $0{,}999 \times 0{,}998 \times 0{,}995 \times 0{,}990 =
0{,}9821$, soit \textbf{98,21~\%}~: environ 12,9~heures d'interruption
attendues par mois, alors que l'engagement de NETPORT SERVICES sur les
seuls serveurs (99,5~\%) est tenu.

\medskip
\textbf{Avec un second opérateur mobile en secours.} Le réseau passe à
$1 - (0{,}01 \times 0{,}01) = 0{,}9999$, et le service à $0{,}999 \times
0{,}998 \times 0{,}995 \times 0{,}9999 = 0{,}9919$, soit 99,19~\%.

\medskip
\textbf{Analyse} (confiance élevée). Le maillon faible n'est pas le
serveur, sur lequel porte le contrat, mais le réseau mobile, sur lequel
il ne porte pas. Un SLA sur un composant ne garantit pas le service.
\end{boxscenario}

% ============================================================================
\section{Les Pénalités}
\label{sec:c15-penalites}
% ============================================================================

\subsection{Un mécanisme contractuel}

Les pénalités (souvent appelées « crédits de service ») sanctionnent le
non-respect des engagements. Elles ont trois fonctions~: inciter le
fournisseur, compenser partiellement le client, et documenter les
manquements. Elles ne compensent presque jamais le préjudice réel.

\begin{boxscenario}{Opération KRIBI. Le calcul du mois de juin}
\textbf{Clause du nouveau contrat.} Redevance mensuelle~: 25~millions
FCFA. Engagement~: 99,5~\% de disponibilité du service de suivi, mesurée
de bout en bout. Pénalité~: 4~\% de la redevance par tranche complète de
0,1~point en dessous de l'engagement, plafonnée à 30~\%.

\medskip
\textbf{Faits du mois.} Trois interruptions de 2~heures, 4,5~heures et
1,5~heure, soit 8~heures sur 720.

\medskip
\textbf{Calcul.} Disponibilité~: $(720-8)/720 = 98{,}89~\%$. Écart~: $99{,}5
- 98{,}89 = 0{,}61$~point, soit 6~tranches complètes de 0,1~point.
Pénalité~: $6 \times 4~\% = 24~\%$, sous le plafond, soit $0{,}24 \times 25 =
6$~millions FCFA.

\medskip
\textbf{Mise en perspective.} L'interruption de 4,5~heures a immobilisé
un navire une demi-journée~: environ 6~millions FCFA de surestaries à elle
seule. La pénalité du mois couvre une seule interruption sur trois.
\end{boxscenario}

\begin{boiteRemarque}{Faut-il des pénalités plus lourdes~?}
Pas nécessairement. Des pénalités très lourdes renchérissent le contrat,
poussent le fournisseur à contester les mesures et ne remplacent pas un
service fiable. La protection réelle du client vient d'ailleurs~: des
engagements de résultat, une architecture redondante, un plan de
continuité, un droit de résiliation en cas de manquements répétés et une
supervision effective du fournisseur. La pénalité est un signal, pas une
assurance.
\end{boiteRemarque}

% ============================================================================
\section{La Continuité de Service}
\label{sec:c15-continuite}
% ============================================================================

\subsection{L'analyse d'impact sur l'activité}

La continuité vise à maintenir les services à un niveau acceptable en cas
de sinistre majeur~: incendie du centre de données, coupure prolongée
d'électricité, cyberattaque destructrice. Elle commence par une analyse
d'impact sur l'activité.

\begin{boxdef}
\termecle{Analyse d'impact sur l'activité}\index{analyse d'impact sur
l'activité} (\textit{business impact analysis}, \sigle{BIA})~: analyse qui
identifie les activités essentielles, mesure l'effet de leur interruption
dans le temps, et en déduit les objectifs de reprise.

\medskip
\termecle{Durée maximale d'interruption admissible}\index{durée maximale
d'interruption admissible} (\sigle{DMIA})~: durée au-delà de laquelle
l'interruption d'une activité devient inacceptable pour l'organisation.

\medskip
\termecle{RTO}\index{RTO} (\textit{recovery time objective}, durée
maximale de reprise)~: délai cible dans lequel le service doit être
rétabli après un sinistre. Il doit être inférieur à la DMIA.

\medskip
\termecle{RPO}\index{RPO} (\textit{recovery point objective}, perte de
données maximale admissible)~: quantité maximale de données, exprimée en
durée, que l'on accepte de perdre. Un RPO de 15~minutes impose une
réplication ou une sauvegarde au moins toutes les 15~minutes.
\end{boxdef}

\begin{tableaucours}{L{3.6cm}L{2.0cm}L{1.6cm}L{1.6cm}L{3.0cm}}{%
    \tableauHeader{Service} &
    \tableauHeader{DMIA} &
    \tableauHeader{RTO} &
    \tableauHeader{RPO} &
    \tableauHeader{Dispositif}
}
Système central (SANAGA) & 8 heures & 4 h & 15 min & Site de secours
répliqué \\
Paie des fonctionnaires (ANPP) & 48 heures & 24 h & 1 h & Sauvegarde
horaire, site de secours \\
Suivi des conteneurs (KRIBILOG) & 4 heures & 2 h & 15 min & Réplication,
second opérateur \\
Messagerie (KRIBILOG) & 3 jours & 24 h & 24 h & Sauvegarde quotidienne \\
\end{tableaucours}
\vspace{-0.8em}
\begin{center}{\small\itshape Tableau~15.3. Objectifs de reprise de
quelques services du cours}\end{center}

\subsection{Le plan et ses tests}

Le plan de continuité ne vaut que s'il est \emph{testé}. COBIT y consacre
l'objectif DSS04~; la norme ISO~22301 en formule les exigences
(Chapitre~\ref{chap:normes}).

\begin{boxobjet}{DSS04 · Gérer la continuité}
\textbf{Finalité~:} s'adapter rapidement, poursuivre les activités et
maintenir la disponibilité des ressources et de l'information à un niveau
acceptable en cas de perturbation majeure.

\textbf{Objectifs liés~:} EG01, EG02, EG06, EG08~; AG05, AG07.

\textbf{Pratiques~:} DSS04.01 définir la politique, les objectifs et le
périmètre de la continuité~; DSS04.02 maintenir la résilience~; DSS04.03
élaborer et mettre en œuvre la réponse~; DSS04.04 exercer, tester et revoir
les plans~; DSS04.05 revoir, maintenir et améliorer les plans~; DSS04.06
former~; DSS04.07 gérer les sauvegardes~; DSS04.08 conduire une revue après
la reprise.

\textbf{Indicateurs~:} part des services critiques dont le plan a été testé
dans l'année~; écart entre RTO visé et durée de reprise constatée lors des
tests.
\end{boxobjet}
{\small\itshape Finalité, objectifs liés et pratiques résumés et traduits
par nos soins à partir de COBIT~2019~\cite{isaca2018cobitobjectifs}.}

\begin{tableaucours}{L{3.4cm}L{4.6cm}L{4.2cm}}{%
    \tableauHeader{Type de test} &
    \tableauHeader{Déroulement} &
    \tableauHeader{Ce qu'il prouve}
}
Exercice sur table & Les responsables déroulent le plan sur un scénario,
autour d'une table & Que chacun connaît son rôle \\
Simulation & Les équipes exécutent les procédures sans basculer
réellement & Que les procédures sont praticables \\
Bascule réelle & Le service est réellement repris sur le site de secours
& Que le RTO et le RPO sont tenus \\
\end{tableaucours}
\vspace{-0.8em}
\begin{center}{\small\itshape Tableau~15.4. Types de tests de continuité}\end{center}

\begin{boxalert}
\textbf{Violation typique.} La BANQUE SANAGA disposait d'un plan de
continuité rédigé en 2022. Il n'avait jamais été testé, et le site de
secours qu'il désignait hébergeait encore l'ancienne version du système
central. Le jour de la panne, personne n'a pensé à l'ouvrir.

\smallskip
\textbf{Règle~:} un plan de continuité se teste au moins une fois par an
pour chaque service critique, et se met à jour à chaque changement majeur
(Chapitre~\ref{chap:itil-changements}).
\end{boxalert}

% ============================================================================
\section{Piloter les Fournisseurs}
\label{sec:c15-fournisseurs}
% ============================================================================

\subsection{L'externalisation ne transfère pas la responsabilité}

Le Chapitre~\ref{chap:defaillances} l'a montré avec TSB, le
Chapitre~\ref{chap:normes} avec l'ISO/IEC~20000-1~: confier un service à
un tiers ne décharge pas l'organisation de sa responsabilité. Le droit
camerounais le dit aussi.

\begin{boxjuris}
\textbf{\loicam{n°~2010/012}, article~26, alinéa~1} (extrait, version
anglaise officielle)~:

\medskip
\textit{« Operators of information systems shall take every technical and
administrative measure to ensure the security of services offered. To
this end, they shall have standardized systems enabling them to at all
times identify, assess, process or manage any risk relating to the
security of the information systems of the services provided directly or
indirectly. »}

\medskip
\textbf{Portée.} L'obligation de gérer les risques de sécurité porte sur
les services fournis « directement ou indirectement ». Une organisation qui
exploite ses systèmes par l'intermédiaire d'un infogérant reste tenue
d'identifier, d'évaluer et de traiter les risques de sécurité de ces
services~: elle doit donc s'en donner les moyens dans le contrat (droit
d'audit, information sur les incidents, exigences de sécurité).
\end{boxjuris}

\subsection{Le cycle de la relation fournisseur}

\begin{boxobjet}{APO10 · Gérer les fournisseurs}
\textbf{Finalité~:} obtenir des fournisseurs les services convenus, en
maîtrisant leurs risques, et en entretenant des relations efficaces.

\textbf{Objectifs liés~:} EG01, EG08~; AG05.

\textbf{Pratiques~:} APO10.01 identifier et évaluer les relations et les
contrats~; APO10.02 sélectionner les fournisseurs~; APO10.03 gérer les
relations et les contrats~; APO10.04 gérer les risques liés aux
fournisseurs~; APO10.05 surveiller la performance et la conformité des
fournisseurs.

\textbf{RACI (extrait, KRIBILOG)~:} APO10.05~: \emph{A} au responsable des
services informatiques, \emph{R} au gestionnaire du contrat, \emph{C} à la
direction des achats et au RSSI.

\textbf{Indicateurs~:} part des fournisseurs critiques évalués dans
l'année, nombre de manquements contractuels, délai de traitement des
écarts.
\end{boxobjet}
{\small\itshape Pratiques et objectifs liés traduits par nos soins à partir
de COBIT~2019~\cite{isaca2018cobitobjectifs}~; finalité reformulée.}

\subsection{Les clauses essentielles}

\begin{boxprocedure}
\textbf{Le contrat d'un fournisseur critique.} Pour qu'un auditeur puisse
établir que l'organisation maîtrise un service externalisé, le contrat
doit comporter au minimum~:
\begin{itemize}
    \item des engagements de service mesurables, dont au moins un
          indicateur de résultat, et la méthode de mesure~;
    \item des pénalités et un droit de résiliation pour manquements
          répétés~;
    \item l'obligation de signaler les incidents et les failles de sécurité
          dans un délai fixé~;
    \item un droit d'audit, par l'organisation ou par un tiers qu'elle
          désigne~;
    \item la localisation des données et les conditions de tout transfert
          (loi n°~2024/017, Chapitre~\ref{chap:droit})~;
    \item les conditions de la sous-traitance en cascade~;
    \item un plan de continuité du fournisseur, testé~;
    \item une clause de réversibilité~: restitution des données et
          assistance au transfert vers un autre prestataire.
\end{itemize}
\end{boxprocedure}

\begin{boxmetier}{Responsable des services IT}
\textbf{Le comité de service trimestriel avec NETPORT SERVICES}
\begin{enumerate}
    \item Revoir les indicateurs du trimestre, mesurés par KRIBILOG et non
          par le seul fournisseur.
    \item Examiner chaque incident P1 et P2, et les problèmes ouverts
          imputables au fournisseur.
    \item Arrêter les pénalités et les plans d'action.
    \item Revoir les risques du fournisseur (dépendance à une personne,
          santé financière, sous-traitants).
    \item Consigner les décisions au procès-verbal, transmis au comité
          stratégique des SI.
\end{enumerate}
\end{boxmetier}

\begin{boxmetier}{Auditeur SI}
\textbf{Auditer un fournisseur critique}
\begin{enumerate}
    \item Vérifier que le contrat contient les clauses essentielles
          ci-dessus.
    \item Recalculer, sur trois mois, la disponibilité et les pénalités à
          partir des données de l'organisation, et les comparer au rapport
          du fournisseur.
    \item Vérifier que les pénalités dues ont été effectivement appliquées.
    \item Vérifier le dernier test du plan de continuité du fournisseur.
\end{enumerate}
\end{boxmetier}

% ============================================================================
\separateur
% ============================================================================

\begin{boxresume}
\textbf{Ce qu'il faut retenir de ce chapitre~:}
\begin{itemize}
    \item \textbf{SLA, OLA, contrat de sous-traitance}~: on ne promet pas au
          client plus que ce que garantissent les accords qui soutiennent le
          service.
    \item \textbf{Un SLA utile} mesure au moins un résultat~; sinon
          apparaît l'effet pastèque.
    \item \textbf{La disponibilité se calcule}~: 99,5~\% sur 720~heures
          autorisent 3,6~heures d'interruption. En série, les
          disponibilités se multiplient~; en parallèle, les secours se
          composent.
    \item \textbf{Les pénalités} sont un signal, pas une assurance~: elles
          couvrent rarement le préjudice réel.
    \item \textbf{La continuité} part de l'analyse d'impact (DMIA, RTO,
          RPO) et ne vaut que si le plan est testé.
    \item \textbf{L'externalisation ne transfère pas la responsabilité}~:
          la loi n°~2010/012 vise les services fournis directement ou
          indirectement. Le contrat doit donner les moyens du contrôle.
\end{itemize}
\end{boxresume}

\bigskip

\begin{boxexo}[title={\ding{45}\quad Exercice~15.1. Calculs de disponibilité \niveaubase}]
\begin{enumerate}[(a)]
    \item Combien d'heures d'interruption par mois de 30~jours autorise une
          disponibilité de 99,8~\%~?
    \item Un service a connu 5~heures d'interruption en avril (30~jours).
          Quelle est sa disponibilité~?
    \item Un service dépend de trois composants en série disponibles à
          99,9~\%, 99,7~\% et 99,5~\%. Quelle est sa disponibilité~?
    \item On double le composant le plus faible par un secours identique.
          Quelle est la nouvelle disponibilité~?
\end{enumerate}
\end{boxexo}

\begin{boxexo}[title={\ding{45}\quad Exercice~15.2. Les pénalités de juillet \niveaumoyen}]
Avec la clause du scénario de la section~\ref{sec:c15-penalites}, le mois
de juillet (31~jours) a connu quatre interruptions de 3~heures,
2~heures, 6~heures et 1~heure.
\begin{enumerate}[(a)]
    \item Calculez la disponibilité du mois.
    \item Calculez la pénalité en FCFA.
    \item NETPORT SERVICES conteste l'interruption de 6~heures, qu'elle
          attribue à l'opérateur mobile. Que dit le contrat, et que
          devrait-il dire~?
\end{enumerate}
\end{boxexo}

\begin{boxexo}[title={\ding{45}\quad Exercice~15.3. La continuité de l'ANPP \niveauexpert}]
L'ANPP verse chaque mois les salaires de plus de 300~000 agents publics.
\begin{enumerate}[(a)]
    \item Conduisez une analyse d'impact simplifiée de ses trois services
          principaux (paie, paiement des fournisseurs, portail
          d'information), et fixez pour chacun DMIA, RTO et RPO.
    \item Proposez un dispositif de continuité proportionné, et le programme
          de tests de l'année.
    \item Identifiez les clauses à imposer à l'hébergeur de l'agence.
    \item Indiquez à quelle instance de gouvernance les résultats des tests
          doivent être présentés, et pourquoi.
\end{enumerate}
\end{boxexo}

\bigskip

\begin{boxinfo}
\textbf{Transition.} Les Chapitres~\ref{chap:itil-svs} à~\ref{chap:itil-niveaux}
ont présenté les pratiques d'ITIL~4 une par une. Le
Chapitre~\ref{chap:processus} apprend à les assembler~: modéliser un
processus, en attribuer les responsabilités par une matrice RACI, le
mesurer, et établir la correspondance entre COBIT, ITIL et
l'ISO/IEC~20000-1.
\end{boxinfo}

        % ============================================================================
% Fichier  : partie4/P04_C16_processus_raci_mesure.tex
% Titre    : Processus : Modéliser, Responsabiliser, Mesurer
% Auteur   : MINKA MI NGUIDJOÏ Thierry Emmanuel
% Version  : 1.0 -- octobre 2026
% ============================================================================
% Sources : BPMN 2.0 (OMG, ISO/IEC 19510) ; COBIT 2019 (dossier Normes) ;
% ITIL 4 (sources publiques). La correspondance avec les articles de
% l'ISO/IEC 20000-1:2018 est établie par nos soins (norme absente du
% dossier Normes) : A VERIFIER dans le texte de la norme.
% Laboratoire associé : GCI-LAB-09 (matrice RACI et correspondance COBIT-ITIL
% d'un processus de NOUN PAY).
% ============================================================================

\chapter{Processus : Modéliser, Responsabiliser, Mesurer}
\label{chap:processus}
\chaptermark{Processus : modéliser, responsabiliser}

\epigraph{%
    « Quand tout le monde est responsable, personne ne l'est. »%
}{Adage du management}

\niveauChapitre{Expert}{%
    Maîtrise des Chapitres~\ref{chap:cobit-objectifs} (les objectifs COBIT)
    et~\ref{chap:itil-niveaux} (niveaux de service et fournisseurs).%
}

\begin{boxobjectifs}
\begin{itemize}
    \item Modéliser un processus en BPMN.
    \item Construire une matrice RACI sans ambiguïté.
    \item Mettre en correspondance COBIT, ITIL et ISO/IEC~20000-1.
    \item Définir les indicateurs (KPI, KGI) et un tableau de bord.
    \item Désigner les propriétaires de processus et organiser
          l'amélioration continue.
\end{itemize}
\end{boxobjectifs}

\bigskip

% ============================================================================
\section*{Question directrice : qui fait quoi, et comment le sait-on~?}
% ============================================================================

Les Chapitres~\ref{chap:itil-svs} à~\ref{chap:itil-niveaux} ont présenté
les pratiques une par une. Mais une organisation ne vit pas des pratiques
isolées~: elle vit des enchaînements d'activités, qui traversent les
équipes, les fournisseurs et les outils. Chez NOUN PAY, à Bafoussam, la
mise à jour du modèle de notation de crédit passe par cinq personnes et un
fournisseur de nuage~; personne ne sait qui l'autorise. Chez KRIBILOG,
NETPORT SERVICES et le centre de services se renvoient les incidents.

Ce chapitre apprend à assembler les pratiques en processus maîtrisés~: les
modéliser (section~\ref{sec:c16-bpmn}), répartir les responsabilités
(section~\ref{sec:c16-raci}), établir la correspondance entre les
référentiels (section~\ref{sec:c16-correspondance}), mesurer
(section~\ref{sec:c16-mesure}) et améliorer (section~\ref{sec:c16-proprietaires}).
Le laboratoire GCI-LAB-09 fait construire la matrice RACI et la
correspondance d'un processus de NOUN PAY.

\separateur

% ============================================================================
\section{Processus et Pratiques}
\label{sec:c16-notions}
% ============================================================================

\begin{boxdef}
\termecle{Processus}\index{processus}~: ensemble d'activités corrélées qui
transforme des entrées en sorties, en vue d'un objectif. Il a un début, une
fin, des acteurs et des résultats mesurables.

\medskip
\termecle{Propriétaire de processus}\index{propriétaire de processus}
(\textit{process owner})~: personne qui répond de la conception, de la
performance et de l'amélioration d'un processus dans toute
l'organisation. Elle ne l'exécute pas nécessairement elle-même.

\medskip
\termecle{Gestionnaire de processus}\index{gestionnaire de processus}
(\textit{process manager})~: personne qui pilote l'exécution quotidienne du
processus, sur un périmètre donné.
\end{boxdef}

Un processus n'est pas une pratique. Une \emph{pratique} ITIL ou un
\emph{objectif} COBIT décrivent un domaine de maîtrise, avec toutes ses
composantes (Chapitre~\ref{chap:cobit-archi}). Un \emph{processus} est
l'enchaînement concret par lequel une organisation réalise une partie de
ce domaine. La pratique « gestion des incidents » peut donner lieu à
plusieurs processus~: incidents ordinaires, incidents majeurs, incidents
de sécurité.

% ============================================================================
\section{Modéliser en BPMN}
\label{sec:c16-bpmn}
% ============================================================================

\subsection{Les éléments du langage}

Le langage BPMN (\textit{Business Process Model and Notation}), publié par
l'Object Management Group dans sa version~2.0 et repris par la norme
ISO/IEC~19510, est la notation de référence pour modéliser les
processus. Il suffit de quelques éléments pour en tirer l'essentiel.

\begin{tableaucours}{L{3.0cm}L{3.6cm}L{5.6cm}}{%
    \tableauHeader{Élément} &
    \tableauHeader{Représentation} &
    \tableauHeader{Usage}
}
Événement & Cercle (début, intermédiaire, fin) & Ce qui déclenche, ponctue
ou termine le processus \\
Activité & Rectangle arrondi & Une tâche ou un sous-processus \\
Passerelle & Losange & Un choix (exclusif), une séparation en parallèle,
une jonction \\
Flux de séquence & Flèche pleine & L'ordre des activités \\
Flux de messages & Flèche en pointillé & Un échange entre deux acteurs \\
Couloir & Bande horizontale & Le rôle ou l'organisation qui exécute les
activités \\
\end{tableaucours}
\vspace{-0.8em}
\begin{center}{\small\itshape Tableau~16.1. Les éléments essentiels de
BPMN~2.0}\end{center}

\subsection{Un exemple}

\noindent\begin{minipage}{\linewidth}
\begin{center}
\begin{tikzpicture}[
    tache/.style = {
        draw=CouleurBordDef, fill=CouleurDef, rounded corners=3pt,
        text width=1.8cm, minimum height=0.9cm, align=center,
        font=\scriptsize\color{CouleurPrimaire}
    },
    evt/.style = {circle, draw=CouleurPrimaire, thick, minimum size=0.5cm},
    evtfin/.style = {circle, draw=CouleurPrimaire, line width=2pt,
        minimum size=0.5cm},
    pass/.style = {diamond, draw=CouleurAccent, fill=CouleurAccent!15,
        minimum size=0.8cm, font=\scriptsize\bfseries\color{CouleurPrimaire},
        inner sep=0pt},
    fleche/.style = {-{Stealth[length=4pt]}, thick, color=CouleurSecondaire},
    couloir/.style = {font=\scriptsize\bfseries\color{CouleurPrimaire},
        rotate=90, anchor=south}
]
% Couloirs
\draw[CouleurPrimaire] (-0.6,0.9) rectangle (13.0,-1.1);
\draw[CouleurPrimaire] (-0.6,-1.1) rectangle (13.0,-3.1);
\node[anchor=north west, font=\tiny\bfseries\color{CouleurPrimaire}]
    at (-0.6,0.9) {Centre de services (KRIBILOG)};
\node[anchor=north west, font=\tiny\bfseries\color{CouleurPrimaire}]
    at (-0.6,-1.1) {NETPORT SERVICES (infogérant)};
% Centre de services
\node[evt] (d) at (0.2,-0.1) {};
\node[tache] (e) at (1.8,-0.1) {Enregistrer et prioriser};
\node[pass] (p) at (3.7,-0.1) {X};
\node[tache] (r1) at (5.5,-0.1) {Résoudre au niveau~1};
\node[tache] (v) at (10.2,-0.1) {Vérifier avec l'utilisateur};
\node[evtfin] (f) at (12.2,-0.1) {};
% NETPORT
\node[tache] (dg) at (5.5,-2.1) {Diagnostiquer};
\node[tache] (r2) at (8.0,-2.1) {Résoudre ou contourner};
% Flux
\draw[fleche] (d) -- (e);
\draw[fleche] (e) -- (p);
\draw[fleche] (p) -- node[above, font=\tiny]{connu} (r1);
\draw[fleche] (p) |- node[pos=0.2, right, font=\tiny]{inconnu} (dg);
\draw[fleche] (dg) -- (r2);
\draw[fleche] (r2) -| (v);
\draw[fleche] (r1) -- (v);
\draw[fleche] (v) -- (f);
\end{tikzpicture}
\end{center}
\vspace{-0.4em}
\begin{center}{\small\itshape Figure~16.1. Processus de traitement d'un
incident chez KRIBILOG (BPMN simplifié)}\end{center}
\end{minipage}

\medskip

La figure fait apparaître ce que le texte des procédures masquait~: le
passage d'un couloir à l'autre. C'est à ces frontières que se perdent les
incidents. Chez KRIBILOG, la flèche de retour de NETPORT SERVICES vers la
vérification n'existait pas~: l'infogérant clôturait lui-même les tickets
(Chapitre~\ref{chap:itil-incidents}).

\begin{boxalert}
\textbf{Violation typique.} Une organisation modélise ses processus tels
qu'ils \emph{devraient} être, à partir des référentiels, sans observer ce
qui se passe réellement. Les schémas sont beaux, mais personne ne les
applique.

\smallskip
\textbf{Règle~:} modéliser d'abord le processus \emph{tel qu'il est}, à
partir d'entretiens et des traces de l'outil (principe « commencer là où
l'on est », Chapitre~\ref{chap:itil-svs}), puis le processus cible, et
mesurer l'écart.
\end{boxalert}

% ============================================================================
\section{Responsabiliser : la Matrice RACI}
\label{sec:c16-raci}
% ============================================================================

\subsection{Quatre rôles}

\begin{boxdef}
\termecle{Matrice RACI}\index{matrice RACI}~: tableau qui croise les
activités d'un processus et les rôles de l'organisation, et indique pour
chaque case le rôle tenu~:
\begin{itemize}
    \item \textbf{R} (\textit{Responsible})~: réalise l'activité~;
    \item \textbf{A} (\textit{Accountable})~: en répond et l'approuve~;
    \item \textbf{C} (\textit{Consulted})~: est consulté avant~;
    \item \textbf{I} (\textit{Informed})~: est informé après.
\end{itemize}
\end{boxdef}

COBIT~2019 fournit pour chaque pratique une matrice RACI type, établie
avec des rôles génériques (Chapitre~\ref{chap:cobit-objectifs})~; il
revient à chaque organisation de la traduire avec ses propres fonctions.

\subsection{Cinq règles}

\begin{enumerate}
    \item \textbf{Un seul A par activité.} Deux personnes qui « répondent »
          d'une même activité, c'est n'en avoir aucune.
    \item \textbf{Au moins un R par activité.} Une activité sans exécutant
          ne sera pas faite.
    \item \textbf{Le A au bon niveau.} Celui qui répond doit avoir
          l'autorité de décider et les moyens d'agir.
    \item \textbf{Peu de C.} Trop de consultations ralentissent tout~;
          consulter, c'est s'engager à attendre la réponse.
    \item \textbf{Des noms de fonctions, pas de personnes.} La matrice
          survit aux départs.
\end{enumerate}

\begin{boxscenario}{Opération NOUN. Qui autorise~?}
\textbf{Faits.} Chez NOUN PAY, une nouvelle version du modèle de notation
de crédit est préparée par l'équipe de science des données, testée par le
directeur technique, déployée par l'équipe DevOps sur l'infrastructure du
fournisseur de nuage. La fondatrice est « tenue au courant ». Le délégué à
la protection des données n'est pas associé.

\medskip
\textbf{Matrice constatée.}

\medskip
\begin{tableaucours}{L{3.4cm}ccccc}{%
    \tableauHeader{Activité} &
    \tableauHeader{DG} &
    \tableauHeader{Dir. tech.} &
    \tableauHeader{Données} &
    \tableauHeader{DevOps} &
    \tableauHeader{DPD}
}
Préparer la version & I & A & R & · & · \\
Tester & · & A/R & C & · & · \\
Évaluer l'impact sur les clients & · & · & · & · & · \\
Autoriser la mise en production & I & A & · & A & · \\
Déployer & · & · & · & A/R & · \\
\end{tableaucours}

\medskip
\textbf{Constatation.} L'activité « évaluer l'impact sur les clients » n'a
ni R ni A. L'autorisation compte deux A. Le délégué à la protection des
données n'apparaît nulle part.

\medskip
\textbf{Analyse} (confiance élevée). Trois règles sur cinq sont violées.
C'est l'explication directe de la demande de l'Autorité de protection des
données (Chapitre~\ref{chap:droit})~: personne n'était chargé de vérifier
qu'une nouvelle version ne pénalisait pas une catégorie de clients.
\end{boxscenario}

\begin{boxprocedure}
\textbf{Une matrice RACI auditable.} Pour qu'une matrice RACI soit
opposable, elle doit~:
\begin{itemize}
    \item porter sur des activités précises, rattachées à une pratique ou à
          un objectif identifié~;
    \item utiliser les fonctions réelles de l'organisation, telles qu'elles
          figurent à l'organigramme~;
    \item respecter les règles « un seul A » et « au moins un R »~;
    \item être approuvée par le propriétaire du processus et datée~;
    \item être cohérente avec les fiches de poste et les délégations de
          pouvoir.
\end{itemize}
\end{boxprocedure}

% ============================================================================
\section{COBIT, ITIL et ISO/IEC 20000-1}
\label{sec:c16-correspondance}
% ============================================================================

\subsection{Pourquoi une correspondance}

Une organisation qui utilise COBIT pour sa gouvernance, ITIL pour ses
pratiques et vise l'ISO/IEC~20000-1 doit éviter de construire trois
dispositifs parallèles. La correspondance permet de n'avoir qu'un
processus, documenté une fois, qui répond aux trois référentiels.

\begin{tableaucours}{L{3.0cm}L{4.6cm}L{4.6cm}}{%
    \tableauHeader{COBIT 2019} &
    \tableauHeader{Pratiques ITIL 4} &
    \tableauHeader{ISO/IEC 20000-1:2018}
}
APO08 Relations & Gestion des relations & 8.3.2 Relations avec les
clients \\
APO09 Accords de service & Gestion des niveaux de service, du catalogue &
8.3.3 Niveaux de service~; 8.2.4 Catalogue \\
APO10 Fournisseurs & Gestion des fournisseurs & 8.3.4 Fournisseurs \\
APO06 Budget et coûts & Gestion financière des services & 8.4.1
Budgétisation et comptabilité \\
BAI04 Disponibilité et capacité & Gestion de la disponibilité, de la
capacité et des performances & 8.7.1 Disponibilité~; 8.4.3 Capacité \\
BAI06 Changements & Habilitation des changements & 8.5.1 Changements \\
BAI07 Transition & Mises en production, déploiements, validation et tests
& 8.5.2 Conception et transition~; 8.5.3 Mises en production et
déploiement \\
BAI09 Actifs & Gestion des actifs informatiques & 8.2.5 Actifs \\
BAI10 Configuration & Gestion de la configuration des services & 8.2.6
Configuration \\
DSS02 Demandes et incidents & Centre de services, incidents, demandes de
service & 8.6.1 Incidents~; 8.6.2 Demandes de service \\
DSS03 Problèmes & Gestion des problèmes & 8.6.3 Problèmes \\
DSS04 Continuité & Gestion de la continuité des services & 8.7.2
Continuité \\
APO13, DSS05 Sécurité & Gestion de la sécurité de l'information & 8.7.3
Sécurité de l'information \\
MEA01 Performance & Mesure et compte rendu & 9.1 Surveillance et mesure \\
\end{tableaucours}
\vspace{-0.8em}
\begin{center}{\small\itshape Tableau~16.2. Correspondance entre COBIT~2019,
ITIL~4 et l'ISO/IEC~20000-1:2018 (établie par nos soins~; numéros
d'articles de la norme à vérifier dans son texte)}\end{center}

\subsection{Les limites d'une correspondance}

Une correspondance n'est jamais une équivalence. Trois précautions
s'imposent.
\begin{itemize}
    \item \textbf{Les périmètres diffèrent.} L'objectif BAI07 de COBIT couvre
          la conversion des données et la revue après mise en œuvre, que la
          pratique ITIL de gestion des mises en production ne traite pas de
          la même manière.
    \item \textbf{Les natures diffèrent.} COBIT fixe un objectif, ITIL décrit
          une pratique, la norme formule une exigence~: satisfaire l'une ne
          garantit pas de satisfaire les autres.
    \item \textbf{Les éditions changent.} Une correspondance se date et se
          révise à chaque nouvelle édition d'un des référentiels.
\end{itemize}

\begin{boxmetier}{Responsable GRC}
\textbf{Construire un référentiel unique de processus}
\begin{enumerate}
    \item Partir des objectifs COBIT prioritaires
          (Chapitre~\ref{chap:cobit-conception}).
    \item Pour chacun, désigner le processus de l'organisation qui le
          réalise, et la pratique ITIL qui guide sa conception.
    \item Repérer les exigences de l'ISO/IEC~20000-1 si une certification
          est visée, et les obligations légales (Chapitre~\ref{chap:droit}).
    \item Documenter le processus une seule fois, avec un tableau de
          correspondance en annexe.
    \item Faire approuver l'ensemble par le comité stratégique des SI.
\end{enumerate}
\end{boxmetier}

% ============================================================================
\section{Mesurer : KPI, KGI et Tableau de Bord}
\label{sec:c16-mesure}
% ============================================================================

\subsection{Deux familles d'indicateurs}

\begin{boxdef}
\termecle{Indicateur clé d'objectif}\index{indicateur clé d'objectif}
(\textit{key goal indicator}, \sigle{KGI})~: indicateur qui dit si un
objectif est atteint. Il mesure un résultat, après coup.

\medskip
\termecle{Indicateur clé de performance}\index{indicateur clé de
performance} (\textit{key performance indicator}, \sigle{KPI})~:
indicateur qui dit si un processus fonctionne bien et a des chances
d'atteindre l'objectif. Il mesure une performance, en cours de route.
\end{boxdef}

Cette distinction, mise en avant par COBIT~4.1~\cite{itgi2007cobit41}, reste
utile~: COBIT~2019 propose pour chaque objectif d'entreprise, chaque
objectif d'alignement et chaque pratique des « exemples
d'indicateurs », qui relèvent de l'une ou l'autre famille.

\begin{tableaucours}{L{3.0cm}L{4.6cm}L{4.6cm}}{%
    \tableauHeader{Niveau} &
    \tableauHeader{Indicateur (KRIBILOG)} &
    \tableauHeader{Destinataire}
}
Objectif d'entreprise (KGI) & Nombre de navires retardés par une panne du
suivi & Conseil d'administration \\
Objectif d'alignement (KGI) & Disponibilité de bout en bout du service de
suivi & Comité stratégique des SI \\
Processus (KPI) & Taux de réouverture des incidents~; âge moyen des
problèmes ouverts & Responsable des services informatiques \\
Activité (KPI) & Part des tickets enregistrés dans les 5~minutes &
Responsable du centre de services \\
\end{tableaucours}
\vspace{-0.8em}
\begin{center}{\small\itshape Tableau~16.3. La cascade des indicateurs}\end{center}

\subsection{Un bon indicateur}

Un indicateur utile est \textbf{défini} (formule, source, fréquence),
\textbf{assorti d'une cible}, \textbf{attribué} à un responsable,
\textbf{difficile à manipuler} (accompagné de son indicateur de
contrôle, Chapitre~\ref{chap:itil-incidents}) et \textbf{suivi d'action}~:
un indicateur dont l'écart ne déclenche jamais rien est un indicateur
inutile.

\begin{boxobjet}{MEA01 · Gérer la surveillance de la performance et de la conformité}
\textbf{Finalité~:} assurer la transparence de la performance et de la
conformité, et piloter l'atteinte des objectifs.

\textbf{Contenu (résumé)~:} fixer l'approche de surveillance, définir les
cibles, collecter et traiter les données, analyser et rendre compte de la
performance, assurer la mise en œuvre des actions correctives.

\textbf{Sorties~:} tableau de bord, rapports de performance, plans d'action.

\textbf{RACI (extrait, KRIBILOG)~:} \emph{A} au directeur général,
\emph{R} au responsable des services informatiques et au responsable GRC,
\emph{I} au conseil.

\textbf{Indicateurs~:} part des indicateurs dotés d'une cible et d'un
responsable~; délai entre l'écart constaté et l'action décidée.
\end{boxobjet}
{\small\itshape Résumé établi par nos soins à partir de
COBIT~2019~\cite{isaca2018cobitobjectifs}.}

% ============================================================================
\section{Propriétaires et Amélioration Continue}
\label{sec:c16-proprietaires}
% ============================================================================

Un processus sans propriétaire se dégrade. Le propriétaire répond de
quatre choses~: la conception du processus et sa documentation, la
définition de ses indicateurs, sa performance dans toute l'organisation,
et son amélioration. Il conduit, au moins une fois par an, une revue du
processus~: indicateurs, incidents, constats d'audit, retours des
utilisateurs, et alimente le registre des améliorations
(Chapitre~\ref{chap:itil-svs}).

\begin{boxmetier}{DSI}
\textbf{Désigner les propriétaires de processus}
\begin{enumerate}
    \item Un propriétaire par processus prioritaire, de rang suffisant pour
          imposer ses décisions aux équipes concernées.
    \item Pour les processus qui traversent un fournisseur, un propriétaire
          \emph{interne}~: on n'externalise pas la propriété d'un processus.
    \item Une lettre de mission, qui précise les pouvoirs, les indicateurs et
          le rythme de la revue.
    \item Une présentation annuelle de chaque propriétaire au comité
          stratégique des SI.
\end{enumerate}
\end{boxmetier}

\begin{boiteRemarque}{Le processus n'est-il pas l'ennemi de l'agilité~?}
Un processus lourd, imposé d'en haut et jamais révisé, l'est sans doute.
Mais l'absence de processus n'est pas l'agilité~: c'est la dépendance aux
personnes. Les organisations agiles ont des processus courts, visibles,
automatisés quand c'est possible, et révisés souvent~; elles appliquent les
principes d'ITIL « rester simple et pratique » et « progresser par
itérations ». Le Chapitre~\ref{chap:agile-cloud} y reviendra.
\end{boiteRemarque}

% ============================================================================
\separateur
% ============================================================================

\begin{boxresume}
\textbf{Ce qu'il faut retenir de ce chapitre~:}
\begin{itemize}
    \item \textbf{Processus et pratique}~: la pratique décrit un domaine de
          maîtrise, le processus l'enchaînement concret par lequel
          l'organisation le réalise.
    \item \textbf{BPMN} rend visibles les frontières entre acteurs, là où se
          perdent les dossiers. Modéliser d'abord l'existant.
    \item \textbf{RACI}~: un seul A, au moins un R, le A au bon niveau, peu
          de C, des fonctions et non des personnes.
    \item \textbf{La correspondance COBIT, ITIL, ISO/IEC~20000-1} permet un
          seul processus pour trois référentiels~; ce n'est jamais une
          équivalence.
    \item \textbf{KGI et KPI}~: le premier dit si l'objectif est atteint, le
          second si le processus a des chances de l'atteindre. Les
          indicateurs se déclinent du conseil à l'activité.
    \item \textbf{Un propriétaire par processus}, interne même quand le
          processus traverse un fournisseur, qui en conduit la revue
          annuelle.
\end{itemize}
\end{boxresume}

\bigskip

\begin{boxexo}[title={\ding{45}\quad Exercice~16.1. Corriger une matrice RACI \niveaubase}]
Reprenez la matrice constatée de NOUN PAY (section~\ref{sec:c16-raci}).
\begin{enumerate}[(a)]
    \item Listez toutes les violations des cinq règles.
    \item Proposez une matrice corrigée, en ajoutant les activités qui
          manquent.
    \item Indiquez qui doit approuver la matrice corrigée.
\end{enumerate}
\end{boxexo}

\begin{boxexo}[title={\ding{45}\quad Exercice~16.2. Modéliser le processus de changement de l'ANPP \niveaumoyen}]
Modélisez en BPMN le processus de changement normal de l'ANPP, depuis la
demande d'une direction métier jusqu'à la clôture, avec quatre couloirs~:
direction métier, division informatique, comité des changements,
prestataire.
\begin{enumerate}[(a)]
    \item Faites figurer au moins une passerelle exclusive (autorisé ou
          refusé) et un flux de messages.
    \item Repérez les frontières où le dossier risque de se perdre.
    \item Établissez la matrice RACI correspondante.
    \item Rattachez chaque activité à une pratique de BAI06 et à l'article
          correspondant de l'ISO/IEC~20000-1.
\end{enumerate}
\end{boxexo}

\begin{boxexo}[title={\ding{45}\quad Exercice~16.3. Le tableau de bord de la BANQUE SANAGA \niveauexpert}]
Construisez le tableau de bord de la gouvernance des SI de la BANQUE
SANAGA pour l'année qui suit la panne.
\begin{enumerate}[(a)]
    \item Proposez deux KGI par niveau (objectifs d'entreprise prioritaires,
          objectifs d'alignement) et deux KPI pour chacun des dix objectifs
          prioritaires (Chapitre~\ref{chap:cobit-objectifs}).
    \item Pour chaque indicateur, précisez la formule, la source, la
          fréquence, la cible et le responsable.
    \item Associez à chaque indicateur susceptible d'être manipulé son
          indicateur de contrôle.
    \item Indiquez ce qui sera présenté au conseil, au comité stratégique
          des SI et au comité d'audit.
\end{enumerate}
\end{boxexo}

\bigskip

\begin{boxinfo}
\textbf{Transition.} Le cours a désormais présenté toutes les pièces~: la
gouvernance, COBIT, le cadre normatif et juridique, ITIL et les
processus. Le Chapitre~\ref{chap:cas-integrateurs} les assemble sur les
quatre cas fil rouge, en études de cas intégratrices.
\end{boxinfo}

        % ============================================================================
% Fichier  : partie4/P04_C17_cas_integrateurs.tex
% Titre    : Études de Cas Intégrateurs
% Auteur   : MINKA MI NGUIDJOÏ Thierry Emmanuel
% Version  : 1.0 -- octobre 2026
% ============================================================================
% Chapitre-cas. Chaque cas suit le plan : brief initial, investigation
% guidée, révélations progressives, synthèse et rapport. Le chapitre se
% termine par une grille d'erreurs et une simulation d'examen.
% Laboratoire associé : GCI-LAB-10 (diagnostic et feuille de route chiffrée).
% Tous les cas sont fictifs (noms, dates, montants) ; les mécanismes sont
% exacts.
% ============================================================================

\chapter{Études de Cas Intégrateurs}
\label{chap:cas-integrateurs}
\chaptermark{Études de cas intégrateurs}

\epigraph{%
    « On ne comprend vraiment un référentiel que le jour où il faut
    l'appliquer à une organisation qui ne lui ressemble pas. »%
}{Principe de l'apprentissage par les cas}

\niveauChapitre{Expert}{%
    Maîtrise des Chapitres~\ref{chap:cobit-capacite}
    (capacité et feuille de route), \ref{chap:audit} (audit et constats) et
    \ref{chap:processus} (processus, RACI, indicateurs).%
}

\begin{boxobjectifs}
\begin{itemize}
    \item Mobiliser l'ensemble du cours sur les quatre cas fil rouge.
    \item Produire un diagnostic de gouvernance complet.
    \item Bâtir une feuille de route COBIT--ITIL chiffrée.
    \item Présenter et défendre ses recommandations.
    \item Se préparer à l'examen (grille d'erreurs, simulation).
\end{itemize}
\end{boxobjectifs}

\bigskip

% ============================================================================
\section*{Mode d'emploi du chapitre}
% ============================================================================

Ce chapitre ne présente aucune notion nouvelle. Il reprend les quatre
organisations du cours, un an après les événements qui les ont fait
connaître, et pose à chaque fois la même question~: \emph{la gouvernance
des SI a-t-elle réellement progressé, et que faut-il faire maintenant~?}

Chaque cas suit quatre temps. Le \textbf{brief} donne la situation et la
demande. L'\textbf{investigation guidée} propose les questions à se poser,
avec le chapitre qui permet d'y répondre. Les \textbf{révélations} apportent,
par étapes, des faits nouveaux qui obligent à réviser le premier
jugement. La \textbf{synthèse} donne une réponse de référence, à comparer
avec la vôtre. Travaillez chaque cas \emph{avant} de lire sa synthèse.

La section~\ref{sec:c17-methode} fixe la méthode commune. Les
sections~\ref{sec:c17-sanaga} à~\ref{sec:c17-noun} traitent les quatre
cas. Les sections~\ref{sec:c17-erreurs} et~\ref{sec:c17-examen}
proposent une grille d'erreurs et une simulation d'examen. Le laboratoire
GCI-LAB-10 de NDAP SUKULU reprend la démarche sur les quatre dossiers.

\separateur

% ============================================================================
\section{La Méthode du Diagnostic Intégré}
\label{sec:c17-methode}
% ============================================================================

\subsection{Six volets}

Un diagnostic de gouvernance des SI couvre six volets. Chacun mobilise une
partie du cours.

\begin{tableaucours}{L{3.0cm}L{5.2cm}L{4.0cm}}{%
    \tableauHeader{Volet} &
    \tableauHeader{Questions} &
    \tableauHeader{Chapitres}
}
1. Organes et décisions & Qui décide de quoi~? Les organes fonctionnent-ils~?
& \ref{chap:gouv-entreprise}, \ref{chap:alignement} \\
2. Valeur et portefeuille & Les investissements sont-ils arbitrés et leurs
bénéfices suivis~? & \ref{chap:alignement} \\
3. Système COBIT & Objectifs prioritaires, capacité constatée et cible,
écarts & \ref{chap:cobit-objectifs} à~\ref{chap:cobit-capacite} \\
4. Conformité & Obligations légales et réglementaires, preuves &
\ref{chap:normes}, \ref{chap:droit} \\
5. Services & Pratiques ITIL, niveaux de service, indicateurs &
\ref{chap:itil-svs} à~\ref{chap:processus} \\
6. Fournisseurs et continuité & Contrats, supervision, plans testés &
\ref{chap:itil-niveaux} \\
\end{tableaucours}
\vspace{-0.8em}
\begin{center}{\small\itshape Tableau~17.1. Les six volets du diagnostic
intégré}\end{center}

\subsection{Le livrable}

\begin{boxprocedure}
\textbf{Le rapport de diagnostic.} Le rapport remis à l'organe de
gouvernance comporte~:
\begin{itemize}
    \item une synthèse d'une page, avec trois messages et trois décisions
          demandées~;
    \item le périmètre, les sources et la méthode (entretiens, pièces,
          tests)~;
    \item pour chaque volet, les constats en cinq éléments
          (Chapitre~\ref{chap:audit}), chacun avec son niveau de
          confiance~;
    \item le tableau des écarts de capacité (Chapitre~\ref{chap:cobit-capacite})~;
    \item la feuille de route chiffrée, avec ses propriétaires et ses
          échéances~;
    \item les risques de la feuille de route elle-même.
\end{itemize}
\end{boxprocedure}

% ============================================================================
\section{Cas 1 : la BANQUE SANAGA, un An Après}
\label{sec:c17-sanaga}
% ============================================================================

\subsection{Brief initial}

\begin{boxscenario}{Opération SANAGA. Mars 2027}
Un an après la panne, la COBAC annonce une mission de contrôle sur place
consacrée au risque informatique. Le président du conseil
d'administration vous demande, comme consultant indépendant, un
diagnostic préalable~: « Avons-nous fait ce que nous avions annoncé à la
COBAC, et sommes-nous prêts~? »

\medskip
\textbf{Pièces remises.} Le rapport circonstancié à la COBAC (avril~2026),
la feuille de route à dix-huit mois (Chapitre~\ref{chap:cobit-capacite}),
les procès-verbaux du comité stratégique des SI, le registre des
changements de l'année, le contrat renégocié avec DIGICORE AFRIQUE, le
compte rendu de l'agence pilote de la nouvelle bascule.
\end{boxscenario}

\subsection{Investigation guidée}

\begin{enumerate}
    \item Le comité stratégique des SI fonctionne-t-il selon les critères du
          Chapitre~\ref{chap:gouv-entreprise}~? Reçoit-il une information
          indépendante de celle du directeur général~?
    \item Où en est chaque ligne de la feuille de route~? Quelles capacités
          ont été réévaluées, et par qui (Chapitre~\ref{chap:cobit-capacite})~?
    \item Le registre des changements montre-t-il une gestion conforme à la
          grille de risque (Chapitre~\ref{chap:itil-changements})~?
    \item Le contrat de DIGICORE AFRIQUE contient-il les clauses essentielles
          (Chapitre~\ref{chap:itil-niveaux}), et sont-elles appliquées~?
    \item Les obligations légales et réglementaires sont-elles tenues et
          prouvées (Chapitre~\ref{chap:droit})~?
\end{enumerate}

\subsection{Révélations progressives}

\begin{boxscenario}{Opération SANAGA. Première révélation}
Le comité stratégique des SI s'est réuni quatre fois, comme prévu. Mais
ses procès-verbaux montrent qu'il n'a examiné que des présentations du
directeur général~; la DSI et l'audit interne n'ont jamais été entendus
séparément. Le pouvoir d'accès direct prévu par la charte n'a pas été
exercé.
\end{boxscenario}

\begin{boxscenario}{Opération SANAGA. Deuxième révélation}
Le registre compte 412~changements dans l'année, dont 91 qualifiés
d'urgents (22~\%). La procédure des changements urgents a bien été
rédigée, mais 64 de ces 91~changements n'ont pas fait l'objet de la revue
a posteriori prévue.
\end{boxscenario}

\begin{boxscenario}{Opération SANAGA. Troisième révélation}
Le contrat renégocié avec DIGICORE AFRIQUE contient désormais des
engagements de service, des pénalités et un droit d'audit. Les pénalités
dues pour deux manquements en 2026 (évaluées à 38~millions FCFA) n'ont pas
été réclamées~: la direction des achats « ne voulait pas dégrader la
relation pendant la nouvelle bascule ». Le droit d'audit n'a jamais été
exercé.
\end{boxscenario}

\subsection{Synthèse et rapport}

\begin{tableaucours}{L{3.2cm}L{5.2cm}cL{2.4cm}}{%
    \tableauHeader{Volet} &
    \tableauHeader{Constat principal} &
    \tableauHeader{Conf.} &
    \tableauHeader{Objectif}
}
Organes & Le comité existe, mais reproduit l'asymétrie d'information
qu'il devait réduire & Élevée & EDM01 \\
COBIT & Gains rapides réalisés (DSS04, MEA02)~; chantiers de fond en retard
(EDM03, APO12) & Moyenne & EDM03, APO12 \\
Changements & Taux d'urgence élevé, revues a posteriori non faites (70~\%)
& Élevée & BAI06 \\
Fournisseurs & Clauses présentes mais non appliquées & Élevée & APO10 \\
Conformité & Audit de sécurité légal réalisé en 2026 & Élevée & MEA03 \\
\end{tableaucours}
\vspace{-0.8em}
\begin{center}{\small\itshape Tableau~17.2. Synthèse du diagnostic de la
BANQUE SANAGA}\end{center}

\textbf{Message principal au conseil.} La banque a mis en place les
\emph{structures} et les \emph{documents} annoncés, mais pas encore les
\emph{comportements}~: le comité n'écoute qu'une voix, les contrôles sont
contournés quand ils dérangent, les clauses ne sont pas appliquées. C'est
exactement ce que la grille de De Haes et Van Grembergen permet de
diagnostiquer (Chapitre~\ref{chap:gouv-entreprise})~: des structures sans
mécanismes relationnels. La COBAC le verra.

\begin{boxscenario}{Opération SANAGA. Feuille de route}
\begin{tableaucours}{L{5.6cm}L{2.4cm}rL{1.6cm}}{%
    \tableauHeader{Action} &
    \tableauHeader{Propriétaire} &
    \tableauHeader{Coût (M FCFA)} &
    \tableauHeader{Échéance}
}
Audition séparée de la DSI et de l'audit interne à chaque comité &
Président du comité & 0 & Immédiat \\
Revue a posteriori obligatoire des changements urgents, indicateur
trimestriel & Responsable des changements & 15 & 3 mois \\
Réclamation des pénalités 2026 et premier audit de DIGICORE AFRIQUE &
Direction des achats, RSSI & 25 & 3 mois \\
Appétence au risque approuvée par le conseil, registre des risques SI relié
au registre de la banque & Directeur des risques & 60 & 6 mois \\
Réévaluation indépendante de la capacité des dix objectifs & Audit interne
et cabinet externe & 45 & 9 mois \\
Formation des administrateurs à la gouvernance des SI & Secrétaire du
conseil & 12 & 6 mois \\
\textbf{Total} & & \textbf{157} & \\
\end{tableaucours}
\medskip
Le coût est à rapporter aux 1,9~milliard FCFA de la panne de 2026~: moins
de 9~\%.
\end{boxscenario}

% ============================================================================
\section{Cas 2 : KRIBILOG, la Certification en Jeu}
\label{sec:c17-kribi}
% ============================================================================

\subsection{Brief initial}

\begin{boxscenario}{Opération KRIBI. Mois 12 sur 18}
KRIBILOG SA doit obtenir la certification ISO/IEC~20000-1 dans six mois
pour conserver le contrat de l'armateur (Chapitre~\ref{chap:normes}).
M.~MBARGA vous demande un avis~: « Pouvons-nous solliciter l'organisme de
certification comme prévu au mois~16~? »

\medskip
\textbf{Pièces remises.} Le plan de certification, l'avenant au contrat de
NETPORT SERVICES, les indicateurs mensuels des six derniers mois, le
registre des problèmes, le rapport d'audit interne du mois~11.
\end{boxscenario}

\subsection{Investigation guidée}

\begin{enumerate}
    \item Les pratiques des articles~8.5 et~8.6 de la norme
          fonctionnent-elles depuis assez longtemps pour produire des
          enregistrements probants (Chapitre~\ref{chap:normes})~?
    \item Les indicateurs ont-ils progressé, et sont-ils mesurés par
          KRIBILOG elle-même (Chapitres~\ref{chap:itil-incidents}
          et~\ref{chap:itil-niveaux})~?
    \item Le correctif de l'éditeur a-t-il été installé, et le problème
          clos (Chapitre~\ref{chap:itil-changements})~?
    \item L'avenant donne-t-il à KRIBILOG la maîtrise d'un service exploité
          par un tiers~?
\end{enumerate}

\subsection{Révélations progressives}

\begin{boxscenario}{Opération KRIBI. Première révélation}
Les progrès sont réels. Le taux de réouverture est passé de 30~\% à 12~\%,
le volume mensuel de tickets de 4~800 à 3~100, grâce au libre-service et à
l'installation du correctif. La disponibilité de bout en bout du service
de suivi atteint 99,3~\% en moyenne sur six mois, avec un second
opérateur mobile depuis le mois~8.
\end{boxscenario}

\begin{boxscenario}{Opération KRIBI. Deuxième révélation}
L'audit interne du mois~11 relève que NETPORT SERVICES a confié, sans en
informer KRIBILOG, l'hébergement des sauvegardes à un sous-traitant situé
à Casablanca. L'avenant ne traite ni de la sous-traitance en cascade, ni de
la localisation des données.
\end{boxscenario}

\begin{boxscenario}{Opération KRIBI. Troisième révélation}
La revue de direction prévue au mois~15 par la norme n'est pas encore
programmée~; le directeur général « n'en voit pas l'intérêt, puisque les
indicateurs sont bons ».
\end{boxscenario}

\subsection{Synthèse et rapport}

\textbf{Avis.} Le calendrier peut être tenu, \emph{sous deux conditions}. La
première tient au fournisseur~: un nouvel avenant doit traiter la
sous-traitance et la localisation des données avant l'audit, faute de
quoi l'organisme de certification relèvera une non-conformité majeure sur
la maîtrise des processus confiés à des tiers, et KRIBILOG s'exposera au
titre de la loi n°~2024/017 (Chapitre~\ref{chap:droit}). La seconde tient à
la direction~: la revue de direction est une exigence de la norme
(article~9 de la structure harmonisée), et non une formalité~; sans elle,
pas de certification.

\begin{boxmetier}{Responsable des services IT}
\textbf{Présenter l'avis au directeur général}
\begin{enumerate}
    \item Commencer par les progrès, chiffrés~: ils sont réels et
          mesurés.
    \item Présenter les deux conditions comme des risques pour le contrat
          de l'armateur, en termes d'affaires et non de norme.
    \item Proposer la date de la revue de direction et son ordre du jour.
    \item Demander une décision écrite sur l'avenant avec NETPORT
          SERVICES.
\end{enumerate}
\end{boxmetier}

% ============================================================================
\section{Cas 3 : l'ANPP face à la Chambre des Comptes}
\label{sec:c17-mefou}
% ============================================================================

\subsection{Brief initial}

\begin{boxscenario}{Opération MEFOU. Le rapport provisoire}
La Chambre des comptes a transmis à l'ANPP son rapport provisoire d'audit
des SI (Chapitre~\ref{chap:audit}). Le directeur général dispose de deux
mois pour répondre. Il vous demande de préparer la réponse et le plan
d'action.

\medskip
\textbf{Constats du rapport provisoire.} Absence de comité stratégique des
SI jusqu'en 2026~; projets lancés sans dossier d'investissement~; trois
projets ayant consommé 71~\% du budget du schéma directeur sans être
achevés~; 14~comptes d'agents partis actifs sur le système de paiement~;
séparation des tâches défaillante sur 6~comptes sur 25.
\end{boxscenario}

\subsection{Investigation guidée}

\begin{enumerate}
    \item Quels constats sont fondés, et lesquels peuvent être nuancés par
          des pièces que l'auditeur n'a pas vues (Chapitre~\ref{chap:audit})~?
    \item Comment réarbitrer le portefeuille restant (Chapitre~\ref{chap:alignement})~?
    \item Comment rattacher les projets aux programmes budgétaires
          (loi n°~2018/012)~?
    \item Quelles actions immédiates sur les accès (Chapitres~\ref{chap:itil-incidents}
          et~\ref{chap:audit})~?
\end{enumerate}

\subsection{Révélations progressives}

\begin{boxscenario}{Opération MEFOU. Première révélation}
Le projet de paie dématérialisée, présenté comme « inachevé » dans le
rapport, est en réalité en service depuis janvier~2027. Mais aucun suivi
des bénéfices n'a été fait~: personne ne sait si les délais de paiement
ont diminué.
\end{boxscenario}

\begin{boxscenario}{Opération MEFOU. Deuxième révélation}
Les 14~comptes actifs d'agents partis ont été désactivés le lendemain du
passage de l'auditeur. Mais la cause (l'absence de notification des
départs par les ressources humaines) n'a pas été traitée~: trois nouveaux
départs, depuis, n'ont pas été signalés.
\end{boxscenario}

\subsection{Synthèse et rapport}

\textbf{La réponse à la Chambre des comptes} admet les constats fondés,
apporte les pièces qui nuancent le constat sur le projet de paie, et
présente un plan d'action qui traite les \emph{causes}~:
\begin{itemize}
    \item un comité stratégique des SI, créé par décision du conseil, avec
          une charte inspirée de l'ISO/IEC~38500 (Chapitre~\ref{chap:normes})~;
    \item un réarbitrage du portefeuille restant par score multicritère
          (Chapitre~\ref{chap:alignement}), et l'arrêt d'un projet~;
    \item un suivi des bénéfices du projet de paie, avec un indicateur
          rattaché au projet de performance annuel du programme~;
    \item une notification automatique des départs et une revue trimestrielle
          des accès (DSS05, DSS06).
\end{itemize}

\begin{boxalert}
\textbf{Violation typique.} Corriger le symptôme le lendemain de l'audit
(désactiver les 14~comptes) et présenter la correction comme une réponse
au constat. L'auditeur vérifiera la cause, et constatera la récidive.

\smallskip
\textbf{Règle~:} une réponse à un rapport d'audit traite la cause, désigne
un responsable et une échéance, et indique la pièce qui prouvera la mise
en œuvre.
\end{boxalert}

% ============================================================================
\section{Cas 4 : NOUN PAY, Croître sans Perdre la Confiance}
\label{sec:c17-noun}
% ============================================================================

\subsection{Brief initial}

\begin{boxscenario}{Opération NOUN. La levée de fonds}
NOUN PAY prépare une levée de fonds de 2~milliards FCFA. Les investisseurs
exigent un diagnostic de gouvernance des SI avant de signer. Au même
moment, l'Autorité de protection des données poursuit son enquête sur le
modèle de notation (Chapitre~\ref{chap:droit}).

\medskip
\textbf{Pièces remises.} La matrice RACI du processus de mise à jour du
modèle (Chapitre~\ref{chap:processus}), le contrat type du fournisseur de
nuage, la réponse provisoire à l'Autorité, les statistiques de refus de
crédit par région.
\end{boxscenario}

\subsection{Investigation guidée}

\begin{enumerate}
    \item Quels facteurs de conception caractérisent NOUN PAY, et quel
          système de gouvernance proportionné en découle
          (Chapitre~\ref{chap:cobit-conception})~?
    \item Le contrat du fournisseur de nuage permet-il de maîtriser la
          localisation des données, la réversibilité et l'audit
          (Chapitre~\ref{chap:itil-niveaux})~?
    \item Qui répond, dans l'entreprise, des effets du modèle sur les
          clients (Chapitre~\ref{chap:processus})~?
\end{enumerate}

\subsection{Révélations progressives}

\begin{boxscenario}{Opération NOUN. Première révélation}
Les statistiques montrent que le taux de refus est de 61~\% pour les
demandeurs d'une région, contre 24~\% en moyenne nationale. Le modèle
utilise, parmi ses variables, l'opérateur mobile du demandeur, très
corrélé à la région.
\end{boxscenario}

\begin{boxscenario}{Opération NOUN. Deuxième révélation}
Le contrat type du fournisseur de nuage ne permet ni de choisir la
localisation des données, ni d'auditer le fournisseur~; il renvoie à ses
propres certifications. Un contrat d'entreprise négocié existe chez le
même fournisseur, pour un surcoût annuel de 18~millions FCFA.
\end{boxscenario}

\subsection{Synthèse et rapport}

\textbf{Message aux investisseurs.} NOUN PAY n'a pas besoin d'un système
COBIT complet, mais de six objectifs bien tenus (Chapitre~\ref{chap:cobit-conception})~:
EDM03 (appétence au risque), APO10 (fournisseur de nuage), APO12
(risques), APO14 (données), BAI06 (changements du modèle) et MEA03
(conformité). Deux décisions sont urgentes~: souscrire le contrat
d'entreprise du fournisseur de nuage, et suspendre l'usage de la variable
« opérateur mobile » en attendant une analyse d'impact du modèle. Le
Chapitre~\ref{chap:donnees-ia} reprendra ce cas sous l'angle de la
gouvernance de l'intelligence artificielle.

% ============================================================================
\section{Grille des Erreurs Transversales}
\label{sec:c17-erreurs}
% ============================================================================

Les quatre cas font réapparaître, sous des formes différentes, les mêmes
erreurs. La grille du Chapitre~\ref{chap:defaillances} s'enrichit de
celles que révèle la mise en œuvre.

\begin{tableaucours}{L{4.4cm}L{2.4cm}L{5.4cm}}{%
    \tableauHeader{Erreur} &
    \tableauHeader{Cas} &
    \tableauHeader{Antidote}
}
Des structures sans comportements & SANAGA & Évaluer le fonctionnement des
organes, pas leur existence (Chapitre~\ref{chap:cobit-capacite}) \\
Des contrôles contournés quand ils gênent & SANAGA & Indicateur de
contournement (taux d'urgence, revues non faites) \\
Des clauses non appliquées & SANAGA & Comité de service, réclamation
systématique des pénalités \\
Une sous-traitance en cascade ignorée & KRIBI & Clause de sous-traitance et
de localisation, droit d'audit \\
Un indicateur bon qui dispense de la revue & KRIBI & La revue de
direction est une exigence, pas une option \\
Un symptôme corrigé, une cause intacte & MEFOU & Traiter la cause, prouver
par une pièce \\
Un projet livré sans suivi des bénéfices & MEFOU & EDM02, indicateur
rattaché au programme \\
Une décision automatisée sans responsable & NOUN & Un A pour l'évaluation
de l'impact sur les clients \\
\end{tableaucours}
\vspace{-0.8em}
\begin{center}{\small\itshape Tableau~17.3. Erreurs transversales de mise en
œuvre}\end{center}

% ============================================================================
\section{Simulation d'Examen}
\label{sec:c17-examen}
% ============================================================================

\begin{boxinfo}
\textbf{Sujet type.} \emph{Durée~: 3~heures. Documents autorisés~: aucun,
sauf le tableau des 40~objectifs COBIT fourni.}

\medskip
\textbf{Cas} \textit{(fictif, construit pour l'examen)}. La SOCIÉTÉ D'EAU
DU NORD, entreprise publique
installée à Garoua, distribue l'eau potable à 180~000~abonnés. Son système
de facturation, exploité par un prestataire, a émis en septembre des
factures erronées à 30~000~abonnés. Le conseil d'administration, qui se
réunit deux fois par an, l'a appris par la presse. La direction
informatique compte huit personnes. Aucun comité ne traite du numérique.

\medskip
\begin{enumerate}
    \item \textbf{(4~points)} Distinguez la cause technique probable et les
          causes de gouvernance, en appliquant les trois questions réflexes.
    \item \textbf{(4~points)} Proposez les organes de gouvernance des SI
          adaptés à une entreprise publique de cette taille, en tenant
          compte de la tutelle.
    \item \textbf{(4~points)} Renseignez les facteurs de conception et
          retenez six objectifs COBIT prioritaires, avec leur capacité
          cible.
    \item \textbf{(4~points)} Rédigez deux constats d'audit en cinq éléments,
          l'un sur la gestion des changements, l'autre sur le contrat du
          prestataire.
    \item \textbf{(4~points)} Présentez une feuille de route à douze mois,
          chiffrée en FCFA, en distinguant gains rapides et chantiers de
          fond.
\end{enumerate}
\end{boxinfo}

\begin{boxalert}
\textbf{Ce que le correcteur attend.} Des constatations séparées des
analyses~; des références précises (objectifs COBIT, pratiques ITIL,
textes de loi vérifiés)~; des recommandations qui traitent des causes~; une
gouvernance \emph{proportionnée} à une organisation de huit informaticiens.
Une copie qui recommande de « déployer COBIT et ITIL » sans choisir
n'obtient pas la moyenne.
\end{boxalert}

% ============================================================================
\separateur
% ============================================================================

\begin{boxresume}
\textbf{Ce qu'il faut retenir de ce chapitre~:}
\begin{itemize}
    \item \textbf{Six volets}~: organes et décisions, valeur, système COBIT,
          conformité, services, fournisseurs et continuité. Un diagnostic
          les couvre tous.
    \item \textbf{Les structures ne suffisent pas}~: un an après, la BANQUE
          SANAGA a ses comités et ses documents, mais pas encore les
          comportements.
    \item \textbf{Les bons indicateurs ne dispensent de rien}~: KRIBILOG
          progresse, mais une sous-traitance cachée et une revue de
          direction négligée menacent sa certification.
    \item \textbf{Traiter les causes}~: l'ANPP corrige des symptômes le
          lendemain de l'audit et voit les écarts revenir.
    \item \textbf{La proportion}~: NOUN PAY n'a besoin que de six objectifs
          bien tenus, et de deux décisions urgentes.
    \item \textbf{Une feuille de route se chiffre}~: à rapporter au coût des
          défaillances qu'elle prévient.
\end{itemize}
\end{boxresume}

\bigskip

\begin{boxexo}[title={\ding{45}\quad Exercice~17.1. Le rapport de la BANQUE SANAGA \niveaubase}]
Rédigez la synthèse d'une page du diagnostic de la BANQUE SANAGA~: trois
messages, trois décisions demandées au conseil, et le coût de la feuille
de route rapporté au coût de la panne.
\end{boxexo}

\begin{boxexo}[title={\ding{45}\quad Exercice~17.2. L'avenant de KRIBILOG \niveaumoyen}]
Rédigez les clauses de l'avenant au contrat de NETPORT SERVICES qui
traitent la sous-traitance en cascade, la localisation des sauvegardes et
le droit d'audit. Pour chacune, indiquez le risque traité et la pièce qui
permettra de vérifier son application.
\end{boxexo}

\begin{boxexo}[title={\ding{45}\quad Exercice~17.3. Traiter le sujet d'examen \niveauexpert}]
Traitez en trois heures le sujet de la section~\ref{sec:c17-examen}, puis
comparez votre copie avec les attentes du correcteur. Pour chaque
question, indiquez les chapitres du cours que vous avez mobilisés.
\end{boxexo}

\bigskip

\begin{boxinfo}
\textbf{Transition.} La Partie~\ref{part:pratique} s'achève. Les quatre cas
ont montré que la gouvernance des SI se joue sur des terrains qui
évoluent vite~: l'informatique en nuage, les méthodes agiles, les données
et l'intelligence artificielle. La Partie~\ref{part:avance} les aborde,
en commençant par le Chapitre~\ref{chap:agile-cloud}~: gouverner à grande
vitesse.
\end{boxinfo}


    % ======================================================================
    % PARTIE V — DÉFIS AVANCÉS ET ÉMERGENTS
    % Standards : ISO/IEC 38507, ISO/IEC 42001.
    % Niveau : expert
    % ======================================================================
    \part{Défis Avancés et Émergents}
    \label{part:avance}

        % ============================================================================
% Fichier  : partie5/P05_C18_agilite_devops_cloud.tex
% Titre    : Agilité, DevOps et Cloud : Gouverner à Grande Vitesse
% Auteur   : MINKA MI NGUIDJOÏ Thierry Emmanuel
% Version  : 1.0 -- octobre 2026
% ============================================================================
% Sources : Manifeste agile (2001) ; Forsgren, Humble et Kim (2018) ; NIST
% SP 800-145 (2011) ; ISO/IEC 27017 et 27018 ; loi n° 2010/012 (texte
% anglais officiel) ; loi n° 2024/017 (présentations publiées). Le tableau
% comparatif africain ne cite que l'intitulé et l'année des lois ; aucun
% article n'est cité.
% ============================================================================

\chapter{Agilité, DevOps et Cloud : Gouverner à Grande Vitesse}
\label{chap:agile-cloud}
\chaptermark{Agilité, DevOps et cloud}

\epigraph{%
    « La vitesse n'est pas l'ennemie du contrôle.\\
    C'est le contrôle manuel qui est l'ennemi de la vitesse. »%
}{Principe des équipes DevOps}

\niveauChapitre{Expert}{%
    Maîtrise des Chapitres~\ref{chap:itil-changements} (gestion des
    changements) et~\ref{chap:itil-niveaux} (niveaux de service et
    fournisseurs).%
}

\begin{boxobjectifs}
\begin{itemize}
    \item Concilier contrôle et agilité (DevOps, gestion automatisée des
          changements).
    \item Gouverner l'informatique en nuage (responsabilité partagée,
          réversibilité, localisation des données).
    \item Évaluer un contrat d'informatique en nuage au regard du droit
          camerounais.
    \item Comparer les approches de gouvernance numérique de plusieurs pays
          africains.
\end{itemize}
\end{boxobjectifs}

\bigskip

% ============================================================================
\section*{Question directrice : peut-on gouverner ce qui change chaque jour~?}
% ============================================================================

La nouvelle équipe mobile de la BANQUE SANAGA déploie une version de
l'application chaque semaine, et voudrait passer à un déploiement par
jour. Le comité des changements, qui se réunit le mardi, ne peut pas
suivre. À Bafoussam, NOUN PAY n'a pas un seul serveur à elle~: tout son
système d'information est loué, à la minute, à un fournisseur de nuage
étranger.

Les Parties~\ref{part:cobit} à~\ref{part:pratique} ont construit une
gouvernance faite de comités, de procédures et d'autorisations. La
Partie~\ref{part:avance}, qui s'ouvre ici, la confronte aux pratiques
qui en bousculent les rythmes et les frontières. Ce chapitre traite de
l'agilité et de DevOps (sections~\ref{sec:c18-agile}
et~\ref{sec:c18-devops}), de l'informatique en nuage
(section~\ref{sec:c18-cloud}), du contrat de nuage au regard du droit
camerounais (section~\ref{sec:c18-contrat}), puis des approches
africaines de la gouvernance numérique (section~\ref{sec:c18-afrique}).

\separateur

% ============================================================================
\section{L'Agilité}
\label{sec:c18-agile}
% ============================================================================

\subsection{Le manifeste et ses conséquences}

Le \emph{Manifeste pour le développement agile de logiciels}, publié en
2001 par dix-sept praticiens, pose quatre préférences~: les individus et
leurs interactions plutôt que les processus et les outils, des logiciels
opérationnels plutôt qu'une documentation exhaustive, la collaboration
avec les clients plutôt que la négociation contractuelle, l'adaptation au
changement plutôt que le suivi d'un plan~\cite{agile2001manifeste}. Il
précise que les éléments de droite gardent leur valeur, mais que ceux de
gauche en ont davantage.

\begin{boxalert}
\textbf{Violation typique.} Une équipe invoque l'agilité pour ne rien
documenter, ne rien faire autoriser et ne rendre compte à personne. Lors
d'un incident, personne ne sait quelle version est en production ni qui
l'a validée.

\smallskip
\textbf{Règle~:} l'agilité change la \emph{forme} du contrôle, pas son
existence. Le manifeste ne supprime ni la documentation, ni les contrats,
ni les plans~: il en change le poids relatif.
\end{boxalert}

\subsection{Ce qui change pour la gouvernance}

\begin{tableaucours}{L{3.2cm}L{4.4cm}L{4.6cm}}{%
    \tableauHeader{Sujet} &
    \tableauHeader{Approche traditionnelle} &
    \tableauHeader{Approche agile}
}
Décision d'investissement & Un projet, un budget, une date de fin & Un
produit, une équipe stable, un financement par période \\
Pilotage & Respect du plan, du budget, des délais & Valeur livrée à chaque
itération \\
Exigences & Définies au départ & Priorisées en continu par un responsable
de produit \\
Contrôle & Jalons et validations formelles & Démonstrations régulières,
critères de « fini » \\
\end{tableaucours}
\vspace{-0.8em}
\begin{center}{\small\itshape Tableau~18.1. Gouvernance traditionnelle et
gouvernance agile}\end{center}

La gouvernance ne disparaît pas~: elle se déplace. Le comité stratégique
n'approuve plus un cahier des charges, il fixe les objectifs d'un produit,
l'enveloppe d'une équipe et les indicateurs de valeur
(Chapitre~\ref{chap:alignement}), puis revoit ces décisions à intervalles
réguliers.

% ============================================================================
\section{DevOps}
\label{sec:c18-devops}
% ============================================================================

\subsection{Rapprocher le développement et l'exploitation}

\begin{boxdef}
\termecle{DevOps}\index{DevOps}~: ensemble de pratiques culturelles et
techniques qui rapproche les équipes de développement et d'exploitation,
pour livrer des changements fréquents, fiables et réversibles, grâce à
l'automatisation de la construction, des tests et du déploiement.

\medskip
\termecle{Chaîne de déploiement}\index{chaîne de déploiement}
(\textit{pipeline})~: suite automatisée des étapes par lesquelles un
changement de code passe, de sa validation par un pair à sa mise en
production.
\end{boxdef}

\subsection{Mesurer la performance de la livraison}

Les travaux de Forsgren, Humble et Kim~\cite{forsgren2018accelerate} ont
popularisé quatre indicateurs de la performance de la livraison
logicielle. Ils ont une propriété précieuse pour la gouvernance~: ils
mesurent à la fois la vitesse \emph{et} la stabilité.

\begin{tableaucours}{L{3.8cm}L{4.4cm}L{4.0cm}}{%
    \tableauHeader{Indicateur} &
    \tableauHeader{Mesure} &
    \tableauHeader{Dimension}
}
Fréquence de déploiement & Nombre de mises en production par période &
Vitesse \\
Délai de livraison des changements & Temps entre la validation du code et
sa mise en production & Vitesse \\
Taux d'échec des changements & Part des déploiements qui provoquent un
incident ou un retour arrière & Stabilité \\
Délai de rétablissement & Temps pour rétablir le service après un échec &
Stabilité \\
\end{tableaucours}
\vspace{-0.8em}
\begin{center}{\small\itshape Tableau~18.2. Les quatre indicateurs de la
performance de livraison}\end{center}

Le taux d'échec des changements et le délai de rétablissement sont les
indicateurs de BAI06 et de DSS02 (Chapitres~\ref{chap:itil-changements}
et~\ref{chap:itil-incidents})~: DevOps ne remplace pas la gestion des
changements, il en mesure le résultat.

\subsection{Déplacer le contrôle dans la chaîne}

\begin{tableaucours}{L{4.0cm}L{8.2cm}}{%
    \tableauHeader{Contrôle traditionnel} &
    \tableauHeader{Équivalent dans une chaîne DevOps}
}
Autorisation par le comité des changements & Changement standard
pré-autorisé, si la chaîne valide toutes les étapes~; comité réservé aux
changements hors norme \\
Séparation des tâches entre développeur et exploitant & Revue obligatoire
du code par un pair~; le développeur ne peut pas fusionner seul son code \\
Tests de recette manuels & Tests automatisés obligatoires, avec seuil de
couverture bloquant \\
Plan de retour arrière & Déploiement progressif et retour arrière
automatique en cas d'erreur \\
Procès-verbal d'autorisation & Journal inaltérable de la chaîne (qui, quoi,
quand, avec quels résultats de tests) \\
\end{tableaucours}
\vspace{-0.8em}
\begin{center}{\small\itshape Tableau~18.3. Les contrôles de la gestion des
changements dans une chaîne DevOps}\end{center}

\begin{boxscenario}{Opération SANAGA. L'équipe mobile}
\textbf{Demande.} L'équipe mobile veut déployer chaque jour, sans passer
par le comité des changements.

\medskip
\textbf{Décision du comité stratégique.} Les déploiements de l'application
mobile sont reconnus comme \emph{changements standard} à quatre conditions,
vérifiées par la chaîne~: revue du code par un pair, tests automatisés
couvrant au moins 90~\% des fonctions de paiement, déploiement progressif
(5~\% des clients pendant une heure), retour arrière automatique. Tout
changement touchant l'interface avec le système central reste un
changement normal, soumis au comité.

\medskip
\textbf{Contrôle.} L'auditrice interne examine chaque trimestre un
échantillon de déploiements dans le journal de la chaîne~: elle teste
désormais un contrôle \emph{automatisé}, et non plus une signature.
\end{boxscenario}

\begin{boxmetier}{Auditeur SI}
\textbf{Auditer une chaîne DevOps}
\begin{enumerate}
    \item Obtenir la description de la chaîne et des contrôles qu'elle
          impose, puis vérifier qu'ils sont bloquants (un test échoué
          arrête le déploiement).
    \item Vérifier qui peut modifier la chaîne elle-même~: c'est le
          nouveau point de contrôle critique.
    \item Rapprocher les déploiements en production du journal de la
          chaîne~: un déploiement hors chaîne est un changement non
          autorisé.
    \item Suivre les quatre indicateurs de performance de livraison, et
          notamment le taux d'échec.
\end{enumerate}
\end{boxmetier}

% ============================================================================
\section{L'Informatique en Nuage}
\label{sec:c18-cloud}
% ============================================================================

\subsection{Définitions}

La définition de référence est celle du NIST~\cite{nist2011cloud}~:
l'informatique en nuage permet un accès à la demande, par le réseau, à un
ensemble partagé de ressources informatiques configurables, rapidement
mises à disposition et libérées avec un effort de gestion minimal. Elle
présente cinq caractéristiques~: libre-service à la demande, large accès
par le réseau, mutualisation des ressources, élasticité rapide, service
mesuré.

\begin{tableaucours}{L{2.8cm}L{4.8cm}L{4.6cm}}{%
    \tableauHeader{Modèle} &
    \tableauHeader{Ce que fournit le prestataire} &
    \tableauHeader{Exemple}
}
IaaS (infrastructure) & Serveurs, stockage, réseau virtuels & L'hébergement
des serveurs de NOUN PAY \\
PaaS (plateforme) & Environnement d'exécution, bases de données gérées &
La base de données gérée du modèle de notation \\
SaaS (logiciel) & Application complète & La messagerie de la BANQUE SANAGA \\
\end{tableaucours}
\vspace{-0.8em}
\begin{center}{\small\itshape Tableau~18.4. Les trois modèles de service du
NIST}\end{center}

\subsection{La responsabilité partagée}

Dans le nuage, la responsabilité de la sécurité et de l'exploitation est
\emph{partagée} entre le client et le fournisseur, selon le modèle de
service. Quel que soit le modèle, \textbf{le client reste responsable de
ses données, de ses utilisateurs et de ses droits d'accès}.

\begin{tableaucours}{L{4.0cm}ccc}{%
    \tableauHeader{Domaine} &
    \tableauHeader{IaaS} &
    \tableauHeader{PaaS} &
    \tableauHeader{SaaS}
}
Données et leur classification & Client & Client & Client \\
Identités et droits d'accès & Client & Client & Client \\
Application & Client & Client & Fournisseur \\
Système d'exploitation, intergiciels & Client & Fournisseur & Fournisseur \\
Infrastructure physique, réseau & Fournisseur & Fournisseur & Fournisseur \\
\end{tableaucours}
\vspace{-0.8em}
\begin{center}{\small\itshape Tableau~18.5. Le modèle de responsabilité
partagée (simplifié)}\end{center}

\begin{boxalert}
\textbf{Violation typique.} « Nos données sont chez un grand fournisseur
mondial, elles sont donc en sécurité. » Le stockage de NOUN PAY était
configuré en accès public par erreur~; le fournisseur, dont la plateforme
fonctionnait parfaitement, n'y était pour rien.

\smallskip
\textbf{Règle~:} la sécurité du nuage est celle du fournisseur~; la
sécurité \emph{dans} le nuage est celle du client.
\end{boxalert}

\subsection{Gouverner le nuage}

Le nuage déplace quatre risques que la gouvernance doit arbitrer
explicitement~:
\begin{itemize}
    \item \textbf{la dépendance}~: changer de fournisseur peut être long et
          coûteux, surtout en PaaS~;
    \item \textbf{la localisation}~: les données peuvent être stockées
          hors du territoire, ce que le droit encadre
          (section~\ref{sec:c18-contrat})~;
    \item \textbf{l'informatique fantôme}~: n'importe quel service peut
          souscrire un abonnement avec une carte bancaire, hors de tout
          contrôle~;
    \item \textbf{le coût}~: payé à l'usage, il dérive vite sans pilotage.
\end{itemize}
Deux normes aident à traiter la sécurité~: l'ISO/IEC~27017, qui donne des
mesures de sécurité propres aux services en nuage, et l'ISO/IEC~27018, qui
traite de la protection des données personnelles dans le nuage public.

% ============================================================================
\section{Le Contrat de Nuage et le Droit Camerounais}
\label{sec:c18-contrat}
% ============================================================================

\begin{boxjuris}
\textbf{\loicam{n°~2010/012}, articles~26 et~29} (version anglaise
officielle, extraits)~:

\medskip
\textit{« [Operators of information systems] shall have standardized
systems enabling them to at all times identify, assess, process or manage
any risk relating to the security of the information systems of the
services provided directly or indirectly. »} (art.~26, al.~1)

\medskip
\textit{« Operators of information systems shall be bound to conserve the
connection and traffic data of their information systems for a period of
10 (ten) years. »} (art.~29, al.~1)

\medskip
\textbf{\loicam{n°~2024/017}} (d'après les présentations publiées). La loi
encadre strictement le transfert des données personnelles hors du
territoire national (Chapitre~\ref{chap:droit}).

\medskip
\textbf{Portée pour un contrat de nuage.} Une organisation qui confie son
système d'information à un fournisseur de nuage reste tenue de maîtriser
les risques de sécurité des services fournis « indirectement », de
conserver ses données de connexion pendant dix ans, et de faire examiner
tout transfert de données personnelles hors du Cameroun au regard de la
loi n°~2024/017. Le contrat type d'un fournisseur mondial n'a pas été
écrit pour ces obligations~: il faut en vérifier chaque point.
\end{boxjuris}
% A VERIFIER : conditions exactes des transferts hors du territoire dans le
% texte publié de la loi n° 2024/017 et ses décrets d'application.

\begin{boxprocedure}
\textbf{La grille d'évaluation d'un contrat de nuage.} Avant de signer, ou
lors de l'audit, chaque point est noté conforme, partiellement conforme
ou non conforme, avec la clause correspondante~:
\begin{itemize}
    \item localisation des données et des sauvegardes, et liste des pays~;
    \item conditions de transfert hors du Cameroun~;
    \item conservation et accès aux journaux pendant la durée légale~;
    \item droit d'audit, ou rapports d'audit indépendants et leur
          périmètre~;
    \item notification des incidents de sécurité, dans un délai fixé~;
    \item engagements de service et pénalités~;
    \item sous-traitants du fournisseur et information du client~;
    \item réversibilité~: formats, délais, assistance, coûts~;
    \item droit applicable et juridiction compétente.
\end{itemize}
\end{boxprocedure}

\begin{boxscenario}{Opération NOUN. Deux contrats comparés}
\begin{tableaucours}{L{4.2cm}L{3.8cm}L{3.8cm}}{%
    \tableauHeader{Point de la grille} &
    \tableauHeader{Contrat type} &
    \tableauHeader{Contrat d'entreprise}
}
Localisation des données & Non conforme (non choisie) & Conforme (région
choisie, liste des pays) \\
Droit d'audit & Non conforme (renvoi aux certifications) & Partiel
(rapports d'audit indépendants) \\
Notification des incidents & Partiel (« dans les meilleurs délais ») &
Conforme (72~heures) \\
Réversibilité & Non conforme & Partiel (exports, sans assistance) \\
Journaux & Partiel (90~jours) & Partiel (configurable, coût de stockage) \\
\end{tableaucours}

\medskip
\textbf{Analyse} (confiance moyenne). Le contrat d'entreprise, pour un
surcoût de 18~millions FCFA par an (Chapitre~\ref{chap:cas-integrateurs}),
corrige les points les plus graves. Il ne règle pas tout~: la conservation
des journaux pendant dix ans suppose un archivage complémentaire, et la
réversibilité doit être préparée par NOUN PAY elle-même (formats ouverts,
exercice de restitution annuel).
\end{boxscenario}

\begin{boxmetier}{RSSI}
\textbf{Gouverner les abonnements en nuage}
\begin{enumerate}
    \item Tenir l'inventaire des services en nuage utilisés, y compris ceux
          payés par carte bancaire par les métiers.
    \item Classer les données selon leur sensibilité, et interdire les
          plus sensibles sur les services non évalués.
    \item Vérifier la configuration des services (accès publics, chiffrement,
          comptes à privilèges) au moins chaque mois.
    \item Présenter chaque année au comité des risques la carte des
          dépendances aux fournisseurs de nuage.
\end{enumerate}
\end{boxmetier}

% ============================================================================
\section{La Gouvernance Numérique en Afrique}
\label{sec:c18-afrique}
% ============================================================================

\subsection{Un cadre continental}

L'Union africaine a adopté en 2020 une \emph{Stratégie de transformation
numérique pour l'Afrique 2020--2030}, puis en 2022 un \emph{Cadre
stratégique sur les données}, qui encourage l'harmonisation des règles de
circulation des données sur le continent. La Convention de Malabo, entrée
en vigueur en 2023, en constitue le volet juridique
(Chapitre~\ref{chap:droit}). L'initiative Smart Africa, lancée en 2013,
réunit des États autour de projets numériques communs.

\subsection{Des lois nationales de protection des données}

\begin{tableaucours}{L{3.0cm}L{6.2cm}L{3.0cm}}{%
    \tableauHeader{Pays} &
    \tableauHeader{Texte} &
    \tableauHeader{Année}
}
Sénégal & Loi sur la protection des données à caractère personnel & 2008 \\
Maroc & Loi n°~09-08 relative à la protection des personnes physiques à
l'égard du traitement des données personnelles & 2009 \\
Ghana & Data Protection Act & 2012 \\
Côte d'Ivoire & Loi relative à la protection des données à caractère
personnel & 2013 \\
Kenya & Data Protection Act & 2019 \\
Rwanda & Loi sur la protection des données personnelles et de la vie
privée & 2021 \\
Nigeria & Nigeria Data Protection Act & 2023 \\
Cameroun & Loi n°~2024/017 & 2024 \\
\end{tableaucours}
\vspace{-0.8em}
\begin{center}{\small\itshape Tableau~18.6. Quelques lois africaines de
protection des données (intitulés abrégés)}\end{center}

\subsection{Trois enseignements}

\begin{enumerate}
    \item \textbf{La localisation des données est un choix politique.}
          Certains États privilégient la libre circulation pour attirer les
          services, d'autres la conservation sur le territoire pour des
          raisons de souveraineté. Une organisation présente dans plusieurs
          pays doit en tenir un registre (Chapitre~\ref{chap:droit}).
    \item \textbf{Les infrastructures nationales se développent.} Plusieurs
          États ont investi dans des centres de données nationaux, qui
          offrent une alternative au nuage étranger~; leur niveau de
          service et de sécurité doit être évalué comme celui de tout
          fournisseur.
    \item \textbf{Le droit précède souvent les autorités.} Une loi n'a
          d'effet que si l'autorité qu'elle crée dispose de moyens.
          L'efficacité réelle d'un régime se mesure à ses décisions
          publiées, encore rares (Chapitre~\ref{chap:defaillances}).
\end{enumerate}

\begin{boiteRemarque}{Souveraineté numérique ou isolement~?}
Imposer que toutes les données restent au Cameroun protège contre
certains risques (accès par des autorités étrangères, coupure de
connectivité internationale), mais en crée d'autres~: dépendance à un
nombre réduit de fournisseurs locaux, coûts plus élevés, moindre
résilience si un seul centre de données héberge tout. La bonne réponse
de gouvernance n'est ni l'interdiction ni l'indifférence, mais la
\emph{classification}~: localiser ce qui doit l'être, et choisir
librement pour le reste, contrat par contrat.
\end{boiteRemarque}

% ============================================================================
\separateur
% ============================================================================

\begin{boxresume}
\textbf{Ce qu'il faut retenir de ce chapitre~:}
\begin{itemize}
    \item \textbf{L'agilité change la forme du contrôle, pas son
          existence}~: la gouvernance finance des produits et des équipes,
          et revoit ses décisions à intervalles réguliers.
    \item \textbf{DevOps déplace le contrôle dans la chaîne}~: revue par un
          pair, tests bloquants, déploiement progressif, retour arrière
          automatique, journal inaltérable. Le point critique devient la
          chaîne elle-même.
    \item \textbf{Quatre indicateurs} mesurent ensemble la vitesse et la
          stabilité de la livraison.
    \item \textbf{Dans le nuage, la responsabilité est partagée}~: le client
          reste responsable de ses données, de ses utilisateurs et de ses
          accès.
    \item \textbf{Le contrat de nuage s'évalue au regard du droit
          camerounais}~: risques des services fournis indirectement,
          conservation des journaux pendant dix ans, transferts hors du
          territoire.
    \item \textbf{L'Afrique se dote de règles}~: stratégie continentale,
          lois nationales de protection des données, centres de données
          nationaux. La réponse de gouvernance est la classification des
          données.
\end{itemize}
\end{boxresume}

\bigskip

\begin{boxexo}[title={\ding{45}\quad Exercice~18.1. Qui est responsable~? \niveaubase}]
Pour chacune des situations suivantes, indiquez si la responsabilité
incombe au client ou au fournisseur de nuage, selon le Tableau~18.5.
\begin{enumerate}[(a)]
    \item Un ancien salarié de NOUN PAY se connecte encore à la console
          d'administration.
    \item Une panne électrique arrête le centre de données du fournisseur.
    \item La messagerie en SaaS de la BANQUE SANAGA présente une faille
          logicielle.
    \item Une base de données gérée (PaaS) est accessible depuis Internet
          sans mot de passe.
\end{enumerate}
\end{boxexo}

\begin{boxexo}[title={\ding{45}\quad Exercice~18.2. Calculer les indicateurs de livraison \niveaumoyen}]
En un mois, l'équipe mobile de la BANQUE SANAGA a réalisé 22~déploiements.
Trois ont provoqué un incident~; les délais de rétablissement ont été de
20~minutes, 45~minutes et 3~heures. Le délai moyen entre la validation du
code et la mise en production est de 2~jours.
\begin{enumerate}[(a)]
    \item Calculez le taux d'échec des changements et le délai moyen de
          rétablissement.
    \item L'incident de 3~heures a touché le paiement. Qu'en déduisez-vous
          sur le périmètre des changements standard~?
    \item Proposez une cible pour chacun des quatre indicateurs.
\end{enumerate}
\end{boxexo}

\begin{boxexo}[title={\ding{45}\quad Exercice~18.3. Migrer l'ANPP vers le nuage~? \niveauexpert}]
Un fournisseur propose à l'ANPP d'héberger son système de paiement dans
son nuage, hors du Cameroun, pour un coût annuel inférieur de 30~\% à celui
de son centre de données actuel.
\begin{enumerate}[(a)]
    \item Classez les données de l'agence selon leur sensibilité.
    \item Appliquez la grille d'évaluation du contrat de nuage à l'offre,
          en indiquant les informations qui vous manquent.
    \item Identifiez les obligations légales en jeu, en citant uniquement
          les articles vérifiés, et signalez ce qui reste à vérifier.
    \item Rédigez la recommandation au conseil d'administration, en
          présentant au moins deux options.
\end{enumerate}
\end{boxexo}

\bigskip

\begin{boxinfo}
\textbf{Transition.} Le nuage concentre les données~; l'intelligence
artificielle les exploite pour décider. Le Chapitre~\ref{chap:donnees-ia}
traite de la gouvernance des données et de l'intelligence artificielle, à
partir du cas de la notation de crédit de NOUN PAY.
\end{boxinfo}

        % ============================================================================
% Fichier  : partie5/P05_C19_gouvernance_donnees_ia.tex
% Titre    : Gouvernance des Données et de l'Intelligence Artificielle
% Auteur   : MINKA MI NGUIDJOÏ Thierry Emmanuel
% Version  : 1.0 -- octobre 2026
% ============================================================================
% Sources : COBIT 2019, fiche APO14 (dossier Normes) ; ISO/IEC 38507:2022,
% ISO/IEC 42001:2023, ISO/IEC 23894:2023 (présentation d'après leur objet
% public) ; NIST AI RMF 1.0 (2023) ; Recommandation de l'UNESCO sur
% l'éthique de l'IA (2021) ; règlement (UE) 2024/1689. Loi n° 2024/017 :
% aucun article cité ; le régime des décisions automatisées est à vérifier
% dans le texte publié.
% ============================================================================

\chapter{Gouvernance des Données et de l'Intelligence Artificielle}
\label{chap:donnees-ia}
\chaptermark{Données et intelligence artificielle}

\epigraph{%
    « Un algorithme ne décide jamais seul.\\
    Quelqu'un a décidé de le laisser décider. »%
}{Principe de la gouvernance de l'IA}

\niveauChapitre{Expert}{%
    Maîtrise des Chapitres~\ref{chap:droit} (protection des données
    personnelles) et~\ref{chap:agile-cloud} (informatique en nuage).%
}

\begin{boxobjectifs}
\begin{itemize}
    \item Définir la gouvernance des données (propriété, qualité, cycle de
          vie).
    \item Présenter l'ISO/IEC~38507 et l'ISO/IEC~42001.
    \item Évaluer les risques d'un système d'IA (biais, explicabilité,
          responsabilité).
    \item Appliquer la loi de 2024 sur les données personnelles à un cas de
          notation par IA.
\end{itemize}
\end{boxobjectifs}

\bigskip

% ============================================================================
\section*{Question directrice : qui répond d'une décision prise par un modèle~?}
% ============================================================================

Le Chapitre~\ref{chap:cas-integrateurs} l'a établi~: le modèle de notation
de crédit de NOUN PAY refuse 61~\% des demandes d'une région, contre
24~\% en moyenne nationale. Il utilise l'opérateur mobile du demandeur,
très corrélé à la région. Personne, dans l'entreprise, n'était chargé
d'évaluer les effets du modèle sur les clients (Chapitre~\ref{chap:processus}).
L'Autorité de protection des données demande des comptes.

Ce cas réunit les deux sujets de ce chapitre. La \textbf{gouvernance des
données} (sections~\ref{sec:c19-donnees} et~\ref{sec:c19-qualite}) traite
de ce que l'organisation possède~: qui en répond, avec quelle qualité, pour
combien de temps. La \textbf{gouvernance de l'intelligence artificielle}
(sections~\ref{sec:c19-ia} à~\ref{sec:c19-noun}) traite de ce que
l'organisation en fait quand elle confie des décisions à des modèles.

\separateur

% ============================================================================
\section{La Gouvernance des Données}
\label{sec:c19-donnees}
% ============================================================================

\subsection{Les données, un actif}

\begin{boxdef}
\termecle{Gouvernance des données}\index{gouvernance des données}~:
exercice de l'autorité et du contrôle sur la gestion des données~:
qui décide de leur définition, de leur usage, de leur qualité, de leur
protection et de leur conservation, selon quelles règles, et qui en rend
compte.
\end{boxdef}

COBIT~2019 consacre aux données un objectif de management à part entière,
APO14, qui couvre tout le cycle de vie des données, de leur création à
leur archivage. Ses liens principaux sont la qualité de l'information
financière (EG04) et de l'information de gestion (EG07, AG10)~: les
données sont d'abord ce sur quoi l'on décide.

\begin{boxobjet}{APO14 · Gérer les données}
\textbf{Description~:} gérer efficacement les actifs de données de
l'organisation sur tout leur cycle de vie, de la création à la diffusion,
à la maintenance et à l'archivage.

\textbf{Finalité~:} assurer l'utilisation efficace des données critiques
pour atteindre les objectifs de l'organisation.

\textbf{Objectifs liés~:} EG04, EG07~; AG10.

\textbf{Pratiques~:} APO14.01 définir la stratégie de gestion des données
et les rôles~; APO14.02 tenir un glossaire métier~; APO14.03 gérer les
métadonnées~; APO14.04 définir une stratégie de qualité des données~;
APO14.05 profiler les données~; APO14.06 évaluer leur qualité~; APO14.07
définir leur nettoyage~; APO14.08 gérer le cycle de vie des actifs de
données~; APO14.09 organiser l'archivage et la conservation~; APO14.10
gérer les sauvegardes et les restaurations.

\textbf{Indicateurs~:} part des données critiques dotées d'un
propriétaire, taux d'anomalies de qualité, part des données conservées
au-delà de leur durée prévue.
\end{boxobjet}
{\small\itshape Description, finalité, objectifs liés et pratiques résumés
et traduits par nos soins à partir de COBIT~2019~\cite{isaca2018cobitobjectifs}.}

\subsection{Les rôles}

\begin{tableaucours}{L{3.4cm}L{5.0cm}L{3.8cm}}{%
    \tableauHeader{Rôle} &
    \tableauHeader{Responsabilité} &
    \tableauHeader{Chez NOUN PAY}
}
Propriétaire des données & Répond d'un domaine de données~: usage
autorisé, niveau de qualité, accès & Directeur commercial pour les données
clients \\
Intendant des données (\textit{data steward}) & Veille au quotidien à la
définition et à la qualité & Analyste du crédit \\
Dépositaire & Assure le stockage, la sécurité et la sauvegarde &
Directeur technique et fournisseur de nuage \\
Délégué à la protection des données & Veille au respect du droit des
données personnelles & M.~DJOUMESSI \\
\end{tableaucours}
\vspace{-0.8em}
\begin{center}{\small\itshape Tableau~19.1. Les rôles de la gouvernance des
données}\end{center}

La règle cardinale~: \textbf{le propriétaire des données est un métier,
pas l'informatique}. La DSI conserve et protège~; c'est le métier qui
décide de ce que les données signifient et de l'usage qui en est fait.

% ============================================================================
\section{Qualité et Cycle de Vie}
\label{sec:c19-qualite}
% ============================================================================

\subsection{Les dimensions de la qualité}

\begin{tableaucours}{L{3.0cm}L{5.0cm}L{4.2cm}}{%
    \tableauHeader{Dimension} &
    \tableauHeader{Question} &
    \tableauHeader{Exemple}
}
Exactitude & La donnée reflète-t-elle la réalité~? & Le solde converti
est-il le bon (SANAGA)~? \\
Complétude & Toutes les données requises sont-elles présentes~? &
Revenus manquants dans 18~\% des dossiers de crédit \\
Cohérence & La même donnée a-t-elle la même valeur partout~? & L'adresse
du client dans deux systèmes \\
Actualité & La donnée est-elle à jour~? & Position d'un conteneur vieille
de six heures (KRIBILOG) \\
Validité & Respecte-t-elle le format et les règles~? & Numéro de
téléphone au bon format \\
Unicité & Chaque entité n'apparaît-elle qu'une fois~? & Un fournisseur
enregistré trois fois à l'ANPP \\
\end{tableaucours}
\vspace{-0.8em}
\begin{center}{\small\itshape Tableau~19.2. Les dimensions de la qualité
des données}\end{center}

\subsection{Le cycle de vie}

Une donnée se crée, s'utilise, se partage, s'archive et se détruit. La
gouvernance fixe, pour chaque catégorie, sa \emph{durée de conservation}~:
assez longue pour répondre aux obligations (dix ans pour les données de
connexion selon l'article~29 de la loi n°~2010/012, Chapitre~\ref{chap:droit}),
pas plus longue que nécessaire pour les données personnelles. Conserver
« au cas où » n'est pas une politique, c'est un risque.

\begin{boxprocedure}
\textbf{Le registre des données critiques.} Pour qu'un auditeur puisse
établir que les données critiques sont gouvernées, l'organisation tient
un registre qui indique, pour chaque domaine de données~:
\begin{itemize}
    \item le propriétaire métier et l'intendant~;
    \item la définition, dans le glossaire métier~;
    \item la classification (publique, interne, confidentielle,
          personnelle, sensible)~;
    \item la localisation (systèmes, pays)~;
    \item les règles de qualité et leur dernier contrôle~;
    \item la durée de conservation et son fondement.
\end{itemize}
\end{boxprocedure}

% ============================================================================
\section{Gouverner l'Intelligence Artificielle}
\label{sec:c19-ia}
% ============================================================================

\subsection{Ce que l'IA change}

Un système d'intelligence artificielle n'est pas un logiciel comme un
autre, pour trois raisons qui intéressent directement la gouvernance~:
\begin{itemize}
    \item \textbf{il apprend des données}~: ses erreurs viennent souvent de
          biais présents dans les données d'entraînement, que personne n'a
          voulus~;
    \item \textbf{il est difficile à expliquer}~: certains modèles
          produisent un résultat sans que l'on puisse dire simplement
          pourquoi~;
    \item \textbf{il évolue}~: un modèle réentraîné est un nouveau modèle,
          et ses effets peuvent dériver avec le temps.
\end{itemize}

\subsection{Les normes et cadres de référence}

\begin{tableaucours}{L{3.4cm}L{2.0cm}L{6.8cm}}{%
    \tableauHeader{Référence} &
    \tableauHeader{Nature} &
    \tableauHeader{Objet}
}
ISO/IEC 38507 (2022) & Lignes directrices & Implications, pour l'organe de
gouvernance, de l'usage de l'IA par l'organisation~; prolonge
l'ISO/IEC~38500 \\
ISO/IEC 42001 (2023) & Norme certifiable & Système de management de l'IA,
selon la structure harmonisée (Chapitre~\ref{chap:normes}) \\
ISO/IEC 23894 (2023) & Lignes directrices & Management des risques liés à
l'IA, dans la logique de l'ISO~31000 \\
NIST AI RMF 1.0 (2023) & Cadre & Quatre fonctions~: gouverner,
cartographier, mesurer, gérer~\cite{nist2023airmf} \\
Recommandation de l'UNESCO (2021) & Instrument normatif international &
Valeurs et principes éthiques de l'IA, adoptés par les États
membres~\cite{unesco2021ethique} \\
\end{tableaucours}
\vspace{-0.8em}
\begin{center}{\small\itshape Tableau~19.3. Normes et cadres de la
gouvernance de l'IA}\end{center}

La Recommandation de l'UNESCO sur l'éthique de l'intelligence
artificielle, adoptée en novembre~2021, est le premier instrument normatif
mondial sur le sujet. Elle insiste sur la proportionnalité, la sécurité,
l'équité et la non-discrimination, la transparence et l'explicabilité, la
responsabilité et la surveillance humaine. Pour une formation dispensée par
une Chaire UNESCO, elle est le point de départ naturel.

\begin{boxinfo}
\textbf{Le contexte africain et international.} L'Union africaine a adopté
en 2024 une stratégie continentale sur l'intelligence artificielle.
L'Union européenne a adopté le règlement (UE)~2024/1689 sur l'intelligence
artificielle, qui classe l'évaluation de la solvabilité des personnes
physiques parmi les systèmes à haut risque~\cite{ue2024aiact}. Ce règlement
ne s'applique pas au Cameroun en tant que tel, mais il s'impose aux
entreprises qui y servent des clients, et sert déjà de référence aux
investisseurs.
\end{boxinfo}

% ============================================================================
\section{Évaluer les Risques d'un Système d'IA}
\label{sec:c19-risques}
% ============================================================================

\subsection{Une grille en six risques}

\begin{tableaucours}{L{3.0cm}L{5.0cm}L{4.2cm}}{%
    \tableauHeader{Risque} &
    \tableauHeader{Question} &
    \tableauHeader{Mesure de maîtrise}
}
Biais et discrimination & Le modèle traite-t-il défavorablement un groupe~?
& Mesure des écarts entre groupes, revue des variables \\
Explicabilité & Peut-on expliquer une décision à la personne concernée~?
& Motifs principaux de chaque décision, documentation du modèle \\
Responsabilité & Qui répond d'une décision du modèle~? & Un A dans la
matrice RACI (Chapitre~\ref{chap:processus}) \\
Robustesse & Le modèle résiste-t-il aux données inhabituelles et aux
manipulations~? & Tests de robustesse, surveillance des entrées \\
Dérive & Ses performances se dégradent-elles avec le temps~? & Indicateurs
suivis en production, réentraînement contrôlé \\
Conformité & Les données sont-elles traitées licitement~? & Analyse au
regard de la loi n°~2024/017 \\
\end{tableaucours}
\vspace{-0.8em}
\begin{center}{\small\itshape Tableau~19.4. Les risques d'un système d'IA}\end{center}

\subsection{Mesurer un biais}

Une mesure simple de l'écart entre groupes est le \termecle{ratio d'impact
disproportionné}\index{ratio d'impact disproportionné}~: le taux de
décisions favorables du groupe le moins favorisé, divisé par celui du
groupe le plus favorisé. Une pratique répandue, issue du droit américain
du travail, considère comme suspect un ratio inférieur à~0,8 (la « règle
des quatre cinquièmes »). Ce seuil n'a pas de valeur juridique au
Cameroun~; il sert de signal d'alerte.

\begin{boxscenario}{Opération NOUN. Le calcul du biais}
\textbf{Données.} Taux de refus de 61~\% dans la région concernée, soit un
taux d'acceptation de 39~\%. Taux de refus moyen national de 24~\%, soit un
taux d'acceptation de 76~\%.

\medskip
\textbf{Calcul.} Ratio d'impact $= 0{,}39 / 0{,}76 \approx 0{,}51$, très en
dessous du seuil d'alerte de 0,8.

\medskip
\textbf{Constatation.} Le ratio est de 0,51~; la variable « opérateur
mobile », corrélée à la région, figure parmi les variables du modèle.

\medskip
\textbf{Analyse} (confiance moyenne). L'écart peut s'expliquer en partie
par des différences réelles de risque, mais l'usage d'une variable
corrélée à la région fait peser un risque sérieux de discrimination
indirecte. Seule une analyse du modèle sans cette variable, à risque égal,
permettra de conclure. En attendant, la variable doit être suspendue
(Chapitre~\ref{chap:cas-integrateurs}).
\end{boxscenario}

% ============================================================================
\section{La Notation de NOUN PAY au Regard du Droit}
\label{sec:c19-noun}
% ============================================================================

\subsection{Les questions à instruire}

La loi n°~2024/017 (Chapitre~\ref{chap:droit}) s'applique aux données
personnelles utilisées par le modèle. Sans préjuger de la lettre de ses
articles, que l'on vérifiera dans le texte publié, six questions doivent
être instruites~:
\begin{enumerate}
    \item \textbf{La base juridique}~: sur quel fondement le traitement des
          données des demandeurs est-il licite~?
    \item \textbf{La finalité et la proportionnalité}~: les données utilisées
          sont-elles nécessaires à l'évaluation du crédit~? L'opérateur
          mobile l'est-il~?
    \item \textbf{L'information des personnes}~: les demandeurs savent-ils
          que leur demande est évaluée par un modèle~?
    \item \textbf{Les décisions automatisées}~: de nombreuses lois de
          protection des données de cette génération encadrent les
          décisions fondées exclusivement sur un traitement automatisé et
          ouvrent un droit à l'intervention humaine~; il faut vérifier si
          et comment la loi camerounaise le prévoit.
    \item \textbf{Le transfert hors du territoire}~: les données sont
          hébergées hors du Cameroun (Chapitre~\ref{chap:agile-cloud}).
    \item \textbf{La sécurité et la conservation}~: qui accède aux données,
          et combien de temps sont-elles gardées~?
\end{enumerate}

\begin{boxalert}
\textbf{Violation typique.} Une entreprise supprime de son modèle la
variable « région » pour éviter toute discrimination, mais conserve des
variables qui lui sont étroitement corrélées (opérateur mobile, préfixe
téléphonique, agence). Le biais demeure, par une autre voie.

\smallskip
\textbf{Règle~:} on ne prouve pas l'absence de biais en retirant une
variable~; on la prouve en \emph{mesurant} les résultats par groupe.
\end{boxalert}

\subsection{Le dispositif de gouvernance}

\begin{boxmetier}{Responsable GRC}
\textbf{Doter NOUN PAY d'une gouvernance de l'IA proportionnée}
\begin{enumerate}
    \item \textbf{Un inventaire}~: lister les modèles en production, leur
          usage, les données qu'ils emploient et leurs effets sur les
          clients.
    \item \textbf{Une classification}~: la notation de crédit, qui décide
          de l'accès d'une personne à un service financier, est à risque
          élevé.
    \item \textbf{Un responsable}~: désigner un A pour chaque modèle à
          risque élevé, au niveau de la direction (Chapitre~\ref{chap:processus}).
    \item \textbf{Une évaluation d'impact avant chaque version}~: mesure
          des écarts entre groupes, avis du délégué à la protection des
          données, décision écrite.
    \item \textbf{Une surveillance en production}~: ratio d'impact et taux
          d'erreur suivis chaque mois, avec un seuil qui déclenche la
          suspension.
    \item \textbf{Un recours humain}~: toute personne refusée peut demander
          un réexamen par un analyste.
\end{enumerate}
\end{boxmetier}

\begin{boxmetier}{Auditeur SI}
\textbf{Auditer un modèle de décision}
\begin{enumerate}
    \item Demander la documentation du modèle~: finalité, données
          d'entraînement, variables, performances, limites connues.
    \item Vérifier que chaque mise à jour a fait l'objet d'une évaluation
          d'impact et d'une autorisation (BAI06).
    \item Recalculer, sur un échantillon de décisions récentes, les taux
          par groupe et le ratio d'impact.
    \item Vérifier l'existence et l'usage effectif du recours humain.
\end{enumerate}
\end{boxmetier}

\begin{boiteRemarque}{L'humain dans la boucle suffit-il~?}
Prévoir qu'un analyste « valide » chaque décision du modèle est souvent
présenté comme la garantie ultime. L'expérience montre pourtant qu'un
analyste qui traite deux cents dossiers par jour finit par entériner les
propositions du modèle. La surveillance humaine n'est réelle que si
l'analyste dispose du temps, de l'information (les motifs de la
proposition) et du pouvoir de s'écarter du modèle sans avoir à s'en
justifier davantage que s'il le suivait.
\end{boiteRemarque}

% ============================================================================
\separateur
% ============================================================================

\begin{boxresume}
\textbf{Ce qu'il faut retenir de ce chapitre~:}
\begin{itemize}
    \item \textbf{Les données sont un actif gouverné}~: un propriétaire
          métier, un intendant, un dépositaire~; COBIT y consacre
          l'objectif APO14.
    \item \textbf{Six dimensions de qualité} (exactitude, complétude,
          cohérence, actualité, validité, unicité) et un cycle de vie avec
          des durées de conservation fondées.
    \item \textbf{L'IA apprend, s'explique mal et dérive}~: elle appelle une
          gouvernance propre.
    \item \textbf{Les références}~: ISO/IEC~38507 pour le conseil,
          ISO/IEC~42001 pour le système de management, ISO/IEC~23894 pour
          les risques, le NIST AI RMF, et la Recommandation de l'UNESCO
          pour les principes éthiques.
    \item \textbf{Mesurer le biais}~: le ratio d'impact de NOUN PAY (0,51)
          est un signal d'alerte~; retirer une variable ne prouve rien.
    \item \textbf{Un modèle à risque élevé} a un responsable, une évaluation
          d'impact à chaque version, une surveillance en production et un
          recours humain réel.
\end{itemize}
\end{boxresume}

\bigskip

\begin{boxexo}[title={\ding{45}\quad Exercice~19.1. Quelle dimension de qualité~? \niveaubase}]
Rattachez chacune des anomalies suivantes à une dimension de qualité
(Tableau~19.2).
\begin{enumerate}[(a)]
    \item Un agent payé deux fois par l'ANPP parce qu'il figure sous deux
          matricules.
    \item Un client de la BANQUE SANAGA dont la date de naissance est le
          31~février.
    \item Un conteneur localisé à Kribi alors qu'il est à Douala depuis la
          veille.
    \item Un dossier de crédit sans les revenus du demandeur.
\end{enumerate}
\end{boxexo}

\begin{boxexo}[title={\ding{45}\quad Exercice~19.2. Le registre des données de la BANQUE SANAGA \niveaumoyen}]
Établissez les cinq premières lignes du registre des données critiques de
la BANQUE SANAGA (exigence de contrôle de la section~\ref{sec:c19-qualite}),
en commençant par les soldes des comptes.
\begin{enumerate}[(a)]
    \item Pour chaque domaine, désignez le propriétaire métier et
          justifiez.
    \item Fixez les règles de qualité et leur fréquence de contrôle.
    \item Indiquez la durée de conservation et son fondement, en signalant
          ce qui reste à vérifier.
\end{enumerate}
\end{boxexo}

\begin{boxexo}[title={\ding{45}\quad Exercice~19.3. La réponse à l'Autorité de protection des données \niveauexpert}]
Rédigez, pour Mme~KAMGA, la réponse de NOUN PAY à l'Autorité de
protection des données.
\begin{enumerate}[(a)]
    \item Exposez les faits, en séparant constatations et analyses.
    \item Répondez aux six questions de la section~\ref{sec:c19-noun}, en
          indiquant ce qui suppose une vérification du texte de la loi.
    \item Présentez les mesures prises et prévues, avec leur échéance.
    \item Expliquez en quoi la Recommandation de l'UNESCO sur l'éthique de
          l'IA éclaire les engagements pris.
\end{enumerate}
\end{boxexo}

\bigskip

\begin{boxinfo}
\textbf{Transition.} La Partie~\ref{part:avance} a confronté la gouvernance
des SI aux défis de la vitesse, du nuage, des données et de
l'intelligence artificielle. La Partie~\ref{part:recherche} prend du
recul~: le Chapitre~\ref{chap:recherche-mesure} présente l'état de la
recherche sur la mesure de la gouvernance des SI.
\end{boxinfo}


    % ======================================================================
    % PARTIE VI — RECHERCHE ET CONTRIBUTIONS
    % Niveau : expert — recherche
    % ======================================================================
    \part{Recherche et Contributions}
    \label{part:recherche}

        % ============================================================================
% Fichier  : partie6/P06_C20_mesurer_gouvernance_recherche.tex
% Titre    : Mesurer la Gouvernance des SI : État de la Recherche
% Auteur   : MINKA MI NGUIDJOÏ Thierry Emmanuel
% Version  : 1.0 -- octobre 2026
% ============================================================================
% Chapitre de recherche. Les références citées sont des publications
% académiques identifiées ; leurs résultats sont résumés, non reproduits.
% Les chiffres repris d'une étude sont attribués à leurs auteurs.
% ============================================================================

\chapter{Mesurer la Gouvernance des SI : État de la Recherche}
\label{chap:recherche-mesure}
\chaptermark{État de la recherche}

\epigraph{%
    « Ce qui compte ne peut pas toujours être compté,\\
    et ce qui peut être compté ne compte pas toujours. »%
}{Maxime attribuée, sous diverses formes, à plusieurs auteurs}

\niveauChapitre{Recherche}{%
    Maîtrise du Chapitre~\ref{chap:cobit-capacite} (capacité et maturité)
    et du Chapitre~\ref{chap:alignement} (alignement et valeur).%
}

\begin{boxobjectifs}
\begin{itemize}
    \item Synthétiser les travaux de référence sur la gouvernance des SI.
    \item Discuter la validité des modèles de maturité.
    \item Lire de façon critique une étude empirique.
    \item Formuler une question de recherche sur la gouvernance des SI en
          Afrique centrale.
\end{itemize}
\end{boxobjectifs}

\bigskip

% ============================================================================
\section*{Question directrice : la gouvernance des SI sert-elle à quelque chose~?}
% ============================================================================

Tout ce cours repose sur une conviction~: une meilleure gouvernance des SI
produit de meilleurs résultats. Le praticien l'admet volontiers~; le
chercheur demande des preuves. Qu'entend-on par « meilleure
gouvernance »~? Comment la mesure-t-on~? Et si l'on observe que les
organisations bien gouvernées réussissent mieux, est-ce la gouvernance qui
produit le succès, ou le succès qui permet de financer la gouvernance~?

La Partie~\ref{part:recherche}, qui s'ouvre ici, prend ce recul. Ce
chapitre présente les grandes étapes de la recherche
(section~\ref{sec:c20-courants}), ce qu'elle a établi
(section~\ref{sec:c20-resultats}), le débat sur les modèles de maturité
(section~\ref{sec:c20-maturite}), une méthode de lecture critique
(section~\ref{sec:c20-lecture}), puis les pistes de recherche ouvertes en
Afrique centrale (section~\ref{sec:c20-afrique}). Il s'adresse en
particulier aux étudiants qui préparent un mémoire.

\begin{boxinfo}
\textbf{Règle de lecture.} Les travaux cités dans ce chapitre sont des
publications académiques identifiées dans la bibliographie. Leurs
résultats sont résumés et attribués à leurs auteurs. Un résumé n'est pas
une lecture~: avant de citer une étude dans un mémoire, lisez-la.
\end{boxinfo}

\separateur

% ============================================================================
\section{Trois Courants de Recherche}
\label{sec:c20-courants}
% ============================================================================

\subsection{Les structures de décision}

Le premier courant s'est demandé \emph{où} se prennent les décisions
informatiques. Dans les années 1990, la recherche oppose les modèles
centralisés, décentralisés et fédéraux. Sambamurthy et Zmud montrent que
le bon arrangement dépend de plusieurs facteurs de contingence qui
agissent ensemble~: la taille, la stratégie, la structure de
l'organisation~\cite{sambamurthy1999arrangements}. Weill et Ross
prolongent ce travail avec les archétypes et la matrice des droits de
décision (Chapitre~\ref{chap:alignement})~\cite{weill2004itgovernance}.

\subsection{Les mécanismes}

Le deuxième courant passe de la question « qui décide~? » à la question
« par quels moyens~? ». Peterson distingue les structures, les processus
et les capacités relationnelles~\cite{peterson2004crafting}~; De Haes et
Van Grembergen reprennent et outillent cette distinction en structures,
processus et mécanismes relationnels (Chapitre~\ref{chap:gouv-entreprise}),
et étudient leur effet sur l'alignement~\cite{dehaes2009exploratory,dehaes2015egit}.

\subsection{L'effet sur la performance}

Le troisième courant cherche à établir le lien entre gouvernance des SI,
alignement et performance de l'organisation. Brown et Grant en ont dressé
une revue structurante au milieu des années 2000~\cite{brown2005framing}.
Les études quantitatives se sont depuis multipliées, avec des méthodes de
plus en plus exigeantes, comme celle de Wu, Straub et Liang, qui croise
les réponses de dirigeants métiers et de responsables informatiques des
mêmes entreprises~\cite{wu2015mechanisms}.

\begin{tableaucours}{L{2.6cm}L{4.8cm}L{4.8cm}}{%
    \tableauHeader{Courant} &
    \tableauHeader{Question} &
    \tableauHeader{Travaux de référence}
}
Structures & Où se prennent les décisions~? & Sambamurthy et Zmud (1999)~;
Weill et Ross (2004) \\
Mécanismes & Par quels moyens la gouvernance s'exerce-t-elle~? & Peterson
(2004)~; De Haes et Van Grembergen (2009, 2015) \\
Performance & La gouvernance améliore-t-elle les résultats~? & Brown et
Grant (2005)~; Wu, Straub et Liang (2015) \\
\end{tableaucours}
\vspace{-0.8em}
\begin{center}{\small\itshape Tableau~20.1. Trois courants de la recherche
sur la gouvernance des SI}\end{center}

% ============================================================================
\section{Ce que la Recherche a Établi}
\label{sec:c20-resultats}
% ============================================================================

\subsection{Des résultats convergents}

Avec les précautions qui s'imposent, trois résultats reviennent avec
régularité dans la littérature.

\begin{enumerate}
    \item \textbf{La gouvernance compte.} Weill et Ross ont rapporté, sur
          un large échantillon d'entreprises, que celles dont la
          gouvernance des SI était la plus efficace obtenaient, à
          stratégie comparable, une rentabilité nettement supérieure à
          celles dont la gouvernance était la plus faible, de l'ordre de
          20~\% selon leurs mesures~\cite{weill2004itgovernance}.
    \item \textbf{Les mécanismes relationnels comptent autant que les
          structures.} Plusieurs études observent qu'à structures égales,
          ce sont la compréhension mutuelle et la collaboration entre
          métiers et informatique qui font la différence
          (Chapitre~\ref{chap:cas-integrateurs}~: le cas de la BANQUE SANAGA un an
          après).
    \item \textbf{L'effet passe par l'alignement.} La gouvernance agit sur
          la performance surtout par l'intermédiaire de l'alignement
          stratégique~: elle aide à faire les bons choix plus qu'elle ne
          produit directement des résultats~\cite{wu2015mechanisms}.
\end{enumerate}

\subsection{Ce qui reste débattu}

\begin{tableaucours}{L{4.6cm}L{7.6cm}}{%
    \tableauHeader{Question} &
    \tableauHeader{État du débat}
}
Le sens de la causalité & Les organisations performantes ont-elles une
bonne gouvernance parce qu'elles en ont les moyens~? Les études
transversales ne permettent pas de trancher. \\
La mesure de la gouvernance & Elle repose souvent sur l'autoévaluation
d'un seul répondant, ce qui fragilise les résultats. \\
La transposition des résultats & La plupart des études portent sur des
organisations d'Amérique du Nord, d'Europe et d'Asie~; leur validité
ailleurs n'est pas démontrée. \\
L'effet des référentiels & L'adoption de COBIT ou d'ITIL est fréquente,
mais leur effet propre, toutes choses égales par ailleurs, reste mal
établi. \\
\end{tableaucours}
\vspace{-0.8em}
\begin{center}{\small\itshape Tableau~20.2. Questions débattues}\end{center}

\begin{boxalert}
\textbf{Violation typique.} Un mémoire affirme que « COBIT améliore la
performance des entreprises de 20~\% », en citant Weill et Ross. Or Weill
et Ross n'étudient pas COBIT, et leur écart de rentabilité est une
association mesurée sur leur échantillon, pas un effet causal
transposable.

\smallskip
\textbf{Règle~:} citer un résultat, c'est citer aussi ce qu'il mesure, sur
quel échantillon, et avec quelle méthode.
\end{boxalert}

% ============================================================================
\section{Le Débat sur les Modèles de Maturité}
\label{sec:c20-maturite}
% ============================================================================

\subsection{Un succès considérable}

Depuis le modèle de maturité du logiciel de l'université Carnegie Mellon,
les modèles de maturité se sont multipliés~: COBIT en utilise un depuis
sa version~3, et le modèle de capacité de COBIT~2019 s'inspire du CMMI
(Chapitre~\ref{chap:cobit-capacite}). Leur succès tient à leur
simplicité~: une échelle de 0 à~5, que tout dirigeant comprend.

\subsection{Quatre critiques}

La recherche en a pourtant relevé les limites~\cite{becker2009maturity,poppelbuss2011maturity}.

\begin{enumerate}
    \item \textbf{La linéarité.} Les modèles supposent un chemin unique,
          du niveau~1 au niveau~5. Or une organisation peut avoir de bonnes
          raisons de ne pas viser le niveau~5, ou de progresser autrement.
    \item \textbf{Le fondement théorique.} Beaucoup de modèles sont construits
          à partir de l'expérience de leurs auteurs, sans validation
          empirique de leurs niveaux ni de leurs critères.
    \item \textbf{La fiabilité de la mesure.} Deux évaluateurs n'attribuent
          pas toujours le même niveau au même processus, surtout en
          autoévaluation.
    \item \textbf{Le lien avec la performance.} Qu'un niveau de maturité plus
          élevé produise de meilleurs résultats est postulé plus souvent que
          démontré. Certaines études l'observent~\cite{simonsson2010maturity},
          d'autres non, et le lien dépend du contexte.
\end{enumerate}

\begin{boiteRemarque}{Faut-il renoncer aux modèles de maturité~?}
Non, mais il faut les utiliser pour ce qu'ils sont~: des outils de
\emph{diagnostic} et de \emph{dialogue}, pas des instruments de mesure
scientifique. COBIT~2019 en tient compte de deux manières, déjà vues dans
ce cours~: il fait dépendre la capacité \emph{cible} des facteurs de
conception, ce qui répond à la critique de la linéarité
(Chapitre~\ref{chap:cobit-conception}), et il rattache chaque niveau à des
activités précises, ce qui améliore la fiabilité de l'évaluation
(Chapitre~\ref{chap:cobit-capacite}). Pöppelbuß et Röglinger ont proposé des
principes de conception qui permettent de juger si un modèle de maturité
est utile~\cite{poppelbuss2011maturity}.
\end{boiteRemarque}

% ============================================================================
\section{Lire une Étude Empirique}
\label{sec:c20-lecture}
% ============================================================================

\subsection{Une grille de lecture}

\begin{tableaucours}{L{3.2cm}L{9.0cm}}{%
    \tableauHeader{Élément} &
    \tableauHeader{Questions à se poser}
}
Question et théorie & Quelle question précise~? Quelle théorie justifie
les hypothèses (agence, alignement, contingence)~? \\
Échantillon & Combien d'organisations, de quels pays, secteurs, tailles~?
Comment ont-elles été choisies~? \\
Mesures & Comment la gouvernance et la performance sont-elles mesurées~?
Par qui~? Les échelles sont-elles validées~? \\
Biais de méthode & Le même répondant décrit-il à la fois la gouvernance et
la performance~? \\
Causalité & L'étude est-elle transversale (une seule date) ou
longitudinale~? Peut-elle établir le sens de l'effet~? \\
Généralisation & Les résultats valent-ils au-delà de l'échantillon~? Pour
une organisation camerounaise~? \\
Statut & Article évalué par des pairs, communication, prépublication,
rapport de cabinet~? \\
\end{tableaucours}
\vspace{-0.8em}
\begin{center}{\small\itshape Tableau~20.3. Grille de lecture critique d'une
étude empirique}\end{center}

\subsection{Le statut des sources}

\begin{boxdef}
\termecle{Évaluation par les pairs}\index{évaluation par les pairs}
(\textit{peer review})~: examen d'un article par des chercheurs
indépendants avant sa publication dans une revue ou des actes de
conférence. Elle ne garantit pas que les résultats sont justes, mais
qu'ils ont été soumis à une critique méthodologique.

\medskip
\termecle{Prépublication}\index{prépublication} (\textit{preprint})~:
version d'un travail rendue publique avant son évaluation par les pairs.
Elle doit être citée comme telle.
\end{boxdef}

Les enquêtes des cabinets de conseil et des éditeurs de logiciels sont
utiles pour connaître les pratiques, mais elles poursuivent souvent un
objectif commercial et publient rarement leur méthode. Elles se citent
avec cette réserve.

\begin{boxscenario}{Exercice de lecture. Une étude fictive}
\textbf{Résumé.} « Une enquête menée auprès de 85~DSI d'entreprises
camerounaises montre que les entreprises qui ont adopté COBIT ont une
satisfaction client supérieure de 15~\% à celles qui ne l'ont pas adopté
($p < 0{,}05$). »

\medskip
\textbf{Lecture critique.} L'échantillon est-il représentatif (85~DSI
volontaires, sans doute les plus sensibles au sujet)~? La satisfaction
client est-elle mesurée auprès des clients, ou estimée par le DSI~? Les
entreprises qui adoptent COBIT sont-elles plus grandes, plus riches, mieux
organisées par ailleurs~? Une différence statistiquement significative
n'est pas une preuve d'effet. Conclusion~: résultat intéressant, à
présenter comme une association observée sur un échantillon de
convenance.
\end{boxscenario}

% ============================================================================
\section{Des Pistes de Recherche en Afrique Centrale}
\label{sec:c20-afrique}
% ============================================================================

\subsection{Un champ peu exploré}

La recherche publiée sur la gouvernance des SI porte majoritairement sur
des organisations des pays industrialisés. Les travaux consacrés à
l'Afrique centrale restent peu nombreux, alors que le contexte y présente
des caractéristiques propres~: concentration du pouvoir de décision,
tutelle publique, dépendance aux intégrateurs étrangers, régulation
récente, contraintes d'infrastructure (Chapitre~\ref{chap:gouv-entreprise}).
C'est une raison suffisante pour que les mémoires et les thèses s'en
emparent.

\subsection{Des questions possibles}

\begin{tableaucours}{L{5.6cm}L{3.2cm}L{3.4cm}}{%
    \tableauHeader{Question} &
    \tableauHeader{Méthode adaptée} &
    \tableauHeader{Terrain}
}
Comment les droits de décision informatiques se répartissent-ils dans
les établissements publics camerounais~? & Études de cas multiples & Trois
à cinq établissements \\
Les comités stratégiques des SI créés après une crise fonctionnent-ils
réellement~? & Étude longitudinale & Une banque, sur trois ans \\
Quels facteurs de conception COBIT pèsent le plus dans les organisations
de la zone CEMAC~? & Enquête par questionnaire & Banques et assurances \\
Comment adapter un modèle de capacité aux PME camerounaises~? &
Recherche orientée conception~\cite{hevner2004design} & PME de Douala et
Yaoundé \\
L'entrée en vigueur de la loi n°~2024/017 modifie-t-elle la gouvernance
des données~? & Études de cas avant et après & Entreprises de paiement
mobile \\
\end{tableaucours}
\vspace{-0.8em}
\begin{center}{\small\itshape Tableau~20.4. Exemples de questions de
recherche}\end{center}

\subsection{Les méthodes}

\begin{itemize}
    \item \textbf{L'étude de cas} convient aux questions « comment~? » et
          « pourquoi~? » dans des contextes peu connus~; elle exige un
          accès de confiance au terrain et une rigueur dans la collecte des
          preuves.
    \item \textbf{L'enquête par questionnaire} permet de mesurer des
          associations~; elle exige des échelles validées et un échantillon
          suffisant, souvent difficile à réunir.
    \item \textbf{La recherche orientée conception} construit un artefact
          (un modèle, une méthode, un outil) et l'évalue en situation~; elle
          convient bien aux praticiens qui préparent un mémoire
          professionnel~\cite{hevner2004design}.
\end{itemize}

\begin{boxprocedure}
\textbf{Un projet de mémoire recevable.} Le projet présenté au directeur
de mémoire comporte~:
\begin{itemize}
    \item une question de recherche précise, et sa justification au regard
          de la littérature~;
    \item le cadre théorique mobilisé~;
    \item la méthode, le terrain et l'accès obtenu (accord écrit de
          l'organisation)~;
    \item les règles de confidentialité et d'anonymisation des données
          recueillies~;
    \item le calendrier et les livrables~;
    \item une bibliographie d'au moins quinze références, dont une majorité
          évaluées par les pairs.
\end{itemize}
\end{boxprocedure}

\begin{boxmetier}{Auditeur SI}
\textbf{Faire de sa pratique un terrain de recherche}
\begin{enumerate}
    \item Les missions d'audit produisent des données rares sur la
          gouvernance réelle des organisations~: écarts, causes,
          capacités.
    \item Elles ne peuvent servir à la recherche qu'avec l'accord des
          organisations et une anonymisation rigoureuse, dans le respect de
          la confidentialité (Chapitre~\ref{chap:audit}).
    \item Agrégées sur plusieurs missions, elles permettent de répondre à
          des questions qu'aucune enquête par questionnaire n'atteint.
\end{enumerate}
\end{boxmetier}

% ============================================================================
\separateur
% ============================================================================

\begin{boxresume}
\textbf{Ce qu'il faut retenir de ce chapitre~:}
\begin{itemize}
    \item \textbf{Trois courants}~: les structures de décision, les
          mécanismes, l'effet sur la performance.
    \item \textbf{Des résultats convergents}~: la gouvernance compte, les
          mécanismes relationnels comptent autant que les structures, et
          l'effet passe par l'alignement.
    \item \textbf{Des questions ouvertes}~: le sens de la causalité, la
          fiabilité des mesures, la transposition des résultats, l'effet
          propre des référentiels.
    \item \textbf{Les modèles de maturité} sont des outils de diagnostic et
          de dialogue, pas des instruments scientifiques~: linéarité,
          fondement, fiabilité et lien avec la performance sont discutés.
    \item \textbf{Lire une étude}, c'est en examiner l'échantillon, les
          mesures, les biais, la causalité et le statut.
    \item \textbf{L'Afrique centrale} est un terrain peu exploré, qui
          appelle études de cas, enquêtes et recherche orientée conception.
\end{itemize}
\end{boxresume}

\bigskip

\begin{boxexo}[title={\ding{45}\quad Exercice~20.1. Association ou causalité~? \niveaubase}]
Pour chacune des affirmations suivantes, indiquez si elle décrit une
association ou un effet causal, et ce qu'il faudrait pour établir la
causalité.
\begin{enumerate}[(a)]
    \item « Les banques dotées d'un comité stratégique des SI ont moins
          d'incidents majeurs. »
    \item « Après la création de son comité, la BANQUE SANAGA a réduit ses
          incidents majeurs de moitié. »
    \item « Les entreprises certifiées ISO/IEC~27001 subissent moins de
          fuites de données. »
\end{enumerate}
\end{boxexo}

\begin{boxexo}[title={\ding{45}\quad Exercice~20.2. Lire un article \niveaumoyen}]
Choisissez l'un des articles cités dans ce chapitre et appliquez-lui la
grille du Tableau~20.3. Rédigez une note de lecture de deux pages~:
question, théorie, méthode, résultats, limites, intérêt pour une
organisation camerounaise.
\end{boxexo}

\begin{boxexo}[title={\ding{45}\quad Exercice~20.3. Un projet de mémoire \niveauexpert}]
Rédigez un projet de mémoire sur l'une des questions du Tableau~20.4, ou
sur une question de votre choix, en respectant l'exigence de contrôle de
la section~\ref{sec:c20-afrique}. Justifiez en particulier le choix de la
méthode et la manière d'obtenir l'accès au terrain.
\end{boxexo}

\bigskip

\begin{boxinfo}
\textbf{Transition.} Ce chapitre a présenté l'état de la recherche et les
questions qu'elle laisse ouvertes. Le Chapitre~\ref{chap:contributions}
présente des contributions et des pistes de recherche propres à
l'auteur et à son laboratoire.
\end{boxinfo}

        % ============================================================================
% Fichier  : partie6/P06_C21_contributions_recherche.tex
% Titre    : Contributions et Pistes de Recherche
% Auteur   : MINKA MI NGUIDJOÏ Thierry Emmanuel
% Version  : 0.1 -- octobre 2026 (squelette)
% ============================================================================

\chapter{Contributions et Pistes de Recherche}
\label{chap:contributions}

\niveauChapitre{Recherche}{%
    Maîtrise du Chapitre~\ref{chap:recherche-mesure}%
}

\begin{boxobjectifs}
\begin{itemize}
    \item Objectifs à définir avec l'auteur selon les travaux présentés.
\end{itemize}
\end{boxobjectifs}

\begin{boxalert}
Ce chapitre est en cours de rédaction.
\end{boxalert}


% --------------------------------------------------------------------------
% ANNEXES
% --------------------------------------------------------------------------
\appendix

    % ============================================================================
% Fichier : annexes/Ann_A_glossaire.tex
% Titre   : Glossaire
% Version : 1.0 -- novembre 2026
% ============================================================================
% Définitions reprises, en abrégé, des encadrés « Définition » des chapitres
% (termes marqués \termecle) ; renvoi au chapitre de référence.
% ============================================================================

\chapter{Glossaire}
\label{ann:glossaire}

Ce glossaire reprend, sous une forme abrégée, les définitions données dans
les chapitres. Le renvoi indique où le terme est défini et illustré.

\begin{description}
    \item[Accord de niveau de service (SLA)] Accord documenté entre un
          fournisseur et son client, qui précise le service et les niveaux
          attendus (Ch.~\ref{chap:itil-niveaux}).
    \item[Accord de niveau opérationnel (OLA)] Accord entre le fournisseur et
          une autre équipe interne qui soutient le service
          (Ch.~\ref{chap:itil-niveaux}).
    \item[Activité (COBIT)] Action concrète qui met en œuvre une pratique~;
          rattachée à un niveau de capacité de~2 à~5
          (Ch.~\ref{chap:cobit-objectifs}).
    \item[Alignement stratégique] Degré de cohérence entre la stratégie de
          l'organisation, celle des technologies de l'information, et les
          infrastructures et processus qui les mettent en œuvre
          (Ch.~\ref{chap:alignement}).
    \item[Alignement volontaire] Choix de s'inspirer d'un référentiel ou
          d'une norme sans viser de certificat (Ch.~\ref{chap:normes}).
    \item[Analyse d'impact sur l'activité (BIA)] Analyse qui identifie les
          activités essentielles, mesure l'effet de leur interruption dans
          le temps et en déduit les objectifs de reprise
          (Ch.~\ref{chap:itil-niveaux}).
    \item[Autorité de changement] Personne ou groupe qui autorise un
          changement~; l'autorisation est déléguée selon le risque
          (Ch.~\ref{chap:itil-changements}).
    \item[Base de données de gestion de la configuration (CMDB)] Base qui
          enregistre les éléments de configuration, leurs attributs et
          leurs relations (Ch.~\ref{chap:itil-changements}).
    \item[Cadre de référence] Ensemble structuré de principes et
          d'objectifs qui guide un domaine sans imposer d'exigences
          vérifiables par un tiers~; par exemple COBIT, COSO
          (Ch.~\ref{chap:referentiels}).
    \item[Capacité d'un processus] Mesure de la manière dont un processus
          est mis en œuvre et fonctionne, de~0 à~5 dans COBIT~2019
          (Ch.~\ref{chap:cobit-capacite}).
    \item[Cascade des objectifs] Chaîne qui relie les besoins des parties
          prenantes, les objectifs d'entreprise, les objectifs d'alignement
          et les objectifs de gouvernance et de management
          (Ch.~\ref{chap:cobit-objectifs}).
    \item[Certification] Attestation, par un organisme tiers accrédité,
          qu'un système de management est conforme à une norme, sur un
          périmètre et pour une durée donnés (Ch.~\ref{chap:normes}).
    \item[Chaîne de déploiement] Suite automatisée des étapes par lesquelles
          un changement de code passe jusqu'à la mise en production
          (Ch.~\ref{chap:agile-cloud}).
    \item[Changement] Ajout, modification ou retrait de tout élément
          susceptible d'avoir un effet sur un service
          (Ch.~\ref{chap:itil-changements}).
    \item[Chronologie de défaillance] Suite datée des faits, décisions et
          absences de décision autour d'une défaillance, chaque entrée
          rattachée à une pièce (Ch.~\ref{chap:prologue}).
    \item[Co-création de valeur] Idée selon laquelle la valeur naît de la
          collaboration entre fournisseur et consommateur du service
          (Ch.~\ref{chap:itil-svs}).
    \item[Comité stratégique, comité de pilotage des SI] Le premier assiste
          le conseil sur les orientations, la valeur et les risques du
          numérique~; le second, au niveau de la direction générale,
          priorise les projets et arbitre les ressources
          (Ch.~\ref{chap:gouv-entreprise}).
    \item[Composante] Élément qui contribue au fonctionnement du système de
          gouvernance~; COBIT~2019 en distingue sept types
          (Ch.~\ref{chap:cobit-archi}).
    \item[Conformité] Respect d'exigences obligatoires (lois, règlements,
          contrats)~; elle se démontre (Ch.~\ref{chap:normes}).
    \item[Contrat de sous-traitance] Contrat entre le fournisseur et un tiers
          extérieur dont il dépend pour fournir le service
          (Ch.~\ref{chap:itil-niveaux}).
    \item[Contrôles d'application] Contrôles intégrés à un système, qui
          agissent sur le traitement des transactions (Ch.~\ref{chap:audit}).
    \item[Contrôles généraux informatiques] Contrôles qui s'appliquent à tout
          l'environnement informatique~: accès, changements, exploitation,
          développement et acquisition (Ch.~\ref{chap:gouv-entreprise},
          \ref{chap:audit}).
    \item[Coûts d'agence] Dépenses de surveillance du principal, dépenses
          d'engagement de l'agent, et perte résiduelle due aux divergences
          d'intérêts (Ch.~\ref{chap:gouv-entreprise}).
    \item[Coût total de possession (TCO)] Somme des coûts sur la durée de vie
          d'une solution, de l'acquisition à la sortie
          (Ch.~\ref{chap:alignement}).
    \item[Déclaration d'applicabilité] Liste des mesures de l'annexe~A de
          l'ISO/IEC~27001, retenues ou exclues, avec leur justification
          (Ch.~\ref{chap:normes}).
    \item[Défaillance de gouvernance] Absence, inadéquation ou contournement
          d'un mécanisme de décision, de contrôle ou de reddition de
          comptes, qui laisse un risque connu se réaliser
          (Ch.~\ref{chap:prologue}).
    \item[Demande de service] Demande d'un utilisateur portant sur une
          prestation déjà prévue dans le service (Ch.~\ref{chap:itil-incidents}).
    \item[DevOps] Pratiques qui rapprochent développement et exploitation,
          pour livrer des changements fréquents, fiables et réversibles
          (Ch.~\ref{chap:agile-cloud}).
    \item[Disponibilité] Aptitude d'un service à remplir sa fonction quand on
          en a besoin~; se mesure sur le temps de service convenu
          (Ch.~\ref{chap:itil-niveaux}).
    \item[Domaine d'intérêt] Thème de gouvernance traité par un ensemble
          d'objectifs et de composantes (Ch.~\ref{chap:cobit-archi}).
    \item[Droits de décision] Répartition formelle du pouvoir de décider sur
          une question donnée (Ch.~\ref{chap:prologue},
          \ref{chap:alignement}).
    \item[Durée maximale d'interruption admissible (DMIA)] Durée au-delà de
          laquelle l'interruption d'une activité devient inacceptable
          (Ch.~\ref{chap:itil-niveaux}).
    \item[Élément de configuration] Tout composant qui doit être géré pour
          fournir un service (Ch.~\ref{chap:itil-changements}).
    \item[Erreur connue] Problème analysé, dont la cause est identifiée, mais
          pas encore résolu (Ch.~\ref{chap:itil-incidents}).
    \item[Évaluation par les pairs] Examen d'un article par des chercheurs
          indépendants avant sa publication (Ch.~\ref{chap:recherche-mesure}).
    \item[Facteur de conception] Facteur qui influence la conception du
          système de gouvernance~; COBIT~2019 en retient onze
          (Ch.~\ref{chap:cobit-archi}).
    \item[Garantie] Assurance que le service répondra aux exigences
          convenues~: disponibilité, capacité, sécurité, continuité
          (Ch.~\ref{chap:itil-svs}).
    \item[Gestionnaire de processus] Personne qui pilote l'exécution
          quotidienne d'un processus (Ch.~\ref{chap:processus}).
    \item[Gouvernance d'entreprise] Relations, règles et mécanismes par
          lesquels une société est dirigée et contrôlée
          (Ch.~\ref{chap:gouv-entreprise}).
    \item[Gouvernance des données] Exercice de l'autorité et du contrôle sur
          la gestion des données (Ch.~\ref{chap:donnees-ia}).
    \item[Gouvernance des SI] Activités par lesquelles l'organe de
          gouvernance évalue, dirige et surveille l'usage de l'information
          et des technologies, pour qu'il crée de la valeur
          (Ch.~\ref{chap:prologue}, \ref{chap:gouv-entreprise}).
    \item[Habilitation des changements] Pratique qui maximise le nombre de
          changements réussis, en évaluant les risques, en autorisant et en
          maîtrisant le calendrier (Ch.~\ref{chap:itil-changements}).
    \item[Incident] Interruption non planifiée d'un service, ou dégradation
          de sa qualité (Ch.~\ref{chap:itil-incidents}).
    \item[Incident majeur] Incident dont l'impact exige une réponse
          dépassant le traitement ordinaire (Ch.~\ref{chap:itil-incidents}).
    \item[Indicateur clé d'objectif (KGI), de performance (KPI)] Le premier
          dit si un objectif est atteint, après coup~; le second, si un
          processus a des chances de l'atteindre, en cours de route
          (Ch.~\ref{chap:processus}).
    \item[Management des SI] Activités qui planifient, bâtissent, exploitent
          et surveillent les SI, selon la direction fixée par la gouvernance
          (Ch.~\ref{chap:prologue}).
    \item[Matrice RACI] Tableau qui croise activités et rôles~: qui réalise
          (R), qui en répond (A, un seul), qui est consulté (C), qui est
          informé (I) (Ch.~\ref{chap:processus}).
    \item[Maturité] Mesure de performance d'un domaine d'intérêt~: il atteint
          un niveau quand tous ses processus l'atteignent
          (Ch.~\ref{chap:cobit-capacite}).
    \item[Modèle des trois lignes] Répartition des rôles entre l'organe de
          gouvernance, le management et l'audit interne
          (Ch.~\ref{chap:audit}).
    \item[Niveau de capacité cible] Niveau de~0 à~5 que l'organisation décide
          d'atteindre pour un processus (Ch.~\ref{chap:cobit-conception}).
    \item[Norme certifiable] Norme qui formule des exigences dont le respect
          peut être audité par un organisme de certification
          (Ch.~\ref{chap:referentiels}).
    \item[Objectif d'alignement, d'entreprise] Le second est un but de
          l'organisation~; le premier, un but que l'informatique doit
          atteindre pour le soutenir (Ch.~\ref{chap:cobit-objectifs}).
    \item[Objectif de gouvernance ou de management] Ce qu'il faut atteindre
          pour que l'information et les technologies contribuent aux
          objectifs~; COBIT~2019 en compte quarante
          (Ch.~\ref{chap:cobit-archi}).
    \item[Optimisation des ressources, des risques] Disposer des ressources
          suffisantes et bien utilisées~; maintenir les risques dans
          l'appétence fixée par la gouvernance (Ch.~\ref{chap:alignement}).
    \item[Partie prenante] Groupe ou individu qui affecte les objectifs d'une
          organisation ou en est affecté (Ch.~\ref{chap:gouv-entreprise}).
    \item[Plan de retour arrière] Actions qui ramènent à l'état antérieur si
          un changement échoue, avec leurs critères et leur délai
          (Ch.~\ref{chap:itil-changements}).
    \item[Portefeuille] Ensemble des projets, programmes et services, géré
          comme un tout pour maximiser la valeur (Ch.~\ref{chap:alignement}).
    \item[Pratique (COBIT)] Action clé pour atteindre un objectif~; codée,
          par exemple DSS02.04 (Ch.~\ref{chap:cobit-objectifs}).
    \item[Prépublication] Version d'un travail rendue publique avant son
          évaluation par les pairs (Ch.~\ref{chap:recherche-mesure}).
    \item[Problème] Cause, réelle ou potentielle, d'un ou de plusieurs
          incidents (Ch.~\ref{chap:itil-incidents}).
    \item[Processus] Ensemble d'activités corrélées qui transforme des
          entrées en sorties, en vue d'un objectif (Ch.~\ref{chap:processus}).
    \item[Propriétaire de processus] Personne qui répond de la conception, de
          la performance et de l'amélioration d'un processus
          (Ch.~\ref{chap:processus}).
    \item[Ratio d'impact disproportionné] Taux de décisions favorables du
          groupe le moins favorisé divisé par celui du plus favorisé~; un
          ratio inférieur à~0,8 est un signal d'alerte
          (Ch.~\ref{chap:donnees-ia}).
    \item[Réalisation des bénéfices] Créer de la valeur grâce aux
          technologies, et éliminer ce qui n'en crée pas assez
          (Ch.~\ref{chap:alignement}).
    \item[Recueil de bonnes pratiques] Pratiques reconnues, présentées comme
          recommandations, que l'organisation adapte~; par exemple ITIL
          (Ch.~\ref{chap:referentiels}).
    \item[Redevabilité] Obligation de rendre compte d'une responsabilité
          reçue, et d'en assumer les conséquences (Ch.~\ref{chap:prologue}).
    \item[Relation d'agence] Contrat par lequel un principal délègue à un
          agent un pouvoir de décision (Ch.~\ref{chap:gouv-entreprise}).
    \item[RTO, RPO] Durée maximale de reprise d'un service~; perte de
          données maximale admissible, exprimée en durée
          (Ch.~\ref{chap:itil-niveaux}).
    \item[Service] Moyen de permettre la co-création de valeur, en
          facilitant les résultats que veut le client
          (Ch.~\ref{chap:itil-svs}).
    \item[Solution de contournement] Moyen de réduire l'effet d'un incident
          ou d'un problème sans en traiter la cause
          (Ch.~\ref{chap:itil-incidents}).
    \item[Source de première main] Document produit par une partie qui a eu
          accès direct aux faits et engage sa responsabilité sur leur
          exactitude (Ch.~\ref{chap:defaillances}).
    \item[Système de management] Ensemble d'éléments en interaction qui sert
          à établir des politiques et des objectifs, et à les atteindre
          (Ch.~\ref{chap:normes}).
    \item[Système de valeur des services] Selon ITIL~4, la manière dont les
          composantes et activités d'une organisation transforment une
          demande en valeur (Ch.~\ref{chap:itil-svs}).
    \item[Utilité] Ce que fait le service pour répondre à un besoin
          (Ch.~\ref{chap:itil-svs}).
    \item[Valeur] Bénéfices, utilité et importance perçus~; elle est perçue
          par celui qui la reçoit (Ch.~\ref{chap:itil-svs}).
    \item[Valeur actuelle nette (VAN)] Somme des flux d'un projet ramenés à
          aujourd'hui par un taux d'actualisation (Ch.~\ref{chap:alignement}).
    \item[Vérification des systèmes d'information] Selon le GUID~5100,
          examen des contrôles liés aux SI pour déceler des écarts à des
          critères (Ch.~\ref{chap:audit}).
\end{description}

    % ============================================================================
% Fichier : annexes/Ann_B_outils.tex
% Titre   : Référentiel des Outils
% Version : 1.0 -- novembre 2026
% ============================================================================
% Outils mobilisés dans les chapitres et les labs ; les noms commerciaux sont
% des exemples, à vérifier au moment du choix.
% ============================================================================

\chapter{Référentiel des Outils}
\label{ann:outils}

Cette annexe rassemble les outils de la gouvernance et de la gestion des
services mobilisés dans le cours et dans les labs de NDAP SUKULU. Les noms
de logiciels sont des exemples~: les offres évoluent, et un outil ne
remplace ni une décision ni une procédure.

\section*{Gouvernance et conception}

\begin{tableaucours}{L{3.6cm}L{5.4cm}L{3.2cm}}{%
    \tableauHeader{Outil} &
    \tableauHeader{Usage} &
    \tableauHeader{Chapitre, lab}
}
Cascade des objectifs (tableur) & Relier besoins des parties prenantes,
objectifs d'entreprise, d'alignement et objectifs COBIT &
Ch.~\ref{chap:cobit-objectifs}, GCI-LAB-03 \\
Canevas de conception & Appliquer les onze facteurs de conception et fixer
les capacités cibles & Ch.~\ref{chap:cobit-conception}, GCI-LAB-02 \\
Grille d'évaluation de capacité & Coter les activités d'un processus,
niveau par niveau & Ch.~\ref{chap:cobit-capacite}, GCI-LAB-04 \\
Tableau de portefeuille & VAN, coût total de possession, risque, alignement
& Ch.~\ref{chap:alignement}, GCI-LAB-01 \\
Matrice des droits de décision & Qui est consulté, qui décide, pour chaque
type de décision & Ch.~\ref{chap:alignement} \\
\end{tableaucours}

\section*{Gestion des services}

\begin{tableaucours}{L{3.6cm}L{5.4cm}L{3.2cm}}{%
    \tableauHeader{Outil} &
    \tableauHeader{Usage} &
    \tableauHeader{Chapitre, lab}
}
Outil de gestion des services & Tickets d'incidents et de demandes,
priorités, escalade, problèmes & Ch.~\ref{chap:itil-incidents}, GCI-LAB-06 \\
Registre des changements & Demande, évaluation du risque, autorisation,
calendrier, retour arrière & Ch.~\ref{chap:itil-changements}, GCI-LAB-07 \\
CMDB & Éléments de configuration et leurs relations, pour l'analyse
d'impact & Ch.~\ref{chap:itil-changements} \\
Supervision & Mesure de la disponibilité, alertes &
Ch.~\ref{chap:itil-niveaux}, GCI-LAB-08 \\
Matrice RACI & Un seul A par activité~; rôles et non personnes &
Ch.~\ref{chap:processus}, GCI-LAB-09 \\
\end{tableaucours}

\section*{Architecture, livraison et données}

\begin{tableaucours}{L{3.6cm}L{5.4cm}L{3.2cm}}{%
    \tableauHeader{Outil} &
    \tableauHeader{Usage} &
    \tableauHeader{Chapitre}
}
Cadre d'architecture d'entreprise (par exemple TOGAF) & Décrire processus,
données, applications et technologies, et leur évolution &
Ch.~\ref{chap:referentiels} \\
Chaîne de déploiement & Construire, tester et déployer automatiquement,
avec traçabilité & Ch.~\ref{chap:agile-cloud} \\
Catalogue et dictionnaire de données & Propriétaires, définitions,
qualité, conservation & Ch.~\ref{chap:donnees-ia} \\
\end{tableaucours}

\section*{Contrôle et audit}

\begin{tableaucours}{L{3.6cm}L{5.4cm}L{3.2cm}}{%
    \tableauHeader{Outil} &
    \tableauHeader{Usage} &
    \tableauHeader{Chapitre, lab}
}
Extractions et empreintes & Tester les contrôles généraux sur des données
tracées & Ch.~\ref{chap:audit}, GCI-LAB-05 \\
Tableur ou SQL & Rapprochements (départs et comptes actifs, déploiements et
changements autorisés) & Ch.~\ref{chap:audit},
\ref{chap:itil-changements} \\
\end{tableaucours}

\section*{Le bureau de NDAP SUKULU}

Les labs s'exécutent dans un bureau d'analyste virtuel. Les pièces sont
remises dans \texttt{\textasciitilde/preuves}, en lecture seule, avec leur
manifeste d'empreintes~; une copie de travail se trouve dans
\texttt{\textasciitilde/travail}. Le bureau propose un terminal et une suite
bureautique~; l'activité dans le terminal est journalisée.

    % ============================================================================
% Fichier : annexes/Ann_C_templates.tex
% Titre   : Modèles de Documents Opérationnels
% Version : 1.0 -- novembre 2026
% ============================================================================
% Modèles annoncés dans le mode d'emploi (charte, RACI, SLA, grille de
% capacité), complétés de la demande de changement, de la fiche
% d'investissement, du constat en cinq éléments et de la chronologie.
% Contenus repris des chapitres 1, 3, 8, 11, 14, 15 et 16.
% ============================================================================

\chapter{Modèles de Documents Opérationnels}
\label{ann:templates}

Chaque modèle reprend les rubriques présentées dans un chapitre. Ils sont
à adapter à la taille de l'organisation~: un modèle complet mal tenu vaut
moins qu'un modèle simple tenu à jour.

\section*{1. Charte d'un comité stratégique des SI (Chapitre~\ref{chap:gouv-entreprise})}

\begin{tableaucours}{L{3.2cm}L{9.0cm}}{%
    \tableauHeader{Rubrique} &
    \tableauHeader{Contenu}
}
Mission & Assister le conseil dans l'évaluation de la stratégie numérique,
de la valeur des investissements et des risques liés au SI \\
Composition & Au moins deux administrateurs, dont le président~; DG et DSI
membres~; audit interne invité permanent \\
Pouvoirs & Accès direct à la DSI et à l'audit interne~; droit de demander
une expertise externe \\
Fréquence & Au moins trimestrielle~; sur convocation pour un projet majeur
en difficulté \\
Quorum et décisions & Règle de quorum~; décisions inscrites au procès-verbal
avec un responsable et une échéance \\
Reddition de comptes & Rapport écrit au conseil après chaque séance \\
\end{tableaucours}

\section*{2. Fiche d'investissement (Chapitre~\ref{chap:alignement})}

\begin{tableaucours}{L{3.2cm}L{9.0cm}}{%
    \tableauHeader{Rubrique} &
    \tableauHeader{Contenu}
}
Objectif & Objectif d'entreprise servi, et objectif d'alignement \\
Coût & Coût total de possession sur la durée de vie, par poste \\
Bénéfices & Bénéfices attendus, datés, avec leur indicateur et leur
propriétaire \\
Rentabilité & Valeur actuelle nette, taux d'actualisation retenu \\
Risques & Risques du projet et risques de ne pas le faire \\
Décision & Organe qui décide, date, conditions~; revue des bénéfices
prévue \\
\end{tableaucours}

\section*{3. Grille d'évaluation de la capacité d'un processus}

\begin{tableaucours}{L{2.8cm}cL{4.4cm}L{3.4cm}}{%
    \tableauHeader{Pratique} &
    \tableauHeader{Niv.} &
    \tableauHeader{Activité} &
    \tableauHeader{Preuve et cote}
}
BAI06.01 & 2 & \ldots & Pièce~; note~: complètement, largement,
partiellement, pas \\
BAI06.01 & 3 & \ldots & \ldots \\
\ldots & \ldots & \ldots & \ldots \\
\end{tableaucours}

\noindent Notes~: complètement (plus de 85~\%), largement (50 à 85~\%),
partiellement (15 à 50~\%), pas (moins de 15~\%). Un niveau est atteint
s'il est noté au moins « largement » et si tous les niveaux inférieurs sont
notés « complètement »~; la capacité retenue est celle que les preuves
établissent, non celle que l'on déclare (Chapitre~\ref{chap:cobit-capacite}).

\section*{4. Demande de changement (Chapitre~\ref{chap:itil-changements})}

\begin{tableaucours}{L{3.2cm}L{9.0cm}}{%
    \tableauHeader{Rubrique} &
    \tableauHeader{Contenu}
}
Description & Ce qui change, pourquoi, éléments de configuration concernés \\
Type & Standard, normal ou urgent \\
Risque & Impact et probabilité~; services et clients exposés (analyse
d'impact sur la CMDB) \\
Tests & Recette réalisée, taux, anomalies bloquantes closes \\
Mise en production & Date dans le calendrier, période de gel, équipes
mobilisées \\
Retour arrière & Plan testé, critères de déclenchement, heure limite \\
Autorisation & Autorité de changement compétente, décision motivée \\
Revue & Résultat constaté après la mise en production \\
\end{tableaucours}

\section*{5. Accord de niveau de service (Chapitre~\ref{chap:itil-niveaux})}

\begin{tableaucours}{L{3.2cm}L{9.0cm}}{%
    \tableauHeader{Rubrique} &
    \tableauHeader{Contenu}
}
Description du service & Ce que le service permet au client d'obtenir \\
Horaires & Périodes où les engagements s'appliquent \\
Disponibilité & Part du temps de service convenu où le service est
utilisable \\
Délais de rétablissement & Par priorité \\
Indicateur de résultat & Ce que le client vit, défini avec lui \\
Mesure et compte rendu & Source des données, fréquence \\
Pénalités & Conséquences du non-respect, et leur application \\
Revue & Fréquence de révision, comité de service \\
\end{tableaucours}

\section*{6. Matrice RACI}

\begin{tableaucours}{L{4.0cm}ccccc}{%
    \tableauHeader{Activité} &
    \tableauHeader{Rôle 1} & \tableauHeader{Rôle 2} &
    \tableauHeader{Rôle 3} & \tableauHeader{Rôle 4} &
    \tableauHeader{Rôle 5}
}
Activité 1 & A & R & C & I & \\
Activité 2 & & A & R & C & I \\
\ldots & & & & & \\
\end{tableaucours}

\noindent Règles (Chapitre~\ref{chap:processus})~: un seul A par activité~;
au moins un R~; le A au bon niveau~; peu de C~; des fonctions et non des
personnes.

\section*{7. Constat d'audit en cinq éléments (Chapitre~\ref{chap:audit})}

\begin{tableaucours}{L{3.2cm}L{9.0cm}}{%
    \tableauHeader{Élément} &
    \tableauHeader{Question}
}
Constatation & Que se passe-t-il réellement~? \\
Critère & Que devrait-il se passer~? Avec sa source \\
Cause & Pourquoi l'écart existe-t-il~? \\
Conséquence & Quel risque, ou quel effet~? \\
Recommandation & Que faut-il faire, par qui, pour quand~? \\
\end{tableaucours}

\section*{8. Chronologie de défaillance}

\begin{tableaucours}{L{2.6cm}L{6.4cm}L{3.2cm}}{%
    \tableauHeader{Date} &
    \tableauHeader{Fait établi} &
    \tableauHeader{Pièce}
}
\ldots & Fait, décision, ou absence de décision & Procès-verbal, courriel,
journal \\
\end{tableaucours}

\noindent La chronologie ne contient que des constatations~; l'analyse vient
ensuite, avec son niveau de confiance (Chapitre~\ref{chap:prologue}).

    % ============================================================================
% Fichier : annexes/Ann_D_textes.tex
% Titre   : Textes Législatifs et Réglementaires
% Version : 1.0 -- novembre 2026
% ============================================================================
% Extraits identiques à ceux des chapitres (prologue, C01, C03, C04, C07,
% C10). Loi n° 2010/012 : version anglaise officielle. COBAC : relevé en
% ligne, à collationner. AUSCGIE, loi n° 2018/012, loi n° 2024/017 :
% présentés d'après des sources secondaires, à vérifier.
% ============================================================================

\chapter{Textes Législatifs et Réglementaires}
\label{ann:textes}

Cette annexe rassemble les textes cités dans le cours, avec leur degré de
vérification. Seuls les extraits lus dans leur version officielle sont
reproduits entre guillemets~; les autres sont présentés sans citation.

\section*{Degré de vérification}

\begin{tableaucours}{L{4.6cm}L{3.8cm}L{3.8cm}}{%
    \tableauHeader{Texte} &
    \tableauHeader{Vérification} &
    \tableauHeader{Chapitres}
}
Loi n°~2010/012 (cybersécurité) & Version anglaise officielle lue &
Prologue, \ref{chap:referentiels}, \ref{chap:cobit-conception},
\ref{chap:droit} \\
Règlement COBAC R-2016/04 & Art.~1 et~2 relevés en ligne~; à collationner &
\ref{chap:droit} \\
AUSCGIE révisé (2014), art.~829-1 & D'après la présentation de la réforme
par l'OHADA & \ref{chap:gouv-entreprise}, \ref{chap:droit} \\
Loi n°~2018/012 (régime financier), art.~69 & D'après la présentation du
MINFI~; à vérifier au Journal officiel & \ref{chap:alignement} \\
Loi n°~2024/017 (données personnelles) & D'après les présentations
publiées~; à vérifier au Journal officiel & \ref{chap:cobit-conception},
\ref{chap:droit} \\
Convention de Malabo (UA, 2014) & Entrée en vigueur vérifiée~; statut du
Cameroun non établi & \ref{chap:droit} \\
\end{tableaucours}

\section*{Loi n° 2010/012 relative à la cybersécurité et à la cybercriminalité}

\begin{boxjuris}
\textbf{Article~7, alinéa~2}~: \textit{« The Agency \textelp{} shall be
responsible for the regulation, control and monitoring of activities related
to the security of electronic communication networks, information systems,
and electronic certification on behalf of the State. »}

\medskip
\textbf{Article~11} (extrait)~: \textit{« The following activities may be
subject to authorization: \textelp{} the provision of security auditing,
security programs editing, and other authorized security services to the
public. »}

\medskip
\textbf{Article~26, alinéa~1}~: \textit{« Operators of information systems
shall take every technical and administrative measure to ensure the security
of services offered. To this end, they shall have standardized systems
enabling them to at all times identify, assess, process or manage any risk
relating to the security of the information systems of the services
provided directly or indirectly. »}

\medskip
\textbf{Article~29, alinéa~1}~: \textit{« Operators of information systems
shall be bound to conserve the connection and traffic data of their
information systems for a period of 10 (ten) years. »}

\medskip
\textbf{Article~30, alinéa~1}~: \textit{« Operators of information systems
shall assess and revise their security systems and, where necessary, make
the appropriate modifications to their security practices, measures and
techniques according to technological change. »}

\medskip
\textbf{Article~31, alinéa~1}~: \textit{« Electronic communication networks
and information systems content providers shall be bound to ensure the
availability of material, as well as the data stored in their
installations. »}

\medskip
\textbf{Article~32, alinéas~1 et~2}~: \textit{« (1) Electronic
communication networks and information systems shall be subject to a regime
of compulsory and periodic auditing of their security systems by the
Agency. (2) Security audit and severity scale rating shall be undertaken
each year or as required by the prevailing circumstances. »}

\medskip
\textbf{Articles~41 et~42}~: \textit{« 41. Every individual shall have the
right to the protection of their privacy. \textelp{} 42. The confidentiality
of information channelled through electronic communication and information
systems networks, including traffic data, shall be ensured by operators of
electronic communication and networks information systems. »}

\medskip
\textbf{Article~64, alinéas~1 à~3}~: \textit{« (1) Corporate bodies shall be
criminally liable for offences committed on their account by their
management structures. (2) The criminal liability of corporate bodies shall
not preclude that of natural persons who commit such offences or are
accomplices. (3) The penalties to be meted out on defaulting corporate
bodies shall be fines of from 5~000~000 (five million) to 50~000~000 (fifty
million) CFA francs. »}
\end{boxjuris}

L'article~13 soumet en outre les systèmes d'information des opérateurs à un
audit de sécurité obligatoire, dont les modalités sont fixées par voie
réglementaire~; il est présenté sans citation (Prologue).

\section*{Règlement COBAC R-2016/04}

\begin{boxjuris}
\textbf{Articles~1 et~2} (texte relevé sur une base réglementaire en ligne,
à collationner avec la version du SG-COBAC)~: article~1~: \textit{« Le
présent règlement fixe les exigences minimales relatives au système de
contrôle interne que les établissements assujettis doivent mettre en
œuvre. »}~; article~2, risque opérationnel~: \textit{« risque de pertes
résultant de l'inadéquation ou de la défaillance des processus internes, des
personnes et des systèmes, ou d'évènements externes. »}
\end{boxjuris}

\section*{Textes présentés sans citation}

\begin{itemize}
    \item \textbf{AUSCGIE révisé, article~829-1}~: selon la présentation de
          la réforme, comité d'audit obligatoire dans les sociétés cotées,
          chargé notamment du suivi de l'efficacité du contrôle interne et
          de la gestion des risques (Chapitre~\ref{chap:gouv-entreprise}).
          Responsabilité des dirigeants~: articles~161 et suivants, 740,
          d'après la doctrine (Chapitre~\ref{chap:droit}).
    \item \textbf{Loi n°~2018/012, article~69}~: selon le MINFI, rôle du
          responsable de programme dans le budget par programmes
          (Chapitre~\ref{chap:alignement}).
    \item \textbf{Loi n°~2024/017}~: protection des données à caractère
          personnel~; autorité de protection, encadrement des transferts,
          délai de mise en conformité de dix-huit mois, d'après les
          présentations publiées (Chapitres~\ref{chap:cobit-conception}
          et~\ref{chap:droit}).
    \item \textbf{Convention de Malabo}~: cybersécurité et protection des
          données, en vigueur depuis le 8~juin~2023~; ratification par le
          Cameroun non établie (Chapitre~\ref{chap:droit}).
\end{itemize}

    % ============================================================================
% Fichier : annexes/Ann_E_bibliographie.tex
% Titre   : Bibliographie Commentée
% Version : 1.0 -- novembre 2026
% ============================================================================
% Uniquement des entrées de bibliographie/GCI_references.bib.
% ============================================================================

\chapter{Bibliographie Commentée}
\label{ann:biblio}

Pour chaque partie, quelques lectures, avec ce qu'elles apportent. Les
références complètes figurent dans la bibliographie du livre. Les
publications d'ISACA et d'AXELOS sont paraphrasées dans le cours~; leur
lecture directe reste indispensable.

\section*{Prologue et Partie I · Fondements}

\begin{itemize}
    \item Weill et Ross~\cite{weill2004itgovernance}~: les droits de décision
          et les archétypes de gouvernance~; le livre fondateur.
    \item Jensen et Meckling~\cite{jensen1976agency}, Freeman~\cite{freeman1984stakeholder},
          Davis et al.~\cite{davis1997stewardship}~: agence, parties
          prenantes, intendance~; les théories sous la gouvernance.
    \item Cadbury~\cite{cadbury1992}, principes de l'OCDE~\cite{oecd2023principles},
          King~IV~\cite{iod2016king4}~: les codes de gouvernance d'entreprise.
    \item Henderson et Venkatraman~\cite{henderson1993alignment}~: le modèle
          d'alignement stratégique.
    \item Les sources de première main sur les défaillances~: SEC sur Knight
          Capital~\cite{sec2013knight}, FCA et PRA sur TSB~\cite{fca2022tsb},
          analyses de l'incident CrowdStrike~\cite{crowdstrike2024rca,microsoft2024crowdstrike}.
    \item Thorp~\cite{thorp1998paradox} et Weill et Broadbent~\cite{weill1998infrastructure}~:
          valeur et portefeuille.
\end{itemize}

\section*{Partie II · COBIT 2019}

\begin{itemize}
    \item ISACA, \emph{Introduction and Methodology}~\cite{isaca2018cobitintro}~:
          principes, composantes, facteurs de conception, capacité.
    \item ISACA, \emph{Governance and Management Objectives}~\cite{isaca2018cobitobjectifs}~:
          les quarante objectifs, à consulter objectif par objectif.
    \item COBIT~4.1~\cite{itgi2007cobit41}~: encore cité dans beaucoup de
          dossiers~; les modèles de maturité.
    \item De Haes et Van Grembergen~\cite{dehaes2015egit}~: la gouvernance
          des SI d'entreprise, entre recherche et pratique.
\end{itemize}

\section*{Partie III · Cadre normatif, juridique et déontologique}

\begin{itemize}
    \item ISO/IEC~38500~\cite{iso38500,iso38500_2024} et
          ISO/IEC~27000~\cite{iso27000_2018}.
    \item COSO ERM~\cite{coso2004ermfr,coso2017erm}~: la gestion des risques
          de l'entreprise.
    \item IIA~: CRIPP~\cite{iia2017cripp}, Normes mondiales~\cite{iia2024gias},
          modèle des trois lignes~\cite{iia2020troislignes}~; INTOSAI,
          GUID~5100~\cite{intosai2019guid5100}.
    \item Textes~: loi n°~2010/012~\cite{loi2010012}, COBAC
          R-2016/04~\cite{cobac2016r04}, AUSCGIE~\cite{ohada2014auscgie},
          régime financier de l'État~\cite{minfi2023regimefinancier},
          Sarbanes-Oxley~\cite{sox2002}~; voir l'Annexe~\ref{ann:textes}.
\end{itemize}

\section*{Partie IV · ITIL 4 et pratique opérationnelle}

\begin{itemize}
    \item AXELOS, ITIL~4~\cite{axelos2019itil4}~: le système de valeur des
          services et les pratiques.
    \item PeopleCert, ITIL (2026)~\cite{peoplecert2026itil5}~: l'évolution
          du référentiel, à suivre.
\end{itemize}

\section*{Partie V · Défis avancés}

\begin{itemize}
    \item Manifeste agile~\cite{agile2001manifeste} et Forsgren et
          al.~\cite{forsgren2018accelerate}~: agilité et performance de la
          livraison.
    \item NIST, définition du nuage~\cite{nist2011cloud}.
    \item NIST AI RMF~\cite{nist2023airmf}, Recommandation de
          l'UNESCO~\cite{unesco2021ethique}, règlement européen sur
          l'IA~\cite{ue2024aiact}~: gouverner l'IA.
\end{itemize}

\section*{Partie VI · Recherche}

\begin{itemize}
    \item Sambamurthy et Zmud~\cite{sambamurthy1999arrangements}, Peterson~\cite{peterson2004crafting},
          De Haes et Van Grembergen~\cite{dehaes2009exploratory}, Brown et
          Grant~\cite{brown2005framing}, Wu et al.~\cite{wu2015mechanisms}~:
          les arrangements et mécanismes de gouvernance.
    \item Becker et al.~\cite{becker2009maturity}, Pöppelbuß et
          Röglinger~\cite{poppelbuss2011maturity}, Simonsson et
          al.~\cite{simonsson2010maturity}~: concevoir et éprouver un modèle
          de maturité.
    \item Hevner et al.~\cite{hevner2004design}~: la science du design en
          systèmes d'information.
\end{itemize}

    % ============================================================================
% Fichier : annexes/Ann_F_contacts.tex
% Titre   : Contacts et Ressources
% Version : 1.0 -- novembre 2026
% ============================================================================
% Adresses de sites institutionnels données à titre indicatif, à vérifier
% avant diffusion. Contact de NDAP SUKULU : celui fixé par l'auteur.
% ============================================================================

\chapter{Contacts et Ressources}
\label{ann:contacts}

Les adresses ci-dessous sont celles des sites institutionnels à la date de
rédaction~; elles peuvent changer et sont à vérifier avant usage.

\section*{Autorités camerounaises et régionales}

\begin{tableaucours}{L{3.6cm}L{4.6cm}L{4.0cm}}{%
    \tableauHeader{Institution} &
    \tableauHeader{Rôle dans le cours} &
    \tableauHeader{Site}
}
ANTIC & Sécurité des SI, audits de sécurité (loi n°~2010/012) &
\texttt{www.antic.cm} \\
ART & Régulation des communications électroniques & \texttt{www.art.cm} \\
MINPOSTEL & Politique du secteur~; destinataire des rapports d'audit de
sécurité & \texttt{www.minpostel.gov.cm} \\
COBAC (Secrétariat général) & Contrôle interne et risques des
établissements de crédit de la CEMAC & \texttt{www.sgcobac.org} \\
OHADA & Droit des sociétés, responsabilité des dirigeants, comité
d'audit & \texttt{www.ohada.org} \\
Autorité de protection des données & Loi n°~2024/017 & · \\
Chambre des comptes de la Cour suprême & Contrôle des finances publiques
& · \\
Union africaine & Convention de Malabo & \texttt{au.int} \\
\end{tableaucours}

\section*{Organisations professionnelles et normatives}

\begin{tableaucours}{L{3.6cm}L{4.6cm}L{4.0cm}}{%
    \tableauHeader{Organisation} &
    \tableauHeader{Ressources} &
    \tableauHeader{Site}
}
ISACA & COBIT, certifications (CISA, CGEIT) & \texttt{www.isaca.org} \\
PeopleCert (AXELOS) & ITIL et ses certifications & \texttt{www.peoplecert.org} \\
ISO & ISO/IEC~38500, 20000-1, 27001, 27002, 31000 & \texttt{www.iso.org} \\
IIA & Normes mondiales d'audit interne, modèle des trois lignes &
\texttt{www.theiia.org} \\
COSO & Contrôle interne, gestion des risques & \texttt{www.coso.org} \\
INTOSAI & ISSAI, GUID~5100 & \texttt{www.intosai.org} \\
NIST & Définition du nuage, AI RMF & \texttt{www.nist.gov} \\
UNESCO & Recommandation sur l'éthique de l'IA & \texttt{www.unesco.org} \\
\end{tableaucours}

\section*{NDAP SUKULU}

Les fiches d'auto-évaluation et les dix labs du cours sont disponibles sur
NDAP SUKULU, sous le code \codeUE. Leur présentation et le tableau des labs
figurent au début du livre, dans « NDAP SUKULU, la salle de classe compagnon ».
NDAP SUKULU est une propriété de MaletYon, mise à la disposition des
étudiants et des praticiens souhaitant se perfectionner~:
\texttt{maletyon@proton.me}.


% --------------------------------------------------------------------------
% ÉLÉMENTS FINAUX
% --------------------------------------------------------------------------
\backmatter

    \nocite{*}
    \printbibliography[heading=bibintoc, title={Références bibliographiques}]

    \printindex

% ==============================================================================
% 4ÈME DE COUVERTURE
% ==============================================================================
\clearpage
\thispagestyle{empty}
\pagecolor{CouleurPrimaire}

\begin{tikzpicture}[remember picture, overlay]
    \fill[CouleurAccent]
        (current page.south west) rectangle
        ([yshift=1.8cm] current page.south east);
    \fill[CouleurAccent]
        ([yshift=-1.4cm] current page.north west) rectangle
        ([yshift=-1.7cm] current page.north east);
\end{tikzpicture}

\vspace*{2.5cm}

\begin{center}
    {\color{CouleurAccent}\rule{0.6\linewidth}{1pt}}\\[1.5ex]
    {\color{white}\Large\bfseries À propos de ce cours}\\[1.5ex]
    {\color{CouleurAccent}\rule{0.6\linewidth}{1pt}}
\end{center}

\vspace{1cm}

\begin{center}
\begin{minipage}{0.75\textwidth}
\color{white!85!gray}\setstretch{1.4}
Quand un système d'information s'arrête, la question décisive n'est pas
technique : elle porte sur qui avait décidé, et qui devait rendre compte.
Ce cours apprend à y répondre avant la panne.

\medskip

Il progresse en six parties, précédées d'un prologue : les fondements de la
gouvernance des systèmes d'information, le référentiel COBIT~2019, le cadre
normatif et juridique, ITIL~4 et les processus de gestion des services, les
défis de l'agilité, du nuage et de l'intelligence artificielle, et enfin la
recherche.

\medskip

Chaque notion est ancrée à la fois dans les standards internationaux (ISO,
ISACA, IIA, INTOSAI) et dans le droit camerounais et communautaire (CEMAC,
OHADA). Le cours s'adresse aux étudiants du Master Professionnel en
Management des Systèmes d'Information et aux professionnels qui exercent comme
DSI, responsables GRC, auditeurs des SI, responsables des services
informatiques ou RSSI.
\end{minipage}
\end{center}

\vfill

\begin{center}
    {\color{CouleurAccent}\rule{0.7\linewidth}{1.2pt}}\\[1.2ex]
    {\color{white}\large\bfseries MINKA MI NGUIDJOÏ Thierry Emmanuel}\\[0.4ex]
    {\color{white!80!gray}\small
        CISA~·~CISM~·~CGEIT~·~CRISC~·~CDPSE~·~ISO 27001~·~ISO 22301~·~ISO 9001}\\[0.3ex]
    {\color{white!70!gray}\footnotesize\itshape
        Chercheur au laboratoire d'Ingénierie Mathématiques et des Systèmes d'Information}\\[0.3ex]
    {\color{white!70!gray}\footnotesize\itshape
        Chercheur~|~Cryptologue~|~Auditeur Senior~|~Expert Judiciaire~|~GRC \& Sécurité de l'Information}\\[1.2ex]
    {\color{CouleurAccent}\rule{0.7\linewidth}{1.2pt}}
\end{center}

\vspace{1cm}

\afterpage{\nopagecolor}
\end{document}
